\documentclass{report}%
\usepackage[T1]{fontenc}%
\usepackage[utf8]{inputenc}%
\usepackage{lmodern}%
\usepackage{textcomp}%
\usepackage{lastpage}%
\usepackage{geometry}%
\geometry{a4paper=True,margin=25mm}%
\usepackage{amsmath}%
\usepackage{amssymb}%
\usepackage{amsfonts}%
\usepackage{mathtools}%
\usepackage{bm}%
\usepackage{physics}%
\usepackage{graphicx}%
\usepackage{microtype}%
\usepackage{iftex}%
\usepackage{listings}%
\usepackage{listingsutf8}%
\usepackage{xcolor}%
\usepackage[strings]{underscore}%
\usepackage{url}%
\usepackage{xurl}%
\usepackage{hyperref}%
\usepackage{bookmark}%
\usepackage[square,numbers]{natbib}%
%
\ifLuaTeX\usepackage{fontspec}\usepackage{unicode-math}\else\usepackage[utf8]{inputenc}\usepackage[T1]{fontenc}\fi\usepackage[czech]{babel}%
\graphicspath{{figures/}}%
\renewcommand{\bibname}{Použitá literatura}%
\renewcommand{\bibsection}{\chapter*{\bibname}\addcontentsline{toc}{chapter}{\bibname}}%
% Robust LaTeX settings for AutoGenBook (LuaLaTeX-friendly).
\ProvidesFile{robustness.tex}[2026/01/01 Robust LaTeX settings]

\RequirePackage{xparse}
\RequirePackage{fvextra}
\RequirePackage{enumitem}
\RequirePackage{adjustbox}
\RequirePackage{tabularx}
\RequirePackage{ltablex}
\keepXColumns

% Fallback for unknown environments from conversions.
\ProvideDocumentEnvironment{para}{+b}{\par\smallskip\noindent#1\par\smallskip}{}

% Prompt/code block (breaks anywhere to avoid overfull boxes).
% Prompt/code block with safe line breaking.
\DefineVerbatimEnvironment{PromptBlock}{Verbatim}{
  formatcom=\ttfamily\small,
  breaklines=true,
  breakanywhere=true,
  breaksymbolleft=,
  breaksymbolright=
}

% Make plain verbatim blocks breakable as well.
\DefineVerbatimEnvironment{verbatim}{Verbatim}{
  formatcom=\ttfamily\small,
  breaklines=true,
  breakanywhere=true,
  breaksymbolleft=,
  breaksymbolright=
}

% Listings defaults (robust line breaking).
\lstset{
  basicstyle=\ttfamily\small,
  breaklines=true,
  breakatwhitespace=false,
  columns=fullflexible,
  keepspaces=true,
  showstringspaces=false
}

% Description list labels should wrap to next line.
\setlist[description]{style=nextline,leftmargin=0pt,labelsep=0.5em}

% Safer line breaking for prose.
\emergencystretch=3em
\tolerance=4000
\pretolerance=1000
\hfuzz=0.2pt

% Better URL breaks.
\Urlmuskip=0mu plus 2mu

% Images never exceed linewidth.
\makeatletter
\def\maxwidth{%
  \ifdim\Gin@nat@width>\linewidth\linewidth\else\Gin@nat@width\fi}
\setkeys{Gin}{width=\maxwidth,keepaspectratio}
\makeatother

% Fit-width wrapper for wide tables.
\NewDocumentEnvironment{fitwidth}{+b}{\begin{adjustbox}{max width=\linewidth}#1\end{adjustbox}}{}

% LuaLaTeX fonts and Unicode fallbacks.
\ifLuaTeX
  \defaultfontfeatures{Ligatures=TeX}
  \newcommand{\AutogenMainFont}{Latin Modern Roman}
  \IfFontExistsTF{Latin Modern Roman}{}{\renewcommand{\AutogenMainFont}{TeX Gyre Termes}}
  \newcommand{\AutogenSansFont}{Latin Modern Sans}
  \IfFontExistsTF{Latin Modern Sans}{}{\renewcommand{\AutogenSansFont}{TeX Gyre Heros}}
  \newcommand{\AutogenMonoFont}{DejaVu Sans Mono}
  \newcommand{\AutogenMonoBold}{DejaVu Sans Mono Bold}
  \newcommand{\AutogenMonoItalic}{DejaVu Sans Mono Oblique}
  \newcommand{\AutogenMonoBoldItalic}{DejaVu Sans Mono Bold Oblique}
  \IfFontExistsTF{DejaVu Sans Mono}{}{
    \renewcommand{\AutogenMonoFont}{Latin Modern Mono}
    \renewcommand{\AutogenMonoBold}{Latin Modern Mono Bold}
    \renewcommand{\AutogenMonoItalic}{Latin Modern Mono Slanted}
    \renewcommand{\AutogenMonoBoldItalic}{Latin Modern Mono Bold Slanted}
  }

  \setmainfont{\AutogenMainFont}
  \setsansfont{\AutogenSansFont}
  \IfFontExistsTF{\AutogenMonoBold}{
    \setmonofont{\AutogenMonoFont}[
      BoldFont=\AutogenMonoBold,
      ItalicFont=\AutogenMonoItalic,
      BoldItalicFont=\AutogenMonoBoldItalic,
      Scale=MatchLowercase
    ]
  }{
    \setmonofont{\AutogenMonoFont}[Scale=MatchLowercase]
  }
\fi
%
\title{METODICKÁ PŘÍRUČKA "AI VE VÝUCE VUT"}%
\author{Petr Koňas, OpenAI{-}GPT{-}5{-}mini}%
\date{\today}%
%
\begin{document}%
\normalsize%
\lstset{backgroundcolor=\color[gray]{0.95},basicstyle=\ttfamily\small,breaklines=true,breakatwhitespace=false,columns=fullflexible,keepspaces=true,frame=single,numbers=left,numberstyle=\tiny,tabsize=4,inputencoding=utf8,extendedchars=true,showstringspaces=false}%
\hypersetup{pageanchor=false}%
\pagenumbering{roman}%
\maketitle%
\tableofcontents%
\clearpage%
\hypersetup{pageanchor=true}%
\pagenumbering{arabic}%
\chapter{ÚVOD DO AI PRO PEDAGOGY}%
Kapitola uvádí pedagogické publikum do základů generativní AI: definuje umělou inteligenci a generativní AI, popisuje rozdíl mezi tradičními AI přístupy a moderními LLM a vysvětluje základní principy fungování (tokens, embeddings, training data). Nabízí praktické příklady toho, co může pedagog od nástrojů reálně očekávat, a jak bezpečně začít (halucinace, bias, cut‑off dat). Obsahuje krátké technické vysvětlení LLM, rychlý start s nástroji (ChatGPT, Claude, Gemini, Microsoft Copilot) a návrhy cvičení první interakce ve třídě. Zahrnuje odkazy na tutoriály a další studijní zdroje.%
\section{Co je AI a generativní AI}%
Stručné pedagogické uvedení: základní definice umělé inteligence a generativní AI, proč je téma relevantní pro výuku a jak ho rámovat vůči studentům. Zaměření na srozumitelné vysvětlení bez hloubkových technických detailů.%


%
Krátké uvedení (co si pedagogové potřebují pamatovat)

Obecná definice
- Umělá inteligence (AI) jsou metody a nástroje, které umožňují strojům vykonávat úlohy spojené s lidskou inteligencí (rozpoznávání vzorů, rozhodování, generování textu).
- Generativní AI vytváří nový obsah — text, obrázky, kód nebo multimodální výstupy — podle instrukcí nebo vzorů.
- Omezení: modely odrážejí tréninková data, mohou opakovat chyby či zkreslení a neznamenají „porozumění“.
- Ve výuce rozlišujte mezi simulací chování a skutečným učením z dat; mnoho demonstračních aktivit ilustruje principy, nikoli autonomní učení.

Proč je téma relevantní pro výuku
- Generativní AI zrychluje přípravu materiálů, testů a vizuálů, usnadňuje diferenciaci a interaktivní výuku \cite{kb_ai_101_tipu_pdf_04e4a115_page_36_1}.
- Hry a vzdělávací světy (např. Minecraft Education) umožňují žákům experimentovat s principy ML a AI bez náročné teorie \cite{kb_ai_101_tipu_pdf_04e4a115_page_15_2}.
- Praktické přínosy: úspora času, rychlá personalizace úloh a tvorba variant.
- Rizika: nutnost ověřování faktů, ochrana osobních údajů, prevence plagiátorství a zachování pedagogického záměru úloh.

Praktické příklady pro rychlý start ve třídě

1) Demonstrace principů AI pomocí Minecraft Education
- Požadavky: počítače/tablety s Minecraft Education; dostupné Hour of Code lekce ilustrují řízení agentů a principy učení \cite{kb_ai_101_tipu_pdf_04e4a115_page_15_2}.
- Krátký postup:
  1. Připravte zadání (např. „navrhněte agenta, který rozpozná překážky“).
  2. Spusťte lekci Hour of Code, nechte žáky upravovat pravidla a pozorovat chování agenta.
  3. Diskutujte rozdíl mezi pravidly/simulací a učením z dat.
- Doba: 45–90 minut podle věku a zkušeností.
- Cíle: pochopení senzory/akce, základní představa o učení přes zpětnou vazbu.
- Pozn.: některé světy jsou součástí školních verzí Microsoft 365; část Hour of Code je zdarma \cite{kb_ai_101_tipu_pdf_04e4a115_page_15_2}.

2) Rychlé vytvoření kvízu pomocí Microsoft 365 Copilot + Forms
- Požadavky: přístup k Microsoft 365 s Copilotem nebo Copilot Chat (některé funkce i ve web. Office 365) \cite{kb_ai_101_tipu_pdf_04e4a115_page_3_1}.
- Workflow:
  1. Otevřete Microsoft Forms a aktivujte Copilot.
  2. Zadejte prompt (např. „Vytvoř test 10 otázek — 4x ABC, 4x pravda/nepravda, 2 otevřené“).
  3. Zkontrolujte a upravte otázky, ověřte správnost a úroveň před sdílením \cite{kb_ai_101_tipu_pdf_04e4a115_page_36_1}.
- Tipy: přidejte instrukce (úroveň třídy, čas, styl), validujte obsah.
- Čas: vytvoření a kontrola obvykle 10–30 minut.

Krátké pedagogické zásady (praktické minimum)
- Transparentnost: informujte studenty, když používáte AI k tvorbě materiálů nebo hodnocení.
- Validace: vždy ověřte výstupy lidskou kontrolou — modely mohou být nepřesné.
- Didaktika: používejte AI jako pomocníka (varianty úloh, nápověda), ne jako náhradu studentské práce.
- Ochrana dat: anonymizujte školní údaje při použití online služeb a dodržujte GDPR a školní pravidla.
- Hodnocení: navrhujte úkoly zaměřené na proces a porozumění, aby se omezilo povrchní kopírování generovaných odpovědí.

Rychlé odkazy a další kroky
- Vyzkoušejte okamžitě: stáhněte Minecraft Education a najděte Hour of Code aktivity; pro kvízy použijte Copilot + Forms s ukázkovým promptem \cite{kb_ai_101_tipu_pdf_04e4a115_page_15_2, kb_ai_101_tipu_pdf_04e4a115_page_36_1}.
- Začněte pilotně: vyzkoušejte jednu lekci, vyhodnoťte zpětnou vazbu a upravte postupy před širším nasazením.
- Podrobnosti o ověřování, etice, GDPR a integraci do sylabů jsou rozpracovány v dalších kapitolách této příručky.%

%
\section{Tradiční přístupy v AI versus moderní LLM}%
Srovnání klasických přístupů (např. pravidlové systémy, klasické machine learningové modely) a moderních velkých jazykových modelů; pedagogické dopady rozdílů (jaké typy úloh který přístup zvládá).%


%
Tato sekce krátce porovnává tři rodiny přístupů v AI, které pedagogům pomohou zvolit vhodnou didaktiku a aktivity: pravidlové systémy, klasické strojové učení a moderní velké jazykové modely (LLM). Následující technické odlišení a pedagogické důsledky jsou uvedeny stručně tak, aby je bylo možné přímo aplikovat při návrhu výuky.

Krátké technické rozlišení (general background knowledge)
\begin{itemize}
  \item Pravidlové systémy (rule‑based): chovají se podle explicitně definovaných pravidel a rozhodovacích stromů; jsou relativně předvídatelné a snadno vysvětlitelné, ale křehké vůči nepředvídaným vstupům.
  \item Klasické strojové učení (supervised/feature‑based): modely se učí z označených dat, jejich chování závisí na zvolené reprezentaci (features) a dostupnosti trénovacích příkladů; vhodné pro úlohy s dobře strukturovanými vstupy a jasnými cíli.
  \item Velké jazykové modely (LLM, generativní modely): škálované modely předtrénované na rozsáhlých textech, které provádějí predikci tokenů a generují plynulý text či jiné multimodální výstupy; silné v generování a zobecnění vzorů, zároveň méně přímo vysvětlitelné. (general background knowledge)
\end{itemize}

Pedagogické důsledky — co použít k čemu
\begin{itemize}
  \item Kdy preferovat pravidlové systémy: pokud chcete demonstrovat deterministické chování, kontrolovatelné rozhodovací cesty nebo stavový stroj (např. logika herní mechaniky, kontrola vstupů), pravidla jsou vhodná pro názorné ukázky a cvičení zaměřená na interpretaci.
  \item Kdy preferovat klasické ML: při úlohách, kde existují dobře očíslovaná data a jasná metrika (např. klasifikace senzorických dat, předpovědi založené na měřeních), je vhodné ukázat proces sběru dat, tvorbu featureů a evaluaci modelu.
  \item Kdy využít LLM a generativní AI: pro tvorbu textových materiálů, návrh variant úloh, generování vysvětlení v různých úrovních složitosti nebo rychlé vytvoření zadání a kvízů — LLM zvyšují produktivitu přípravy a umožňují experimenty s adaptivními odpověďmi. Konkrétně: rychlé vytvoření kvízu s pomocí Microsoft 365 Copilot + Forms je praktický příklad využití AI pro tvorbu testů a úloh \cite{kb_ai_101_tipu_pdf_04e4a115_page_8_1}.
\end{itemize}

Praktické ukázky a aktivita do výuky
\begin{enumerate}
  \item Demonstrace: pravidla vs. učení v Minecraft Education (podpora konceptů)
    \begin{itemize}
      \item Cíl: studentům názorně ukázat rozdíl mezi explicitní logikou (pravidly) a chováním navrženým učením či adaptací.
      \item Postup (stupeň připravenosti: nízký až střední): stáhněte Minecraft Education a spusťte jednu z volně dostupných hodin z Hour of Code, které jsou zaměřené na AI; nechte studenty upravovat jednoduchá pravidla chování agenta a diskutujte, jak by se chování lišilo, kdyby agent „učil“ své reakce z dat \cite{kb_ai_101_tipu_pdf_04e4a115_page_15_2}.
      \item Diskuse: porovnejte předvídatelnost, opakovatelnost a možnosti vysvětlení chování v obou přístupech.
    \end{itemize}

  \item Rychlý experiment: vytvoření a validace kvízu pomocí Copilot + Forms
    \begin{itemize}
      \item Cíl: prakticky ukázat pedagogům, jak generativní AI urychlí tvorbu a jaký je potřeba krok lidské validace.
      \item Postup (10–30 minut): v Microsoft Forms aktivujte Copilot, zadejte parametrizovaný prompt (např. úroveň kurzu, počet otázek, formát), vygenerované otázky pečlivě zkontrolujte a upravte podle didaktických cílů \cite{kb_ai_101_tipu_pdf_04e4a115_page_8_1}.
      \item Ověření: přidejte krok, kde studenti nahrávají postup řešení (pokud je to podstatné), nebo nastavte časově omezené in‑class verze úloh (viz kapitola o hodnocení).
    \end{itemize}

  \item Vizuální rozdíly: generování grafiky v Designer jako ilustrace generativních přístupů
    \begin{itemize}
      \item Použití: ukázat, že generativní nástroje v oblasti grafiky (např. Designer) mohou rychle vytvořit různé varianty vizuálů a přizpůsobit poměry stran či styly; vhodné pro tvorbu materiálů a vizuální diferenciaci výukových listů \cite{kb_ai_101_tipu_pdf_04e4a115_page_51_1}.
      \item Didaktická poznámka: vždy zkontrolujte autorské a licenční podmínky pro použití v kurzu a upravte vygenerované výstupy z hlediska pedagogického záměru.
    \end{itemize}
\end{enumerate}

Doporučení pro učitele (general background knowledge)
\begin{itemize}
  \item Volte technologii podle cíle učení: pro vysvětlení deterministických pravidel použijte pravidlové systémy; pro ukázku procesu učení a práce s daty zvolte klasické ML; pro tvorbu textů, adaptivních vysvětlení a rychlé přípravy materiálů využijte LLM.
  \item Validace a transparentnost: při použití generativní AI vždy plánujte krok lidské kontroly a informujte studenty o využití AI v materiálech (viz kapitola o etice a GDPR).
  \item Aktivně zapojte reflexi: dejte studentům úkol porovnat výstupy více přístupů (pravidla vs. ML vs. generativní model) a popsat rozdíly ve spolehlivosti, interpretovatelnosti a vhodnosti pro konkrétní úlohy.
\end{itemize}%

%
\section{Jak fungují LLM: tokens, embeddings a tréninková data}%
Koncentrované technické jádro: vysvětlení pojmů tokens a embeddings, role tréninkových dat a pojem cut‑off dat; názorné grafické schéma (bez hluboké matematiky) vhodné pro pedagogické použití.%
\subsection{Co jsou tokens a jak je chápat v pedagogickém kontextu}%
Krátké, srozumitelné vysvětlení pojmu token (jednotka textu pro LLM), rozdíl mezi tokenem a písmenem/slovem, ukázka rozdělení věty na tokeny a jednoduché pedagogické metafory pro vysvětlení studentům. Praktické důsledky pro promptování (dlouhé vs. krátké promptování, limit kontextu).%


%
% (Tělo sekce: Co jsou tokens a jak je chápat v pedagogickém kontextu)

% Stručné vysvětlení (označeno jako obecné znalosti)
\noindent\textbf{Co je token?} \emph{(general background knowledge)} Token je v kontextu velkých jazykových modelů základní jednotka textu, kterou model zpracovává. Token nemusí odpovídat přesně slovu nebo písmenku — může to být celé slovo, jeho část nebo symbol; závisí to na použitém tokenizéru.\\

% Jednoduchá pedagogická metafora
\noindent\textbf{Metafora pro studenty} \emph{(general background knowledge)}: Představte si tokeny jako stavební cihly věty. Některé cihly jsou malé (krátké části slov), jiné větší (celá slova). Stejně jako při stavění modelu z cihel, počet cihel ovlivní, kolik materiálu (textu) lze najednou postavit — model má limit, kolik tokenů zmůže v jednom „posloupnosti“.\\

% Praktický příklad (označeno jako obecné znalosti)
\noindent\textbf{Krátký příklad rozdělení věty} \emph{(general background knowledge)}: Věta „Dobrý den, jak se máte?“ může být tokenizována do několika tokenů, přičemž interpunkce a koncovky mohou tvořit samostatné tokeny u některých tokenizérů. (Pozn.: přesné rozdělení závisí na konkrétním modelu/tokenizéru.)\\

% Důsledky pro promptování — praktické tipy (označeno jako obecné znalosti)
\noindent\textbf{Praktické důsledky pro pedagoga a promptování} \emph{(general background knowledge)}:
\begin{itemize}
  \item Kratší a cílené prompty — jasně formulované instrukce a příklady vedou obvykle k přesnějším odpovědím; dlouhé nepřehledné prompty mohou obsahovat irelevantní kontext a zvýšit riziko chyb nebo „odbočení“ od úkolu.
  \item Limit kontextu — modely mají omezení počtu tokenů, která mohou zpracovat najednou; při velmi dlouhých učebních materiálech je vhodné použít postupy jako rozdělení na menší části nebo RAG (retrieval‑augmented generation). 
  \item Validace výstupu — vždy zkontrolujte fakta a didaktickou vhodnost vygenerovaného textu (lidská kontrola jako povinný krok při použití AI v přípravě materiálů).
\end{itemize}

% Ukázka reálného použití promptu (podpořené KB)
\noindent\textbf{Krátký ukázkový workflow} (použití v praxi): zadejte do nástroje jasný prompt s požadovaným formátem (např. „Vytvoř test s 10 otázkami, 4× ABC, 4× pravda/nepravda, 2× otevřené“); vygenerovaný obsah pečlivě zkontrolujte a upravte podle didaktických cílů — tento postup je ilustrován u Microsoft 365 Copilot, který dokáže na základě takového promptu vytvořit kompletní test včetně označení správných odpovědí a vysvětlení \cite{kb_ai_101_tipu_pdf_04e4a115_page_36_1}.\\

% Krátké doporučení pro výuku
\noindent\textbf{Doporučení pro výuku}:
\begin{itemize}
  \item Předveďte studentům jednoduchou ukázku tokenizace jako aktivitu (ukázat různé výsledky u různých nástrojů) a diskutujte, proč to ovlivňuje chování modelu. \emph{(general background knowledge)}
  \item Vysvětlete omezení kontextu a nechte studenty navrhnout, jak dělit rozsáhlý text pro spolehlivou práci s modelem. \emph{(general background knowledge)}
\end{itemize}%

%
\subsection{Embeddings: co to jsou a jak je používat ve výuce}%
Přehled, co jsou vektory embeddingů a jak zachycují význam slov/vět; demonstrace jednoduchého příkladu (vizualizace 2D/3D) a konkrétní pedagogické scénáře: vyhledávání podobných materiálů, clustering studijních odpovědí, základ RAG (retrieval‑augmented generation).%


%
\noindent\textbf{Co jsou embeddings?} \emph{(general background knowledge)} Vektorové embeddingy jsou numerické reprezentace slov, vět nebo dokumentů v‑dimenzionálním prostoru tak, aby podobné položky měly blízké vektory. Můžete si je představit jako souřadnice slov na mapě významů: slova se podobným významem leží blízko u sebe.\\

\noindent\textbf{Pedagogická metafora a vizualizace} \emph{(general background knowledge)}: Představte si sada bodů v rovině (2D), kde body reprezentují slova nebo krátké texty. Studenti mohou v malém cvičení promítnout více‑rozměrné embeddingy do 2D pomocí PCA nebo t‑SNE a vizuálně hledat shluky témat — to pomáhá pochopit, že „podobnost“ je vektorová blízkost.\\

\noindent\textbf{Praktické scénáře použití ve výuce} \emph{(general background knowledge)}:
\begin{itemize}
  \item Vyhledávání podobných materiálů: pro dotaz (např. krátký popis tématu) se spočtou embeddingy dotazu i dokumentů a vrátí se nejbližší dokumenty — rychlý způsob, jak studentům nabídnout relevantní doplňkové zdroje.
  \item Clustering studentských odpovědí: embeddingy krátkých studentských textů lze seskupit (clustering) a identifikovat opakující se koncepty, nejasnosti nebo typické chyby ve třídě.
  \item Základ pro RAG (retrieval‑augmented generation): indexace interních materiálů pomocí embeddingů umožní modelu nejprve „vyhledat“ relevantní pasáže a teprve pak generovat odpověď, čímž se zvyšuje relevance a kontrola nad zdroji.
\end{itemize}

\noindent\textbf{Krátký ukázkový workflow pro vyučující (praktické nasazení)} \emph{(general background knowledge)}:
\begin{enumerate}
  \item Připravte sadu krátkých dokumentů nebo odpovědí (např. abstrakty, výpisy z přednášek).
  \item Vygenerujte embeddingy pro každý dokument a pro ukázkové dotazy (použijte dostupnou službu/model podle pravidel instituce).
  \item Pro vizualizaci použijte redukci dimenzí (PCA/t‑SNE/UMAP) a ukažte výsledky studentům; nechte je interpretovat shluky.
  \item Pro vyhledávání nastavte jednoduché rozhraní: při dotazu vypište top‑k relevantních dokumentů podle podobnosti.
  \item Diskutujte omezení: embeddingy odrážejí data a model, se kterým byly vytvořeny — vždy ověřujte relevanci a přesnost návrhů.
\end{enumerate}

\noindent\textbf{Krátké didaktické tipy} \emph{(general background knowledge)}:
\begin{itemize}
  \item Začněte s malým datasetem (10–100 textů), aby studenti viděli jasné shluky a vztahy.
  \item Použijte vizualizaci jako diskusní podnět — požádejte studenty, aby navrhli, proč se některé odpovědi seskupily dohromady.
  \item Při zavádění RAG zdůrazněte transparentnost: vždy ukažte, které zdroje byly vybrány pro generovanou odpověď.
\end{itemize}%

%
\subsection{Tréninková data, cut‑off dat a názorné schéma pro výuku}%
Role tréninkových dat u LLM, co znamená cut‑off dat (hranice, do kdy model ‚ví‘ fakta), dopady na aktualitu informací, základní otázky kolem kvality dat, zkreslení a ochrany soukromí (stručně pro pedagogické posouzení). Součástí je jedno přehledné grafické schéma (návrh obsahu schématu a tipy pro zařazení do výuky) a krátký návod, jak toto schéma prezentovat bez hluboké matematiky.%


%
\noindent\textbf{Role tréninkových dat (stručně)} \emph{(general background knowledge)}: Tréninková data určují vzory, fakta a možná zkreslení v odpovědích; kvalita, rozmanitost a časové pokrytí ovlivňují přesnost a aktuálnost.\\[6pt]
\noindent\textbf{Co znamená cut‑off dat} \emph{(general background knowledge)}: Cut‑off je poslední časový bod zahrnutí tréninkových dat — model nemusí znát události po tomto datu; v učebním kontextu zdůrazněte ověřování z externích zdrojů.\\[6pt]
\noindent\textbf{Klíčové pedagogické body k zahrnutí do schématu} \begin{itemize}
  \item Zdroj dat: weby, knihy, interní dokumenty, anotace; kvalita a bias: reprezentativnost, chyby, systematická zkreslení.
  \item Časová osa / cut‑off: vizuální osa znázorňující „do kdy“ model zná informace.
  \item Soukromí a právní aspekty: osobní údaje, DPIA; Retrieval/RAG: jak indexy a retrieval vrstva mění chování.
  \item Didaktika: kdy vyžadovat ověření a ukázat zdroje.
\end{itemize}
\noindent\textbf{Návrh jednoduchého schématu pro slide (1–2 snímky)} \emph{(general background knowledge)}: \begin{enumerate}
  \item Snímek 1 — „Mapa tréninkových dat“: centrální kruh „Model“ obklopený bublinami „Web“, „Knihy“, „Interní dokumenty“; pod tím osa s cut‑off.
  \item Snímek 2 — „Dopady a mitigace“: tři sloupce: rizika, příznaky v odpovědích, opatření (ověřit zdroj, DPIA, lidská validace).
\end{enumerate}
\noindent\textbf{Jak toto schéma prezentovat bez matematiky} \emph{(general background knowledge)}: \begin{itemize}
  \item Metafora „knihovny“ a krátké demonstrace: porovnejte otázky o událostech před a po cut‑off.
  \item Praktická pravidla: vždy ukázat zdroj a provést lidskou kontrolu u kritických tvrzení.
\end{itemize}
\noindent\textbf{Krátký workshop / aktivita ve třídě (praktický příklad)}: \begin{enumerate}
  \item Fáze A (bez AI): studenti vypracují krátký text.
  \item Fáze B (s AI): studenti použijí model k rozšíření textu a porovnají přínosy a chyby; diskutujte odkud model čerpá informace \cite{kb_malinka_chatgpt_ve_vyuce_dva_roky_pote_2024_pdf_90bc3aaf_page_12_1}.
  \item Reflexe a alternativa: studenti popíší ověřování tvrzení; jako alternativa použijte didaktický svět (např. Minecraft Education) k názornému zobrazení vlivu tréninkových dat \cite{kb_ai_101_tipu_pdf_04e4a115_page_15_2}.
\end{enumerate}
\noindent\textbf{Krátké bezpečnostní a etické připomenutí} \cite{kb_malinka_chatgpt_ve_vyuce_dva_roky_pote_2024_pdf_90bc3aaf_page_12_1}: zvažte dopad na ochranu soukromí a zařaďte analýzu rizik (DPIA) při práci s osobními nebo interními daty.\\[6pt]
\noindent\textbf{Doporučené diskusní otázky pro studenty} \emph{(general background knowledge)}: \begin{itemize}
  \item Jaké zdroje byste považovali za důvěryhodné pro ověření výroku z modelu?
  \item Co by model mohl „nevědět“ kvůli cut‑off datu a jak to rozpoznat?
  \item Jak mitigovat riziko, že model předá citlivé informace z tréninkových dat?
\end{itemize}%

%
\section{Rizika a bezpečné používání ve výuce}%
Přehled hlavních rizik relevantních pro výuku: halucinace (nesprávné výstupy), bias, otázky ochrany osobních údajů a akademické integrity; doporučení prvních bezpečnostních opatření a transparentní praxe vůči studentům.%


%
Krátké shrnutí rizik (co mít na paměti)
\begin{itemize}
  \item Soukromí a právní rizika: nasazení generativní AI může zahrnovat zpracování osobních údajů studentů (konverzace, odevzdané soubory, logy). Před pilotem proveďte posouzení dopadů (DPIA) a zvažte minimální sběr dat a pseudonymizaci — toto by mělo být součástí plánování nasazení \cite{kb_malinka_chatgpt_ve_vyuce_dva_roky_pote_2024_pdf_90bc3aaf_page_12_1}. 
  \item Nesprávné nebo vymyšlené výstupy („halucinace“): modely mohou generovat přesvědčivě znějící, leč nepravdivé informace. (Tento jev stručně označujeme jako halucinace; níže najdete doporučení k mitigaci.) \emph{(general background knowledge)}
  \item Zkreslení (bias): odpovědi modelu odrážejí vzory v tréninkových datech, což může vést k systematickému zkreslení obsahu nebo reprezentace určitých skupin. \emph{(general background knowledge)}
  \item Dopad na akademickou integritu: dostupnost generativních nástrojů mění motivace i postupy studentů; je potřeba upravit zadání a hodnoticí postupy tak, aby měřily žádané kompetence.
  \item Provozní rizika a bezpečnost: úniky dat do třetích stran, nevhodné sdílení interních materiálů (indexace interních dokumentů bez kontroly), chybné konfigurace enterprise nastavení.
\end{itemize}

Doporučená první bezpečnostní opatření (praktický návod — krok za krokem)
\begin{enumerate}
  \item Plánování a DPIA: před nasazením popište rozsah zpracování (co se bude posílat do modelu, kde se ukládá, kdo k tomu má přístup) a proveďte DPIA; pro výuku využijte vzor DPIA z přílohy (Příloha 9.4) jako kontrolní seznam. Dokumentujte také očekávané přínosy a zmírňující opatření (pseudonymizace, retence, šifrování) \cite{kb_malinka_chatgpt_ve_vyuce_dva_roky_pote_2024_pdf_90bc3aaf_page_12_1}.
  \item Začněte malým pilotem se zřejmými hranicemi: navrhněte dvoufázové aktivity — fáze A bez AI (nezávislé řešení), fáze B s AI pro rozšíření a reflexi — a porovnejte výsledky a bezpečnostní rizika před škálováním \cite{kb_malinka_chatgpt_ve_vyuce_dva_roky_pote_2024_pdf_90bc3aaf_page_12_1}.
  \item Minimalizace a lidská kontrola: omezte, které typy dat se posílají k externím poskytovatelům; pro automatizované hodnocení zachovejte konečnou lidskou kontrolu — učitel by měl finálně potvrdit nebo upravit návrhy AI (příklad v nástroji pro matematiku, kde AI předpřipraví hodnocení a učitel je výsledně finalizuje) \cite{kb_ai_101_tipu_pdf_04e4a115_page_8_1}.
  \item Transparentnost vůči studentům: předem oznamte, kde a jak AI využíváte (v sylabu, viz Příloha 9.3), jaké typy dat se budou zpracovávat a jaká opatření byla přijata (DPIA, anonymizace). Doporučené texty a šablony najdete v přílohách.
  \item Audit a záznamy: zapněte logování a auditní stopy tam, kde to lze (kdo co dotazoval, jaké verze modelu byly použity), a průběžně revidujte chování systému během pilotu.
\end{enumerate}

Praktické příklady mitigace (konkrétní scénáře)
\begin{itemize}
  \item Cvičení s automatickou nápovědou: nastavte aktivitu tak, aby studenti nejprve samostatně vytvořili řešení; teprve poté mohou použít AI k rozšíření nebo kontroverznímu rozboru nápadů — takto vidíte rozdíl ve schopnostech a zároveň omezíte přímé kopírování výstupů AI \cite{kb_malinka_chatgpt_ve_vyuce_dva_roky_pote_2024_pdf_90bc3aaf_page_12_1}.
  \item Testy a kvízy generované Copilot/Forms: použijte nástroje typu Microsoft 365 Copilot pro návrh otázek, ale ponechte kontrolu nad finální sadou otázek a klíčem učiteli; systém může velmi rychle vytvořit návrh testu, učitel však musí zkontrolovat přesnost a vhodnost otázek \cite{kb_ai_101_tipu_pdf_04e4a115_page_36_1}.
  \item Vysvětlení AI principů studentům: pro praktickou demonstraci principu tréninku a limitů modelu lze využít vzdělávací světy (např. Minecraft Education) jako názorný doplněk při diskusi o tom, odkud model „čerpá“ informace \cite{kb_ai_101_tipu_pdf_04e4a115_page_15_2}.
\end{itemize}

Krátký checklist před spuštěním aktivity s AI (pro vyučujícího)
\begin{itemize}
  \item Mám popsaný účel zpracování a provedenou DPIA (Příloha 9.4)? \\
  \item Jsou citlivá data (osobní údaje, fotografie, zdravotní údaje) vyloučena nebo pseudonymizována? \\
  \item Je zřejmá role člověka v procesu (kdo finalizuje hodnocení / schvaluje výstupy)? \cite{kb_ai_101_tipu_pdf_04e4a115_page_8_1} \\
  \item Jsou studenti informováni o použití AI a existují jasné požadavky na deklaraci využití nástroje (šablona v Příloze 9.3)? \\
  \item Bude pilot monitorován a jaké metriky/incidenty hlásíme IT/GDPR týmu (viz online repozitář ai.vut.cz pro proces hlášení a další šablony)? 
\end{itemize}

Poznámky k hodnocení rizik a další kroky
\begin{itemize}
  \item Nasazení AI není binary problém — zaměřte se na řízené experimenty s jasným plánem monitorování rizik a s možností rychlého stažení řešení, pokud se objeví nepřijatelná rizika \cite{kb_malinka_chatgpt_ve_vyuce_dva_roky_pote_2024_pdf_90bc3aaf_page_12_1}.
  \item Využijte interní šablony a repozitář (ai.vut.cz) pro rychlý přístup k DPIA/DPA vzorům, šablonám do sylabů a video tutoriálům — to zkracuje dobu potřebnou k bezpečnému pilotu a usnadňuje dokumentaci.
  \item U citlivých nebo klíčových rozhodnutí (pracovní toky zahrnující osobní údaje, rozsáhlé sdílení interních materiálů) konzultujte IT a právní/GDPR oddělení fakulty před spuštěním.
\end{itemize}

Závěrem: cílem je nezákazovat nástroje, ale nasazovat je s rozmyslem — plánujte, dokumentujte, minimalizujte data a ponechte lidské rozhodování jako poslední instanci. \emph{(Praktické šablony a kontrolní seznamy najdete v Přílohách 9.3–9.5 a v online repozitáři.)}%

%
\section{Co od nástrojů reálně očekávat — doporučení pro pedagogy}%
Praktický přehled schopností a omezení nástrojů generativní AI z pohledu výuky a návrh ‚bezpečného startu‘: jak začít s malými piloty, co měřit a jak nastavit očekávání studentů i garanta kurzu.%


%
% Text section body for: Co od nástrojů reálně očekávat — doporučení pro pedagogy

V této části shrnujeme, co lze od generativních AI nástrojů reálně očekávat ve výuce, jak definovat „bezpečný start“ pilotu a které metriky a artefakty sledovat. Cílem je dát garantům a učitelům praktický checklist a jednoduchý workflow pro první piloty s jasnými hranicemi odpovědnosti.

Krátké očekávání schopností a omezení
\begin{itemize}
  \item Nástroje dobře pomáhají s opakovatelnými, šablonovatelnými úlohami: generování návrhů otázek, řešení příkladů, návrhů slidů nebo prvotních revizí textů. U nástrojů pro výuku matematiky například funkce „Pokrok v matematice“ generuje příklady, pomáhá s vyhodnocením postupů a připravuje návrh hodnocení, které učitel finálně schválí \cite{kb_ai_101_tipu_pdf_04e4a115_page_8_1}.
  \item Nástroje lze použít i k názorné demonstraci principů AI (např. vzdělávací světy v Minecraft Education), což je užitečné při zavádění tématu studentům \cite{kb_ai_101_tipu_pdf_04e4a115_page_15_2}.
  \item Limitace: modely mohou generovat přesvědčivě znějící, ale nepřesné nebo nepravdivé informace (tzv. „halucinace“) — toto je stručné obecné upozornění pro pedagogickou praxi (general background knowledge).
\end{itemize}

Doporučený postup „bezpečného startu“ (praktický workflow)
\begin{enumerate}
  \item Definujte rozsah a účel pilotu: konkrétní předmět, typ úloh (např. tvorba pracovních listů, formativní zpětná vazba, generování testů), časový horizont a očekávané přínosy.
  \item Proveďte DPIA a zdokumentujte datový tok: co se bude posílat do externích služeb, jaká data jsou citlivá a jak budou pseudonymizována či vyloučena; využijte vzor DPIA z příloh jako kontrolní seznam \cite{kb_malinka_chatgpt_ve_vyuce_dva_roky_pote_2024_pdf_90bc3aaf_page_12_1}.
  \item Začněte dvoufázově: fáze A bez AI (nezávislý výsledek od studentů), fáze B s AI (rozšíření, reflexe). Tím získáte srovnávací data o přínosu a rizicích nasazení \cite{kb_malinka_chatgpt_ve_vyuce_dva_roky_pote_2024_pdf_90bc3aaf_page_12_1}.
  \item Určete „human‑in‑the‑loop“ pravidla: kdo finálně schvaluje obsah, jak probíhá kontrola automaticky generovaných odpovědí a jaká je eskalace chyb.
  \item Transparentnost vůči studentům: v sylabu uveďte, kde a jak AI používáte (šablona v Příloze 9.3), a požadujte deklaraci použití tam, kde je to relevantní \cite{kb_malinka_chatgpt_ve_vyuce_dva_roky_pote_2024_pdf_90bc3aaf_page_12_1}.
\end{enumerate}

Co měřit během pilotu (praktické metriky)
\begin{itemize}
  \item Pedagogické výsledky: srovnání kvality výstupů a porozumění mezi fázemi A a B (např. hodnocení řešení, počet opravných koleček).
  \item Procesní důkazy: artefakty studentů (drafty, postupy, komentáře) a záznamy interakcí s AI (logy dotazů/odpovědí) k ověření autenticity a sledování chyb.
  \item Efektivita práce učitele: časová úspora při přípravě materiálů nebo hodnocení (např. čas na finalizaci návrhů generovaných AI), včetně kvality „předpřipravených“ návrhů, které učitel upravuje \cite{kb_ai_101_tipu_pdf_04e4a115_page_8_1}.
  \item Bezpečnostní a právní ukazatele: typy dat zasílaných do externích služeb, incidennce nežádoucích sdílení, výsledky konzultací s IT/GDPR a splnění požadavků z DPIA \cite{kb_malinka_chatgpt_ve_vyuce_dva_roky_pote_2024_pdf_90bc3aaf_page_12_1}.
\end{itemize}

Krátký checklist pro pedagoga před spuštěním pilotu
\begin{itemize}
  \item Popsaný účel pilotu a provedená DPIA (Příloha 9.4)? \cite{kb_malinka_chatgpt_ve_vyuce_dva_roky_pote_2024_pdf_90bc3aaf_page_12_1}
  \item Jsou citlivá data vyloučena nebo pseudonymizována? 
  \item Je jasně definováno, kdo finalizuje výstupy AI (human‑in‑the‑loop)? \cite{kb_ai_101_tipu_pdf_04e4a115_page_8_1}
  \item Jsou studenti informováni a je ve sylabu viditelná politika využití AI (šablona v Příloze 9.3)?
  \item Budou logovány interakce a naplánován pravidelný audit chování systému během pilotu?
\end{itemize}

Krátké příklady nasazení pro první pilot (inspirace)
\begin{itemize}
  \item Podpora přípravy: nechte AI navrhnout sadu cvičných úloh a vy jako učitel vyberete a doladíte konečný klíč — rychlý způsob, jak otestovat kvalitu návrhů bez ztráty kontroly \cite{kb_ai_101_tipu_pdf_04e4a115_page_8_1}.
  \item Didaktická demonstrace: využijte připravené světy v Minecraft Education k názornému vysvětlení, odkud model „čerpá“ a jak fungují demonstrační scénáře \cite{kb_ai_101_tipu_pdf_04e4a115_page_15_2}.
\end{itemize}

Poznámka o dalším kroku: výstupy pilotu zdokumentujte a po skončení proveďte krátké vyhodnocení pro rozhodnutí o škálování (co fungovalo, kolik práce vyžadovalo dohled, jaké úpravy v DPIA či workflow jsou nutné). Pro šablony DPIA, text do sylabu a checklisty využijte přílohy 9.3–9.5 a online repozitář ai.vut.cz jako living resource.%

%
\section{Rychlý start s nástroji (ChatGPT, Claude, Gemini, Microsoft Copilot)}%
Krok‑za‑krokem návod pro první použití běžných nástrojů v praxi výuky včetně poznámky o dostupnosti Microsoft 365 Copilot pro školní prostředí (integrační příležitosti do Office/Teams).%


%
Krátké úvodní poznámky
\begin{itemize}
  \item Tato část nabízí praktický „rychlý start“ pro pedagoga, který chce poprvé vyzkoušet generativní AI v přípravě nebo v hodině. Obsahuje minimální bezpečnostní kroky, doporučený postup pilotu a konkrétní ukázku (Microsoft 365 Copilot + Forms). Informace o dostupnosti školních verzí Copilota jsou uvedeny níže \cite{kb_ai_101_tipu_pdf_04e4a115_page_3_1, kb_ai_101_tipu_pdf_04e4a115_page_36_1}.
  \item Poznámka: Seznam konkrétních služeb (např. ChatGPT, Claude, Gemini) uvádíme jako orientační — tyto stručné zmínky jsou označeny jako general background knowledge a nevyjadřují doporučení konkrétního produktu v dané instituci (general background knowledge).
\end{itemize}

Rychlý 6‑krokový workflow pro první použití
\begin{enumerate}
  \item Definujte cíl a rozsah pilotu: konkretizujte předmět, typ úlohy (např. generování cvičných testů, návrh slidů, krátká formativní zpětná vazba) a délku pilotu. Zároveň rozhodněte, které typy dat nebudou nikdy odesílány do externích služeb (osobní údaje, studentské práce s identifikátory apod.) \cite{kb_malinka_chatgpt_ve_vyuce_dva_roky_pote_2024_pdf_90bc3aaf_page_12_1}.
  \item Proveďte zjednodušené DPIA a datový plán: zdokumentujte, jaký datový tok bude v pilotu (co se posílá, kde se logy ukládají, kdo má k nim přístup) a uložte tuto dokumentaci jako součást pilotu \cite{kb_malinka_chatgpt_ve_vyuce_dva_roky_pote_2024_pdf_90bc3aaf_page_12_1}.
  \item Zvolte nástroj podle potřeby a licence: ověřte, zda vaše fakulta/univerzita nabízí enterprise nebo školní variantu (např. Office 365 pro školy/A1 obsahuje webovou verzi Copilot Chat; plná integrace Copilota do Word/PowerPoint/Forms je k dispozici v placené verzi Microsoft 365 Copilot) \cite{kb_ai_101_tipu_pdf_04e4a115_page_3_1}. Ujistěte se, že licence a DPA jsou v souladu s požadavky vašeho pracoviště.
  \item Začněte malým dvoufázovým pilotem (bez AI → s AI): porovnejte kontrolní skupinu bez AI a skupinu s AI, zaznamenejte artefakty a procesy (drafty, interakce s AI) pro pozdější vyhodnocení \cite{kb_malinka_chatgpt_ve_vyuce_dva_roky_pote_2024_pdf_90bc3aaf_page_12_1}.
  \item Nastavte pravidla „human‑in‑the‑loop“ a transparentnost: určete, kdo finálně schvaluje výstupy AI, jak probíhá eskalace chyb, a informujte studenty o použití AI (text do sylabu/oznámení) \cite{kb_malinka_chatgpt_ve_vyuce_dva_roky_pote_2024_pdf_90bc3aaf_page_12_1}.
  \item Vyhodnoťte pilot podle předem definovaných metrik: pedagogická kvalita výstupů, časová úspora, bezpečnostní incidenty, dopad na akademickou integritu; na základě výsledků upravte rozsah a procesy.
\end{enumerate}

Konkrétní příklad: rychlé vytvoření kvízu v Microsoft Forms pomocí Copilota
\begin{itemize}
  \item Kontext: Microsoft 365 Copilot je dostupný ve více variantách; školní Office 365 může obsahovat webovou verzi Copilot Chat a placená varianta Copilot nabízí integrace přímo do aplikací (Word, PowerPoint, Forms apod.) \cite{kb_ai_101_tipu_pdf_04e4a115_page_3_1}.
  \item Praktický postup:
    \begin{enumerate}
      \item Ověřte, že vaše školní/ fakultní účetní nastavení zahrnuje Copilot Chat nebo Microsoft 365 Copilot (kontaktujte IT/elearning pro potvrzení licence).
      \item Otevřete Microsoft Forms, aktivujte Copilot (pokud je dostupný) a použijte jednoduchý prompt pro vytvoření testu.
      \item Upravený výsledek zkontrolujte a doplňte manuálně (správné odpovědi, jmenné formátování, úprava úrovně obtížnosti).
    \end{enumerate}
  \item Ukázkový prompt (převzatý z praktických příkladů): 
\begin{verbatim}
Vygeneruj mi test na hodinu zeměpisu na téma Austrálie, test bude obsahovat otázky na základní geografické oblasti, bude pro žáky 8. ročníku základní školy v České republice, test bude mít 10 otázek, 4 otázky budou typu ABC, 4 otázky typu pravda/nepravda a 2 otázky otevřené.
\end{verbatim}
Tento způsob generování vrátí sadu otázek včetně označení správných odpovědí a často i krátkých vysvětlení — vždy ale prověřte přesnost a přiměřenost pro váš kontext před použitím v hodnocení \cite{kb_ai_101_tipu_pdf_04e4a115_page_36_1}.
\end{itemize}

Krátký bezpečnostní checklist před prvním použitím v hodině
\begin{itemize}
  \item Máte záznam DPIA nebo minimální datový plán? \cite{kb_malinka_chatgpt_ve_vyuce_dva_roky_pote_2024_pdf_90bc3aaf_page_12_1}
  \item Jsou citlivá data (jména, osobní identifikátory, studentské nahrávky) vyloučena nebo pseudonymizována? \cite{kb_malinka_chatgpt_ve_vyuce_dva_roky_pote_2024_pdf_90bc3aaf_page_12_1}
  \item Je jasně určeno, kdo finálně validuje výstupy AI (human‑in‑the‑loop)? \cite{kb_malinka_chatgpt_ve_vyuce_dva_roky_pote_2024_pdf_90bc3aaf_page_12_1}
  \item Máte potvrzení licence / dostupnosti Copilota nebo jiného nástroje u interní IT podpory (pokud plánujete používat školní Office 365/Copilot)? \cite{kb_ai_101_tipu_pdf_04e4a115_page_3_1}
  \item Plánujete krátké vyhodnocení pilotu (srovnání bez AI vs. s AI) a sběr zpětné vazby od studentů?
\end{itemize}

Drobné tipy pro rychlé nasazení
\begin{itemize}
  \item Začněte s nehodnoticím obsahem (příprava cvičných listů, návrhy slidů), kde případné chyby lze rychle opravit — to sníží riziko dopadu na studenty (general background knowledge).
  \item Při generování testů nebo úloh vždy proveďte rychlou kontrolu faktů a jazykové úpravy před sdílením se studenty; u automaticky generovaných řešení doplňte pedagogické komentáře.
  \item Pokud máte přístup k Microsoft 365 Copilot v rámci školního Office, využijte integrace s Forms pro rychlou tvorbu kvízů a s Word/PowerPoint pro převod sylabu do slidů — ale vždy zachovejte lidskou kontrolu výsledků \cite{kb_ai_101_tipu_pdf_04e4a115_page_3_1, kb_ai_101_tipu_pdf_04e4a115_page_36_1}.
\end{itemize}

Odkazy na další kroky v příručce
\begin{itemize}
  \item Pro podrobnější postupy, DPIA šablonu a texty do sylabu využijte přílohy 9.3–9.5 a online repozitář ai.vut.cz (living document) — doporučené šablony a checklisty jsou připravené pro rychlé zkopírování a přizpůsobení.
\end{itemize}%

%
\section{Návrhy cvičení první interakce ve třídě}%
Sada konkrétních krátkých aktivit pro první vyzkoušení s žáky/studenty: včetně ukázkových zadání pro diskusi a hands‑on aktivity. Obsahuje praktické příklady z praxe zmíněné v KB, např. použití Minecraft Education nebo nástrojů pro ‚pokrok v matematice‘.%


%
Krátké praktické aktivity nízkého rizika, které učiteli umožní rychle vyzkoušet interakci studentů s generativními nástroji a zároveň otestovat provozní postupy (DPIA, pravidla pro sdílení dat, human‑in‑the‑loop). Před každou aktivitou si ověřte základní bezpečnostní kroky: provedení minimálního DPIA / datového plánu, vyloučení citlivých identifikátorů z odesílaných dat a jasné určení, kdo schvaluje výstupy AI (human‑in‑the‑loop) \cite{kb_malinka_chatgpt_ve_vyuce_dva_roky_pote_2024_pdf_90bc3aaf_page_12_1}. 

Přehled navrhovaných krátkých aktivit
\begin{enumerate}
  \item Icebreaker: „Co je AI?“ — Minecraft Education (30–45 min; general background suggestion)
    \begin{enumerate}
      \item Cíl: ukázat studentům konkrétní, interaktivní prostředí, kde lze demonstrovat principy učení agentů a jednoduché automatizace.
      \item Postup: krátká 10–15min demonstrační ukázka učitelem (využijte jedno z připravených světů zaměřených na AI / Hodinu kódu), poté rozdělte třídu do malých skupin, každá skupina si vyzkouší připravenou „Hour of Code“ aktivitu v Minecraft Education a společně diskutují, co v ní naznačuje učení agentů a jaké máme obavy/omezení. Minecraft Education obsahuje stovky připravených světů a některé aktivity pro Hodinu kódu jsou volně dostupné; zbytek obsahu je součástí školních licencí Microsoft 365 A3/A5 \cite{kb_ai_101_tipu_pdf_04e4a115_page_15_2}.
      \item Reflexe: rychlý záznam jedné otázky/odpovědi od každé skupiny do sdíleného dokumentu; krátká plenární diskuse o tom, co AI „umí“ a co ne.
    \end{enumerate}

  \item Praktické cvičení: „Pokrok v matematice“ — demonstration + hands‑on (45–60 min; general background suggestion)
    \begin{enumerate}
      \item Cíl: ukázat, jak nástroj může generovat cvičné příklady, sbírat postupy řešení a poskytovat předběžnou analýzu silných/slabých stránek studenta.
      \item Postup (učitel připraví zadání):
        \begin{enumerate}
          \item Učitel vybere tematickou oblast (např. zlomky, lineární rovnice) a vygeneruje sadu příkladů v nástroji typu "Pokrok v matematice".
          \item Studenti vypracují příklady a nahrají i postup řešení — nástroj předopraví odpovědi a označí oblasti, kde mohl student udělat chybu; učitel provede finální kontrolu a doplní pedagogický komentář \cite{kb_ai_101_tipu_pdf_04e4a115_page_8_1}.
          \item Diskuse: porovnejte, kde se nástroj „mýlí“ a jaká je přidaná pedagogická hodnota (diagnostika chyb, přizpůsobení obtížnosti).
        \end{enumerate}
      \item Doporučení k bezpečnosti: nepožadujte zasílání identifikovatelných osobních údajů nebo fotografií studentů do externích služeb; pokud takové data vzniknou, pseudonymizujte je před odesláním \cite{kb_malinka_chatgpt_ve_vyuce_dva_roky_pote_2024_pdf_90bc3aaf_page_12_1}.
    \end{enumerate}

  \item Dvoufázový workshop: „bez AI → s AI“ (60–90 min)
    \begin{enumerate}
      \item Cíl: pomoci studentům porozumět rozdílům v produktech lidské práce a AI‑generovaných rozšířeních, a procvičit ověřování faktů a zdrojů.
      \item Postup:
        \begin{enumerate}
          \item Fáze A (bez AI): studenti samostatně napíší krátký text (např. 300–400 slov na odborné téma nebo reflexi) během omezeného času.
          \item Fáze B (s AI): poté stejná skupina použije model (ve schválené a bezpečné instanci) k rozšíření nebo přepracování textu. Studenti dokumentují kroky, které AI provedla (prompt, varianty odpovědí, úpravy).
          \item Porovnání a reflexe: skupiny porovnají oba artefakty, hledají chyby, „halucinace“ a diskutují, jak model doplňoval informace mimo své cut‑off (ověřování aktuálnosti). Aktivita je doporučena jako způsob, jak „bezpečně startovat“ piloty a vést reflexi o ověřování a zdrojích \cite{kb_malinka_chatgpt_ve_vyuce_dva_roky_pote_2024_pdf_90bc3aaf_page_12_1, kb_ai_101_tipu_pdf_04e4a115_page_15_2}.
        \end{enumerate}
      \item Výstup: krátká skupinová prezentace o tom, kde AI přidala hodnotu a kde vznikla potřeba lidské korekce.
    \end{enumerate}

  \item Rychlá kreativní aktivita: generování výukové vizualizace (20–30 min; general background suggestion)
    \begin{enumerate}
      \item Cíl: ukázat, jak lze rychle vytvořit vizuální materiál pro hodinu (obrázek, plakát, monogram).
      \item Postup: učitel nebo skupiny studentů zformulují krátký prompt a nechají nástroj typu Designer vygenerovat návrh; následná společná úprava a sledování variant. Designer umí generovat novou grafiku, upravovat existující obrázky a umožňuje volbu poměru stran (např. širokoúhlé obrázky); mezi praktickými výstupy jsou např. monogramy, omalovánky či pozvánky \cite{kb_ai_101_tipu_pdf_04e4a115_page_51_1}.
      \item Diskuse: etické aspekty užití generované grafiky (licence, autorská práva) — krátké zamyšlení vedené učitelem.
    \end{enumerate}

  \item Krátká reflexe a deklarace použití AI (10–15 min)
    \begin{enumerate}
      \item Po každé aktivitě požádejte studenty o krátkou sebereflexi (co použili, proč, jak model měnil jejich postup) a jednoduchou deklaraci, že případné AI výstupy byly ověřeny. Transparentní oznámení a dokumentace použití AI jsou doporučenou praxí při pilotních nasazeních \cite{kb_malinka_chatgpt_ve_vyuce_dva_roky_pote_2024_pdf_90bc3aaf_page_12_1}.
      \item Vzorová formulace pro studentskou deklaraci (možno vložit do LMS):
\begin{verbatim}
Použil(a) jsem při vypracování následujících částí AI: [uveďte nástroj a stručně, jaký vstup/úpravu poskytl]. Ověřil(a) jsem fakta a upravil(a) jsem jazyk výstupu. Podpis: ...
\end{verbatim}
    \end{enumerate}
\end{enumerate}

Krátký checklist před aktivitou
\begin{itemize}
  \item Máte záznam minimálního DPIA / datového plánu pro aktivitu? (ano / ne) \cite{kb_malinka_chatgpt_ve_vyuce_dva_roky_pote_2024_pdf_90bc3aaf_page_12_1}
  \item Jsou vyloučena osobní identifikovatelná data z odesílaného obsahu? (ano / ne) \cite{kb_malinka_chatgpt_ve_vyuce_dva_roky_pote_2024_pdf_90bc3aaf_page_12_1}
  \item Je jasně určeno, kdo finálně kontroluje a schvaluje výstupy AI (human‑in‑the‑loop)? (ano / ne) \cite{kb_malinka_chatgpt_ve_vyuce_dva_roky_pote_2024_pdf_90bc3aaf_page_12_1}
  \item Má učitel připravenou krátkou reflexní otázku pro studenty (záznam o ověření výstupů)? (ano / ne)
\end{itemize}

Poznámky a doporučení
\begin{itemize}
  \item Začněte s nehodnoticím obsahem (demonstrační aktivity, cvičné listy, vizuály), kde případné chyby lze snadno opravit — to snižuje riziko negativního dopadu na studenty (general background knowledge).
  \item Po ukončení pilotu zaznamenejte artefakty (studentské výstupy, prompty, záznamy o opravách) a krátce vyhodnoťte pedagogický přínos a provozní rizika; výsledky slouží pro další úpravy a případné širší nasazení.
\end{itemize}%

%
\section{Odkazy na tutoriály, zdroje a doporučená literatura}%
Kurátorský seznam online tutoriálů, vstupních zdrojů pro pedagogy a doporučené kroky pro další vzdělávání; odkazy na repozitář s šablonami a videonávody jako další krok po pilotu.%


%
Krátký kurátorský přehled zdrojů, tutoriálů a doporučených kroků pro další vzdělávání pedagogů při pilotním nasazení generativní AI ve výuce. Níže najdete ověřené praktické vstupy (nástroje a ukázky aktivit) s odkazy na doporučené postupy a krátký akční plán „co dělat dál“.

Doporučené online tutoriály a vstupní zdroje
\begin{itemize}
  \item Minecraft Education — úvodní interaktivní světy a Hour of Code jako jednoduchý hands‑on způsob, jak studentům přiblížit principy strojového učení a agentního chování \cite{kb_ai_101_tipu_pdf_04e4a115_page_15_2}. (Praktické: stáhnout stránky Minecraft Education a vyzkoušet jednu z volně dostupných hodin.)
  \item Pokrok v matematice (Microsoft educational accelerators) — demonstrace generování sady příkladů, sběru postupů a předběžného opravení studentských odpovědí nástrojem; užitečné pro cvičení a diagnostiku slabých míst ve třídě \cite{kb_ai_101_tipu_pdf_04e4a115_page_8_1}.
  \item Microsoft 365 Copilot — pro úlohy „deep reasoning“ doporučen agent Výzkumník (Researcher), který v rámci Copilota zpracovává rozsáhlejší souvztažnosti a podklady; vhodné při přípravě strategických dokumentů a komplexnějších výukových materiálů \cite{kb_ai_101_tipu_pdf_04e4a115_page_48_1}.
  \item Repozitář VUT s šablonami a video‑tutoriály — živý repozitář (ai.vut.cz) obsahující šablony promptů, checklisty pro piloty a krátké video‑návody; doporučený další krok po ukončení lokálního pilotu (kontaktujte metodické oddělení pro přístup a podporu). (Organizační informace vyplývají z interních doporučení příručky — viz úvodní část příručky.)
\end{itemize}

Rychlý akční plán pro pedagoga — 5 kroků k prvnímu pilotu
\begin{enumerate}
  \item Projděte si krátké video‑tutoriály v repozitáři ai.vut.cz a vyberte jednu nízkorizikovou aktivitu (např. Minecraft Education nebo Pokrok v matematice). (repozitář jako další krok po pilotu; viz výše)
  \item Před spuštěním aktivity vypracujte minimální DPIA / datový plán a ověřte, že se do služby neodesílají osobně identifikovatelné údaje; zajistěte pseudonymizaci, pokud je to nutné \cite{kb_malinka_chatgpt_ve_vyuce_dva_roky_pote_2024_pdf_90bc3aaf_page_12_1}.
  \item Spusťte krátký pilot (30–90 min) s jasně definovanými cíli a metrikami (pedagogický cíl, očekávaný výstup, způsob validace). Jako první piloty doporučujeme aktivity popsané výše: Minecraft Education (Hour of Code) a Pokrok v matematice \cite{kb_ai_101_tipu_pdf_04e4a115_page_15_2, kb_ai_101_tipu_pdf_04e4a115_page_8_1}.
  \item Dokumentujte artefakty: prompty, varianty odpovědí, záznamy o opravách a krátkou reflexi (co fungovalo, kde AI selhala). Ponechte vždy human‑in‑the‑loop — vyučující kontroluje a finalizuje výstupy \cite{kb_malinka_chatgpt_ve_vyuce_dva_roky_pote_2024_pdf_90bc3aaf_page_12_1}.
  \item Pokud projekt vyžaduje hlubší rešerši nebo strategickou přípravu obsahu, vyzkoušejte použití Microsoft 365 Copilot s agentem Výzkumník pro agregaci a strukturování podkladů \cite{kb_ai_101_tipu_pdf_04e4a115_page_48_1}.
\end{enumerate}

Krátký checklist před pilotem (rychlé bezpečnostní body)
\begin{itemize}
  \item Existuje minimální DPIA / datový plán pro aktivitu? (ano / ne) \cite{kb_malinka_chatgpt_ve_vyuce_dva_roky_pote_2024_pdf_90bc3aaf_page_12_1}
  \item Jsou vyloučena osobně identifikovatelná data z odesílaného obsahu? (ano / ne) \cite{kb_malinka_chatgpt_ve_vyuce_dva_roky_pote_2024_pdf_90bc3aaf_page_12_1}
  \item Je vymezeno, kdo finalizuje výstupy a zajišťuje validaci (human‑in‑the‑loop)? (ano / ne) \cite{kb_malinka_chatgpt_ve_vyuce_dva_roky_pote_2024_pdf_90bc3aaf_page_12_1}
  \item Máte zdokumentované prompty a krátký plán pro evaluaci pedagogického přínosu? (ano / ne)
\end{itemize}

Kde najít další materiály
\begin{itemize}
  \item VUT repozitář (ai.vut.cz) — šablony promptů, checklisty pro DPIA a krátké video‑návody (living document; kvartální aktualizace).
  \item Ověřené ukázky aktivit v dokumentaci Microsoft Education (Minecraft Education, Pokrok v matematice) pro rychlé hands‑on nasazení \cite{kb_ai_101_tipu_pdf_04e4a115_page_15_2, kb_ai_101_tipu_pdf_04e4a115_page_8_1}.
  \item Pro hlubší rešerši a tvorbu strategických materiálů doporučujeme prozkoumat funkce agenta Výzkumník v Microsoft 365 Copilot \cite{kb_ai_101_tipu_pdf_04e4a115_page_48_1}.
\end{itemize}

Poznámka (general background knowledge): krátké tutoriály a praktická cvičení s jasně vymezenými bezpečnostními pravidly jsou nejrychlejší cestou, jak si vybudovat zkušenosti a získat podklady pro širší zavedení AI na úrovni katedry.%

%
\section{Krátký slovník pojmů pro pedagogy}%
Stručný glosář klíčových termínů (tokens, embeddings, hallucination, bias, RAG apod.) vhodný pro rychlou orientaci při přípravě materiálů a komunikaci se studenty.%


%
Krátký glosář — rychlé definice a praktické tipy pro pedagogy. Uvedené záznamy kombinují stručné (general background knowledge) definice klíčových pojmů s konkrétními vzdělávacími příklady a odkazy na nástroje doporučené v této příručce. Cílem je poskytnout lehce použitelné „kapesní“ poznámky, které lze rychle adaptovat do výuky nebo pilotních aktivit.
\begin{itemize}
\item \textbf{Tokeny} (general background knowledge) --- základní jednotka textu, se kterou pracují velké jazykové modely; může to být část slova, celé slovo nebo interpunkce. (Použijte tuto intuici při tvorbě promptů a při odhadování délky požadavků na model.) Praktická rada: při zadávání úloh omezte délku vstupu, pokud chcete předvídatelnější výstupy.
\item \textbf{Embeddings} (general background knowledge) --- číselné (vektorové) reprezentace textu nebo dokumentů, užitečné pro vyhledávání podobných dokumentů, clustering nebo pro RAG‑workflow. Praktický nápad: předveďte studentům jednoduché hledání „nejpodobnějších“ odstavců z učebních materiálů.
\item \textbf{Halucinace (hallucination)} (general background knowledge) --- situace, kdy model generuje fakticky nesprávné nebo nepodložené informace; vždy ověřujte kritická fakta z nezávislých zdrojů a učte studenty metody verifikace. V hodině lze ukázat příklady správných a chybných odpovědí a diskutovat, jak chybu odhalit.
\item \textbf{Bias / předpojatost} (general background knowledge) --- systematické zkreslení výstupů modelu vyplývající z tréninkových dat nebo návrhu modelu; při výuce upozorněte studenty na možné zdroje zkreslení a zvažte kontrolní otázky. Aktivita: nechte studenty navrhnout prototesty, které odhalí možné zkreslení v konkrétních datech.
\item \textbf{RAG — Retrieval‑Augmented Generation} (general background knowledge) --- vzor, kdy se výstup modelu doplňuje o relevantní dokumenty získané z indexu/embeddingů před generováním; vhodné pro zajištění podložených odpovědí při práci s interními materiály. Ukažte rozdíl mezi odpovědí s a bez retrievalu na konkrétním odborném dotazu.
\item \textbf{Pokrok v matematice (Microsoft)} --- nástroj z rodiny „Vzdělávací akcelerátory“, který generuje sady příkladů, sbírá postupy studentů a nabízí předběžné opravy; užitečný příklad, jak AI v praxi zdůrazňuje sběr postupu nad pouhým výsledkem \cite{kb_ai_101_tipu_pdf_04e4a115_page_8_1}. Doporučení: použijte krátké cvičení na sběr postupu a následnou reflexi chyb.
\item \textbf{Minecraft Education} --- edukační prostředí s předpřipravenými světy (včetně aktivit „Hour of Code\"), vhodné pro názorné ukázky principů AI ve výuce a hands‑on demonstrace strojového učení \cite{kb_ai_101_tipu_pdf_04e4a115_page_15_2}. Praktické využití: 10min demonstrační aktivita ilustrující agentní chování nebo jednoduché senzory.
\item \textbf{Microsoft 365 Copilot — agent \"Výzkumník\"} --- doporučený pracovní režim Copilota pro hlubší rešerše a agregaci podkladů při přípravě komplexních výukových materiálů \cite{kb_ai_101_tipu_pdf_04e4a115_page_48_1}. Tip: použijte jej jako pomoc při sběru zdrojů, nikoli jako jediný ověřený zdroj.
\item \textbf{Interní repozitář VUT (ai.vut.cz)} --- living document s šablonami promptů, checklisty a video‑tutoriály; první místo, kam sáhnout po ověřených šablonách a krátkých návodech pro pilotní aktivity (kontaktujte metodické oddělení pro přístup a podporu). Obsah repozitáře pravidelně aktualizujte a sdílejte vlastní osvědčené postupy.
\end{itemize}
Rychlá praktická poznámka pro učitele: ke každému klíčovému pojmu mějte připravenou jednu krátkou ukázku/aktivitu (např. 10min demonstrace v Minecraft Education pro koncept AI; nebo generování dvou‑tří příkladů v Pokrok v matematice pro ilustraci významu sběru postupu) — tyto konkrétní nástroje jsou uvedeny výše a v repozitáři ai.vut.cz \cite{kb_ai_101_tipu_pdf_04e4a115_page_15_2, kb_ai_101_tipu_pdf_04e4a115_page_8_1, kb_ai_101_tipu_pdf_04e4a115_page_48_1}. Doporučujeme také krátké reflexní úlohy po každé aktivitě, které pomohou studentům kriticky zhodnotit výsledky.
Poznámka: stručné definice označené jako (general background knowledge) jsou záměrně minimální; pro výuku doporučujeme rozšířit je o krátké praktické ukázky a ověřovací úkoly z repozitáře. Přizpůsobte aktivity věkové skupině a cílovým kompetencím třídy a pravidelně iterujte podle zpětné vazby.%

%
\chapter{ETICKÉ ZÁSADY POUŽITÍ AI}%
Kapitola definuje klíčové etické principy použití AI ve výuce: transparentnost, ochrana soukromí, spravedlnost/equita, integrita a zodpovědnost. Nabízí praktické postupy pro GDPR‑kompatibilní práci s daty studentů (anonymizace, minimalizace dat, DPIA, uzavření DPA s poskytovateli), doporučení pro využití enterprise či lokálních řešení, a stručný přehled evropského AI Act (2024) a jeho dopadů na vzdělávání. Rozpracovává akademickou integritu v éře AI, navrhuje kategorizaci úloh (zakázáno / povoleno s deklarací / vyžadováno), šablony do sylabů a vzory citací AI v různých citačních stylech. Obsahuje i doporučení, proč se nespoléhat na detektory AI a jak místo toho navrhovat úkoly a hodnoticí kritéria.%
\section{Přehled etických principů použití AI ve výuce}%
Vymezení klíčových principů: transparentnost, ochrana soukromí, spravedlnost (equita), akademická integrita a zodpovědnost. Stručné vysvětlení, proč jsou tyto principy relevantní pro pedagogickou praxi a jak ovlivňují rozhodování při nasazení AI v předmětech.%


%
\paragraph{}
Tato sekce shrnuje pět klíčových etických principů, které by měly řídit nasazení generativní AI ve výuce: transparentnost, ochrana soukromí, spravedlnost (equita), akademická integrita a zodpovědnost. Následující stručné definice a praktické poznámky slouží jako pracovní rámec pro garanty předmětů, vyučující a metodiky. Každý princip obsahuje krátké vysvětlení, konkrétní doporučení pro praxi a upozornění na běžné rizikové situace, které se v pedagogickém prostředí mohou vyskytnout.

\subparagraph{Transparentnost}
- Definice (general background knowledge): jasné informování studentů o tom, kde a jaký nástroj AI je používán, jaké jsou limity jeho výstupů a jak mají studenti deklarovat své využití AI. Transparentnost zahrnuje i sdílení informací o tom, zda byly použity předpřipravené šablony, externí dataset nebo modely třetích stran.
- Praktická doporučení: vložte krátké oznámení o použití AI do sylabu, uveďte verzi/nastavení asistenta a pravidla pro citování AI‑výstupů (viz příslušné šablony v přílohách). Uveďte, jak budou výsledky AI ověřovány a jaký vliv může mít použití AI na hodnocení. Doporučujeme také informovat studenty o možných chybách modelu a poskytnout krátký návod, jak ověřit faktické tvrzení v generovaném textu. (general background knowledge)

\subparagraph{Ochrana soukromí}
- Definice (general background knowledge): minimalizace sběru osobních údajů, pseudonymizace tam, kde je to možné, a předchozí posouzení dopadů (DPIA) pro nové piloty. Zahrnuje také posouzení přenosu dat mezi institucí a poskytovateli cloudových služeb a dodržování relevantních právních předpisů.
- Praktická implikace: při používání edukačních nástrojů, které sbírají postupy a výsledky žáků (např. nástroje typu „Pokrok v matematice“), ověřte, jaká data jsou odesílána a kde jsou uložena \cite{kb_ai_101_tipu_pdf_04e4a115_page_8_1}. Doporučujeme omezit nahrávání citlivých nebo přímo identifikovatelných údajů do externích služeb. Před spuštěním pilotu dokumentujte, jaká metadata se ukládají, kdo má k nim přístup, a stanovte dobu uchovávání dat. (general background knowledge)

\subparagraph{Spravedlnost / equita}
- Definice (general background knowledge): návrh aktivit a podpory tak, aby přístup k AI ani jeho výsledky nezvýhodňovaly systematicky některé skupiny studentů. Princip se vztahuje jak na sémantickou stránku výstupů (obsah a jazyk), tak na technický přístup (přístup k hardware a síťovým zdrojům).
- Praktické kroky: zajistěte alternativní způsoby přístupu (offline materiály, asistivní technologie) a pravidelně kontrolujte, zda generované materiály nepřevádějí nebo nezesilují existující biasy. Při demonstracích využijte ověřené edukační prostředí (např. Minecraft Education pro názorné ukázky principů) s ohledem na dostupnost a licencování \cite{kb_ai_101_tipu_pdf_04e4a115_page_15_2}. Provádějte testování materiálů na reprezentativním vzorku studentů a zaznamenávejte zjištěné disproporce pro následné úpravy obsahů.

\subparagraph{Akademická integrita}
- Definice (general background knowledge): jasné normy toho, co je povoleno, povinnosti deklarace použití AI a hodnocení procesu místo pouhého produktu. Dále sem patří pravidla pro spoluvlastnictví výsledků a uznání přínosu nástrojů AI v tvůrčích úkolech.
- Praktická doporučení: přepracujte zadání tak, aby hodnotila proces (drafty, záznamy postupu, interview), a zvažte integraci polí pro deklaraci AI do odevzdávacích šablon. Pokud povolujete AI nástroje, požadujte reflexi o tom, jak byl nástroj použit. Doporučuje se také používat techniky ověřování originality a kontrolních otázek, které odhalí, zda student rozumí obsahu, který prezentuje. (general background knowledge)

\subparagraph{Zodpovědnost}
- Definice (general background knowledge): jasné určení, kdo nese odpovědnost za validaci výstupů AI (vyučující / garant / IT) a jak budou řešeny chybné či škodlivé výstupy. Zahrnuje také procedury pro hlášení incidentů a postupy nápravy v případě škodlivého či nevhodného obsahu.
- Doporučení pro provoz: vždy fungujte s lidským validátorem (human‑in‑the‑loop) pro kritické rozhodnutí; dokumentujte prompty a verze modelů pro případ auditů. Pokud používáte enterprise nástroje pro přípravu testů či materiálů (např. Microsoft 365 Copilot pro generování testů ve Forms), definujte roli lidského přezkumu při konečné editaci \cite{kb_ai_101_tipu_pdf_04e4a115_page_36_1}. Vypracujte jasný eskalační scénář pro situace, kdy výstup AI obsahuje chybná data nebo eticky problematický obsah, a určete kontaktní osoby zodpovědné za rychlou nápravu.

\paragraph{Krátký praktický checklist před pilotem (shrnutí, general background knowledge)}
\begin{itemize}
  \item Vymezte pedagogický cíl pilotu a očekávaný přínos.
  \item Proveďte základní posouzení rizik a DPIA (pokud jsou zpracovávána osobní data).
  \item Ověřte licenční a datové podmínky nástroje (kde jsou data uložena, kdo je zpracovává).
  \item Zajistěte human‑in‑the‑loop pro finální validaci citlivých výstupů.
  \item Připravte krátkou formulaci pro sylabus/informaci pro studenty o pravidlech použití AI.
  \item Zdokumentujte použité prompty, verze modelů a testovací scénáře pro budoucí revizi.
  \item Naplánujte monitorování pilotu (sběr zpětné vazby od studentů, metriky úspěšnosti).
  \item Určete odpovědné osoby a kontaktní postup pro případ bezpečnostních nebo etických incidentů.
\end{itemize}

\paragraph{}
Pro rychlé praktické příklady a šablony (označené texty do sylabů, checklisty DPIA a prompt‑šablony) využijte centrální repozitář VUT (ai.vut.cz), kde budou dostupné i video‑tutoriály a aktualizované materiály pro pedagogy \cite{kb_ai_101_tipu_pdf_04e4a115_page_15_2}. Repozitář obsahuje rovněž doporučené formulace pro informování studentů, vzory pro záznam použitých promptů a příklady jednoduchých DPIA formulářů, které lze upravit pro konkrétní piloty. Doporučujeme pravidelně kontrolovat aktualizace repozitáře před každým větším nasazením.%

%
\section{Ochrana soukromí a GDPR‑kompatibilní postupy}%
Praktické postupy pro práci s osobními údaji studentů: princip minimalizace dat, anonymizace/pseudonymizace, provedení DPIA (posouzení dopadů) a dohody o zpracování (DPA) s poskytovateli. Doporučení kroků pro zpracování dat ve výuce a příklady opatření ke snížení rizik.%
\subsection{Zásady zpracování osobních údajů ve výuce}%
Stručné shrnutí hlavních principů, které ovlivňují návrh výuky s AI (princip minimalizace dat, účelové omezení, transparentnost vůči studentům) a stručný přehled, kdy zapojit bezpečnostní/PRIV týmy fakult. Slouží jako orientační rámec před praktickými kroky.%


%
Tato krátká sekce shrnuje zásady pro návrh výuky s generativní AI jako rámec před DPIA, smlouvami a technickým řešením.
Základní principy
\begin{itemize}
  \item Minimalizace dat; nezpracovávejte citlivá či přímo identifikovatelná data v externích službách \cite{kb_ai_101_tipu_pdf_04e4a115_page_8_1}.
  \item Účelové omezení, transparentnost a tam, kde možno, anonymizace/pseudonymizace.
  \item DPIA při plánování pilotů a širších nasazení; DPIA povinná při středních/vysokých rizicích \cite{kb_malinka_chatgpt_ve_vyuce_dva_roky_pote_2024_pdf_90bc3aaf_page_12_1}.
\end{itemize}
Kdy zapojit bezpečnost/PRIV: při zpracování osobních/citlivých dat, ukládání mimo EU nebo při pilotu sbírajícím behaviorální metriky či ovlivňující hodnocení; před spuštěním proveďte checklist (minimalizace PII, DPIA/odůvodnění, ověření DPA, lidská validace).%

%
\subsection{Minimalizace dat a návrh datového toku v kurzech}%
Praktické postupy pro omezení sběru a zpracování (co sbírat, co není nutné), doporučený návrh datového toku v rámci kurzu (sběr, zpracování, uchovávání, mazání) a krátké příklady řešení pro běžné výukové scénáře.%


%
Tato krátká sekce shrnuje kroky pro omezení sběru a bezpečný datový tok; uplatněte minimalizaci, účelové omezení a anonymizaci/pseudonymizaci \cite{kb_ai_101_tipu_pdf_04e4a115_page_8_1}.

\begin{itemize}
  \item Definujte účel polí; sbírejte jen nezbytné údaje, používejte pseudonymy/ID místo jména/fotek \cite{kb_ai_101_tipu_pdf_04e4a115_page_8_1}.
  \item Anonymizujte/pseudonymizujte před sdílením a omezte přístupy; zabraňte nahrávání citlivých souborů do neověřených cloudů, používejte interní nebo enterprise úložiště \cite{kb_ai_101_tipu_pdf_04e4a115_page_8_1, kb_ai_101_tipu_pdf_04e4a115_page_53_2}.
  \item U pilotů prověřte DPIA a DPA s poskytovatelem; stanovte retenci a role human-in-the-loop \cite{kb_malinka_chatgpt_ve_vyuce_dva_roky_pote_2024_pdf_90bc3aaf_page_12_1}.
\end{itemize}

Krátký doporučený datový tok:
\begin{enumerate}
  \item Sběr (minimální pole) → předzpracování (kontrola PII, anonymizace) → šifrované interní/enterprise uložení.
  \item Zpracování: AI pracuje jen s anonymizovanými daty v rámci účelu; retence a trvalé mazání záloh podle plánu.
\end{enumerate}

Praktické příklady: formativní kvíz — jen anonymní skóre a čas; programátorský projekt — školní repozitář; při externích AI službách vyžádejte souhlas a odstranění PII \cite{kb_ai_101_tipu_pdf_04e4a115_page_53_2}.%

%
\subsection{Anonymizace a pseudonymizace: metody a limity}%
Popis rozdílu mezi anonymizací a pseudonymizací, praktické techniky (odstranění identifikátorů, agregace, generalizace, hashing/pseudonymy) a upozornění na omezení a zbytková rizika při použití v učebních datech.%


%
% Anonymizace a pseudonymizace: metody a limity
\textbf{General background knowledge:} Anonymizace = odstranění nebo transformace údajů tak, aby subjekt nebyl identifikovatelný; pseudonymizace = nahrazení identifikátorů stabilním pseudonymem, obnova identity možná jen s klíčem.
Praktická doporučení: princip minimalizace — sbírat jen nezbytná pole; místo jmen/fotek používat pseudonymy/ID \cite{kb_ai_101_tipu_pdf_04e4a115_page_8_1}.
Před sdílením s externími nástroji aplikujte anonymizaci/pseudonymizaci a omezte přístupy; neuploadujte studentské fotografie do neověřených cloudů \cite{kb_ai_101_tipu_pdf_04e4a115_page_8_1, kb_ai_101_tipu_pdf_04e4a115_page_53_2}.
Rychlý pracovní postup:
\begin{enumerate}
\item Definujte účel a uchovávejte jen nezbytná pole; pokud není třeba identifikovat, neukládejte ID \cite{kb_ai_101_tipu_pdf_04e4a115_page_8_1}.
\item Před sdílením zkontrolujte PII, aplikujte anonymizaci/pseudonymizaci a odstraňte citlivá metadata \cite{kb_ai_101_tipu_pdf_04e4a115_page_8_1}.
\item Ukládejte šifrovaně v interním úložišti, nastavte přístupy; v pilotech ověřte DPIA/DPA a retenci \cite{kb_ai_101_tipu_pdf_04e4a115_page_8_1, kb_malinka_chatgpt_ve_vyuce_dva_roky_pote_2024_pdf_90bc3aaf_page_12_1}.
\end{enumerate}
Krátké příklady: kvíz — anonymní skóre nebo interní ID; programátorský projekt — školní repozitář a minimalizace PII v commitech; sdílení s externí AI — sjednat odstranění PII a rozsah v DPA \cite{kb_ai_101_tipu_pdf_04e4a115_page_8_1, kb_ai_101_tipu_pdf_04e4a115_page_53_2, kb_malinka_chatgpt_ve_vyuce_dva_roky_pote_2024_pdf_90bc3aaf_page_12_1}.
Limity a checklist: anonymizace není absolutní — posuzujte reziduální riziko a DPIA; držte mapovací tabulky pseudonymů odděleně; před sdílením ověřte nezbytnost polí, pseudonymizaci, odstranění metadat/fotek a pro piloty DPIA/DPA s lidským dohledem.%

%
\subsection{Posouzení dopadů na ochranu osobních údajů (DPIA) pro AI ve výuce}%
Praktický návod k provedení DPIA zaměřeného na nasazení generativní AI ve výuce: klíčové kroky, co zdokumentovat, ukázková struktura posouzení a kdy je DPIA doporučená/nezbytná z hlediska rizik.%


%
\textbf{Krátké shrnutí a kdy začít DPIA.} Před pilotním nasazením zkontrolujte datový tok; u pilotů a při sdílení studentských dat s externími službami doporučena DPIA/DPA \cite{kb_ai_101_tipu_pdf_04e4a115_page_8_1}.

\textbf{Rychlý pracovní postup pro pedagoga.}
\begin{enumerate}
\item Omezte sběr na nezbytné údaje; neukládejte ID, pokud není nutné \cite{kb_ai_101_tipu_pdf_04e4a115_page_8_1}.
\item Před sdílením odstraňte PII, anonymizujte/pseudonymizujte a vymažte citlivá metadata \cite{kb_ai_101_tipu_pdf_04e4a115_page_8_1}.
\item Šifrujte a omezte přístupy; pro piloty ověřte DPIA/DPA a zajistěte lidský dohled \cite{kb_ai_101_tipu_pdf_04e4a115_page_8_1, kb_malinka_chatgpt_ve_vyuce_dva_roky_pote_2024_pdf_90bc3aaf_page_12_1}.
\end{enumerate}

\textbf{Hlavní položky DPIA a mitigace.}
\begin{itemize}
\item Popis účelu, datový tok, úložiště a přístupy; kategorie údajů a hodnocení rizik; mitigace (minimalizace, pseudonymizace, šifrování, DPA) a plán dohledu.
\end{itemize}

\textbf{Příklady a tip.} „Kvíz v Moodle“: anonymní skóre, externí export jen s DPA; „Projekt s kódem“: odstraňovat metadata a pseudonymizovat \cite{kb_ai_101_tipu_pdf_04e4a115_page_8_1}. Tip: fotografie, citlivá data nebo automatizované rozhodování bez lidského dohledu indikují potřebu DPIA a konzultaci \cite{kb_ai_101_tipu_pdf_04e4a115_page_8_1, kb_malinka_chatgpt_ve_vyuce_dva_roky_pote_2024_pdf_90bc3aaf_page_12_1}.%

%
\subsection{Smlouvy s poskytovateli a dohody o zpracování (DPA)}%
Seznam klíčových bodů do DPA a smluv s poskytovateli AI/nástrojů (role správce/zpracovatele, bezpečnostní požadavky, podmínky uchovávání dat, možnosti auditů a odpovědnost), praktické tipy pro vyjednávání s komerčními službami.%


%
\noindent Před uzavřením dohody o zpracování (DPA) u poskytovatele AI si ověřte výsledky DPIA a datový tok pro daný pilot nebo službu; u pilotů a při sdílení studentských dat je provedení DPIA doporučené \cite{kb_ai_101_tipu_pdf_04e4a115_page_8_1}.

Níže klíčové body, které by DPA měla pokrýt:

\begin{itemize}
  \item Role správců a zpracovatelů; sub‑procesoři a pravidla jejich schvalování.
  \item Bezpečnost (šifrování v přenosu i v klidu, správa klíčů, logging, monitoring) a incident response se SLA notifikací.
  \item Uchovávání a likvidace dat; data residency / lokalita zpracování.
  \item Práva na audit a reporting; odpovědnost, odškodnění a SLA (dostupnost, podpora).
  \item Požadavky na lidský dohled u automatizovaných rozhodnutí.
\end{itemize}

Praktický tip: upřednostňujte enterprise/licencované varianty s jasnou DPA a možností hostingu v přijatelné jurisdikci (např. edukační nabídky Microsoftu) \cite{kb_ai_101_tipu_pdf_04e4a115_page_3_1}; před podpisem konzultujte DPA s právním oddělením a IT.%

%
\subsection{Praktický checklist a opatření ke snížení rizik}%
Jednostránkový/checklistový seznam kroků pro pedagoga před i během použití AI ve výuce: povolení dat, pseudonymizace, retence, přístupová práva, lokální zpracování vs. cloud (např. doporučení: nenahrávat studentské fotky do neznámých služeb), monitorování a incident response.%


%
Krátký praktický checklist — rychlá kontrola před i během použití generativní AI ve výuce.

\begin{itemize}
\item \textbf{Před nasazením:} Proveďte DPIA pokud se zpracovávají studentská data a uveďte ji při vyjednávání DPA;\cite{kb_ai_101_tipu_pdf_04e4a115_page_8_1} zkontrolujte DPA (role, sub‑procesory, retention, audit) a vyžadujte SLA pro incident response;\cite{kb_ai_101_tipu_pdf_04e4a115_page_8_1} upřednostněte enterprise/licencované varianty s možností hostingu v přijatelné jurisdikci;\cite{kb_ai_101_tipu_pdf_04e4a115_page_3_1} minimalizujte sběr osobních údajů a neodesílejte studentské fotografie do neznámých cloudů.\cite{kb_ai_101_tipu_pdf_04e4a115_page_53_2}
\item \textbf{Během provozu:} Pseudonymizujte nebo agregujte data, nastavte role a přístupy s loggingem/monitoringem a sledujte používání a chybové stavy;\cite{kb_ai_101_tipu_pdf_04e4a115_page_8_1} mějte připravený postup incident response (kontakty IT/právo, notifikace podle DPA/SLA).\cite{kb_ai_101_tipu_pdf_04e4a115_page_8_1}
\item \textbf{Krátký „cheat‑sheet“ pro incident:} Detect, Contain, Notify, Remediate.
\end{itemize}%

%
\section{Enterprise a lokální řešení: volba infrastruktury a bezpečnost dat}%
Kritéria pro rozhodování mezi cloudovými enterprise řešeními a lokálními/nebo on‑premise nasazení (včetně RAG), včetně aspektů bezpečnosti, řízení přístupu, logování a smluvních požadavků. Praktické doporučení pro IT spolupráci na VUT.%


%
Volba mezi cloudovým enterprise řešením a lokálním / on‑premise provozem se řídí třemi hlavními kritérii: citlivostí zpracovávaných dat, požadavky na integraci do existujícího ekosystému (SSO, LMS, katalogy uživatelů) a provozní/rozpočtovou dostupností. Před rozhodnutím proveďte krátké posouzení rizik a rozsahu — zmapujte scénáře použití, datové toky a potenciální expozice. Zahrňte do posouzení také otázky provozní resilience, zálohování a požadavky na obnovu po havárii. Pokud v řešení budou zpracovávána osobní nebo citlivá studentská data, je doporučené provést DPIA jako součást plánování pilotu; výsledky DPIA pomohou specifikovat technické i smluvní požadavky vůči poskytovateli a zohlednit právní rizika \cite{kb_malinka_chatgpt_ve_vyuce_dva_roky_pote_2024_pdf_90bc3aaf_page_12_1}. DPIA rovněž slouží jako podklad pro interní rozhodovací procesy a pro komunikaci s vedením školy či fakultou.

Praktická doporučení pro rozhodovací proces (VUT)
\begin{enumerate}
\item Klasifikujte data podle citlivosti a účelu zpracování; vytvořte jednoduchou matici klasifikace (veřejné, interní, citlivé, zvlášť chráněné) a pro scénáře se studentskými údaji naplánujte DPIA. Zahrňte výstupy DPIA do požadavků na bezpečnostní opatření a do smluvního vyjednávání s poskytovatelem \cite{kb_malinka_chatgpt_ve_vyuce_dva_roky_pote_2024_pdf_90bc3aaf_page_12_1}. Taková klasifikace umožní rychlé rozhodování o tom, které datové sady mohou být zpracovávány v cloudu a které musí zůstat lokálně.
\item Preferujte enterprise/licencované varianty pro výukové scénáře, zejména pokud chcete využít integrované školní funkce a centralizovanou správu uživatelů. Enterprise balíčky často nabízejí administrativní nástroje, SLA a integrace (příklady: nástroje Microsoft 365 a školní licence obsahující edukační služby), které usnadní nasazení a podporu \cite{kb_ai_101_tipu_pdf_04e4a115_page_15_2, kb_ai_101_tipu_pdf_04e4a115_page_36_1}. Při posuzování variant porovnejte nejen cenu licence, ale i dostupnost podpory, školení a možnosti rozšíření funkcí.
\item Požadujte jasnou smlouvu o zpracování (DPA) a SLA zahrnující incident response, retenci dat, bezpečnostní standardy a auditní možnosti; výsledky DPIA jsou užitečným podkladem při specifikaci těchto požadavků. Dbejte na definice odpovědnosti za sub‑procesory a na dostupnost logů pro forenzní šetření \cite{kb_ai_101_tipu_pdf_04e4a115_page_8_1, kb_malinka_chatgpt_ve_vyuce_dva_roky_pote_2024_pdf_90bc3aaf_page_12_1}. V dokumentech si vyjasněte postupy oznámení bezpečnostních incidentů a lhůty pro nápravu.
\item Ověřte technické možnosti integrace: SSO/autorizace, centrální správa přístupů a auditní logy — tyto prvky usnadňují provoz a zajišťují sledovatelnost chování systému; zkontrolujte také podporu šifrování v klidu i při přenosu a možnosti exportu dat (general background knowledge). Zjistěte, jak snadno lze z řešení migrovat data či uzavřít službu bez ztráty provozních záznamů a metrik.
\item Pro piloty vyžadujte minimální množství PII v indexovaných zdrojích a tam, kde je to možné, pseudonymizujte nebo agregujte data před nahráním do externích služeb. Uvažujte o testovacích datech a anonymizovaných sadách, abyste snížili riziko úniku citlivých informací \cite{kb_ai_101_tipu_pdf_04e4a115_page_8_1}. Doporučuje se také zavést pravidelné kontroly indexovaných zdrojů a proces pro rychlé odstranění citlivých záznamů.
\end{enumerate}

Konkrétní kroky spolupráce s IT a právním oddělením (postup)
\begin{itemize}
\item Krátký assessment: garant předmětu a IT společně vyplní formulář s typem dat, očekávanými uživateli, účelem zpracování a navrhovanými integracemi; pokud se zpracovávají studentská data, iniciujte DPIA a zapojte právní oddělení dříve než v pozdějších fázích projektu \cite{kb_malinka_chatgpt_ve_vyuce_dva_roky_pote_2024_pdf_90bc3aaf_page_12_1}. Tento krok by měl zahrnovat i odhad nákladů a zdrojů potřebných pro pilotní provoz.
\item Smluvní příprava: na základě výsledků DPIA specifikujte požadavky na DPA/SLA (retence, sub‑procesory, auditní práva, incident response) a v případě potřeby vyžádejte hosting/rezidenci dat v přijatelné jurisdikci. Zahrňte také metriky dostupnosti, sankce za porušení a postupy pro bezpečnostní incidenty (general background knowledge; doporučeno zahrnout do DPA). Jasně definované smluvní podmínky usnadní řešení sporů a zkrátí dobu reakce v případě incidentu.
\item Technická integrace: nastavte SSO a role‑based access, zapněte auditní logování a metriky využití; v pilotní fázi měřte chybovost, výkonnost a bezpečnostní incidenty pro rozhodnutí o dalším rozšíření. Dokumentujte integrační kroky, zálohovací politiky a plán obnovy po havárii. Plánujte také pravidelné revize konfigurace a testy obnovy dat.
\item Testování funkcí a pedagogické ověření: ověřte, že funkce určené pro výuku (např. automatizovaná tvorba testů nebo specializované edukační moduly) fungují v rámci vybraného enterprise balíčku. Proveďte uživatelské testy s garanty předmětů a sbírejte zpětnou vazbu; konkrétním příkladem je Microsoft 365 Copilot ve spojení s Forms, který umí generovat testy a vysvětlení, což ukazuje integraci pedagogických funkcí do enterprise prostředí \cite{kb_ai_101_tipu_pdf_04e4a115_page_36_1}. Zpětnou vazbu dokumentujte a zapracujte do rozhodnutí o nasazení či úpravě integračních nastavení.
\end{itemize}

Poznámka k RAG a lokálnímu indexování
\begin{itemize}
\item Retrieval‑Augmented Generation (RAG) a indexace interních dokumentů mohou snížit potřebu ukládání PII přímo do externích modelů, ale vyžadují přísné řízení přístupů, revizi indexovaných zdrojů a kontrolu, jaké dokumenty jsou indexovány. Implementujte pravidla pro redakci citlivých polí, dočasné cache a minimalizaci kontextu zasílaného do modelu; při návrhu RAG architektury spolupracujte úzce s oddělením IT a právem, aby byly pokryty požadavky vyplývající z DPIA a DPA (general background knowledge). Zvažte rovněž možnosti auditu indexů a pravidelné automatizované kontroly kvality indexovaných dat.
\end{itemize}

Shrnutí pro garanty předmětů: při rozhodování upřednostněte enterprise/licencované řešení tam, kde chcete garantovat školní funkce, centrální správu uživatelů a SLA; tam, kde je kritická kontrola nad daty nebo jsou vyšší právní rizika, zvažte lokální/on‑premise variantu. V obou případech považujte DPIA za výchozí krok plánování pilotu a použijte její výstupy pro smluvní i technické požadavky vůči poskytovateli \cite{kb_malinka_chatgpt_ve_vyuce_dva_roky_pote_2024_pdf_90bc3aaf_page_12_1, kb_ai_101_tipu_pdf_04e4a115_page_15_2}. Dodržení těchto kroků zlepší bezpečnostní postavení projektu, zkrátí dobu nasazení a sníží provozní rizika spojená s nasazením AI nástrojů ve výuce.%

%
\section{Krátký přehled evropského AI Act (2024) a důsledky pro vzdělávání}%
Stručné vymezení relevantních částí AI Act pro akademické prostředí a praktické implikace pro výuku a instituční nasazení AI (rizikové klasifikace, povinnosti transparentnosti a odpovědnosti). Nízkoslovný přenos na interní politiky předmětů.%


%
Toto krátké shrnutí uvádí klíčové body, které by měli pedagogové a garanty předmětů znát při zohlednění evropského AI Act (2024) v kontextu výuky. Vzhledem k rozsahu a rychlému vývoji legislativy následující tvrzení slouží jako stručná orientace; konkrétní právní výklad konzultujte s právním oddělením (general background knowledge).

\begin{itemize}
  \item Hlavní principy AI Act (stručně): legislativa pracuje s rizikově‑orientovaným přístupem k AI systémům, zavádí povinnosti transparentnosti vůči uživatelům a požadavky na odpovědnost poskytovatelů u systémů považovaných za „vysoké riziko“ (general background knowledge).
  \item Povinnosti transparentnosti: u generativních systémů mohou existovat požadavky na informování uživatele, že výstup pochází od AI, a na dokumentaci datových zdrojů a použitých bezpečnostních opatření (general background knowledge).
  \item Dopady na poskytovatele a instituce: pro systémy s vyšším rizikem mohou platit povinnosti provést posouzení souladu, zavést lidský dozor, vést záznamy (logs) pro audit a zajistit mechanismy pro řešení škod či chyb (general background knowledge).
\end{itemize}

Praktické implikace pro výuku a nasazení AI na VUT (doporučené kroky)
\begin{enumerate}
  \item Klasifikujte konrétní edukační scénáře podle rizika: rozlište jednoduché podpůrné nástroje (např. generování prezen­tací) od řešení, která zpracovávají studentská osobní či citlivá data nebo která ovlivňují hodnocení — pro tyto případy uvažujte přísnější režim řízení rizik (general background knowledge).
  \item DPIA při středním či vysokém riziku: pokud je v projektu zpracováváno studentské PII nebo jiné citlivé údaje, zahrňte posouzení dopadů na ochranu osobních údajů (DPIA) již do plánování pilotu; výsledky DPIA by měly ovlivnit technické a smluvní požadavky vůči poskytovateli \cite{kb_malinka_chatgpt_ve_vyuce_dva_roky_pote_2024_pdf_90bc3aaf_page_12_1}.
  \item Preferujte enterprise/licencované varianty tam, kde chcete centrální správu uživatelů, SLA a lepší kontrolu nad daty; enterprise balíčky často nabízejí administrativní nástroje a integrace užitečné pro akademické prostředí \cite{kb_ai_101_tipu_pdf_04e4a115_page_15_2}.
  \item Minimalizace PII v indexech a RAG: při použití retrieval‑augmented řešení omezte množství indexovaných osobních údajů, pseudonymizujte data tam, kde je to možné, a nastavte pravidla pro revizi a mazání indexovaných záznamů (general background knowledge).
  \item Transparentnost do sylabů a komunikace se studenty: do sylabů uveďte, zda a kde je AI používána, jaké jsou limity (ověřování faktů, lidská validace) a jak mají studenti deklarovat použití AI při odevzdávání prací (general background knowledge).
  \item Smlouvy a technické požadavky: při výběru dodavatele požadujte DPA/SLA s ustanoveními o retenci dat, sub‑procesorech, auditních záznamech a incident response; výsledky DPIA použijte jako podklad pro tyto požadavky \cite{kb_malinka_chatgpt_ve_vyuce_dva_roky_pote_2024_pdf_90bc3aaf_page_12_1}.
\end{enumerate}

Krátký checklist pro garanta předmětu (rychlé kroky)
\begin{enumerate}
  \item Zmapovat scénáře použití AI ve výuce a označit jejich potenciální riziko (nízké / střední / vysoké) (general background knowledge).
  \item Pokud je riziko střední či vysoké: spustit DPIA a zapojit právní/PRIV týmy \cite{kb_malinka_chatgpt_ve_vyuce_dva_roky_pote_2024_pdf_90bc3aaf_page_12_1}.
  \item V sylabu explicitně uvést pravidla pro použití AI a požadavek na deklaraci tam, kde je to relevantní (general background knowledge).
  \item Preferovat enterprise řešení pro scénáře vyžadující správu uživatelů, logování a SLA; v jiných případech zvážit lokální/on‑premise variantu (general background knowledge; \cite{kb_ai_101_tipu_pdf_04e4a115_page_15_2}).
  \item Definovat lidský dozor (human‑in‑the‑loop) pro všechny výstupy AI, které mají dopad na hodnocení nebo rozhodování o studentech (general background knowledge).
\end{enumerate}

Poznámka: tento přehled je záměrně nízkoslovný a zaměřený na přenos do interních politik předmětů. Pro detailní právní výklad a závazné požadavky vyhledejte oficiální znění AI Act a konzultaci s právním oddělením VUT (general background knowledge).%

%
\section{Akademická integrita v éře AI}%
Analýza vlivu generativní AI na plagiátorství, autorství a fair play; návrh politiky a postupů pro dodržování akademické integrity. Obsahuje navrhovanou kategorizaci úloh (zakázáno / povoleno s deklarací / vyžadováno využití AI) pro začlenění do předmětů.%


%
Úvodní poznámka. Toto oddílové shrnutí se zaměřuje na praktická doporučení a zásady pro udržení akademické integrity v době dostupných generativních nástrojů. Hlavní doporučení níže jsou uvedena jako praktické návrhy pro garanty předmětů a učitele; tam, kde jde o širší právní či institucionální závazky (DPIA, DPA, smluvní podmínky), odkazujte na příslušné interní postupy a právní konzultace (viz kapitoly 2.2 a 8). Některá doporučení jsou shrnuta jako obecné (general background knowledge).

Cíle oddílu
\begin{itemize}
  \item Pomoci garantům kategorizovat úlohy podle pravidelného nasazení AI a očekávané úrovně deklarace ze strany studentů (praktická politika).
  \item Návrh konkrétních postupů pro navržení úloh a workflow hodnocení, které snižují riziko nevhodného použití AI.
  \item Příklady formulací do sylabů a šablony pro pole deklarace použití AI při odevzdání.
\end{itemize}

Doporučená kategorizace úloh (pro zařazení do sylabu) — návrh (general background knowledge)
\begin{enumerate}
  \item Zakázáno: úlohy, u kterých použití AI by zásadním způsobem podkopalo hodnotu výstupu (např. krátké dozorované písemky, časově omezené in‑class writings). V sylabu uveďte jasné sankce a proces řešení podezření.
  \item Povoleno s deklarací: úlohy, kde je možné AI použít, ale vyžadujeme explicitní deklaraci, popis kroků a připojení procesních důkazů (drafty, verze, poznámky k práci) (general background knowledge).
  \item Vyžadováno využití AI (AI‑enhanced assignments): úlohy, které explicitně integrují AI jako nástroj — student musí popsat, jak AI použil, jaký prompt použil a doložit vlastní přidanou hodnotu (analýza, kritika, adaptace).
\end{enumerate}

Praktické konstrukce zadání a ověřitelnost (general background knowledge)
\begin{itemize}
  \item Preferujte procesně orientované body hodnocení (např. milníky, průběžné odevzdání částí, komentované verze), nikoli pouze konečný produkt.
  \item Požadujte artefakty, které je těžší získat z externího generátoru bez vlastní práce: pracovní poznámky, postupy řešení, lokální logy z vývojového repozitáře, screencasty vývoje nebo sběr původních dat.
  \item U technických úloh (např. matematika, programování) zařaďte požadavky na ukázku postupu — v některých nástrojích lze vyžadovat nahrání postupu řešení jako součást hodnocení (viz příklad z praxe v rámci školních nástrojů: povinnost nahrání postupu v aplikaci „Pokrok v matematice“) \cite{kb_ai_101_tipu_pdf_04e4a115_page_8_1}.
  \item Využívejte živé či polosimulované ústní ověření (spot‑checky, krátké obhajoby), zvláště u závěrečných prací či projektů, kde autenticita výstupu je klíčová (general background knowledge).
\end{itemize}

Formulace do sylabu — ukázky (copy‑\\ \&‑paste připravené)
\begin{itemize}
  \item Variantní text pro povolení s deklarací: ``Použití generativní AI je v tomto předmětu povoleno pouze se souhlasem vyučujícího a za podmínky, že při odevzdání student explicitně uvede: (1) které nástroje použil, (2) konkrétní prompty a verze, (3) vlastní úpravy a reflexi, a (4) přiloží pracovní verze/drafty.'' (general background knowledge)
  \item Variantní text pro zakázání: ``Použití externích generativních nástrojů k vytvoření obsahu zadání je zakázáno. Podezření na porušení integritní politiky bude řešeno podle akademických pravidel VUT.''
  \item Pole deklarace na odevzdání: krátké checkboxy a textové pole se záznamem: nástroj, prompt, popis přidané hodnoty (např. 150–300 znaků).
\end{itemize}

Rubrika hodnocení obsahující AI‑deklaraci — návrh polí (general background knowledge)
\begin{itemize}
  \item Transparentnost použití AI (0–5 bodů): student popíše rozsah použití AI a doloží prompty/konkrétní výstupy.
  \item Originalita a vlastní přínos (0–10 bodů): posuzuje se míra vlastní práce, úpravy a kritické evaluace AI‑výstupů.
  \item Doklad procesu (0–5 bodů): přiložené drafty, commity v repozitáři, screencasty či živá obhajoba.
\end{itemize}

Postupy při podezření z nevhodného použití (general background knowledge)
\begin{enumerate}
  \item Nejprve vyžádejte procesní důkazy (verze, drafty, artefakty) a krátkou obhajobu studenta.
  \item Proveďte krátké interview/ústní ověření, zaměřené na konkrétní kroky řešení.
  \item Pokud podezření přetrvává, zahajte interní proces podle akademických předpisů fakulty (dokumentovat průběh a rozhodnutí).
\end{enumerate}

Poznámka o detektorech a alternativách (general background knowledge)
\begin{itemize}
  \item Automatické „detektory AI“ mohou mít omezení a falešné pozitivy; místo spoléhání se pouze na ně kombinujte procesní ověřování, redesign úloh a požadavek transparentnosti v odevzdání.
\end{itemize}

Krátké doporučení pro implementaci na úrovni předmětu
\begin{enumerate}
  \item Na začátku semestru: aktualizujte sylabus podle jedné z výše uvedených kategorií a komunikujte očekávání studentům.
  \item Upravte rubriky tak, aby hodnotily proces i produkt a zahrnovaly pole pro AI‑deklaraci.
  \item Nasazujte postupy ověřování (drafty, ústní kontrola) u klíčových milníků.
  \item Zapojte metodické oddělení a PR/PRIV týmy při pilotním zavádění nových pravidel (zejména při zpracování PII nebo změnách v systému hodnocení).
\end{enumerate}

Závěrem. Cílem není zákaz technologií, ale udržení férovosti hodnocení a rozvoj dovedností studentů v kritickém používání AI. Implementace politiky by měla být praktická, srozumitelná a komunikovaná již v sylabu; kombinace procesních požadavků, deklarací a místních ověření je nejúčinnějším přístupem pro většinu předmětů (general background knowledge).%

%
\section{Šablony do sylabů a vzory citací AI}%
Hotové texty a formulace, které lze vložit do sylabů (policy o používání AI, požadavek na deklaraci), a doporučené způsoby citování výstupů AI v běžných citačních stylech. Praktické příklady, které pedagog může adaptovat.%


%
Níže jsou zkrácené formulace a šablony pro sylaby a pole deklarace; upravte podle předmětu a místních pravidel.

\bigskip

\textbf{Formulace do sylabu (kopírovat/upravit):}

\begin{quote}
\begin{PromptBlock}


Použití generativní AI je v tomto předmětu povoleno pouze se souhlasem vyučujícího a za podmínky, že při odevzdání student explicitně uvede: (1) které nástroje použil, (2) konkrétní prompty a verze, (3) vlastní úpravy a reflexi, a (4) přiloží pracovní verze/drafty.
\end{PromptBlock}


\end{quote}

\begin{quote}
\begin{PromptBlock}


Použití externích generativních nástrojů k vytvoření obsahu zadání je zakázáno. Podezření na porušení integritní politiky bude řešeno podle akademických pravidel VUT.
\end{PromptBlock}


\end{quote}

\bigskip

\textbf{Pole deklarace (návrh):}
\begin{itemize}
  \item Checkbox: \textit{Použil(a) jsem externí AI nástroj při tvorbě tohoto odevzdání.}
  \item Text (150–300 znaků): nástroj, verze, použité prompty a stručně vlastní přidaná hodnota.
  \item Povinná příloha: pracovní verze / commity / screencast či jiný doklad postupu.
\end{itemize}

\bigskip

\textbf{Hodnocení (návrh):}
\begin{itemize}
  \item Transparentnost použití AI (0–5): rozsah použití a doložení promptů/výstupů.
  \item Originalita a vlastní přínos (0–10): úpravy a kritická evaluace.
  \item Doklad procesu (0–5): drafty, commity, screencasty nebo obhajoba.
\end{itemize}

\bigskip

\textbf{Krátký checklist:}
\begin{enumerate}
  \item Nástroj a verze.
  \item Přesný prompt/postup.
  \item Popis úprav a vlastní analytický přínos.
  \item Procesní důkazy: drafty, commity, screencast, logy.
\end{enumerate}

\bigskip

\textbf{Doporučení k citování AI‑výstupů (stručně):}
\begin{itemize}
  \item APA: \textit{OpenAI. (2024). ChatGPT (verze X) [Large language model]. Výstup použit v úloze: popis a prompt.}
  \item IEEE: \textit{OpenAI, "ChatGPT (verze X)," 2024. Použito pro: [popis úlohy], prompt: "...".}
  \item Interní pole: \textit{Nástroj: ..., Prompt (zkráceně): "...", Vlastní přínos: "..."}
\end{itemize}

\bigskip

Poznámky z praxe:
\begin{itemize}
  \item U integrací v Microsoft 365 (např. Copilot) uveďte možnost generování výstupů do školních šablon \cite{kb_ai_101_tipu_pdf_04e4a115_page_44_2}.
  \item Ukázky AI ve výuce lze doplnit odkazy na Minecraft Education (Hour of Code) \cite{kb_ai_101_tipu_pdf_04e4a115_page_15_2}.
  \item U technických úloh požadujte procesní artefakty — inspirace: "Pokrok v matematice" \cite{kb_ai_101_tipu_pdf_04e4a115_page_8_1}.
\end{itemize}

Pro rychlé nasazení použijte tyto šablony jako výchozí bod a odkažte studenty na living document na ai.vut.cz; upravte text podle kategorie úloh (zakázáno / povoleno s deklarací / vyžadováno využití AI).%

%
\section{Proč se nespoléhat na detektory AI a alternativní přístupy}%
Kritické zhodnocení omezení detektorů AI a doporučení nástrojů a postupů místo nich: návrh úloh odolných vůči zneužití AI, požadavek na procesní důkazy (postupy, drafty) a využití formativního hodnocení.%


%
\noindent
Krátké shrnutí problému (general background knowledge): automatické „detektory AI“ (nástroje, které se snaží rozpoznat, zda text vytvořil člověk nebo model) mají v praxi významná omezení: vysoký podíl falešně pozitivních i falešně negativních výsledků, závislost na konkrétních datech/tréninku detektoru a snadná obcházení úpravami textu. Z tohoto důvodu by pedagogická praxe neměla stavět na detektorech jako na primárním nebo rozhodujícím nástroji při posuzování akademické integrity. (general background knowledge)

\bigskip

\noindent
Praktický přístup doporučený místo spoléhání se na detektory:
\begin{enumerate}
  \item Návrh úloh odolných vůči triviálnímu využití AI
    \begin{itemize}
      \item Preferujte procesně orientované úlohy: postupné milníky, pracovní verze, commit‑history, průběžné přihlášky a záznamy o postupu.
      \item Využívejte lokálně specifická nebo datově závislá zadání (lokální data, vlastní experimenty, pozorování terénu), která je těžké zcela vygenerovat „z gruntu“ externím AI.
      \item Zařazujte časově omezené in‑class writings nebo krátké testy na porozumění v průběhu výuky.
    \end{itemize}

  \item Požadavek na procesní důkazy při odevzdání
    \begin{itemize}
      \item Vyžadujte pracovní verze, poznámky, commity z repozitáře nebo screencasty editace jako součást odevzdání (viz šablony do sylabů).
      \item Do LMS přidejte pole deklarace: checkbox, krátký text s uvedením nástroje/verze/promptu a povinnou přílohu s pracovním materiálem (viz formulace v sekci \emph{Šablony do sylabů a vzory citací AI}).
    \end{itemize}

  \item Hodnocení procesu a transparentnost v rubrikách
    \begin{itemize}
      \item Zahrňte do rubrik explicitní kritéria vztahující se k použití AI a k procesnímu důkazu (např. Transparentnost použití AI 0–5; Originalita a vlastní přínos 0–10; Doklad procesu 0–5). Tato jednoduchá škála lze převzít a upravit podle typu úkolu (viz šablony do sylabů).
    \end{itemize}

  \item Alternativy k automatickému odsouzení: ověření porozumění
    \begin{itemize}
      \item Realizujte krátká ústní ověření nebo follow‑up pohovory k ověření postupu a porozumění (strukturou otázek zaměřených na proces).
      \item Požadujte reflexi: stručný text, kde student popíše rozhodnutí, zdroje a konkrétní kroky, které vedly k výsledku.
    \end{itemize}

  \item Formativní přístupy a prevence
    \begin{itemize}
      \item Využívejte formativní hodnocení a průběžnou zpětnou vazbu, aby studenti rozvíjeli dovednosti psaní a argumentace nezávisle na jednorázovém „výsledku“.
      \item Nabídněte workshopy o etickém použití AI, o tom jak deklarovat výstupy a jak dokumentovat proces (proaktivní transparentnost snižuje pokušení podvádět).
    \end{itemize}
\end{enumerate}

\bigskip

\noindent
Krátký praktický checklist pro vyučující (před řešením podezření)
\begin{enumerate}
  \item Zkontrolujte, zda student odevzdal požadované procesní důkazy (drafty, commity, screencast). (viz šablona pole deklarace)
  \item Projděte timeline práce v repozitáři nebo logy editací; vyžádejte krátké ústní doplnění či obhajobu postupu.
  \item Pokud je podezření, postupujte podle institucionálních pravidel o akademické integritě; nevyvozujte závěry pouze z automatických detektorů. (general background knowledge)
  \item Pokud chcete monitorovat rizika systematicky, navrhněte redesign zadání tak, aby bylo ověřování jednodušší a méně nákladné (milníky, reflexe, lokální data).
\end{enumerate}

\bigskip

\noindent
Závěr: místo spolehnutí se na „černé skříňky“ detektorů doporučujeme kombinaci didaktického designu (odolná zadání), povinné dokumentace procesu, transparentních rubrik a krátkých ověření porozumění. Tím se snižuje motivace ke zneužití AI a současně se posiluje učení a akademická integrita. (general background knowledge)%

%
\section{Hodnocení a rubriky v době AI}%
Navrhované hodnoticí kritérium a rubriky přizpůsobené použití AI (hodnocení procesu vs. produktu, důraz na metody a reflexi), včetně příkladů měřitelných kritérií pro různé typy úloh.%


%
Když v hodnocení zohledníme dostupnost generativní AI, posun je z hodnocení produktu k hodnocení procesu, reflexe a doloženému přínosu studenta. Praktickým jádrem jsou tři prvky, které by každá rubrika měla obsahovat: pole pro deklaraci použití AI, metriky pro transparentnost a procesní důkazy, a disciplinární kritéria (správnost, argumentace, kódování apod.).

\medskip

Návrh základních AI‑aware položek do rubriky (doporučené škály)
\begin{itemize}
\item Transparentnost použití AI (0–5): deklarace nástrojů, promptů a stručné vysvětlení vlastního přínosu.
\item Originalita a vlastní přínos (0–10): úpravy, kritické zhodnocení a přidaná hodnota nad vygenerovaný obsah.
\item Doklad procesu (0–5): drafty, commit‑history, screencast nebo jiný procesní důkaz.
\end{itemize}

Škály vložte do rubriky a přiřaďte váhy podle typu úlohy.

\medskip

Praktické vzory rubrik pro tři typy úloh
\begin{enumerate}
\item Esej / písemná práce (příklad váh: Originalita 40\%, Transparentnost 20\%, Obsah 30\%, Styl 10\%)
  \begin{itemize}
  \item Originalita (0–10): co student upravil a proč; Transparentnost (0–5): nástroj/prompt/verze; Doklad procesu (0–5): drafty nebo historie; Obsah a argumentace (0–30).
  \end{itemize}

\item Programovací úloha / projekt (příklad: Správnost 50\%, Proces 30\%, Kód 20\%)
  \begin{itemize}
  \item Správnost (0–50): funkčnost podle testů; Doklad procesu (0–15): commit‑history, issue tracker, minimálně X commitů nebo screencast; Transparentnost (0–5); Kódová kvalita (0–30).
  \end{itemize}

\item Laboratorní zpráva / praktické cvičení (váhy přizpůsobit)
  \begin{itemize}
  \item Autentičnost dat (0–30): vlastní sběr vs. externí dataset; Procesní důkazy (0–30): protokoly, fotografie, surová data, skripty; Reflexe AI (0–10); Metodika a interpretace (0–30).
  \end{itemize}
\end{enumerate}

Implementační doporučení pro vyučující
\begin{itemize}
\item Do LMS vložte pole deklarace: checkbox o použití AI, krátké pole pro nástroj/verzi/prompt (150–300 znaků) a povinnou přílohu s procesním důkazem — zvyšuje transparentnost.
\item Pro kvantitativní testy využijte nástroje pro tvorbu a vyhodnocení kvízů, což usnadní konzistentní skórování \cite{kb_ai_101_tipu_pdf_page_36_1}.
\item U matematických cvičení nebo automatizovaných modulů vyžadujte nahrání postupu řešení (funkce „Pokrok v matematice“) \cite{kb_ai_101_tipu_pdf_page_8_1}.
\item Zvažte multimodální aktivity pro ověření porozumění a diskusi o AI (např. výukové světy v Minecraft Education) \cite{kb_ai_101_tipu_pdf_page_15_2}.
\end{itemize}

Krátké pokyny pro hodnotitele
\begin{itemize}
\item Před hodnocením zkontrolujte deklaraci a procesní přílohy; chybějící položky vyžádejte doplnit.
\item Při nejasnostech preferujte krátký dialog se studentem k verifikaci postupu a použijte rubriky pro formativní zpětnou vazbu — uveďte konkrétní kroky ke zlepšení.
\end{itemize}

Závěr: Přizpůsobené rubriky kombinují disciplinární kritéria s explicitními položkami pro AI (transparentnost, procesní důkazy, reflexe). Takový rámec snižuje motivaci k neohlášenému využití AI a posiluje kritické ověřování vygenerovaných výstupů.%

%
\section{Praktické checklisty, šablony a odkazy pro učitele}%
Konzizní check‑listy pro provedení DPIA, uzavření DPA, nasazení enterprise řešení a revizi sylabů; odkazy na šablony (smlouvy, deklarace studentů) a doporučené interní kroky pro rychlou implementaci.%


%
Krátké praktické checklisty a odkazy pro rychlou implementaci — stručné pracovní listy určené garantům předmětu a administrativě katedry; níže uvedené checklisty jsou výchozí rámec (detailní šablony v přílohách, viz část Přílohy).

\medskip

\textbf{1) Checklist pro provedení DPIA:} Definovat rozsah a toky dat (účel, typy dat, systémy); identifikovat rizika pro práva subjektů; navrhnout opatření (minimizace, pseudonymizace, šifrování), dokumentovat závěry a naplánovat revize.

\medskip

\textbf{2) Checklist pro uzavření DPA / smlouvy s poskytovatelem AI:} Vymezit role (správce vs. zpracovatel) a odpovědnosti; požadovat bezpečnostní opatření a retenci; zahrnout právo na audit, incident‑response a SLA.

\medskip

\textbf{3) Checklist pro pilotní nasazení enterprise řešení:} Vybrat pilotní kurzy a měřitelné cíle; zajistit SSO, role a auditní logy; ověřit DPA/DPIA, nastavit eskalace k lidskému garantu a školení personálu.

\medskip

\textbf{4) Revize sylabu a oznámení pro studenty:} Jasně uvést povolení/zakázání/požadavek na AI; požadovat deklaraci nástroje, promptu a vlastního přínosu; poskytnout zdroje a šablony.

\medskip

\textbf{Rychlé odkazy a inspirace:} Minecraft Education a Hour of Code pro demonstrace principů \cite{kb_ai_101_tipu_pdf_04e4a115_page_15_2}; „Pokrok v matematice“ pro vyžadování postupu řešení \cite{kb_ai_101_tipu_pdf_04e4a115_page_8_1}; Designer (Microsoft) pro rychlé vizuály \cite{kb_ai_101_tipu_pdf_04e4a115_page_51_1}.

Viz přílohy pro kopírovatelné šablony: AI Policy Template (sylaby), Vzor DPIA, Vzor DPA a Prompt Library (sekce Přílohy).%

%
\chapter{AI V PŘÍPRAVĚ VÝUKY}%
Praktická kapitola věnovaná využití AI při tvorbě osnov, prezentací, cvičení a studijních materiálů. Obsahuje prompt šablony pro generování sylabů, 90minutových přednášek, návrhy aktivit a diferenciaci podle úrovně studentů. Ukazuje nástroje pro tvorbu slidů a vizuálních materiálů (Gamma, Beautiful.ai, DALL·E, Midjourney) a návody krok za krokem. Dále popisuje generování pracovních listů, kvízů a automatických řešení, sumarizaci vědeckých textů (NotebookLM, Elicit, Claude pro dlouhé texty) a ověřování výstupů. Součástí jsou hotové vzory pro redesign předmětů a ukázky z univerzit, kde takové postupy úspěšně zavedli.%
\section{Úvod a přehled kapitol}%
Krátké představení kapitoly: cíle, kompetence, jak využít materiál v přípravě předmětů na VUT a odkazy na etické/GDPR zásady a transparentnost použití AI (odvozeno z kontextu příručky).%


%
Tato kapitola poskytuje krátký, praktický úvod do části „AI v přípravě výuky“ s cílem nabídnout okamžitě použitelné kroky a odkazy pro garanty předmětů a vyučující na VUT. Najdete zde popis klíčových kompetencí, postupů pro rychlý pilot v kurzu a konkrétní inspiraci k nástrojům, které lze vyzkoušet ve výuce.

\medskip

\textbf{Cíle kapitoly.} Po přečtení této kapitoly by měl čtenář být schopen:
\begin{itemize}
  \item identifikovat oblasti předmětu vhodné pro doplnění AI‑podpory (příprava materiálů, cvičení, automatická zpětná vazba),
  \item použít připravené šablony a prompty k rychlé tvorbě sylabu, 90minutové lekce a pracovních listů (viz Přílohy),
  \item navrhnout jednoduchý pilotní workflow včetně základních kroků pro GDPR/etiku a transparentní komunikaci se studenty (viz Kapitola 2).
\end{itemize}
Poznámka: položky kompetencí jsou stručným souhrnem praktických cílů kapitoly (general background knowledge).

\medskip

\textbf{Jak kapitolu využít při přípravě předmětu — doporučený postup (rychlý checklist).}
\begin{enumerate}
  \item Definujte výukové výstupy, u kterých očekáváte přidanou hodnotu od AI (např. více variant cvičení, rychlejší formativní zpětná vazba).
  \item Vyberte jednu konkrétní oblast (např. generování pracovních listů nebo tvorba kvízů) a zvolte nástroj / šablonu z Příloh pro pilot.
  \item Připravte malý pilot: jeden týden / jedna lekce / malá skupina studentů; předem definujte metriky úspěchu (čas úspory, spokojenost studentů, kvalita materiálů).
  \item Ověřte požadavky na ochranu osobních údajů a transparentnost (informujte studenty, jaký nástroj se používá a jak deklarovat použití AI) — viz Kapitola 2 a šablony do sylabu.
  \item Po pilotu proveďte rychlou reflexi a upravte šablony/prompty; pokud pilot uspěje, rozšiřte nasazení postupně.
\end{enumerate}
(Body 3–5 jsou praktické doporučení odvozené z implementačních vzorů v příručce; konkretizace opatření k ochraně dat je rozvedena v Kapitole 2.)

\medskip

\textbf{Krátká inspirace — konkrétní nástroje a příklady k vyzkoušení.}
\begin{itemize}
  \item Pro názorné uvedení principů strojového učení a AI ve výuce lze využít vzdělávací svět Minecraft Education; v praxi jsou dostupné připravené hodiny (Hour of Code), které pomáhají žákům pochopit základní principy AI a programování \cite{kb_ai_101_tipu_pdf_04e4a115_page_15_2}.
  \item Při práci s matematikou a procvičováním postupů doporučujeme vyzkoušet funkce typu „Pokrok v matematice“, které generují příklady a mohou vyžadovat nahrání postupu řešení — užitečné pro hodnocení procesu, nikoli jen výsledku \cite{kb_ai_101_tipu_pdf_04e4a115_page_8_1}.
  \item Pokud chcete sledovat technické novinky a oficiální learning‑materiály k nástrojům (např. Microsoft Copilot, Azure, další), doporučeným zdrojem je portál Microsoft Learn jako systematický průvodce a dokumentace \cite{kb_ai_101_tipu_pdf_04e4a115_page_54_1}.
\end{itemize}

\medskip

\textbf{Krátké pokyny k etice a GDPR (shrnutí).} V této kapitole pouze navádíme, kde etiku a ochranu dat zapracovat do pracovního postupu: vždy zahrňte před pilotem základní informaci studentům o použití AI, zvažte nutnost DPIA pro zpracování osobních údajů a použijte šablony pro DPA/DPIA uvedené v Přílohách. Podrobná doporučení a checklisty jsou k dispozici v Kapitole 2 a v části Přílohy (Vzor DPIA, AI Policy Template).

\medskip

\textbf{Kde najdete šablony a další materiály.} Kopírovatelné prompty, šablony pro sylaby a checklisty najdete v části Přílohy (Prompt Library, AI Policy Template, Vzor DPIA). Online repozitář s video tutoriály a mechanismem kvartálních aktualizací je uveden v závěru příručky (viz informace o repozitáři a aktualizacích).%

%
\section{Tvorba sylabů a osnov pomocí AI}%
Praktické workflow a prompt šablony pro generování sylabů, formulaci učebních výstupů, rozvržení témat a propojení s ECTS/assessment kritérii; doporučení pro ověření a úpravy výstupů.%


%
Tato sekce nabízí praktický workflow a sadu prompt šablon pro rychlé vytvoření návrhu sylabu a osnovy předmětu s pomocí generativní AI. Cílem je získat kvalitní počáteční návrh, který vyučující ověří a upraví tak, aby odpovídal místním požadavkům (ECTS, hodnocení, etika/GDPR) a pedagogickým cílům kurzu.

\medskip
\textbf{Rychlý pracovní postup (workflow).}
\begin{enumerate}
  \item Definujte rámec: cílová skupina, semestr/počet týdnů, kontaktní hodiny, očekávaný workload/ECTS.
  \item Sepište 3--5 měřitelných učebních výstupů.
  \item Pošlete do AI jasný prompt (viz šablony níže) a získejte návrh obsahující: stručný popis kurzu, výstupy, týdenní témata, návrh hodnocení a orientační rozdělení ECTS/workload.
  \item Validujte návrh: faktická správnost, náročnost, soulad s ECTS a lokálními pravidly; upravte jazyk a úroveň.
  \item Doplňte závěry týkající se etiky/GDPR a transparentnosti používání AI; vytvořte doprovodné materiály pomocí specializovaných nástrojů podle potřeby \cite{kb_ai_101_tipu_pdf_04e4a115_page_44_2, kb_ai_101_tipu_pdf_04e4a115_page_51_1}.
\end{enumerate}

\medskip
\textbf{Důležité body validace (kontrolní seznam).}
\begin{itemize}
  \item Soulad výstupů s aktivitami a kritérii hodnocení.
  \item Realističnost workloadu (hodiny vs.\ ECTS).
  \item Jasné rozlišení procesu a produktu v hodnocení (milníky, průběžné odevzdání).
  \item Transparentnost: do sylabu uveďte, co vzniklo s využitím AI (viz Kapitola 2).
  \item Žádná citlivá data studentů nebyla použita bez příslušných DPA/DPIA.
\end{itemize}

\medskip
\textbf{Praktické prompt šablony (kopírovat a upravit).}

\textbf{1) Generování základního sylabu}
\begin{verbatim}
Jsem vyučující [Název], úroveň [bakalář/magistr].
Cílovka: [ročník, předpoklady]; semestr: [týdny]; kontaktní hodiny: [x]; workload: [y h ~= z ECTS].
Úkol: Navrhni sylabus (max 1 str.) s:
 - krátkým popisem (2-3 věty),
 - 4 měřitelnými výstupy,
 - týdenním plánem témat,
 - návrhem hodnocení (milníky + váhy),
 - orientační literaturou.
\end{verbatim}

\textbf{2) Mapování výstupů na ECTS a hodnocení}
\begin{verbatim}
Mám tyto výstupy: [seznam].
Pro každý navrhni:
 - formy ověření (proces/produkt),
 - doporučené váhy ve hodnocení,
 - orientační čas (kontakt vs. samostatně).
Uveď artefakty jako důkazy dosažení výstupu.
\end{verbatim}

\textbf{3) Struktura lekce (90/45/30 min)}
\begin{verbatim}
Pro týden [číslo] připrav osnovu 90min lekce:
 - cíle lekce,
 - časový rozpis (vstup, aktivita, diskuse, závěr),
 - 2 aktivizační úlohy (různá náročnost),
 - domácí úkol / další zdroje.
\end{verbatim}

\medskip
\textbf{Tipy k využití výstupů s nástroji pro prezentace a grafiku.}
\begin{itemize}
  \item Některé nástroje dokážou vytvořit slide deck přímo do školní šablony, což zrychlí přípravu materiálů \cite{kb_ai_101_tipu_pdf_04e4a115_page_44_2}.
  \item Pro vizuální prvky (ikony, obrázky, rozvržení) využijte Designer nebo podobné nástroje, které generují/úpravují grafiku a nastaví poměry stran \cite{kb_ai_101_tipu_pdf_04e4a115_page_51_1}.
  \item Jako názornou ukázku AI v praxi lze vložit lab aktivitu (např. demonstrace v Minecraft Education) pro praktické principy učení AI \cite{kb_ai_101_tipu_pdf_04e4a115_page_15_2}.
\end{itemize}

\medskip
\textbf{Doporučené ověřovací kroky po vygenerování návrhu.}
\begin{enumerate}
  \item Projděte návrh s garantem nebo kolegou (peer review).
  \item Ověřte zdroje a faktická tvrzení.
  \item Přizpůsobte jazyk a úroveň obsahu specifikům kurzu.
  \item Doplňte pravidla pro používání AI studenty (odkazy na AI Policy Template).
\end{enumerate}

\medskip
Hotové prompty a šablony najdete v Přílohách této příručky (Prompt Library) — využijte je jako výchozí bod a vždy proveďte lokální úpravy a validaci. Poznámka: návrhy vygenerované AI urychlí tvorbu; konečná odpovědnost za obsah, pedagogickou kvalitu a soulad s interními pravidly zůstává na garantu předmětu.%

%
\section{Návrh 90minutové přednášky a plán hodin}%
Šablony a konkrétní prompty pro rozfázování 90minutové přednášky (timeboxing, aktivity, přechody), návrhy interaktivních vstupů a varování k validaci faktů vygenerovaných textů.%


%
\medskip
Stručný pracovní postup pro návrh 90minutové přednášky
\begin{enumerate}
\item Definujte vzdělávací cíl (1–2 věty) a max. 2 měřitelné výstupy.
\item Zvolte mix: motivace, krátká expozice, aktivita studentů, formativní kontrola, závěr/artefakt.
\item Timeboxing: bloky a přechody sbírají důkazy učení (draft, krátký quiz, sdílené výstupy).
\item Použijte AI tam, kde přidá hodnotu; ověřte fakta a upravte pro kontext \cite{kb_ai_101_tipu_pdf_04e4a115_page_51_1}.
\end{enumerate}
\medskip
Časové rozvržení (vzor pro 90 min)
\begin{itemize}
\item 0–10 min: Úvod a motivace — příklad, learning goals.
\item 10–25 min: Expozice (max 15 min) + 2–3 kontrolní otázky.
\item 25–50 min: Aktivní fáze I — skupiny, zadání + 20 min, krátký artefakt.
\item 50–65 min: Formativní sdílení — 2–3 prezentace (3–5 min) + diskuse.
\item 65–80 min: Aktivní fáze II — individuální/párová úloha, případně s AI‑nápovědou.
\item 80–88 min: Krátký test porozumění.
\item 88–90 min: Shrnutí a domácí zdroje.
\end{itemize}
Varianty 45/30 min: zachovejte poměr expozice:aktivita:závěr ~=30:50:20; zkraťte bloky, jedno jasné výstupové zadání.
\medskip
Praktické nápady pro interaktivní vstupy
\begin{itemize}
\item Think–pair–share: 2 min přemýšlení, 5 min sdílení, 3 min shrnutí.
\item Live poll + diskuse pro rychlou diagnostiku.
\item Krátké simulace/experimenty (pokud dostupné) \cite{kb_ai_101_tipu_pdf_04e4a115_page_15_2}.
\item Adaptivní cvičení: generování příkladů a sledování postupu žáka \cite{kb_ai_101_tipu_pdf_04e4a115_page_8_1}.
\end{itemize}
\medskip
Zásady bezpečného zapojení AI
\begin{itemize}
\item Definujte roli AI (nápovědy, nikoli hotová řešení).
\item Transparentně informujte o použitých modelech/nástrojích.
\item Ověřte fakta před zveřejněním; chraňte citlivé údaje.
\end{itemize}
\medskip
Prompt šablony (kopírovat a upravit)
\begin{verbatim}
Pro týden [číslo] připrav osnovu 90min lekce:
- cíle (2-3, měřitelné); časový rozpis (0-10,10-25,25-50,50-65,65-80,80-90)
- 2 aktivity (různé náročnosti); návrh artefaktu; doporučené nástroje; kontrolní otázky
Jako alternativa: slide deck prompt — 10 slidů, vizuály, poznámky; uprav jazyk pro [bakalář/magistr]
\end{verbatim}
Pozn.: některé nástroje generují obrázky a přizpůsobí formát slidedecku \cite{kb_ai_101_tipu_pdf_04e4a115_page_51_1}.
\medskip
Krátký checklist po vygenerování materiálu
\begin{enumerate}
\item Ověřte fakta a citace.
\item Zkontrolujte workload a souladu s ECTS; přizpůsobte jazyk a úroveň studentům.
\item Odstraňte zavádějící části a uveďte v sylabu oznámení o využití AI.
\end{enumerate}
\medskip
Rychlá doporučení pro pilot
\begin{itemize}
\item Spusťte jednu 90min lekci jako pilot, testujte s malou skupinou, zaznamenávejte metriky (čas přípravy, spokojenost, kvalita artefaktů); sdílejte výsledky s kolegy a upravte workflow podle zpětné vazby.
\end{itemize}%

%
\section{Tvorba cvičení, pracovních listů a diferenciace}%
Postupy pro generování pracovních listů a cvičení na různé úrovně studentů, strategie diferenciace (lehčí/standardní/náročné verze) a kontrolní mechanismy pro kvalitu úloh a řešení.%


%
\medskip

Tento oddíl popisuje praktické postupy pro rychlé vytvoření cvičení a pracovních listů s podporou generativní AI, základní přístupy k diferenciaci úloh pro různé úrovně studentů a kontrolní mechanismy pro ověřování kvality úloh a řešení.

\subsection*{Rychlý pracovní postup pro generování pracovního listu}
\begin{enumerate}
  \item Definujte cílovou dovednost (1 věta) a požadovaný výstup (např. „vysvětlí postup řešení kvadratické rovnice“).
  \item Vyberte formát úloh (vektor: krátké otázky / postupové úlohy / otevřené problémové úlohy / automaticky opravitelné MCQ).
  \item Vygenerujte sadu variant (např. 8–12 příkladů) pomocí AI nástroje; u matematických úloh zvažte agenta, který k příkladům vygeneruje i očekávaný postup řešení (solution steps) — systém může předpřipravit indikaci, kde studenti často chybují \cite{kb_ai_101_tipu_pdf_04e4a115_page_8_1}.
  \item Přidejte metainformace: obtížnost, klíčová kompetence, doporučený čas na vypracování a kontrolní otázky.
  \item Validujte: zkontrolujte náhodnou podmnožinu příkladů ručně a u matematických úloh ověřte konzistenci kroků (viz kapitola níže).
\end{enumerate}

\subsection*{Specifika pro matematiku a sběr postupů}
U aktivit zaměřených na matematiku je pedagogicky i prakticky velmi výhodné požadovat nahrání postupu řešení, nikoli jen konečného výsledku. Některé vzdělávací AI nástroje (např. nástroje označené jako „Pokrok v matematice“) umožňují učiteli vybrat oblast, generují příklady a zároveň sbírají i postupy žáků; AI pak předvyplní hodnocení a doplní indikaci, v jaké části postupu pravděpodobně došlo k chybě — učitel finálně dokončí hodnocení \cite{kb_ai_101_tipu_pdf_04e4a115_page_8_1}. To zvyšuje zaměření hodnocení na proces místo pouhého výsledku a usnadňuje diagnostiku chyb.

\subsection*{Nástroje a integrace (praktické poznámky)}
\begin{itemize}
  \item OneNote a další nástroje pro převod rukopisu na digitální matematiku výrazně usnadňují sběr postupů a jejich automatické vyhodnocení; tyto funkce mohou rovněž exportovat rovnice a grafy pro další zpracování ve výukových platformách \cite{kb_ai_101_tipu_pdf_04e4a115_page_8_1}.
  \item Webové služby například shromažďují sady úloh a doplňují je o metriky výkonu, které učiteli ukážou, jak si jednotliví studenti vedou v porovnání se třídou \cite{kb_ai_101_tipu_pdf_04e4a115_page_8_1}.
  \item Pro textové materiály a pracovní listy lze využít nástroje s funkcí „shrnutí“ nebo přepisování tónu (např. integrované Copilot funkce v editoru) pro rychlou úpravu jazykové úrovně či formátu instrukcí \cite{kb_ai_101_tipu_pdf_04e4a115_page_14_1}.
\end{itemize}

\subsection*{Strategie diferenciace (lehké / standardní / náročné)}
(Obecné doporučení — označeno jako general background knowledge.)
\begin{itemize}
  \item Lehké varianty: více vodítek v zadání, částečně vyplněné kroky, konkrétní nápovědy (hinty), kratší počet kroků.
  \item Standardní varianty: plné zadání s jasnou strukturou očekávaného postupu; vyžaduje samostatné doplnění všech kroků.
  \item Náročné varianty: otevřené problémy, aplikace konceptu v novém kontextu, úlohy vyžadující argumentaci nebo více mo­dálit (graf/diagram + text).
  \item Praktický režim: při generování nechte AI vytvořit jedno „master“ zadání a poté automaticky generujte varianty s rozdílnou počáteční informací (více/ méně hintů) a s různými úrovněmi číselné složitosti.
\end{itemize}

\subsection*{Kontrolní mechanismy kvality a akademická integrita}
\begin{enumerate}
  \item Požadujte u zadání postupy řešení — umožní to detekovat parcely „kopírovaných“ odpovědí a podpoří hodnocení procesu \cite{kb_ai_101_tipu_pdf_04e4a115_page_8_1}.
  \item Používejte hybridní ověření: automatické testy pro konečné výsledky + ruční kontrola vybraných kroků u každého studenta.
  \item Náhodné variování parametrů úloh s jednoczesným zachováním didaktického cíle snižuje motivaci ke sdílení hotových řešení mezi studenty.
  \item V sylabu uveďte pravidla pro použití AI v domácích úkolech a požadavek na deklaraci, pokud student při vypracování použil nástroj AI (viz kapitola o etice a integritě).
\end{enumerate}

\subsection*{Praktické prompt‑šablony (copy \\ \& adapt) — general background knowledge}
Níže několik kratkých vzorů, které lze zkopírovat a upravit pro konkrétní nástroj:
\begin{verbatim}
1) Vygeneruj 10 příkladů na téma "kvadratické rovnice" s řešením:
 - 4 snadné (kroky vypsané), 4 standardní (kroky stručné), 2 náročné (otevřený problém).
 - U každého příkladu uveď očekávaný čas: 5/10/20 min.
2) Přepiš zadání pro 1. ročník Bc. jazyka: zjednoduš věty, přidej 2 kontrolní otázky.
3) Vytvoř sadu 5 kontrolních otázek (MCQ) s klíčem a krátkým vysvětlením odpovědi.
\end{verbatim}
(Použití šablon je praktická pomůcka; výsledky vždy validujte manuálně.)

\subsection*{Krátký checklist pro nasazení pracovního listu s AI}
\begin{itemize}
  \item Jasně definované vzdělávací cíle a forma odevzdání (výsledek vs. proces).
  \item Vyžádání postupu řešení tam, kde má proces hodnotu — zejména u matematiky \cite{kb_ai_101_tipu_pdf_04e4a115_page_8_1}.
  \item Automatické generování variant a ruční kontrola náhodné části výstupu.
  \item Zahrnutí informací o tom, zda je AI při tvorbě úlohy použita (transparentnost).
  \item Integrace s nástroji pro sběr rukopisu / převod rovnic (OneNote apod.) podle potřeby \cite{kb_ai_101_tipu_pdf_04e4a115_page_8_1}.
\end{itemize}

\medskip
Poznámka: uvedené postupy kombinují příklady z reálných edukačních nástrojů (viz „Pokrok v matematice“ a integrace převodu rukopisu) a praktické pedagogické doporučení; vždy ověřujte generované úlohy před distribucí studentům \cite{kb_ai_101_tipu_pdf_04e4a115_page_8_1, kb_ai_101_tipu_pdf_04e4a115_page_14_1, kb_ai_101_tipu_pdf_04e4a115_page_15_2}.%

%
\section{Generování kvízů a automatické hodnocení}%
Návody pro tvorbu kvízů (MCQ, otevřené úlohy, programovací úkoly), tvorbu item banky a základy automatického hodnocení včetně návrhů na generování řešení a testovacích případů.%


%
\subsection*{Úvod}
Generování kvízů a automatické hodnocení jsou efektivní způsoby, jak škalovat kontrolu porozumění a poskytnout rychlou zpětnou vazbu studentům. Níže je praktický soubor postupů a šablon, které lze přímo použít při přípravě kurzů na VUT; specifická funkčnost pro matematické úlohy (sběr postupů, předoprava hodnocení a diagnostika) je ilustrována na příkladu nástroje „Pokrok v matematice“ \cite{kb_ai_101_tipu_pdf_04e4a115_page_8_1}.

\subsection*{Rychlý workflow pro tvorbu kvízu}
\begin{enumerate}
  \item Definujte cíl kvízu (1–2 věty) a formát (MCQ, krátké otevřené odpovědi, programovací úloha).
  \item Rozhodněte o úrovni automatizace (plně automatické vyhodnocení vs. AI‑asistent + lidská finalizace). U matematických úloh lze využít režim, kdy AI předopracuje hodnocení a učitel finálně dokončí výsledky — tento přístup používají některé vzdělávací nástroje jako „Pokrok v matematice“ \cite{kb_ai_101_tipu_pdf_04e4a115_page_8_1}.
  \item Vygenerujte sadu položek (itemů) pomocí LLM a zároveň automaticky vytvořte klíče/řešení a návrhy testovacích případů (pro programování). Ponechte v šabloně metainformace: obtížnost, odhadovaný čas, typ položky.
  \item Nasaďte pilotní běh na malé skupině (např. 10–20 studentů) a zkombinujte automatické vyhodnocení s ruční kontrolou náhodného výběru položek.
  \item Reflektujte výsledky pilotu a upravte položky (jasnost zadání, vyvážení obtížnosti, chybějící okruhy).
\end{enumerate}

\subsection*{Generování MCQ a item banky (prakticky)}
Postup pro rychlé vytvoření item banky:
\begin{enumerate}
  \item Vstup: seznam témat / dovedností + požadovaný počet položek na téma.
  \item Pro každé téma vygenerujte 4–6 MCQ: správná odpověď + 3 (plausibilní) distraktory + krátké vysvětlení správné odpovědi.
  \item Přiřaďte každé položce metadatové štítky: Bloomova úroveň (rozumění/aplikace/analýza), očekávaný čas, doporučená prerekvizita.
  \item Generujte varianty parametrických položek (proměnné v zadání — čísla, kontext), aby bylo možné snadno produkovat verze pro jednotlivé studenty a snižovat riziko sdílení hotových odpovědí.
\end{enumerate}
Poznámka: samotné generování distraktorů a vysvětlení je obecná doporučená praxe (general background knowledge).

\subsection*{Automatické hodnocení otevřených odpovědí a matematických postupů}
U otevřených textových odpovědí lze použít AI pro předběžné ohodnocení podle definovaných rubrik (např. jasnost argumentace, správnost klíčových bodů). U úloh, kde je důležitý postup (matematika, chemické výpočty), je silně doporučeno požadovat nahrání postupu řešení — nástroje v rodině „Pokrok v matematice“ toto podporují: žák nahrává postup, AI provede předopravu a označí části, kde pravděpodobně došlo k chybě; učitel pak finálně dokončí hodnocení \cite{kb_ai_101_tipu_pdf_04e4a115_page_8_1}. Pro sběr rukopisu a převod na digitální podobu lze využít např. OneNote, který detekuje rovnice a umí je převést a dále zpracovat \cite{kb_ai_101_tipu_pdf_04e4a115_page_8_1}.

\subsection*{Automatické hodnocení programovacích úloh (shrnutí praktik)}
(Toto je obecné doporučení — general background knowledge.)
\begin{itemize}
  \item Požadujte repozitář se studentským kódem (Git, LMS upload) a definujte rozhraní pro spouštění testů (Docker/CI runner nebo školicí sandbox).
  \item Generujte automatické testy (unit/integration) pro ověření funkcionality; součástí by měly být také testovací případy pro hraniční vstupy a chybná data.
  \item Pro detekci plagiátorství využijte kombinaci metrik: porovnání struktury kódu, podobnost v logice a ruční kontroly podezřelých případů.
  \item Doporučený režim: automatické skóre + náhodná ruční revize, případně oral‑exam nebo doplňková reflexe studenta k předloženému kódu.
\end{itemize}

\subsection*{Praktické prompt‑šablony (copy \\ \& adapt) — general background knowledge}
\begin{verbatim}
1) "Vygeneruj 10 MCQ pro téma 'Numerická integrace':
 - uveď otázku, 4 možnosti (označ A-D), správnou odpověď, krátké vysvětlení 1-2 věty,
 - přiřaď Bloom level a odhadovaný čas."
2) "Pro tuto programovací úlohu napiš 8 unit testů (Python pytest): 
 - popiš vstupy a očekávané výstupy, přidej 2 hraniční případy."
3) "Pro zadanou otevřenou otázku navrhni rubriku se 4 kritérii (0-4 body každé) a krátký popis, co se hodnotí."
\end{verbatim}
(Přizpůsobte prompt stylu používaného LLMu a vždy validujte výstupy manuálně.)

\subsection*{Validace kvality a integrita hodnocení}
Checklist pro nasazení automatického hodnocení:
\begin{itemize}
  \item Definované rubriky pro otevřené položky a pravidla pro automatické předohodnocení.
  \item Požadavek na postup (tam, kde má proces hodnotu) — zvyšuje robustnost vůči generickým odpovědím \cite{kb_ai_101_tipu_pdf_04e4a115_page_8_1}.
  \item Náhodné ruční kontroly (min. 5–10\% automaticky ohodnocených záznamů) a jasně definovaný postup eskalace v případě nesouladu.
  \item Logování verzí položek (item bank changes) a auditní záznamy výsledků pro případnou revizi.
\end{itemize}

\subsection*{Etika, GDPR a transparentnost}
Krátce: informujte studenty o způsobu použití AI při vytváření a hodnocení kvízů (transparentnost), zajistěte minimalizaci osobních údajů tam, kde to je možné, a pokud systém shromažďuje data (logy, nahrávky postupů), dodržujte fakultní postupy pro DPIA a DPA (general background knowledge). V případě použití cloudových služeb zvažte enterprise varianty a konzultaci s IT/právním oddělením (viz kapitola o etice a institucionálním nasazení).

\medskip
Poznámka: výstupy generované AI považujte za návrhy — vždy proveďte lidskou validaci přinejmenším u pilotních běhů a u položek s vysokým dopadem na konečné hodnocení; u matematických aktivit lze využít funkcí popsaných v „Pokrok v matematice“ a nástrojích pro převod rukopisu (OneNote) jako podpůrnou technologii \cite{kb_ai_101_tipu_pdf_04e4a115_page_8_1}.%

%
\section{Tvorba výukových prezentací a slide decků}%
Praktické postupy a prompty pro převod sylabu do slidů, návrhy vizuální struktury, integrace multimédií a nástroje pro rychlou tvorbu prezentací (rokování nástrojů uvedených v kapitole).%


%
% Tělo sekce: Tvorba výukových prezentací a slide decků
% (Pozor: název sekce se neuvádí — pouze obsah)

Cílem této části je nabídnout učitelský workflow pro rychlé a bezpečné převedení sylabu a osnovy lekce do srozumitelných slide decků, doporučené vizuální šablony a praktické prompty, které lze okamžitě použít nebo adaptovat. Praktické tipy vycházejí z ověřených možností nástrojů pro generování grafiky a z práce s digitálním rukopisem, které jsou užitečné při tvorbě výukových materiálů.

Úvodní doporučený workflow (shrnutí)
\begin{enumerate}
  \item Sestavte „kostru“ prezentace: 3--5 hlavních výstupů lekce a seznam klíčových bodů pro každý výstup.  (general background knowledge)
  \item Rozdělte obsah na typy slideů: úvodní / cíl, vysvětlení konceptu, příklad / demonstrace, aktivita / úkol pro studenty, shrnutí / takeaway. (general background knowledge)
  \item Pro každý slide napište 1--2 věty „hlavní myšlenky“, které chce učitel předat; ty jsou vstupem pro generování odstavců, bodů a poznámek lektora pomocí LLM. (general background knowledge)
  \item Vytvořte nebo vygenerujte relevantní vizuál k nejméně jednomu klíčovému slideu (diagram, ilustrace, screenshot nebo monogram/ikonka). Pro generování ilustrací lze využít nástroje jako Microsoft Designer, který umí vytvářet nové grafiky z promptu, upravovat stávající grafiku a také generovat obrázky v různých poměrech stran (např. širokoúhlé) \cite{kb_ai_101_tipu_pdf_04e4a115_page_51_1}.
  \item Pokud sbíráte ručně psané vzorce nebo poznámky pro prezentaci, použijte nástroje pro převod rukopisu (např. OneNote), které dokážou převést rukopisné rovnice do digitální podoby a usnadnit jejich vložení do slideů \cite{kb_ai_101_tipu_pdf_04e4a115_page_8_1}.
  \item Validujte každý automaticky vygenerovaný text a vizuál lidskou revizí; doplňte alt text pro obrázky a krátké poznámky k použití při výuce. (general background knowledge)
\end{enumerate}

Praktická šablona vizuální struktury slideu (opakovatelný pattern)  (general background knowledge)
\begin{itemize}
  \item Hlavička (1 věta): název bloku / slide topic.
  \item Klíčová myšlenka (1--2 věty): co si student má odnést.
  \item 3 krátké body (bullety) nebo jedno krátké vysvětlení + příklad (max. 6 řádků textu).
  \item Vizuální prvek (obr. / schéma / screenshot / ikona) — dominantní prvek slideu.
  \item Takeaway / otázka pro studenty (1 věta) nebo krátká aktivita (má‑to‑udělat‑student).
\end{itemize}

Použití Microsoft Designer pro vizuály (konkrétní možnosti)
\begin{itemize}
  \item Designer podporuje tři režimy práce: generování nové grafiky z promptu, úpravu existující grafiky a klasický ruční návrh; je užitečný tam, kde chcete vytvořit rychlé obrázky, ikony nebo stylizované pozvánky a monogramy pro výuku \cite{kb_ai_101_tipu_pdf_04e4a115_page_51_1}.
  \item Designer umí generovat obrázky na základě textového promptu a umožňuje nastavit různé poměry stran (není omezen na 1:1), což usnadňuje tvorbu širokoúhlých slide ilustrací nebo bannerů \cite{kb_ai_101_tipu_pdf_04e4a115_page_51_1}.
  \item Praktické nápady: vytvořte si sadu malých ikon (avatarů, štítků) pro rychlé odlišení typů úloh ve slidech (např. „cvičení“, „domácí úkol“, „checkpoint“). Designer zvládne i varianty pro samolepky/štítky s bílým okrajem, které lze exportovat do prezentací \cite{kb_ai_101_tipu_pdf_04e4a115_page_51_1}.
\end{itemize}

Práce s rukopisem a rovnicemi
\begin{itemize}
  \item Pokud v kurzu používáte ruční zápis (tabule, tablet), převeďte záznamy do digitální podoby pomocí OneNote nebo podobných nástrojů; OneNote dokáže detekovat rovnice a převést je do editovatelné digitální podoby, což usnadní vložení přesných matematických výrazů do slideů \cite{kb_ai_101_tipu_pdf_04e4a115_page_8_1}.
  \item Doporučení (general background knowledge): po převodu rovnice zkontrolujte formátování (font, zarovnání), přidejte krátkou poznámku ke každému kroku a případně animaci, která postupně odhalí kroky řešení během prezentace.
\end{itemize}

Praktické prompt‑šablony pro převod sylabu do slidů (copy \\ \& adapt)  (general background knowledge)
\begin{verbatim}
1) "Vezmi tento odstavec sylabu (vložit text) a rozděl ho na 5 slideů:
 - pro každý slide napiš: název slide (5-7 slov), 2-4 bullet body, jednu krátkou aktivitu pro studenty (1 věta), návrh vizuálu (1-2 slova)."
2) "Pro tento konkrétní učení výsledek: 'Student rozumí základům numerické integrace'
 - vygeneruj 3 slidey: (a) krátké vysvětlení konceptu, (b) jednoduchý příklad krok-po-kroku, (c) návrh krátké in-class aktivty + kontrolní otázka."
3) "Vytvoř návrh 1-snímkového přehledu kurzu: struktura 10 týdnů, u každého týdne 1 věta o cíli a doporučený artefakt (prezentace/úkol/projekt)."
\end{verbatim}

Tipy pro přístupnost a nasazení v LMS (general background knowledge)
\begin{itemize}
  \item Každý obrázek opatřete krátkým alt‑textem (1–2 věty) a v poznámkách lektora uveďte, jak obrázek komentovat při výuce.
  \item Exportujte slide decky do PDF pro studentské materiály a zachovejte samostatný soubor s podrobnými poznámkami lektora (pro asistentky/asistenty).
  \item U větších kurzů připravte „master“ slide šablonu s jednotným stylem a sadou ikon; šablonu uložte do repozitáře katedry pro opakované použití.
\end{itemize}

Rychlý checklist před publikací (general background knowledge)
\begin{itemize}
  \item Text na slidech stručný: max. 6 řádků / 6 slov na řádek.
  \item Vizuál relevantní a čitelný při redukovaném rozměru (thumbnail).
  \item Kontrola správnosti konverze rovnic (pokud byla použita automatika).
  \item Přidat zdroje / odkazy a upozornění na použití AI při generování obsahu (transparentnost).
\end{itemize}

Poznámka k integraci nástrojů a dalšímu zpracování
\begin{itemize}
  \item Pokud používáte generování obrázků z promptů, uložte původní prompty a metadata (styl, rozměr), aby bylo možné vizuály reprodukovat nebo upravit později; u Microsoft Designeru je možné experimentovat s variantami a exportovat různé poměry stran pro použití v slidech \cite{kb_ai_101_tipu_pdf_04e4a115_page_51_1}.
  \item V případech, kdy do slide decku zahrnujete materiály studentů (např. jejich projekty nebo screenshoty), zvažte souhlas ke zveřejnění a pravidla ochrany osobních údajů podle místních směrnic (general background knowledge).
\end{itemize}

Závěrem: kombinujte strukturovaný text, jeden silný vizuál na slide a krátkou aktivitu — to výrazně zvýší srozumitelnost a interaktivitu prezentace. Vizuály lze efektivně generovat a upravovat např. v Microsoft Designeru; ručně psané rovnice převeďte přes OneNote pro přesnost a snadné vložení do slideů \cite{kb_ai_101_tipu_pdf_04e4a115_page_51_1, kb_ai_101_tipu_pdf_04e4a115_page_8_1}.%

%
\section{Generování a úprava obrázků a grafiky}%
Přehled možností generování vizuálů (včetně editace existující grafiky) a ukázky použití nástrojů pro tvorbu grafiky a obrázků ve výuce; zahrnuje praktické tipy pro rozměry a stylizaci (využití funkcí Designer a podobných nástrojů).%


%
% Tělo sekce: Generování a úprava obrázků a grafiky
Tato sekce shrnuje praktické možnosti tvorby i editace vizuálů přímo v rámci nástrojů Microsoft (zejména Designer) a ukazuje doporučené pracovní postupy a ukázkové prompty pro použití ve výuce. Uváděné technické možnosti vycházejí z dostupných funkcí Designeru a z funkcí pro editaci obrázků v Copilot chatu \cite{kb_ai_101_tipu_pdf_04e4a115_page_51_1, kb_ai_101_tipu_pdf_04e4a115_page_52_2}.

Krátké shrnutí schopností Designeru (co využijete ve výuce)
\begin{itemize}
  \item Designer pracuje ve třech režimech: generování nové grafiky z promptu, úprava stávající grafiky pomocí AI a klasický ruční návrh \cite{kb_ai_101_tipu_pdf_04e4a115_page_51_1}.
  \item Umí generovat obrázky na základě textového promptu a umožňuje definovat jiný než čtvercový poměr stran (rychle vytvoříte širokoúhlé bannery pro slide) \cite{kb_ai_101_tipu_pdf_04e4a115_page_51_1}.
  \item Podporuje také specifické kategorie užitečné pro školní použití: monogramy, omalovánky, avatary a nálepky (sticker designy včetně bílého okraje) \cite{kb_ai_101_tipu_pdf_04e4a115_page_51_1}.
  \item Editace části existujícího obrázku (např. změna pozadí, úprava nápisu nebo přidání/odebrání doplňků) je možná prostřednictvím funkcí transformace v Copilot chatu — tato funkce vyžaduje placený Microsoft 365 Copilot \cite{kb_ai_101_tipu_pdf_04e4a115_page_52_2}.
\end{itemize}

Doporučený rychlý workflow pro vytvoření vizuálu k lekci
\begin{enumerate}
  \item Definujte účel obrázku: (a) slide ilustrace, (b) cover pro materiál, (c) ikonka/štítek pro aktivitu.
  \item Zvolte formát a poměr stran podle cílového použití (např. širokoúhlý pro slide 16:9, čtverec pro ikonu, vertikál pro story/banner). (praktické doporučení; viz možnosti Designeru \cite{kb_ai_101_tipu_pdf_04e4a115_page_51_1})
  \item Napište krátký prompt s požadovaným stylem, barvami a elementy (viz šablony níže).
  \item Vygenerujte 3 varianty, vyberte nejvhodnější a případně použijte editaci části obrázku (transformace) pro doladění \cite{kb_ai_101_tipu_pdf_04e4a115_page_52_2}.
  \item Exportujte ve vhodném formátu (PNG/SVG pro ikony, JPG/PDF pro materiály) a uložte originální prompt/metadata pro pozdější reprodukci. (general background knowledge)
\end{enumerate}

Praktické prompt‑šablony (copy \\ \& adapt)
\begin{verbatim}
1) Monogram pro kurz:
"Vytvoř monogram pro kurz 'Fyzika 1' v jednoduchém moderním stylu,
barvy: tmavě modrá + světle zelená, čistý vektorový styl, výstup 1:1,
jednoduchý obrys vhodný pro tisk na samolepku."
2) Motivace na začátek hodiny (wide banner 16:9):
"Motivační banner: učitel jako postavička z modelíny ve stylu 3D
animovaného filmu, pozadí: zelená tabule s nápisem 'Do toho!',
světlé barvy, čitelný nápis uprostřed, vhodné pro slide 16:9."
3) Nálepka/štítek pro úkol:
"Samolepka 'Hotovo' s bílým okrajem, kruhový tvar, barvy: žlutá + černá,
plochý ikonický styl, vysoký kontrast pro tisk na samolepky."
\end{verbatim}

Ukázka postupu pro úpravu části obrázku (transformace)
\begin{enumerate}
  \item Vygenerujte základní obrázek v Designeru nebo nahrajte vlastní fotografii.
  \item V Copilot chatu zvolte režim Vytvořit + Transformace, popište přesnou změnu (např. „vyměň pozadí za kreslenou zelenou tabuli a přidej nápis 'Do toho!'“). Pozn.: tato funkce je dostupná u Microsoft 365 Copilot (placená) \cite{kb_ai_101_tipu_pdf_04e4a115_page_52_2}.
  \item Klepněte na Aktualizovat a opakujte doladění, dokud výsledek neodpovídá potřebám hodiny \cite{kb_ai_101_tipu_pdf_04e4a115_page_52_2}.
\end{enumerate}

Praktické nápady použití ve výuce (rychlé příklady)
\begin{itemize}
  \item Vytvoření sady ikon/štítků pro typy aktivit (cvičení, domácí úkol, checkpoint) pomocí Designeru a export do balíku ikon pro slide šablony \cite{kb_ai_101_tipu_pdf_04e4a115_page_51_1}.
  \item Generování ilustračních bannerů pro začátek lekce (motivace, cíle) a jejich iterativní úprava podle zpětné vazby studentů prostřednictvím transformací \cite{kb_ai_101_tipu_pdf_04e4a115_page_52_2}.
  \item Příprava materiálů pro výtvarnou výuku: monogramy nebo varianty v různých malířských technikách (praktické příklady v Designeru) \cite{kb_ai_101_tipu_pdf_04e4a115_page_51_1}.
\end{itemize}

Checklist při exportu a nasazení vizuálů (stručně; general background knowledge)
\begin{itemize}
  \item Zkontrolujte rozlišení a poměr stran pro cílové médium (projekce, tisk, web).
  \item Přidejte stručný alt‑text k obrázkům a poznámku pro lektora, jak obrázek v hodině komentovat.
  \item Uložte originální prompty a metadata (styl, poměr stran), aby bylo možné vizuály reprodukovat nebo upravovat později. 
  \item Pokud používáte materiály studentů v publikaci, ověřte souhlas a zásady ochrany osobních údajů. 
\end{itemize}

Poznámka k licencím a dostupnosti funkcí
\begin{itemize}
  \item Některé pokročilé editační funkce (např. transformace v Copilot chatu) vyžadují placenou licenci Microsoft 365 Copilot; také některé šablony/obsahy v Designeru mohou být omezeny pro typy účtů \cite{kb_ai_101_tipu_pdf_04e4a115_page_52_2, kb_ai_101_tipu_pdf_04e4a115_page_51_1}.
\end{itemize}%

%
\section{Sumarizace vědeckých textů a ověřování výstupů}%
Postupy a nástroje pro sumarizaci dlouhých odborných textů, extrakci klíčových závěrů a doporučení jak ověřit přesnost a relevanci sumarizací (návrhy workflow pro review a cross-check).%


%
Tato sekce popisuje praktické postupy pro rychlé a bezpečné shrnutí dlouhých odborných textů pomocí nástrojů AI a navrhuje konkrétní kroky pro ověření přesnosti a relevance výsledných sumarizací. Upozornění: níže uvedené doporučené pracovní postupy kombinují dostupné funkce nástrojů (viz příklad Copilota níže) a osvědčené postupy pro lidskou validaci; konkrétní technické možnosti závisí na vámi používaném nástroji a licenci. (general background knowledge)

Automatizované shrnutí — rychlý postup (praktický příklad)
\begin{enumerate}
  \item Sběr vstupních dokumentů: shromážděte původní PDF, kapitoly nebo články, které mají být shrnuty; zaznamenávejte metadata (autor, rok, URL/DOI). (general background knowledge)
  \item Předzpracování: rozdělte dlouhé texty na logické části (abstrakt, úvod, metodika, výsledky, diskuse) a očistěte OCR chyby u skenovaných dokumentů. (general background knowledge)
  \item Automatická agregace a shrnutí: použijte Copilot / lokální LLM nebo jiný vhodný nástroj pro vytvoření konzistentního přehledu z více zdrojů; v prostředí Windows lze např. v Poznámkovém bloku využít funkci „Shrnout“ pro kombinaci textů z různých webů a vytvoření stručného souhrnu \cite{kb_ai_101_tipu_pdf_04e4a115_page_14_1}.
  \item Uložení výstupů a provenance: ke každému shrnutí uložte seznam zdrojů, původní odstavce, použitý prompt a verzi modelu. (general background knowledge)
\end{enumerate}

Tip pro práci s odbornými texty obsahujícími matematiku a rovnice
\begin{itemize}
  \item Pokud text obsahuje rovnice nebo rukopisné poznámky, využijte nástroje, které umějí převádět rukopis do digitální podoby (např. OneNote) a zpracovat rovnice jako součást obsahu před shrnutím \cite{kb_ai_101_tipu_pdf_04e4a115_page_8_1}. To pomáhá zajistit, že číselné výsledky a symbolika nejsou při summarizaci ztraceny nebo zkresleny.
\end{itemize}

Checklist pro rychlé ověření kvality a přesnosti sumarizace (praktické kroky)
\begin{itemize}
  \item Zkontrolujte, zda shrnutí výslovně uvádí hlavní závěry a klíčové metriky (p-hodnoty, intervaly spolehlivosti, numerické výsledky). (general background knowledge)
  \item Porovnejte minimum 2–3 vět z shrnutí s originálními pasážemi, které shrnutí odkazuje; ověřte, že parafráze není překreslením faktů (hallucination). (general background knowledge)
  \item Ověřte citace: pro každý fakticky náročný výrok vygenerujte kontrolní dotaz pro model („Který odstavec v původním článku toto podporuje?“) nebo najděte primární zdroj. (general background knowledge)
  \item Validace metodiky: pokud shrnutí konstatuje metody či závěry, nechte odborníka z oblasti zhodnotit, zda byly metody popsány adekvátně a zda závěry logicky plynou z výsledků. (general background knowledge)
\end{itemize}

Praktické prompty a šablony pro shrnutí (copy-\\ \&-adapt)
\begin{verbatim}
1) Kompaktní shrnutí (1 odstavec):
"Vezmi následující text (vložit) a vytvoř 5-řádkové shrnutí,
kde první věta popisuje hlavní cíl studie, druhá metoda, třetí
klíčový výsledek (číselně pokud je k dispozici), čtvrtá hlavní závěr
a pátá implikace pro praxi/učitele."
2) Extrakce klíčových metrik:
"Najdi a vypiš všechny číselné výsledky a metriky z textu
(výsledky experimentů, p-hodnoty, velikosti efektu, sample size).
Uveď i přesný odstavec/stránku, kde se nacházejí."
3) Porovnání více studií:
"Máme tři články (vložit krátké metadaty a relevantní pasáže).
Vytvoř tabulku: cíl studie | metoda | hlavní výsledek | omezení."
\end{verbatim}

Workflow pro review a cross‑check ve výuce (doporučené rozhraní učitel–AI)
\begin{enumerate}
  \item Student/učitel vygeneruje první automatické shrnutí pomocí zvoleného nástroje (uloží se prompt a seznam zdrojů). (general background knowledge)
  \item Peer‑review: jiný student nebo asistent porovná klíčové výroky se zdroji a označí nesrovnalosti. (general background knowledge)
  \item Odborná validace: vyučující nebo odborný garant zkontroluje vybrané pasáže a potvrdí závěry; v případě sporných bodů se vrátí k původním zdrojům. (general background knowledge)
  \item Zapracování zpětné vazby do finálního shrnutí a explicitní poznámka v materiálech o tom, že shrnutí prošlo validací (transparentnost vůči studentům). (general background knowledge)
\end{enumerate}

Uchování provenance a transparentnost
\begin{itemize}
  \item Ukládejte vždy: vstupní texty (nebo jejich identifikátory), použitý prompt, verzi modelu/nástroje a datum vytvoření. (general background knowledge)
  \item V sylabu nebo k zadání přidejte krátkou poznámku o tom, že některé souhrny mohou být generovány automaticky a jak byl zajištěn review proces. (general background knowledge)
\end{itemize}

Poznámka o nástrojích: Copilot v Poznámkovém bloku lze použít k automatickému vytváření souhrnů z různých zdrojů, což usnadňuje sestavení přehledů z webových textů a poznámek \cite{kb_ai_101_tipu_pdf_04e4a115_page_14_1}. Pro materiály s rovnicemi doporučujeme před shrnutím využít nástroje podporující převod rukopisu/rovnic na digitální formu, např. OneNote, aby se předešlo ztrátě numerické přesnosti \cite{kb_ai_101_tipu_pdf_04e4a115_page_8_1}.%

%
\section{Knihovna promptů a vzorů pro redesign předmětů}%
Sada opakovaně použitelných promptů, checklistů a šablon pro redesign kurikula, včetně ukázkových vstupů a výsledků, které lze adaptovat pro různé typy předmětů.%


%
Tato sekce obsahuje opakovatelné prompty, šablony a krátké checklisty pro rychlý redesign sylabu, návrh hodin a hodnocení v kontextu generativní AI. Prompty jsou copy-\\ \&-adapt: vložte specifika kurzu (cíle, cílová skupina, očekávané výstupy, rozsah, předpoklady) a ověřte výsledky odbornou revizí.

\medskip
\noindent 1) Prompt: Generátor sylabu (rychlá verze)
\begin{verbatim}
"Vezmi následující informace:
- Název, Úroveň (bakalář/magistr), ECTS
- Cíle (3-5), Klíčová témata, Požadavky
Vygeneruj 1-stránkový sylabus se sekcemi:
- Stručný popis (2 věty)
- Učební výstupy (5 bodů)
- Týdenní osnova (10-12 týdnů, 1 věta/týden)
- Formy hodnocení a váhy
- Doporučená literatura/online zdroje
U každé části přidej 1větný důvod volby."
\end{verbatim}
Tip: rozděl výstup na bloky a ověř terminologii a fakta.

\medskip
\noindent 2) Prompt: Převod sylabu na 90min slide deck
\begin{verbatim}
"Vezmi odstavec sylabu (vložit) a rozděl ho na 5 slideů:
Pro každý slide napiš:
- Název (5-7 slov)
- 2-4 bullet body
- Jednu in-class aktivitu (1 věta)
- Návrh vizuálu (1-2 slova)
Výstup: seznam 5 slideů připravených pro export do PPT."
\end{verbatim}
Doporučený pattern: hlavička, klíčová myšlenka, 3 body, vizuál, takeaway/aktivita.

\medskip
\noindent 3) Prompt: Redesign zadání odolného vůči triviálnímu použití AI
\begin{verbatim}
"Navrhni projekt/esej, které:
- Měří autenticitu (lokální/datumované artefakty)
- Mají >=3 milníky (draft, peer-review, final)
- Vyžadují reflexi použití AI (co/ proč/ jak)
- Obsahují hodnotící rubriku (proces, argumentace, originalita, reflexe AI)
Vstup: téma, délka, požadované artefakty."
\end{verbatim}
Příklad rubriky: Proces (0–4), Metodická přesnost (0–4), Originalita (0–4), Dokumentace AI (0–4).

\medskip
\noindent 4) Prompt: Generátor milestone‑based hodnocení a checklistů
\begin{verbatim}
"Pro kurz (název + úroveň) navrhni 3–5 milníků s termíny a krátkými úkoly,
které lze ověřit. U každého milníku uveď: cíl, očekávaný výstup,
hodnoticí kritéria a doporučený způsob kontroly originality.
Vstup: délka semestru (týdny), preferované formy odevzdání."
\end{verbatim}

\medskip
\noindent 5) Prompt: Šablona pro generování formativní zpětné vazby
\begin{verbatim}
"Na základě výstupu studenta (text/kód/slide) vygeneruj:
- 3 silné stránky (1 věta každá)
- 3 konkrétní doporučení ke zlepšení (konkrétní kroky)
- 1 návrh na cvičení nebo zdroj k procvičení"
\end{verbatim}
Použití: hromadné označení návrhů před kontrolou učitelem; vždy doplň lidskou kontrolu.

\medskip
\noindent 6) Prompt: Text do sylabu — zásady použití AI a deklarace
\begin{verbatim}
"Vytvoř 2–3 odstavce do sylabu, které:
- Definují, co je povoleno/zakázáno u generativní AI
- Popisují požadavek na deklaraci použití AI u odevzdání
- Uvádějí důsledky porušení (odkaz na akademickou integritu)
Přizpůsob pro technický kurz, neutrální tón."
\end{verbatim}

\medskip
\noindent 7) Prompt: Návrh inkluzivních a multimodálních úloh
\begin{verbatim}
"Navrhni 3 varianty jednoho úkolu: textová, multimediální (audio/video),
praktická (simulace). Uveď adaptace hodnocení a doporučené nástroje,
s ohledem na podporu neurodiverzních studentů."
\end{verbatim}
Pozn.: při převodu rukopisu/rovnic použijte nástroje pro převod (např. OneNote) před generováním obsahu \cite{kb_ai_101_tipu_pdf_04e4a115_page_8_1}.

\medskip
\noindent 8) Inspirace pro výukové světy a interaktivní ukázky
\begin{itemize}
  \item Zvažte vzdělávací světy (např. Minecraft Education) pro vysvětlení principů AI a hands‑on aktivity s nižší bariérou vstupu \cite{kb_ai_101_tipu_pdf_04e4a115_page_15_2}.
\end{itemize}

\bigskip
\noindent Krátký checklist před použitím promptů
\begin{enumerate}
  \item Uprav prompt pro lokální kontext (jazyk, úroveň studentů).
  \item Otestuj prompt na jednom vzorku a ověř faktickou správnost.
  \item Ulož prompt, vstupy, verzi modelu a datum (provenance).
  \item Zaznamenej do sylabu, jak bude AI používána a hlášena.
\end{enumerate}

\medskip
\noindent Vzorový postup pro adaptaci v katedře
\begin{enumerate}
  \item Vedoucí připraví vstupy (cíle, témata, výstupy).
  \item Použije generátor sylabu a adaptuje návrh s odbornou garancí.
  \item Pilotně otestuje hodnocení a rubriky na malém vzorku studentů.
  \item Upravený sylabus verzuje, dokumentuje a publikuje (living document).
\end{enumerate}

\bigskip
\noindent Poznámka: prompty jsou praktické šablony pro rychlý start. Při použití vždy proveďte odbornou kontrolu výstupů, zajistěte transparentnost vůči studentům a evidenci použitých modelů a promptů.%

%
\section{Ověřené případové studie a inspirace}%
Konkrétní příklady dobré praxe z univerzit a implementací (shrnutí klíčových přístupů a přenosných principů) a doporučení pro pilotní nasazení na úrovni předmětu či katedry.%


%
Tato sekce shrnuje ověřené příklady dobré praxe a přenosné principy, které lze využít při pilotním nasazení generativní AI na úrovni předmětu nebo katedry.
Následují konkrétní příklady z praxe, krátké analýzy klíčových přístupů a praktické doporučení pro pilot (body označené „(general background knowledge)“ nejsou přímo citovány z KB výstupů a slouží jako stručné praktické rady).

\medskip
\textbf{1) Využití vzdělávacích světů: Minecraft Education jako nástroj pro přiblížení AI.}  
Minecraft Education nabízí připravené světy a aktivity, které lze přímo použít k demonstraci principů strojového učení a AI; existují moduly vhodné pro aktivitu „Hour of Code“, které snižují bariéru vstupu a usnadňují praktické ukázky pro žáky a studenty \cite{kb_ai_101_tipu_pdf_04e4a115_page_15_2}.
Použití takových světů se osvědčuje tam, kde je cílem názorně vysvětlit koncepty (datové vstupy, chování agentů, učení na příkladech) a zapojit studenty do hands‑on aktivit.
V praktickém nasazení je vhodné připravit krátké scénáře a návody, které učitelům umožní rychlý start bez nutnosti hlubokých technických znalostí.
Doporučuje se také zahrnout reflexní fázi po aktivitě, kde studenti popíší, co agent „pozoroval“ a jak se jeho chování změnilo.

\medskip
\textbf{2) AI nástroje podpora procvičování a analýzy v matematice.}  
Existují specializované vzdělávací nástroje integrované do rodiny Microsoft aplikací, které generují úlohy, shromažďují postupy řešení a poskytují učiteli analytický přehled o silných/slabých stránkách jednotlivých studentů (funkce „Pokrok v matematice“).
Součástí ekosystému jsou i nástroje pro převod rukopisu na digitální rovnice (OneNote) a webové nástroje pro výuku matematiky (math.microsoft.com) \cite{kb_ai_101_tipu_pdf_04e4a115_page_8_1}.
Tyto nástroje ukazují, že užitečná implementace AI ve výuce často kombinuje několik vrstev funkcionality:
\begin{itemize}
  \item generování diferenciovaných úloh pro procvičování,
  \item sběr dokladů o řešení (postupy, ne pouze výsledky) jako podklad pro hodnocení,
  \item přehledy pro učitele usnadňující cílenou intervenci.
\end{itemize}
Kromě technických funkcí je důležité myslet na didaktickou integraci: jak úkoly zapadají do kurikula a jaké formy zpětné vazby jsou pro studenty nejpřínosnější.
Učitelé ocení možnost nastavovat obtížnost a monitorovat postup studentů na úrovni třídy i jednotlivce.

\medskip
\textbf{3) Nízkoprahové integrační body (příklady z desktopu).}  
Funkce integrované v běžném softwaru (např. Copilot v Poznámkovém bloku — sumarizace a přepis) umožňují rychlé workflow pro učitele při přípravě výukových materiálů a syntéz z více zdrojů \cite{kb_ai_101_tipu_pdf_04e4a115_page_14_1}.
To je praktické zvláště při tvorbě stručných přehledů, konzpektů a rychlých materiálů do výuky.
Nízkoprahové nástroje snižují odpor vůči změně, protože učitelé je mohou začlenit do svého stávajícího pracovního postupu bez větších výdajů na školení.
Doporučuje se pilotovat tyto funkce nejprve interně mezi vyučujícími a sbírat příklady dobré praxe, které lze později sdílet.

\bigskip
\noindent \textbf{Přenosné principy z ověřených příkladů}
\begin{enumerate}
  \item Zaměření na proces, ne jen výsledek — v matematických nástrojích je kladen důraz na nahrávání postupu řešení a jeho vyhodnocení, což snižuje riziko, že AI bude využita pouze pro hotové odpovědi \cite{kb_ai_101_tipu_pdf_04e4a115_page_8_1}.
  \item Používat připravené interaktivní světy pro názornost — světy jako v Minecraft Education usnadňují demonstraci principů AI a zapojení studentů do praktických aktivit \cite{kb_ai_101_tipu_pdf_04e4a115_page_15_2}.
  \item Využívat existující integrační body v běžném software — začněte tam, kde už učitelé pracují (editor, LMS, poznámky), aby bylo nasazení nízkoprahové a rychle přínosné \cite{kb_ai_101_tipu_pdf_04e4a115_page_14_1}.
\end{enumerate}
K těmto principům patří také doporučení pro zapojení stakeholderů: průběžná komunikace s vedením katedry a IT oddělením, aby byla technická podpora a právní aspekty řešeny včas.
Důležitá je rovněž příprava krátkých instruktážních materiálů pro studenty, které vysvětlí cíle pilotu a pravidla využití nástrojů.

\bigskip
\noindent \textbf{Praktický mini‑pilot: krok za krokem (general background knowledge)}
\begin{enumerate}
  \item Vyberte jeden kurz a definujte jasný cíl pilotu (např. zrychlení formativní zpětné vazby u cvičení, nebo zvýšení angažovanosti prostřednictvím interaktivní aktivity).
  \item Zvolte nástroj s nízkou bariérou (např. Minecraft Education pro demonstrační hodiny nebo „Pokrok v matematice“ pro procvičování) a ověřte licenční/přístupové podmínky. (general background knowledge)
  \item Navrhněte 2–3 metriky úspěchu (které lze měřit): účast v aktivitě, počet validovaných postupů řešení, čas potřebný pro poskytnutí formativní zpětné vazby. (general background knowledge)
  \item Pilot proveďte na malém vzorku (1–2 rozstupy výuky), zaznamenejte data o použití a sbírejte zpětnou vazbu od studentů i vyučujících. (general background knowledge)
  \item Vyhodnoťte bezpečnostní a GDPR aspekty (minimalizace dat, anonymizace) před rozšířením. (general background knowledge)
  \item Dokumentujte průběh, ukládejte verze promptů/nastavení a datum nasazení (provenance) a připravte krátký report s doporučeními pro rozšíření. (general background knowledge)
\end{enumerate}
Více detailů v průběhu pilotu zahrnuje plán školení pro vyučující, jednoduché instrukce pro studenty a mechanismus pro rychlé řešení technických problémů.
Doporučuje se také zajistit sadu testovacích scénářů a záložní plán pro případ výpadku služby nebo licenčních omezení.

\bigskip
\noindent \textbf{Krátké doporučení pro přenositelnost na katedru}  
Začněte od malých, dobře definovaných případů užití (laborka, cvičení, doplňkový tutor) a sdílejte šablony, prompt‑knihovnu a rubriky mezi kolegy; využijte existujících nástrojů a připravených výukových světů pro zrychlení adopce \cite{kb_ai_101_tipu_pdf_04e4a115_page_15_2, kb_ai_101_tipu_pdf_04e4a115_page_8_1}. (general background knowledge)
Navrhněte krátké workshopy pro výměnu zkušeností a vytvořte místo na sdílení materiálů (interní repozitář nebo sdílená složka).
Klíčové je také zahrnout měřitelné výstupy a krátké případové studie, které lze prezentovat vedení katedry při žádosti o rozšíření.

\bigskip
\noindent \textbf{Odkazy a další kroky}  
Pro konkrétní šablony, check‑listy a video tutoriály se podívejte do online repozitáře příručky (living document) — tam jsou připravené exportovatelné materiály a vzorové pilotní plány, které lze upravit pro lokální podmínky. (general background knowledge)
Doporučuje se pravidelně aktualizovat repozitář na základě zkušeností z pilotů a sbírat anonymizovaná data o efektivitě, aby se podpořilo systematické zlepšování praxe.
Pokud je to možné, plánujte další iterace pilotu s rozšířeným počtem kurzů a upravenými metrikami podle prvotních zjištění.%

%
\chapter{AI VE VÝUCE (INTERAKTIVNÍ METODY)}%
Kapitola ukazuje, jak integrovat AI do interaktivních metod výuky: virtuální teaching asistenti (custom GPTs), live demonstrace a simulace, facilitace diskusí s AI a gamifikace. Přináší návod na vytvoření vlastního GPT asistenta pro předmět, příklady scénářů pro sokratovský dialog a techniky pro adaptaci vysvětlení podle úrovně studentů (Explain it like I'm 5/15/25). Dále popisuje použití AI pro escape roomy, role‑playing a podporu mezinárodních i neurodivergentních studentů. Opírá se o ověřené příklady z univerzit, kde AI asistenti snižují zátěž office hours a urychlují zpětnou vazbu (inspirováno např. kurzy typu CS50).%
\section{Úvod a pedagogické cíle interaktivních AI metod}%
Stručné vymezení kapitoly: proč integrovat AI do interaktivní výuky, očekávané didaktické přínosy a hlavní rizika k řešení. Nastavení cílů pro předměty (zvýšení angažovanosti, zrychlení zpětné vazby, podpora diferenciace).%


%
\medskip

Integrace interaktivních AI metod do výuky má dvě vzájemně se doplňující roviny: pedagogické cíle (co chceme u studentů podpořit) a praktické nástroje/scénáře, které to umožní. Níže nejprve shrnujeme hlavní očekávané přínosy (stručně jako doporučené cíle) a poté uvádíme konkrétní nízkoprahové příklady, které lze v předmětu otestovat.

\smallskip
\textbf{Doporučené pedagogické cíle (shrnutí, general background knowledge)}  
\begin{itemize}
  \item Zvýšení angažovanosti studentů skrze interaktivní, praktické aktivity a simulace. (general background knowledge)
  \item Zrychlení a škálování formativní zpětné vazby pomocí AI nástrojů, které generují komentáře nebo diagnostiku chyb. (general background knowledge)
  \item Podpora diferenciace učení — adaptivní úlohy a personalizované nápovědy pro různé úrovně studentů. (general background knowledge)
  \item Důraz na proces (dokumentace postupu, milníky) místo pouze konečného výsledku; integrace reflexe o použití AI do hodnocení. (general background knowledge)
\end{itemize}

\smallskip
\textbf{Nízkoprahové praktické příklady — co zkusit v prvním pilotu}  

\begin{enumerate}
  \item Minecraft Education — názorné demonstrace principů AI a „Hour of Code“ aktivity.  
    Minecraft Education obsahuje připravené světy zaměřené i na principy fungování AI a zdarma dostupné hodiny kódu, které lze použít pro rychlou ukázku a hands‑on aktivity ve výuce \cite{kb_ai_101_tipu_pdf_04e4a115_page_15_2}. Praktický krok: stáhněte Minecraft Education a vyzkoušejte se studenty jednu ze zdarma dostupných hodin kódu; po aktivitě zařaďte krátkou reflexi, kde studenti popíšou, co agent „pozoroval“ a jak se chování změnilo \cite{kb_ai_101_tipu_pdf_04e4a115_page_15_2}.
  \item Procvičování s diagnostikou chyb — „Pokrok v matematice“.  
    Funkce „Pokrok v matematice“ umožňuje učiteli vygenerovat sadu příkladů, vyžadovat nahrání postupu řešení a získat přehled o oblastech, kde studenti chybují; nástroj také předopraví a doporučí oblast s chybou, čímž usnadňuje cílenou intervenci \cite{kb_ai_101_tipu_pdf_04e4a115_page_8_1}.
  \item Rychlá tvorba elektronických testů — Microsoft 365 Copilot + Forms.  
    Kombinace Microsoft 365 Copilot a Forms umí automaticky generovat písemky/kvízy (např. mix ABC, pravda/nepravda, otevřené otázky) včetně označení správných odpovědí a krátkých vysvětlení — užitečné pro rychlé, škálovatelné hodnocení v kontrolovaných podmínkách \cite{kb_ai_101_tipu_pdf_04e4a115_page_36_1}.
\end{enumerate}

\smallskip
\textbf{Krátký plán prvního interaktivního sezení (praktický checklist)}  
\begin{enumerate}
  \item Vyberte jeden jasný cíl sezení (např. „ukázat princip učení agenta“ nebo „rychlá diagnostika chyb v řešení rovnic“). (general background knowledge)
  \item Zvolte nástroj podle cíle: Minecraft Education pro demonstrace (viz výše) \cite{kb_ai_101_tipu_pdf_04e4a115_page_15_2}, „Pokrok v matematice“ pro procvičování s požadavkem na postupy řešení \cite{kb_ai_101_tipu_pdf_04e4a115_page_8_1}, nebo Copilot+Forms pro generování kvízu \cite{kb_ai_101_tipu_pdf_04e4a115_page_36_1}.
  \item Připravte krátký instrukční list pro studenty (co se očekává, jak deklarovat použití AI) a jednu reflexní otázku po aktivitě. (general background knowledge)
  \item Po pilotu zaznamenejte data a zpětnou vazbu: co fungovalo, jaké technické/etické problémy přišly, návrh další iterace. (general background knowledge)
\end{enumerate}

\smallskip
Poznámka k rizikům a řízení kvality: konkrétní rizika (ochrana osobních údajů, akademická integrita, nekonzistentní nebo nesprávné výstupy AI) je třeba před pilotem zhodnotit a popsat v krátkém provozním plánu; tato doporučení jsou rozvedena v kapitolách věnovaných rizikům, GDPR a akademické integritě (viz příručka). (general background knowledge)%

%
\section{Vytvoření vlastního GPT asistenta pro předmět}%
Praktický návod krok za krokem: definice rolí a kompetencí asistenta, návrh konverzačních promptů, integrace s LMS/SSO a enterprise nástroji, provozní údržba a monitorování chyb. Zmíněny požadavky na transparentnost a zodpovědnost při nasazení.%
\subsection{Definice role a kompetencí asistenta}%
Vymezení očekávaných rolí (např. tutoring, administrativní podpora), hranic kompetencí a pravidel pro eskalaci k lidskému pracovníkovi; doporučené výstupy pro studenty a učitele.%


%
\begin{itemize}
\item Účel: jasně definujte, zda asistent slouží jako tutor, generuje administrativní odpovědi, provádí předběžnou diagnostiku nebo facilitačně podporuje diskusi.
\item Kompetence: vysvětlovat koncepty v různé hloubce, navrhovat cvičení, kontrolovat formální náležitosti a indikovat možné chyby.
\item Hranice: nesmí poskytovat kompletní řešení hodnocených úloh, měnit hodnocení bez lidské kontroly ani zpracovávat citlivá data bez schváleného workflow.
\item Eskalace: plná řešení směrovat učiteli; při podezření na podvody oznamovat kontaktnímu pracovníkovi/katedře.
\item Eskalace (pokrač.): technické chyby či opakované nesprávnosti hlásit správci a rollbackovat konfiguraci.
\item Výstupy pro studenty: krátký hint (1–3 věty) a diagnostická poznámka „pravděpodobná oblast chyby“ (viz "Pokrok v matematice") \cite{kb_ai_101_tipu_pdf_04e4a115_page_8_1}.
\item Dále: výzva k nahrání mezikroků jako podmínka ověření postupu \cite{kb_ai_101_tipu_pdf_04e4a115_page_8_1}.
\item Výstupy pro učitele: agregovaný přehled problémových oblastí třídy, ukázky typických chyb a log konverzací/configurací (v souladu s GDPR).
\item Dokumentace a školení: doporučená nastavení a trénink pro pedagogy — např. Microsoft Learn \cite{kb_ai_101_tipu_pdf_04e4a115_page_54_1}.
\item Rychlý checklist před spuštěním: definujte role a hranice, zaveďte pravidla eskalace a kontakty, určete kdy požadovat nahrání postupu.
\item Připravte krátké školení pro pedagogy a komunikaci pro studenty o očekáváních a deklaraci využití AI.
\end{itemize}%

%
\subsection{Návrh konverzačních promptů a šablon}%
Praktické šablony promptů pro běžné scénáře (FAQ, vysvětlení úloh, zpětná vazba) a postupy pro testování a ladění konverzačních toků s ohledem na přesnost a bezpečnost.%


%
\begin{quotation}Krátký praktický soubor šablon promptů pro konverzační asistenty (FAQ, vysvětlení úloh, zpětná vazba) a doporučený postup testování a ladění konverzačních toků v kurzu.\end{quotation}

\bigskip

\begin{itemize}
  \item Zásadní: jasně definujte roli asistenta (co smí/nesmí) a u kritických scénářů vyžadujte eskalaci k lidskému vyučujícímu.
\end{itemize}

\medskip

\textbf{Šablony promptů (kopírovatelné)}
\begin{enumerate}
  \item FAQ:
\begin{verbatim}
SYSTEM: Jsi asistent kurzu <NÁZEV>. Stručně, fakticky, odkazuj na materiály.
USER: "Kdy je termín a formát?"
ASSISTANT: [stručná odpověď + odkaz]
\end{verbatim}
  \item Vysvětlení (A/B/C):
\begin{verbatim}
SYSTEM: Nabídni vysvětlení A (ELI5) / B (SŠ) / C (pokročilé).
USER: "Vysvětli gradient descent."
ASSISTANT: [A / B / C]
\end{verbatim}
  \item Zpětná vazba (mezikroky):
\begin{verbatim}
SYSTEM: Hodnoť jen předložené kroky, nepiš hotové řešení.
ASSISTANT: 1) Hint (1–3 věty). 2) Diagnostika. 3) Doporučení eskalace.
\end{verbatim}
Tento vzor vrací hint + diagnostickou poznámku \cite{kb_ai_101_tipu_pdf_04e4a115_page_8_1}.
\end{enumerate}

\textbf{Checklist testování a ladění}: definujte intenty a hranice; sestavte 20–50 testů včetně hraničních příkladů; validujte bezpečnost a diagnostiku mezikroků \cite{kb_ai_101_tipu_pdf_04e4a115_page_8_1}; simulujte pilot (např. Minecraft Education \cite{kb_ai_101_tipu_pdf_04e4a115_page_15_2}).

Krátké tipy: zakazujte kompletní řešení u hodnocených úloh, nastavte délku odpovědí, testujte v pilotu, dokumentujte změny a škálujte postupně.%

%
\subsection{Integrace s LMS, SSO a enterprise nástroji}%
Kroky pro technickou integraci do prostředí VUT (LMS, jednotné přihlášení), zapojení enterprise služeb a RAG patternů; praktické poznámky k licencím a zabezpečení.%


%
% Integrace s LMS, SSO a enterprise nástroji
Cíl: stručně popsat kroky k integraci enterprise AI služeb do výuky (LMS, SSO, interní zdroje).
Kroky integrace — praktický postup
\begin{enumerate}
\item Prověřit licence/DPA/SLA a retenci dat; podepsat DPA před zapojením kurzů.
\item Nastavit autentizaci přes VUT IdP (SAML/OAuth2/OIDC), role, logování a audity.
\item Napojit LMS přes LTI/REST/webhooks; v MS‑ekosystému využít OneNote a převod rukopisu \cite{kb_ai_101_tipu_pdf_04e4a115_page_8_1}.
\item Využít školní enterprise funkce (Minecraft Education, Copilot v appkách) v pilotech \cite{kb_ai_101_tipu_pdf_04e4a115_page_15_2, kb_ai_101_tipu_pdf_04e4a115_page_53_1}.
\item Implementovat RAG/retrieval: indexace autorizovaných kolekcí, embeddings, testy leakage.
\item Spustit mini‑pilot, měřit použitelnost a bezpečnost, dokumentovat prompty a provenance.
\end{enumerate}
Krátký checklist: licence a DPA; SSO a mapování rolí; minimalizovat odesílaná data; testovat Minecraft/OneNote/Copilot; auditování, metriky a rollback; plánovat pilot s fakultním IT a právním oddělením.%

%
\subsection{Provozní údržba a monitorování chyb}%
Postupy pro nasazení, pravidelnou údržbu modelu a promptů, sledování metrik výkonu a chybových stavů, postupy pro rollback a řízení incidentů.%


%
\begin{flushleft}Krátký praktický návod pro provozní údržbu AI asistentů a monitorování chyb v kontextu kurzů a pilotů na VUT; doporučeno krátké piloty a dokumentace promptů/provenance.\end{flushleft}
\begin{itemize}
  \item Pravidelné monitorování: metriky dostupnosti/latence/error rate; vzorkovat kvalitu výstupů (halucinace) a ukládat anotované ukázky; sledovat odesílání citlivých dat.
  \item Verzování promptů a modelů: uchovávat prompty s konfigurací a časovou známkou, vést changelog; verzovat metadata modelu pro možnost rollbacku.
  \item Provenance: anonymizované logy dotazů/odpovědí s kontextem (ID úkolu, verze promptu) a zaznamenat, které dokumenty byly retrieved.
  \item Incidenty a rollback:
    \begin{enumerate}
      \item izolace služby (read‑only / maintenance) a zabránění zápisům obsahujícím citlivá data;
      \item revert na poslední „safe“ verzi promptu/modelu a spuštění smoke tests;
      \item root‑cause analýza, rozhodnutí o dalším postupu a informování vyučujících, IT a GDPR/právního oddělení s dokumentací.
    \end{enumerate}
  \item Testy a metriky: automatizované smoke/acceptance testy pro kritické scénáře a měření uživatelských metrik pilotů (účast, zpětná vazba, spokojenost).
  \item Revize: kvartální kontrola promptů, modelů a DPA/SLA, vyřazení zastaralých promptů/indexů a aktualizace dokumentace pro vyučující.
\end{itemize}
\begin{flushleft}Poznámka: využít enterprise funkce pro lepší audit a řízení přístupu a vždy ověřit smluvní podmínky a retenci dat \cite{kb_ai_101_tipu_pdf_04e4a115_page_3_1}.\end{flushleft}%

%
\subsection{Transparentnost, zodpovědnost a soulad s pravidly}%
Požadavky na transparentní oznámení použití AI, zásady akademické integrity, ochranu osobních údajů a zajištění vysvětlitelnosti odpovědí studentům a garantům předmětu.%


%
Krátké zásady pro transparentní a odpovědné nasazení AI v předmětu, zaměřené na komunikaci se studenty a garanty:

\begin{itemize}
  \item Ve sylabu a na stránce předmětu jasně uveďte, které AI nástroje jsou povoleny/zakázány, jak mají studenti deklarovat použití AI a jak bude použití ovlivňovat hodnocení (viz šablony v kapitole 2.6 a příloha 9.3). % internal cross‑ref
  \item Požadujte jednoduchou provenance: u úloh, kde bylo AI použito, nechte studenty připojit krátké shrnutí co nástroj dělal, jaký prompt použili a které části výsledku upravili (general background knowledge).
  \item Zajistěte lidský dohled pro klíčové situace: kritické hodnocení výstupů AI provádí vyučující nebo asistent, zejména tam, kde výstupy mohou obsahovat chybné nebo citlivé informace.
  \item Logování a audit: pokud používáte enterprise nebo integrované řešení, aktivujte auditní logy a retenci provenance (dotazy/odpovědi, verze promptů/modelů) pro účely vyšetřování incidentů a zpětné kontroly \cite{kb_ai_101_tipu_pdf_04e4a115_page_3_1}.
  \item Informujte studenty o omezeních nástrojů (např. riziko „halucinací“) a o tom, jak ověřovat citace a fakta (general background knowledge). Pro technické integrace do testů a formulářů zvažte konkrétní procesy autorizace nástrojů (např. Copilot ve Forms) a popište je v instrukcích pro kurz \cite{kb_ai_101_tipu_pdf_04e4a115_page_36_1}.
  \item Dohoda o ochraně dat: při sběru dotazů/odpovědí nebo logů uveďte účel zpracování a dobu retention; pokud jde o citlivá data studentů, řešte to přes DPIA a DPA dle kapitoly 2.2. % internal cross‑ref
\end{itemize}

Praktický tip pro vyučující: připravte jednoznačnou větu do sylabu („Používání AI v tomto předmětu je povoleno za těchto podmínek: ...“), a při pilotu zaznamenávejte případy použití a incidenty pro kvartální revizi politiky.%

%
\section{Live demonstrace a simulace (včetně Minecraft Education)}%
Postupy pro živé ukázky a interaktivní simulace v hodině; praktické příklady vedení hodin s využitím virtuálních prostředí (včetně Minecraft Education jako didaktického nástroje pro demonstraci principů AI a Hodina kódu, viz zdroj).%


%
Tato část popisuje praktické postupy pro živé ukázky a interaktivní simulace v hodině se zaměřením na světy typu Minecraft Education a další nástroje, které snadno nasadíte pro názorné předvedení principů agentního chování, učení a diagnostiky chyb.

Krátké shrnutí přínosu
\begin{itemize}
  \item Interaktivní simulace umožňují studentům pozorovat chování „agenta“ v kontrolovaném prostředí a zkoušet dopady jednoduchých změn pravidel či dat (general background knowledge).
  \item Minecraft Education nabízí hotové „světy“ a aktivity, které lze využít k vysvětlení principů strojového učení a k realizaci krátké „Hour of Code“ lekce dostupné zdarma; zbytek obsahu je součástí školní licence Microsoft 365 A3/A5 \cite{kb_ai_101_tipu_pdf_04e4a115_page_15_2}.
\end{itemize}

Praktické scénáře (konkrétní příklady)
\begin{enumerate}
  \item Demonstrace principů agentního chování v Minecraft Education: vyberte předpřipravený svět s aktivitou zaměřenou na AI nebo Hour of Code, spusťte krátkou ukázku (10--20 min), nechte studenty komentovat, proč agent reagoval tak jak reagoval, a diskutujte omezení přístupu. Doporučený start: stáhnout Minecraft Education a vyzkoušet jednu z volně dostupných lekcí Hour of Code \cite{kb_ai_101_tipu_pdf_04e4a115_page_15_2}.
  \item Diagnostika chyb a adaptivní procvičování: využijte nástroj typu „Pokrok v matematice“ jako ukázku, jak systém generuje příklady, sbírá postupy studentů a poskytuje přehledy silných/ slabých stránek třídy; při ukázce vysvětlete, proč je nahrání postupu důležité pro formativní hodnocení \cite{kb_ai_101_tipu_pdf_04e4a115_page_8_1}.
  \item Rychlá tvorba vizuálů pro diskusi: během debriefu si v reálném čase nechte vygenerovat ilustraci nebo schéma (např. pomocí Designer) k probíranému tématu a porovnejte různé stylistické varianty; to pomáhá udržet pozornost a ukazuje limity automatické vizualizace \cite{kb_ai_101_tipu_pdf_04e4a115_page_51_1}.
\end{enumerate}

Před hodinou — kontrolní seznam přípravy
\begin{itemize}
  \item Jasně definujte cíle sezení (co mají studenti pochopit/umět pozorovat). (general background knowledge)
  \item Ověřte technické prostředí: instalace Minecraft Education (nebo alternativy), přístup do potřebných webových služeb, záložní plán pro případ výpadku sítě. Adresa pro Minecraft Education a volné lekce je uvedena v dostupných materiálech produktu \cite{kb_ai_101_tipu_pdf_04e4a115_page_15_2}.
  \item Připravte krátký instrukční list pro studenty s pravidly použití nástrojů a jednou reflexní otázkou (např. „Co by změnila drobná úprava pravidel chování agenta?“). (general background knowledge)
  \item Rozhodněte, zda budou studenti pracovat individuálně, ve dvojicích nebo v malých skupinách; připravte časovou osu aktivity (celkem 30--60 min doporučeno).
\end{itemize}

Průběh ukázky — doporučený časový rámec (45--60 min)
\begin{enumerate}
  \item 5--10 min: úvod a stanovení učebního cíle.
  \item 15--25 min: live demo v Minecraft Education nebo jiném simulátoru (včetně krátkých experimentů s parametry).
  \item 10--15 min: studentské mini-aktivity (např. navrhnout změnu pravidla a předpovědět dopad).
  \item 10--15 min: debrief, diskuse, porovnání pozorování s předpověďmi, závěrečná reflexe.
\end{enumerate}

Facilitace a pedagogické tipy
\begin{itemize}
  \item Moderujte diskusi tak, aby studenti reflektovali nejen „co se stalo“, ale proč se to stalo a jaké jsou limity daného modelu/světa (general background knowledge).
  \item Při použití adaptivních cvičení (např. nástrojů typu Pokrok v matematice) zdůrazněte, že nahraný postup je klíčový pro pedagogickou diagnostiku; ukažte příklady, jak systém identifikuje chybu a navrhuje oblast pro další procvičování \cite{kb_ai_101_tipu_pdf_04e4a115_page_8_1}.
  \item Používejte generované vizuály (Designer apod.) jako doplněk, nikoli jako jediný zdroj vysvětlení; ukažte studentům, jak různé prompty mění výsledek \cite{kb_ai_101_tipu_pdf_04e4a115_page_51_1}.
\end{itemize}

Zhodnocení a reflexe
\begin{itemize}
  \item Na konci sezení požádejte studenty o krátké shrnutí (1--2 věty) o tom, co se naučili a co by změnili v experimentu (může sloužit jako rychlý formativní checkpoint). (general background knowledge)
  \item Zaznamenejte technické a etické problémy pro kvartální revizi pilotu (co fungovalo, kde vznikly nejasnosti, potřeba jiných licencí nebo DPIA pokud sbíráte citlivá data). Použití Minecraft Education a dostupnost Hour of Code materiálů jsou vhodné pro první nízkoprahové piloty \cite{kb_ai_101_tipu_pdf_04e4a115_page_15_2}.
\end{itemize}

Krátké doporučení pro další kroky
\begin{itemize}
  \item Pro první pilot zvolte jednu ukázkovou lekci z Minecraft Education (Hour of Code) a proveďte ji včetně reflexe; výsledky pilotu použijte pro úpravu sylabu a instrukcí pro studenty \cite{kb_ai_101_tipu_pdf_04e4a115_page_15_2}.
  \item Pokud budete demonstrovat adaptivní procvičování (diagnostiku chyb), předveďte i způsob, jak učitel finalizuje automatické návrhy (uvedeno u „Pokrok v matematice“) \cite{kb_ai_101_tipu_pdf_04e4a115_page_8_1}.
\end{itemize}

Poznámka: část praktických a organizačních kroků je uvedena jako obecné pedagogické doporučení; konkrétní právní a GDPR aspekty je vhodné řešit podle kapitol věnovaných ochraně osobních údajů a transparentnosti v této příručce (viz příslušné šablony a checklisty).%

%
\section{Facilitace diskusí s AI: sokratovský dialog a adaptivní vysvětlování}%
Techniky vedení sokratovských dialogů s AI, návrhy scénářů pro postupné zvyšování složitosti vysvětlení (Explain it like I'm 5/15/25) a doporučené postupy pro moderování a zajištění kvality výstupů AI v diskusích.%


%
% (Tělo sekce: Facilitace diskuzí s AI — sokratovský dialog a adaptivní vysvětlování)
Krátké shrnutí a cíle (general background knowledge).
\begin{itemize}
  \item Cíl: vést studenty k hlubšímu porozumění pomocí řízených otázek místo jednoduchého poskytování odpovědí; ukázat, jak lze AI použít jako sokratovského partnera a nástroj pro adaptivní vysvětlování na několika úrovních složitosti.
  \item Výstup: studenti umějí argumentovat vlastními slovy, reflektovat mezikroky řešení a rozpoznat omezení generovaných odpovědí.
\end{itemize}

Praktický workflow před hodinou (general background knowledge).
\begin{enumerate}
  \item Definujte učební cíl diskuse (co chcete, aby studenti vysvětlili/pochopili).
  \item Připravte krátký scénář nebo artefakt (kód, úloha, fragment textu) jako výchozí podnět pro AI a studenty.
  \item Zvolte režim práce (celá třída vs. skupiny) a připravte instrukce pro záznam postupu (poznámky, screenshoty, nahrávky).
  \item Pokud plánujete živou demonstraci s vizualizací nebo simulací, připomeňte, že Minecraft Education nabízí hotové světy pro výuku principů AI včetně obsahu Hour of Code (bezplatně dostupného) — zbytek obsahu je součástí školní licence Microsoft 365 A3/A5 \cite{kb_ai_101_tipu_pdf_04e4a115_page_15_2}. 
\end{enumerate}

Základní role AI v sokratovském dialogu (general background knowledge).
\begin{itemize}
  \item „Sokratovský kouč“: klade otevřené otázky, navrhuje dílčí kroky, vyhodnocuje konzistenci odpovědí.
  \item „Diagnostik“: identifikuje mezery v postupu studenta a navrhuje cvičení či hinty.
  \item „Vizualizátor“: generuje doplňkové schéma nebo obrázek k vysvětlení; k těmto účelům lze v průběhu debriefu použít nástroje jako Designer pro rychlou tvorbu ilustrací a porovnání více stylistických variant \cite{kb_ai_101_tipu_pdf_04e4a115_page_51_1}.
\end{itemize}

Šablony promptů (praktické, copy‑paste ready).
\begin{itemize}
  \item System (role): 
\begin{PromptBlock}


You are a Socratic tutor for university students. Do NOT provide full solutions; instead ask 2--4 guiding questions that reveal the next logical step. Use simple language and request justification for each step.
\end{PromptBlock}


  \item User (student prompt, příklad): 
\begin{PromptBlock}


I have this partial solution to problem X: [vložit stručný popis]. Help me find my next step.
\end{PromptBlock}


  \item Assistant (model output — očekávaný formát): otázky číslované, každá s krátkým vysvětlením proč je relevantní; na konci výzva: \texttt{Which of these do you want to try first?}
\end{itemize}

Adaptivní vysvětlování: ELI5 / ELI15 / ELI25 — návrh scénářů (general background knowledge).
\begin{enumerate}
  \item Představte tři požadavky ve svém promptu: 
\begin{PromptBlock}


Explain this concept (a) like I'm 5, (b) like I'm 15, (c) like I'm 25. For each, keep length under 50/150/300 words respectively and give one short analogy.
\end{PromptBlock}


  \item Použijte ELI varianty k diferenciaci: studenti začnou s ELI5 pro intuitivní uchopení, poté ELI15 pro střední úroveň abstrakce a nakonec ELI25 pro technickou syntézu a kritickou reflexi.
  \item Moderujte: vyzvěte studenty, aby pro každou úroveň napsali jednu kontrolní otázku, kterou by položili modelu — tím trénují kritické čtení odpovědí (general background knowledge).
\end{enumerate}

Příklad hodinové aktivity (45--60 min, adaptováno z doporučených časových rámců).
\begin{enumerate}
  \item 5--10 min: úvod a zadání problému; ukažte artefakt (kód, výňatek textu).
  \item 15--25 min: studenti v malých skupinách vedou dialog s AI podle šablon (sokratovský režim), sbírají otázky a návrhy kroků.
  \item 10--15 min: skupiny prezentují závěry; facilitator vybere jednu odpověď a vsadí ji proti generované vizualizaci nebo simulaci (pokud je relevantní).
  \item 5--10 min: debrief — reflexe, co AI přidala, kde selhala, které otázky byly klíčové.
\end{enumerate}

Moderování a zajištění kvality výstupů (general background knowledge).
\begin{itemize}
  \item Požadujte od modelu, aby vždy uváděl zdůvodnění kroků a případné předpoklady; neakceptujte anonymní tvrzení bez odkazu na postup.
  \item Učitel by měl vzorkovat a anotovat příklady halucinací nebo nevhodných rad pro kvartální revizi pilotu (general background knowledge). Pokud používáte adaptivní cvičení, můžete doplnit o požadavek na nahrání postupu studenta — u podobných nástrojů je právě postup klíčový pro diagnostiku a návrhy dalšího procvičování \cite{kb_ai_101_tipu_pdf_04e4a115_page_8_1}.
  \item Používejte generované vizuály jako pomocný materiál, nikoliv jako jediný zdroj; při porovnání variant ukažte studentům, jak změna promptu mění výsledek \cite{kb_ai_101_tipu_pdf_04e4a115_page_51_1}.
  \item V případě živých simulací/logování, zaznamenávejte hlavní dotazy a verze promptů pro pozdější analýzu (general background knowledge).
\end{itemize}

Rychlé tipy pro učitele (pragmatické).
\begin{itemize}
  \item Nasaďte jasná pravidla: AI smí navrhovat kroky, ale student musí vždy ukázat vlastní mezikroky a rozhodnutí (general background knowledge).
  \item Na první pilot zvažte použít Minecraft Education pro názornou demonstraci principů agentního chování nebo jednoduché „Hour of Code“ lekce — tyto lekce jsou volně dostupné a pomáhají vizualizovat algoritmické chování agentů \cite{kb_ai_101_tipu_pdf_04e4a115_page_15_2}.
  \item Při využití multimédií (obrázky, schémata) mějte záložní plán, pokud generování selže; Designer lze využít pro rychlé generování a porovnání stylů ilustrací během debriefu \cite{kb_ai_101_tipu_pdf_04e4a115_page_51_1}.
\end{itemize}

Poznámka o etice a transparenci (general background knowledge).
\begin{itemize}
  \item Informujte studenty o použití AI při facilitaci diskuse a požadujte krátké sebezprávy o tom, jak studenti AI použili ve své reflexi (deklarace použití).
\end{itemize}%

%
\section{Gamifikace: escape roomy, role‑playing a kvízy s Copilotem}%
Nápady a šablony pro gamifikované aktivity (escape roomy, role‑playing) a praktické využití AI při tvorbě kvízů a testů (vč. Microsoft 365 Copilot + Forms pro generování elektronických písemek, viz zdroj).%


%
Gamifikace nabízí bohaté možnosti, jak zvýšit angažovanost studentů pomocí úkolů typu escape room, role‑playing scénářů nebo interaktivních kvízů. Dvě konkrétní technologická doporučení, která lze na VUT snadno využít, jsou (1) použití Minecraft Education pro názorné, herně založené aktivity a (2) využití Microsoft 365 Copilota v kombinaci s Forms pro rychlé generování elektronických písemek a kvízů. Minecraft Education obsahuje předpřipravené světy a lekce včetně materiálů „Hour of Code“, které usnadňují didaktické nasazení herních aktivit ve výuce \cite{kb_ai_101_tipu_pdf_04e4a115_page_15_2}. Microsoft 365 Copilot lze v prostředí školních Office 365 použít jako integrovaný nástroj a jeho kombinace s Forms umožňuje generování kompletních testů včetně označení správných odpovědí a doplňujících vysvětlení \cite{kb_ai_101_tipu_pdf_04e4a115_page_3_1, kb_ai_101_tipu_pdf_04e4a115_page_36_1}.

Praktické vzory použití
\begin{enumerate}
  \item Escape room (rychlý přehled, general background knowledge):
    \begin{itemize}
      \item Navrhněte 4–6 „stanovišť“ (puzzle), která vyžadují kombinaci logického myšlení, práce s textem/kódem a týmové komunikace.
      \item Pro každé stanoviště připravte primární hádanku a 2–3 postupné nápovědy (hints), které se odhalují podle času nebo podle správných mezikroků.
      \item Použijte Copilot k vygenerování variant hintů a k diagnostice typických chyb (např. rozdílné nápovědy pro začátečníky vs. pokročilé) — tento postup je uveden níže jako šablona promptu (general background knowledge).
    \end{itemize}

  \item Role‑playing (rychlý přehled, general background knowledge):
    \begin{itemize}
      \item Vytvořte scénář s jasně definovanými rolemi (student jako inženýr, klient, etický poradce apod.) a úkolem, který musí tým vyřešit.
      \item Konfigurujte Copilot jako „NPC“ (non‑player character) se zvláštní osobností/odbornostní rolí, aby poskytoval reakce ověřitelné proti zadaným pravidlům (general background knowledge).
    \end{itemize}

  \item Generování kvízů a elektronických písemek s Microsoft 365 Copilot + Forms:
    \begin{enumerate}
      \item Využití: Microsoft 365 Copilot v kombinaci s Forms dokáže vytvořit test obsahující různé typy otázek (ABC, pravda/nepravda, otevřené), označit správné odpovědi a často doplnit vysvětlení volby správné odpovědi \cite{kb_ai_101_tipu_pdf_04e4a115_page_36_1}.
      \item Dostupnost: školní verze Office 365 zahrnuje webového Copilota a některé školy mají přístup k Microsoft 365 Copilot v rámci školních licencí \cite{kb_ai_101_tipu_pdf_04e4a115_page_3_1}.
      \item Příklad promptu (copy‑paste ready, převzato z praktických ukázek):
        \begin{quote}
          
\begin{PromptBlock}


Vygeneruj mi test na hodinu zeměpisu na téma Austrálie, test bude obsahovat otázky na základní geografické oblasti, bude pro žáky 8. ročníku základní školy v České republice, test bude mít 10 otázek, 4 otázky budou typu ABC, 4 otázky typu pravda/nepravda a 2 otázky otevřené.
\end{PromptBlock}


        \end{quote}
        Tento typ promptu ilustruje, že výsledkem je hotový test včetně různých formátů otázek a často i krátkých vysvětlení správných odpovědí \cite{kb_ai_101_tipu_pdf_04e4a115_page_36_1}.
    \end{enumerate}
\end{enumerate}

Šablony promptů a doporučené postupy pro Copilot + Forms
\begin{itemize}
  \item Základní prompt pro rychlý kvíz (čeština, copy‑paste):
    \begin{quote}
      
\begin{PromptBlock}


Vytvoř test s 10 otázkami na téma [TÉMA]. Požadavky: pro [úroveň studentů], 6 otázek typu ABC, 2 pravda/nepravda, 2 otevřené; u každé otázky uveď klíč k odpovědi a 1–2 větami zdůvodni správnou odpověď.
\end{PromptBlock}


    \end{quote}
    (Varianta promptu vychází z použitých příkladů v pramenech; Copilot v prostředí Forms obvykle vrátí strukturovaný soubor otázek a odpovědí, který lze přímo vložit do Forms \cite{kb_ai_101_tipu_pdf_04e4a115_page_36_1}.)
  \item Prompt pro generování nápověd (hints) pro escape room (general background knowledge):
    \begin{quote}
      
\begin{PromptBlock}


Pro hádanku: [stručný popis hádanky] vygeneruj tři postupné nápovědy. 1) velmi nápovědná (pro studenty, kteří jsou ztracení), 2) mírná nápověda, 3) pouze drobný tip, který usměrní myšlení. Každou nápovědu napiš v max. 20–40 slovech.
\end{PromptBlock}


    \end{quote}
  \item Prompt pro nastavení „NPC“ v role‑play (general background knowledge):
    \begin{quote}
      
\begin{PromptBlock}


Chovej se jako [role], tvůj tón je [formální/informální], odpovídej stručně (max. 60 slov). Pokud student dává nesprávné argumenty, polož 2 kontrolní otázky místo opravy.
\end{PromptBlock}


    \end{quote}
\end{itemize}

Bezpečnostní a pedagogické poznámky (stručně)
\begin{itemize}
  \item Transparentnost: informujte studenty předem, že aktivita využívá AI (Copilot nebo jiné nástroje) a jaké výstupy jsou očekávány/akceptovatelné (general background knowledge).
  \item Validace obsahu: vygenerované otázky, nápovědy a vysvětlení vždy ověřte pedagogicky; Copilot může generovat syntakticky korektní, ale obsahově nepřesné formulace — proto kontrola garantem je nezbytná \cite{kb_ai_101_tipu_pdf_04e4a115_page_36_1}.
  \item Licence a přístup: před použitím zkontrolujte dostupnost Copilota v rámci školních licencí Office 365 a nastavení Forms; některé funkce generativní AI jsou dostupné v závislosti na zakoupené verzi produktu \cite{kb_ai_101_tipu_pdf_04e4a115_page_3_1}.
  \item Doporučení pro pilot: začněte s jednou herní aktivitou jako pilotem (1–2 týmy) a sbírejte metriky angažovanosti a zpětnou vazbu studentů pro kvartální revizi materiálů (general background knowledge).
\end{itemize}

Rychlý checklist pro realizaci gamifikované aktivity s Copilotem (general background knowledge)
\begin{enumerate}
  \item Definujte vzdělávací cíl aktivity a očekávané výstupy.
  \item Vyberte formát (escape room / role‑play / kvíz) a naplánujte časovou osu.
  \item Připravte podklady (scénáře, hádanky, datové soubory) a zadejte Copilotu jasné prompty pro generování otázek, hintů nebo NPC odpovědí.
  \item Ověřte a upravte vygenerovaný obsah; prověřte správnost faktů a vhodnost jazyka.
  \item Nasazení: vložte otázky do Forms (nebo jiné platformy), nastavte pravidla zveřejnění nápověd a měření úspěšnosti.
  \item Po aktivitě: vyhodnoťte průběh, zaznamenejte chyby/hallucinace, aktualizujte prompty a materiály pro další běh.
\end{enumerate}

Poznámka o doplnění zdrojů
\begin{itemize}
  \item Pro inspiraci při použití Minecraft Education a Hour of Code viz dostupné předpřipravené světy v Minecraft Education, které usnadňují zapojení herních prvků do výuky \cite{kb_ai_101_tipu_pdf_04e4a115_page_15_2}.
  \item Pro rychlé generování kvízů a elektronických písemek je praktickou možností kombinace Microsoft 365 Copilot a Forms, která umí vytvořit test včetně odpovědí a krátkých vysvětlení \cite{kb_ai_101_tipu_pdf_04e4a115_page_36_1}.
\end{itemize}%

%
\section{Podpora mezinárodních a neurodivergentních studentů}%
Konkrétní adaptace interaktivních AI nástrojů pro podporu jazykové rozmanitosti a neurodiverzity: lokalizace, úprava tempo/komplexity, alternativní modality (text/audio/visual) a zásady inkluzivního designu.%


%
Tato sekce shrnuje postupy pro použití interaktivních AI nástrojů ke podpoře studentů s odlišným jazykovým zázemím a seznámení s neurodiverzitou. Cílem jsou rychlé, opakovatelné šablony a kontrolní body pro výuku i online aktivity.

Praktické principy
\begin{itemize}
  \item Lokalizace a kulturní adaptace: neprovádějte mechanický překlad—ověřte příklady vůči cílové skupině.
  \item Úprava tempa a složitosti: připravte minimálně dvě verze (zjednodušená / standardní) a možnost krokového rozkladu.
  \item Multimodalita: kombinujte text, audio a vizuály; jednoduchý diagram nebo piktogram často zefektivní pochopení. Pro rychlou tvorbu vizuálů použijte Microsoft Designer \cite{kb_ai_101_tipu_pdf_04e4a115_page_51_1}.
  \item Důstojnost a autonomie: AI navrhuje podporu, student má vždy možnost ji přijmout nebo odmítnout.
\end{itemize}

Konkrétní nástroje (ilustrativně)
\begin{description}
  \item[Minecraft Education:] vhodný pro vizuální, interaktivní lekce a připravené světy (např. Hour of Code) pro studenty s omezenou znalostí jazyka \cite{kb_ai_101_tipu_pdf_04e4a115_page_15_2}.
  \item[Microsoft Designer:] rychlé generování piktogramů, diagramů a krokových návodů pro doplnění textu \cite{kb_ai_101_tipu_pdf_04e4a115_page_51_1}.
  \item[OneNote / math.microsoft.com:] převod rukopisu na digitální rovnice a analýza postupu řešení (užitečné v matematice) \cite{kb_ai_101_tipu_pdf_04e4a115_page_8_1}.
\end{description}

Pracovní workflow (pro jedno cvičení)
\begin{enumerate}
  \item Definujte cíl podpory: jazyková pomoc, kognitivní rozklad nebo multimodální materiál.
  \item Připravte základní zadání a dvě upravené verze: zjednodušenou a rozšířenou.
  \item Vytvořte stručný vizuál (diagram/piktogram) a sadu konzistentních ikon (Designer).
  \item Volitelně vložte aktivitu do interaktivního prostředí (Minecraft Education).
  \item Požádejte o nahrání postupu (foto rukopisu, audio/video) a využijte převod rukopisu k analýze.
  \item Poskytněte krátkou personalizovanou zpětnou vazbu a alternativní hint (obrázek nebo 30–60 s audio).
\end{enumerate}

Šablony promptů (kopírovatelně)
\begin{itemize}
  \item Lokalizace / jednoduchý překlad:
    \begin{quote}
      
\begin{PromptBlock}


Přelož následující text do [JAZYK] a vytvoř zjednodušenou verzi pro studenty s nižší jazykovou úrovní (max. 5 vět). Poté navrhni krátký obrázkový návod na 3 krocích.
\end{PromptBlock}


    \end{quote}
  \item Víceúrovňové vysvětlení:
    \begin{quote}
      
\begin{PromptBlock}


Vysvětli [POJEM] ve třech úrovních: 1) velmi jednoduché (1–2 věty), 2) střední s příkladem (3–4 věty), 3) detailní s kroky.
\end{PromptBlock}


    \end{quote}
  \item Generování multimodálního hintu:
    \begin{quote}
      
\begin{PromptBlock}


Pro úlohu [stručný popis] vygeneruj: a) krátké textové vodítko (max. 30 slov), b) návrh diagramu s 4 popisky, c) audio skript pro ~45s.
\end{PromptBlock}


    \end{quote}
\end{itemize}

Checklist pro nasazení
\begin{enumerate}
  \item Zkontrolujte jazyková nastavení a dostupnost překladu.
  \item Připravte minimálně dvě úrovně obtížnosti a multimodální verze materiálu.
  \item Zajistěte možnost nahrávání postupu a workflow pro jeho zpracování (OneNote apod.).
  \item Použijte vizuální pomůcky tam, kde text nestačí.
  \item Informujte studenty o roli AI a podmínkách uchovávání materiálů (souhlas/retence).
\end{enumerate}

Evaluace účinnosti
\begin{itemize}
  \item Po pilotu sbírejte krátké ankety o srozumitelnosti, přístupnosti a vlivu na učení.
  \item Měřte využití modalit (audio vs. vizuál) a upravujte materiály podle výsledků.
  \item Pravidelně aktualizujte šablony a vizuální styly dle zpětné vazby.
\end{itemize}

Poznámka o zdrojích
\begin{itemize}
  \item Minecraft Education nabízí připravené světy pro demonstrace a Hour of Code \cite{kb_ai_101_tipu_pdf_04e4a115_page_15_2}.
  \item Microsoft Designer usnadňuje rychlou tvorbu vizuálů v různých stylech \cite{kb_ai_101_tipu_pdf_04e4a115_page_51_1}.
  \item Nástroje Microsoft pro převod rukopisu podporují práci s postupy v matematice \cite{kb_ai_101_tipu_pdf_04e4a115_page_8_1}.
\end{itemize}%

%
\section{Knihovna promptů a scénářů pro interaktivní hodiny}%
Sada ověřených promptů a konverzačních šablon pro typické scénáře: asistent pro Q\&A, sokratovský kouč, simulace rozhovorů, generování cvičení a automatické opravování návrhů odpovědí.%


%
Tato sekce shrnuje ověřené prompt šablony a scénáře pro interaktivní hodiny (Q\\ \&A, sokratovský kouč, simulace, generování cvičení, automatické opravování). Níže jsou stručné účely, doporučený postup nasazení a kopírovatelné prompt‑bloky.

Pravidla nasazení (krátce)
\begin{itemize}
  \item Definujte roli asistenta a transparentně ji oznamte studentům.
  \item Testujte prompty na malém vzorku otázek a dolaďte styl odpovědí.
  \item U úloh citlivých na podvádění požadujte doložení postupu (fotky rukopisu, meziodevzdání). \cite{kb_ai_101_tipu_pdf_04e4a115_page_8_1}
\end{itemize}

1) Asistent pro Q\\ \&A
\begin{itemize}
  \item Účel: rychlé faktické odpovědi, odkazy na materiály, návrhy dalšího kroku.
  \item Role: stručné odpovědi; žádat doplňující informace namísto kompletních řešení.
\end{itemize}
Prompt:
\begin{quote}
\begin{PromptBlock}


Jsi asistent kurzu [NÁZEV PŘEDMĚTU]. Odpovídej stručně (max. 3 věty) a vždy připoj 1 doporučený zdroj z kurzu (slide/kapitola). Pokud dotaz vyžaduje řešení krok po kroku, nejprve vyžádej studentský postup.
\end{PromptBlock}


\end{quote}

2) Sokratovský kouč
\begin{itemize}
  \item Účel: vést studenta otázkami, rozkládat problém, žádat vysvětlení kroků.
  \item Role: klást otevřené otázky a hinty, neposkytovat kompletní řešení.
\end{itemize}
Prompt:
\begin{quote}
\begin{PromptBlock}


Hraj roli sokratovského kouče. Když student zadá problém, polož 3 postupné otázky, které vedou k rozkladu problému, a na závěr nabídni jeden konkrétní hint (ne řešení). Po každé odpovědi studenta požádej o zdůvodnění kroku.
\end{PromptBlock}


\end{quote}

3) Simulace rozhovorů / ústních zkoušek
\begin{itemize}
  \item Účel: trénink odpovědí a simulace obhajoby; po simulaci generovat zpětnou vazbu podle rubriky.
\end{itemize}
Prompt:
\begin{quote}
\begin{PromptBlock}


Simuluj zkoušejícího na úrovni [bakalář/magistr], polož studentovi 5 otázek k tématu [TÉMA], očekávej odpověď 1–2 minuty a po každé odpovědi poskytněte 3-bodové shrnutí hodnocení (obsah, jasnost, hloubka) podle této rubriky: [vložit kritéria].
\end{PromptBlock}


\end{quote}

4) Generování cvičení a kvízů
\begin{itemize}
  \item Užitečné pro rychlé vytvoření in‑class úloh nebo domácích cvičení; možné integrovat s Forms pro automatickou tvorbu kvízů. \cite{kb_ai_101_tipu_pdf_04e4a115_page_36_1}
\end{itemize}
Prompt:
\begin{quote}
\begin{PromptBlock}


Vytvoř test na téma [TÉMA] pro studenty [ROČNÍK/ÚROVEŇ], 10 otázek: 6x ABC, 2x pravda/nepravda, 2x otevřené (max. 80 slov). U každé ABC otázky uveď správnou odpověď a krátké vysvětlení (1 věta).
\end{PromptBlock}


\end{quote}

5) Automatické opravování postupů (matematika, technické obory)
\begin{itemize}
  \item Workflow: studenti nahrají postup (foto/text); nástroj označí chyby a navrhne cílené cvičení a nápravu. \cite{kb_ai_101_tipu_pdf_04e4a115_page_8_1}
\end{itemize}
Prompt:
\begin{quote}
\begin{PromptBlock}


Analyzuj tento studentský postup (vložený text/obrázek). Vyjmenuj 3 konkrétní chyby nebo nejistoty, u každé navrhni krátké cvičení (max. 2 věty) k odstranění problému. Pokud je postup neúplný, vyžádej chybějící krok.
\end{PromptBlock}


\end{quote}

Mini‑workflow nasazení
\begin{enumerate}
  \item Zvolte scénář (Q\\ \&A, kouč, simulace).
  \item Připravte a otestujte 3–5 promptů mimo třídu.
  \item Ve třídě vysvětlete pravidla interakce a transparentnost AI.
  \item Po interakci shromážděte krátkou reflexi (1–2 otázky) a upravte prompty.
\end{enumerate}

Doplňkové poznámky
\begin{itemize}
  \item Pro vizuální materiály lze rychle generovat piktogramy a diagramy (např. nástroje v ekosystému Microsoft).
  \item Preferujte sběr procesu (drafty, meziodevzdání) před pouhým odevzdáním finální odpovědi pro věrné hodnocení a prevenci zneužití AI. \cite{kb_ai_101_tipu_pdf_04e4a115_page_8_1}
\end{itemize}

Poznámka k licencím a integracím: při plánování plného provozu zvažte enterprise varianty nástrojů a integraci s LMS/SSO (viz kapitola o integraci a enterprise řešeních).%

%
\section{Hodnocení, zpětná vazba a zátěž vyučujících}%
Workflows pro využití AI k rychlejší zpětné vazbě studentům, doplnění office hours a škálování hodnocení; praktické doporučení pro kombinaci automatické a fakultativní lidské validace.%


%
Účel této sekce je nabídnout konkrétní pracovní postupy a doporučení, jak využít nástroje generativní AI pro rychlejší a škálovatelné poskytování zpětné vazby studentům, aniž by se výrazně zvýšila pracovní zátěž vyučujících. Níže jsou uvedeny ověřitelné příklady existujících nástrojů, praktické postupy a krátké šablony, které lze okamžitě vyzkoušet v kurzech.

Praktické příklady nástrojů a jejich schopností
\begin{itemize}
  \item Pokrok v matematice (Microsoft Educational Accelerators): učitel vybere oblast, systém vygeneruje sadu příkladů a umožní žákům nahrát postup řešení; AI provede předběžné opravení a identifikuje oblast, kde došlo k chybě, přičemž učitel provede finální kontrolu a uzavření hodnocení \cite{kb_ai_101_tipu_pdf_04e4a115_page_8_1}. Tento přístup zdůrazňuje hodnotu zachycení postupu (process evidence) před pouhým výsledkem.
  \item OneNote a převod rukopisu: nástroje v rodině Microsoft umí převést rukopisné rovnice na digitální výraz a podpořit zobrazení postupu či návrh opravy — užitečné pro sběr meziodevzdání a automatizovanou analýzu kroků \cite{kb_ai_101_tipu_pdf_04e4a115_page_8_1}.
  \item Microsoft 365 Copilot / Copilot Chat: dostupnost Copilota v prostředí Office 365 umožňuje integrovat generativní AI přímo do běžných nástrojů (Word, OneNote, Forms), což usnadní generování komentářů ke studentským textům i přípravu hromadné formativní zpětné vazby \cite{kb_ai_101_tipu_pdf_04e4a115_page_3_1}.
\end{itemize}

Doporučený workflow pro kombinaci automatické zpětné vazby a lidské validace (general background knowledge)
\begin{enumerate}
  \item Navrhněte úlohu tak, aby studentské artefakty obsahovaly krok‑po‑kroku postup (fotky rukopisu, meziodevzdání, commit historii apod.). Toto umožní smysluplné automatické vyhodnocení a auditovatelnou kontrolu procesu.
  \item Nastavte automatickou předkontrolu: AI provede rozpoznání formátu (OCR/rukopis), základní kontrolu správnosti výsledků a identifikaci typických chyb (např. chybné předpoklady, aritmetika, nesoulad mezi kroky).
  \item Vydefinujte pravidla pro eskalaci: výsledky s nízkou jistotou, rozporné postupy nebo vysoce skórované chyby se propisují do fronty pro lidskou validaci.
  \item Provádějte průběžné spot‑checky (např. 5–10\% odevzdaných prací) a kontrolujte kvalitu AI doporučení; učitelé tím udrží kontrolu nad standardem bez nutnosti opravovat každou práci ručně.
  \item Shromažďujte metriky (čas na opravu, poměr automatických korekcí vs. eskalací, typy chyb) a periodicky dolaďujte prompty a pravidla systému.
\end{enumerate}
(Výše uvedený postup je označen jako general background knowledge — principy je třeba adaptovat na specifika předmětu a smluvní/ bezpečnostní požadavky.)

Praktická pravidla pro snížení zátěže vyučujících (general background knowledge)
\begin{itemize}
  \item Automatizujte šablonové komentáře a opakující se drobné opravy; rezervujte ruční práci pro hlubší, hodnotící komentáře.
  \item Používejte milestone‑based odevzdávání (meziodevzdání) tak, aby AI mohl včas odhalit a nasměrovat studenty, což snižuje počet pozdějších oprav.
  \item Nastavte jasné SLA pro odpovědi AI (např. do 24 hodin) a komunikujte studentům, kdy mohou očekávat lidskou kontrolu.
  \item Zvažte integraci s nástroji, které již používáte (Forms, OneNote, LMS), aby se minimalizovala administrativní režie.
\end{itemize}

Ukázkový jednoduchý prompt pro formativní automatickou zpětnou vazbu (kopírovatelný)
\begin{quote}
\begin{PromptBlock}


Jsi asistent pro zpětnou vazbu. Analyzuj studentský postup a: 1) vypiš maximálně 3 konkrétní chyby nebo nejasnosti; 2) pro každou chybu navrhni krátké cvičení nebo kontrolní otázku (max. 2 věty); 3) pokud je postup neúplný, požádej o chybějící krok. Odpověď stručně označ tagy: [Chyba], [Cvičení], [Požadavek].
\end{PromptBlock}


\end{quote}
Tento typ promtu lze nasadit jako součást předběžného „pre‑gradingu“ a zapojit do pravidla eskalace pro lidskou validaci. Vzorem podobného workflow je právě nástroj „Pokrok v matematice“, který kombinuje generování úloh, sběr postupu a předopravování AI, přičemž učitel finálně dokončí hodnocení \cite{kb_ai_101_tipu_pdf_04e4a115_page_8_1}.

Doporučení pro pilotní nasazení
\begin{itemize}
  \item Začněte na malém vzorku (1–2 semináře) a testujte spolehlivost AI‑feedbacku vs. lidského hodnocení.
  \item Měřte dopad na časovou náročnost vyučujících a spokojenost studentů (krátký dotazník po pilotu).
  \item Dokumentujte pravidla, která využití AI omezují nebo rozšiřují (transparency v sylabu) a informujte studenty, jakým způsobem je AI používána.
\end{itemize}

Krátké shrnutí
\begin{itemize}
  \item Nástroje v ekosystému Microsoft (Pokrok v matematice, OneNote, Copilot) již nabízejí konkrétní funkce, které usnadňují sběr postupu a automatickou předkontrolu, čímž umožňují škálování formativní zpětné vazby \cite{kb_ai_101_tipu_pdf_04e4a115_page_8_1, kb_ai_101_tipu_pdf_04e4a115_page_3_1}.
  \item Klíčové je navrhnout úlohy tak, aby AI mohla smysluplně hodnotit proces (nikoli jen výsledek), a vždy mít definované pravidlo pro lidskou validaci tam, kde je to potřeba (general background knowledge).
\end{itemize}%

%
\section{Implementační checklist a právně‑etické požadavky}%
Kontrolní seznam pro pilotní nasazení interaktivních AI nástrojů: souhlas a transparentnost, GDPR a ochrana osobních dat, akademická integrita, bezpečnostní a provozní opatření.%


%
Tato kontrolní sada slouží jako praktický checklist pro pilotní nasazení interaktivních AI nástrojů ve výuce (GPT‑asistenti, live‑bots, Copilot apod.). U každého bodu je stručně uvedeno, kdo rozhoduje (vyučující / katedra / IT) a zda jde o právně‑etické riziko nebo provozní krok.

\begin{enumerate}
  \item Vymezení rozsahu a zodpovědností
    \begin{itemize}
      \item Určete kurzy/aktivity a typ použití (Q\\ \&A, live demo, auto‑feedback). Odpovědnost: garant + katedra.
      \item Jmenujte kontaktní osobu pro incidenty a eskalaci (IT, GDPR, garant).
    \end{itemize}

  \item Transparentnost a souhlas
    \begin{itemize}
      \item Zaznamenejte do sylabu, kde a jaký nástroj bude použit a jaké typy dat se zpracovávají.
      \item Zajistěte, aby studenti věděli, kdy odpovídá AI a kdy člověk; zvažte informovaný souhlas, pokud se zpracovávají osobní údaje.
    \end{itemize}

  \item Ochrana osobních údajů a DPIA
    \begin{itemize}
      \item Mapujte datové toky: co jde externě vs. co zůstává lokálně; minimalizujte data (pseudonymizace/anonymizace).
      \item Zvažte DPIA při zpracování citlivých nebo identifikovatelných artefaktů.
    \end{itemize}

  \item Smluvní a licenční požadavky
    \begin{itemize}
      \item Ověřte školní/enterprise licence a existenci DPA; některé edukační funkce jsou v Microsoft 365 A3/A5 \cite{kb_ai_101_tipu_pdf_04e4a115_page_15_2}.
      \item Stanovte jasná pravidla retence a zpracování dat u poskytovatele.
    \end{itemize}

  \item Technické a bezpečnostní kontroly
    \begin{itemize}
      \item Zajistěte integraci se školním SSO, řízení přístupu a auditní logy.
      \item Definujte pravidla pro nahrávání médií a ověřte, jak nástroj s nimi pracuje (např. AI funkce v editoru fotek, Copilot) \cite{kb_ai_101_tipu_pdf_04e4a115_page_53_1}.
      \item Nastavte pravidla pro uchovávání a mazání studentských artefaktů.
    \end{itemize}

  \item Akademická integrita
    \begin{itemize}
      \item Formulujte jasná pravidla povoleného a zakázaného využití AI (např. povinná deklarace použití).
      \item Zvažte sběr postupů studentů (fotky, meziodevzdání, commit historie) pro ověření individuálního přínosu a audit výsledků \cite{kb_ai_101_tipu_pdf_04e4a115_page_8_1}.
    \end{itemize}

  \item Eskalace a lidská validace
    \begin{itemize}
      \item Stanovte kritéria pro přepojení na lidskou validaci (nejasné, rozporné nebo nízká jistota).
      \item Zavádějte pravidelné spot‑checky kvality a dokumentujte proces validace.
    \end{itemize}

  \item Provozní připravenost a monitoring
    \begin{itemize}
      \item Měřte metriky pilotu: spolehlivost, poměr eskalací, doba řešení incidentů a spokojenost uživatelů.
      \item Mějte postup pro rollback a záložní lidskou podporu při výpadku služby.
    \end{itemize}

  \item Edukace a komunikace
    \begin{itemize}
      \item Proškolte vyučující a asistenty o fungování nástroje, jeho limitech a interpretaci výstupů; ověřte konkrétní integrace (např. Copilot ve Word/OneNote/Forms) před použitím \cite{kb_ai_101_tipu_pdf_04e4a115_page_8_1}.
      \item Informujte studenty o očekávané úrovni lidské kontroly a o postupu pro stížnosti.
    \end{itemize}

  \item Rychlý kontrolní seznam před spuštěním
    \begin{itemize}
      \item Je jasně definován účel a rozsah použití? Existuje smlouva/DPA nebo školní licence pro požadované funkce?
      \item Jsou mapovány datové toky a provedeny kroky minimalizace dat? Jsou studenti informováni a je nastaven postup pro incidenty/stížnosti?
      \item Jsou nastaveny metriky kvality, pravidla eskalace a záložní mechanismy?
    \end{itemize}
\end{enumerate}

Poznámka: Některé body odkazují na obecné principy (interní postupy). Konkrétní technické a právní kroky konzultujte s interními týmy VUT (IT, právní odd., GDPR/PRIV) a zaznamenejte v dokumentaci pilotu.%

%
\section{Ověřené případové studie a inspirace}%
Krátké přehledy ověřených příkladů z praxe (inspirováno kurzy typu CS50 a dalšími univerzitními nasazeními) s konkrétními poučeními a návrhy, co adaptovat pro VUT.%


%
Krátké, ověřené příklady z praxe slouží jako inspirace pro pilotní nasazení na VUT; níže jsou tři konkrétní příklady‑vzory s praktickými kroky a doporučením, co lze adaptovat lokálně. Některé širší příklady institucionálních nasazení (např. Georgia Tech, Harvard CS50, Arizona State) jsou užitečnou inspirací, jejich popisy zde uvádíme jako obecné pozadí (general background knowledge).

\medskip

\textbf{1) Minecraft Education — výuka principů AI hrou.}  
Minecraft Education nabízí připravené světy a aktivity vhodné pro demonstraci principů strojového učení a pro „Hour of Code“ aktivity; některé z těchto hodin jsou dostupné zdarma, další jsou součástí školních licencí Microsoft 365 A3/A5 \cite{kb_ai_101_tipu_pdf_04e4a115_page_15_2}.  
Praktický pilot (s návrhem kroků):
\begin{enumerate}
  \item Stažení a výběr vhodného světa v Minecraft Education (např. hodina zaměřená na AI nebo Hour of Code) \cite{kb_ai_101_tipu_pdf_04e4a115_page_15_2}.
  \item Příprava krátkého doprovodného materiálu pro studenty: základní pojmy (co je model, trénink, data) a očekávané výstupy z aktivity (diskuse, krátká reflexe).
  \item Vedení hodiny: praktická aktivita ve světě, následná diskuse o možnostech a omezeních AI a závěrečné shrnutí.
\end{enumerate}
Krátké doporučení pro VUT (general background knowledge): vyzkoušet aktivitu v dílčím semináři nebo laboratoři a vyhodnotit pedagogický přínos před širším zavedením; zohlednit potřebu licence Microsoft 365 pro úplnou sadu obsahů.

\medskip

\textbf{2) Automatická tvorba testů: Microsoft 365 Copilot + Forms.}  
Microsoft 365 Copilot ve spojení s Forms umí vygenerovat kompletní test (otázky typu ABC, pravda/nepravda, otevřené), včetně označení správných odpovědí a vysvětlení, proč je odpověď správná \cite{kb_ai_101_tipu_pdf_04e4a115_page_36_1}.  
Praktický pilot (s návrhem kroků):
\begin{enumerate}
  \item Definujte učební cíl a rozsah testu (např. konkrétní kapitola nebo tematický blok).
  \item Vytvořte prompt pro Copilot: počet otázek, typy položek, úroveň (studenti 1.–2. ročníku apod.) a požadavek na klíč a krátké vysvětlení správných odpovědí \cite{kb_ai_101_tipu_pdf_04e4a115_page_36_1}.
  \item Zkontrolujte a upravte vygenerované otázky ručně (ověření faktické přesnosti a vhodnosti pro hodnocení).
  \item Publikujte přes Forms a po odevzdání využijte automatické značení tam, kde je to možné.
\end{enumerate}
Praktické poznámky pro VUT (general background knowledge): Copilot‑workflow zrychlí přípravu písemek a banky úloh; vždy však požadujte lidskou finalizaci otázek a kontrolu vhodnosti pro daný kurs.

\medskip

\textbf{3) Podpora procvičování v matematice — „Pokrok v matematice“ a OneNote.}  
Funkce „Pokrok v matematice“ umožňuje pedagogovi generovat příklady z vybrané oblasti, sbírat odpovědi a (volitelně) postupy žáků; nástroj poskytuje předběžné opravy a diagnostiku oblastí, kde studenti chybují. OneNote má podporu převodu rukopisu na digitální rovnice včetně výpočtu a zobrazení postupu \cite{kb_ai_101_tipu_pdf_04e4a115_page_8_1}.  
Praktický pilot (s návrhem kroků):
\begin{enumerate}
  \item Vyberte téma (např. lineární rovnice) a nastavte v nástroji generování s požadavkem na nahrání postupu řešení od studentů \cite{kb_ai_101_tipu_pdf_04e4a115_page_8_1}.
  \item Pilotně povolte nahrávání fotografií postupu nebo použití OneNote pro převod rukopisu; využijte automatické předopravování jako pomoc učiteli při předkontrole \cite{kb_ai_101_tipu_pdf_04e4a115_page_8_1}.
  \item Definujte pravidla eskalace (kdy poslat práci k lidské revizi) a provádějte spot‑checky kvality výsledků.
\end{enumerate}
Doporučení pro VUT (general background knowledge): integrace do seminářů a cvičení v laboratořích s technickým zaměřením; zaměřit se na využití výstupů AI především jako podpory formativní zpětné vazby.

\medskip

Závěrem: výše uvedené příklady ukazují praktické, nízkonákladové piloty, které lze replikovat na úrovni předmětů nebo laboratoří. Inspirace z větších institucionálních případů (CS50, campus‑licence ChatGPT Enterprise apod.) je užitečná pro strategické plánování, nicméně konkrétní kroky a smluvní ujednání je nutné konzultovat s interními IT a GDPR týmy na VUT (general background knowledge).%

%
\chapter{HODNOCENÍ V ÉŘE AI}%
Kapitola řeší redesign hodnocení v době dostupné AI. Představuje šest principů AI‑aware assessment (autenticita, procesní hodnocení, personalizace, multimodalita, kolaborace, metakognice) a katalog úloh odolných vůči triviálnímu využití AI (ústní zkoušky, in‑class writings, reflexivní portfolia, procesní a lokálně specifické úkoly, AI‑enhanced assignments). Nabízí rubriky, které zahrnují transparentnost použití AI, šablony milestone‑based hodnocení, a doporučení, proč nepoužívat detektory AI. Uvádí nástroje podporující hodnocení (např. Gradescope pro rychlejší a konzistentnější zpětnou vazbu) a konkrétní implementační tipy pro technické obory.%
\section{Úvod a shrnutí kapitoly}%
Stručné uvedení do výzev a cílů hodnocení v éře dostupné generativní AI; jak přistupovat k redesignu hodnocení na úrovni kurzu a katedry, přehled obsahu kapitoly a doporučený pracovní postup pro čtenáře (garanti předmětů, lektoři, asistenti).%


%
Krátké shrnutí: tato kapitola se věnuje výzvám hodnocení v době snadno dostupné generativní AI a nabízí praktický, krok‑za‑krokem přístup k redesignu zadání a hodnoticích postupů na úrovni předmětu i katedry. Cílem je pomoci garantům předmětů, lektorům a asistentům navrhnout hodnocení, které je pedagogicky smysluplné, transparentní vůči studentům a odolné vůči triviálnímu zneužití automatizovaných nástrojů.

Krátké opodstatnění přístupu:
- Namísto snahy „zakázat“ nástroje doporučuje materiál upřednostnit redesign úloh a pilotní nasazení tak, aby hodnocení měřilo procesy, dovednosti a reflektivní metody, nikoli pouze konečný text či kód (general background knowledge).  
- Pro pilotování a plánování strategie lze efektivně využít moderní enterprise agenty (např. Microsoft 365 Copilot) k rychlé tvorbě návrhů, plánů a analýz potřebných kroků; zkušenosti ukazují, že specializované agenty („Výzkumník“) dokážou výrazně zrychlit a zkvalitnit přípravu strategických dokumentů \cite{kb_ai_101_tipu_pdf_04e4a115_page_48_1}.  
- Pilotní nasazení v malém měřítku (seminář, laboratorní cvičení) umožní ověřit pedagogický přínos a doladit postupy před širším rozšířením — analogicky k doporučenému postupu při zavádění aktivit jako Minecraft Education nebo nástrojů pro procvičování matematiky \cite{kb_ai_101_tipu_pdf_04e4a115_page_15_2, kb_ai_101_tipu_pdf_04e4a115_page_8_1}.

Doporučený pracovní postup pro garanty a lektory (rychlý checklist)
\begin{enumerate}
  \item Zmapujte stávající hodnocení: identifikujte typy úloh (výsledkové vs. procesní), body rizika z hlediska automatizovaných nástrojů a které artefakty jsou prokazatelně unikátní (repo‑specifické, lokální data apod.). (general background knowledge)
  \item Rozhodněte o cílech redesignu: které kompetence mají být měřeny (proces, metakognice, aplikace), kde je vhodné povolit AI s povinnou dokumentací a kde je třeba časově/organizačně kontrolované prostředí. (general background knowledge)
  \item Navrhněte a spusťte malý pilot: vyberte jeden modul nebo seminář, připravte nové zadání/rubriky a metodu sběru důkazů (milníky, drafty, portfolia). Vyhodnoťte pilot a sepište lessons‑learned; inspiraci pro pilotní kroky najdete v praktických příkladech (Minecraft Education, nástroje pro „pokrok v matematice“) \cite{kb_ai_101_tipu_pdf_04e4a115_page_15_2, kb_ai_101_tipu_pdf_04e4a115_page_8_1}.
  \item Podpořte práci nástroji pro tvorbu a plánování: při přípravě sylabů, rubrik a checklistů lze využít Copilot‑typ agenta pro generování návrhů a konsolidaci podkladů; vždy ověřte generované výstupy lidskou kontrolou \cite{kb_ai_101_tipu_pdf_04e4a115_page_48_1}.
  \item Formalizujte změny do sylabů a komunikačních šablon: doplňte pole pro deklaraci použití AI, popište pravidla eskalace a spot‑checkové mechanismy pro kontrolu kvality. (general background knowledge)
\end{enumerate}

Co v kapitole najdete (stručný přehled)
- Šest principů AI‑aware assessment a jejich pedagogické důsledky;  
- Katalog úloh odolných vůči triviálnímu využití AI (ústní zkoušky, time‑boxed writings, reflexivní portfolia apod.);  
- Praktické rubriky, milestone‑based hodnocení a šablony pro transparentní deklaraci použití AI;  
- Nástroje pro formativní zpětnou vazbu a tipy pro technické obory;  
- Etické a GDPR aspekty hodnocení, důraz na dokumentaci zpracování a DPIA;  
- Šablony, checklisty a prompt‑knihovna pro rychlé nasazení.

Krátké doporučení na závěr: začněte malým pilotem, využijte dostupné nástroje pro rychlé návrhy a plánování, a výsledky pilotu převeďte do jasně komunikovaných změn v sylabu a rubrikách — tento cyklus „navrhni‑pilotu‑ověř‑škáluj“ je prioritou při adaptaci hodnocení na VUT. \cite{kb_ai_101_tipu_pdf_04e4a115_page_48_1, kb_ai_101_tipu_pdf_04e4a115_page_15_2, kb_ai_101_tipu_pdf_04e4a115_page_8_1}%

%
\section{Šest principů AI‑aware assessment}%
Přehled a krátké definice šesti principů: autenticita (authenticity), procesní hodnocení, personalizace, multimodalita, podpora kolaborace a metakognice; vysvětlení, proč jsou principy základem redesignu úloh a kritérií hodnocení.%
\subsection{Autenticita}%
Krátká definice: úlohy zakotvené v reálných kontextech, které mají smysl mimo zkoušku. Role pro redesign: přehodnocení cílů a kritérií tak, aby úloha měřila reálné schopnosti a motivovala studenty. (general background knowledge)%


%
\paragraph{Definice (stručně)} 
Autenticita: úloha v reálném kontextu měřící schopnost aplikovat znalosti v praktických situacích.

\paragraph{Proč přehodnotit zadání směrem k autenticity}
Autentická zadání motivují studenty a snižují možnost „vyřešení“ jen pomocí generativní AI, protože úspěch závisí na kontextu, procesu a komunikaci výsledků.

\paragraph{Postup redesignu úlohy (praktický návod)}
\begin{enumerate}
  \item Vymezte reálný kontext a adresáta; navrhněte artefakty a milníky (drafty, protokoly, reporty).
  \item Zahrňte rubriku hodnotící proces, relevanci, komunikaci výsledků a reflexi o použití AI (deklarace).
  \item Uskutečnějte pilot v jednom kurzu, ověřte fungování a zapracujte lessons‑learned; při plánování pilotu lze využít Copilot‑typ agenty — vždy s lidskou kontrolou \cite{kb_ai_101_tipu_pdf_04e4a115_page_48_1}.
\end{enumerate}

\paragraph{Krátký příklad} Zadání: analyzujte spotřebu energie labu, navrhněte dvě opatření; artefakty: metodika (milník 1), předběžná analýza (milník 2), finální report + prezentace + reflexe o AI; rubrika: relevance (30\%), kvalita analýzy (30\%), komunikace (20\%), reflexe/AI (20\%).

\paragraph{Rychlý checklist pro vyučující}
\begin{itemize}
  \item Je úloha vázána na adresáta? \item Jsou definovány milníky a artefakty? \item Obsahuje rubrika procesu a reflexi AI? \item Byl naplánován pilot a mechanismus spot‑checků? Pokud ano, zvažte nástroje pro plánování pilotu \cite{kb_ai_101_tipu_pdf_04e4a115_page_48_1}.
\end{itemize}%

%
\subsection{Procesní hodnocení}%
Definice a implikace: zaměření na postupy, kroky a odůvodnění řešení místo pouze konečného výsledku; kritéria hodnotí kvalitu procesu. Praktický důraz v materiálu: u některých nástrojů (např. 'Pokrok v matematice') je vyžadován nahraný postup a tento postup je považován za důležitější než samotný výsledek.%


%
Procesní hodnocení klade důraz na zdokumentovaný postup, kroky a odůvodnění místo jen konečného výsledku; některé nástroje (např. „Pokrok v matematice“) podporují požadavek nahraného postupu \cite{kb_ai_101_tipu_pdf_04e4a115_page_8_1}.

Důsledky pro návrh zadání a rubriky:
\begin{itemize}
  \item Požadujte jasné mezikroky a zdůvodnění (sken/text/video); pokud možno považujte nahraný postup za povinný artefakt \cite{kb_ai_101_tipu_pdf_04e4a115_page_8_1}.
  \item V rubrice uveďte váhy pro kvalitu postupu, správnost metody a reflexi chyb (např. postup 50\%, výsledek 20\%, reflexe 30\%).
  \item Naplánujte milníky (inkrementální odevzdání) a spot‑checky pro formativní zpětnou vazbu.
\end{itemize}

Praktický příklad (matematika):
\begin{enumerate}\item Zadání: řešte modelový problém a nahrajte pracovní postup (fotografie rukopisu nebo export z OneNote) s komentáři k rozhodnutím.\item Hodnocení: nástroj označí pravděpodobné chyby; učitel provede spot‑check 5–10\,\% prací a finálně schválí bodové ohodnocení \cite{kb_ai_101_tipu_pdf_04e4a115_page_8_1}.\end{enumerate}

Rychlý checklist pro vyučující:
\begin{itemize}
  \item Je postup požadován jako samostatný artefakt? Obsahuje rubrika váhy pro procesní kritéria? Jsou naplánovány milníky a spot‑checky?
\end{itemize}%

%
\subsection{Personalizace}%
Definice a implikace: přizpůsobení obtížnosti a obsahu jednotlivým studentům a použité podmínky hodnocení; redesign úloh vede k adaptivním variantám a diferenciovaným kritériím. Podpůrný příklad z materiálu: nástroje ukazují, jaké typy příkladů studentovi jdou a jaké ne, a umožňují nastavovat obtížnost zadání.%


%
Personalizace znamená přizpůsobení obtížnosti a obsahu úloh výkonu studenta; AI sleduje, „jaký typ příkladů studentovi jde a jaký ne“, upravuje výběr/obtížnost a může identifikovat problematická slova a generovat cvičení \cite{kb_ai_101_tipu_pdf_04e4a115_page_8_1, kb_ai_101_tipu_pdf_04e4a115_page_7_1}.

Praktický postup pro pedagoga: \begin{enumerate}
\item Vyberte oblast/kompetenci pro personalizaci (např. algebra).
\item Nastavte parametry nástroje (rozsah obtížnosti, požadavek na postup) — nástroj vygeneruje příklady a úrovně \cite{kb_ai_101_tipu_pdf_04e4a115_page_8_1}.
\item Po odevzdání využijte přehled výkonu k cílenému opakování nebo krátkým intervenčním cvičením na detekované slabiny \cite{kb_ai_101_tipu_pdf_04e4a115_page_7_1}.
\end{enumerate}

Důsledky pro hodnocení: \begin{itemize}
\item Redesign úloh umožní varianty zadání a diferencované hodnoticí kritérium podle úrovní.
\item Požadování artefaktů procesu (postupy, skeny) pomáhá validovat autenticitu pokroku a slouží jako vstup pro další personalizaci \cite{kb_ai_101_tipu_pdf_04e4a115_page_8_1}.
\end{itemize}%

%
\subsection{Multimodalita}%
Definice a implikace: zapojení více modalit (text, obraz, rukopis, interaktivní simulace) do zadání i sběru důkazů u hodnocení. V materiálu ilustrováno např. konverzí rukopisu na digitální rovnice v OneNote a výukovými světy (Minecraft Education) pro praktické aktivity.%


%
Multimodalita znamená zapojení více modalit (text, obraz, rukopis, interaktivní simulace) do zadání i sběru důkazů; umožňuje požadovat artefakty procesu jako doplněk k textovým odpovědím a tím zvyšuje validitu hodnocení \cite{kb_ai_101_tipu_pdf_04e4a115_page_8_1}.
Praktické postupy pro učitele: \begin{itemize}
\item Rukopis→OneNote: studenti píší rukou, pak použijte \textit{Kreslení}→\textit{Rukopis na text} pro strojově čitelný záznam \cite{kb_ai_101_tipu_pdf_04e4a115_page_13_2}.
\item Interaktivní světy (Minecraft Education): použijte připravené světy/lekce z Hour of Code a reflektujte učení diskusí \cite{kb_ai_101_tipu_pdf_04e4a115_page_15_2}.
\item Kombinovaná zadání: požadujte text + screenshot/export z OneNote jako artefakt procesu pro ověření autenticity \cite{kb_ai_101_tipu_pdf_04e4a115_page_8_1, kb_ai_101_tipu_pdf_04e4a115_page_13_2}.
\end{itemize}
Krátký checklist pro návrh multimodálních úloh: \begin{enumerate}
\item Určete modality; \item Dejte technické instrukce (export OneNote, screenshot, nahrání Minecraft artefaktů) \cite{kb_ai_101_tipu_pdf_04e4a115_page_13_2, kb_ai_101_tipu_pdf_04e4a115_page_15_2}; \item Požadujte artefakt procesu + popis workflow \cite{kb_ai_101_tipu_pdf_04e4a115_page_8_1}.
\end{enumerate}%

%
\subsection{Podpora kolaborace}%
Krátká definice: úlohy navržené tak, aby studenty zapojovaly do společné práce; implikace: hodnotit příspěvek, komunikaci a společné řešení úloh. (general background knowledge)%


%
% (bez nadpisu — sekce: Podpora kolaborace)
Krátká definice: úlohy pro společnou práci; hodnocení zahrnuje artefakt, individuální přínos, komunikaci a společné řešení problémů.

Praktická doporučení a checklist:
\begin{itemize}
  \item Definujte role a kritéria pro individuální přínos; zahrňte sebehodnocení, záznamy rozhodnutí a peer review (general background knowledge).
  \item Požadujte sdílené procesní artefakty (OneNote, export rukopisů) \cite{kb_ai_101_tipu_pdf_04e4a115_page_8_1} a využijte Minecraft Education pro týmové simulace \cite{kb_ai_101_tipu_pdf_04e4a115_page_15_2}.
  \item Podpořte kooperaci agenty/Copilotem pro správu rolí, formátování a opakující se úkoly \cite{kb_ai_101_tipu_pdf_04e4a115_page_31_1}.
\end{itemize}

Krátký workflow: \begin{enumerate} \item Rozdělení rolí; vedoucí nahraje průběh do OneNote \cite{kb_ai_101_tipu_pdf_04e4a115_page_8_1}. \item Testy v připraveném světě a export klíčových artefaktů \cite{kb_ai_101_tipu_pdf_04e4a115_page_15_2}. \item Každý vyplní krátkou reflexi; učitel hodnotí artefakt i proces podle rubriky (general background knowledge). \end{enumerate}%

%
\subsection{Metakognice}%
Definice a implikace: podpora sebereflexe, vysvětlování vlastního myšlení a záznamu rozhodnutí; redesign kritérií zahrnuje sebereflexní položky a hodnocení schopnosti zdůvodnit postup. Materiál zdůrazňuje význam dokumentace postupu při automatizované či asistované kontrole úloh.%


%
Krátká definice a účel: metakognice — student popíše vlastní myšlení, odůvodní postup a reflektuje učení; v hodnocení zařaďte pole pro sebereflexi.

Praktická doporučení:
\begin{itemize}
  \item Požadujte řešení + procesní artefakt (fotky/OneNote; OneNote převádí rukopisné rovnice) \cite{kb_ai_101_tipu_pdf_04e4a115_page_8_1}.
  \item Použijte automatizovanou předkontrolu (AI označí pravděpodobné chyby; při nízké jistotě eskalujte k učiteli).
  \item Vložte krátké metakognitivní otázky do zadání (např. „Popište 2 klíčová rozhodnutí a proč“).
\end{itemize}

Krátký workflow:
\begin{enumerate}
  \item Odevzdání řešení + artefakt.
  \item Předkontrola → při nejasnostech učitel.
  \item Student doplní krátkou reflexi; učitel provede spot‑check (5–10\%).
\end{enumerate}

Instrukce pro studenty: 
\begin{PromptBlock}


Nahrát: 1) výsledek, 2) krok-za-krokem (fotky/OneNote), 3) reflexe (max. 200 slov): co bylo klíčové? Kde nejistota?
\end{PromptBlock}

. Pozn.: Zvažte MS 365 Copilot vzhledem k GDPR \cite{kb_ai_101_tipu_pdf_04e4a115_page_3_1}.%

%
\section{Katalog úloh odolných vůči triviálnímu využití AI}%
Sady konkrétních typů úloh s praktickými varian- tami a příklady: ústní zkoušky a pohovory, in‑class writings a time‑boxed psaní, reflexivní portfolia, procesně orientované a lokálně specifické úkoly, AI‑enhanced assignments (úkoly využívající AI jako nástroj s povinnou reflexí).%
\subsection{Ústní zkoušky a pohovory}%
Postupy pro navržení ústních zkoušek a pohovorů odolných vůči triviálnímu využití AI: strukturované otázky, živé dotazování na proces myšlení, scénáře reagující na studentovy odpovědi a ověření porozumění místo pouhého výsledku.%


%
Krátké shrnutí: postupy pro návrh ústních zkoušek zaměřené na ověření porozumění a procesu myšlení, nikoli jen výsledku.

Praktický postup pro návrh a vedení ústní zkoušky
\begin{enumerate}
  \item Příprava otázek: krátké „strukturální“ úkoly s 2–3 doplňujícími variantami; pro každou otázku definujte očekávané kroky řešení.
  \item Požadavek na procesní artefakty: student předloží krátké zápisky/skici/OneNote pro ověření mezikroků \cite{kb_ai_101_tipu_pdf_04e4a115_page_8_1}.
  \item Průběh zkoušky: warm‑up, hlavní úloha, učitel volí 1–3 doplňující otázky zaměřené na odůvodnění a adaptace; kontrola konkrétních mezikroků.
  \item Ověření autenticity: preferujte otázky o rozhodnutích, alternativách a krátkých adaptacích; zaznamenávejte klíčové body (GDPR‑kompatibilně).
\end{enumerate}

Šablony doplňujících otázek (vyberte 1–3): „Proč jste zvolil/a tento postup? Jaké byly alternativy?“, „Jak upravit řešení, pokud se X změní na Y?“, „Který krok je nejrizikovější a proč?“

Krátká rubrika: proces a logika (0–3); reakce/adaptace (0–3); kvalita artefaktů (0–2); metakognice (0–2).

Doporučení: předem informujte studenty o pravidlech a povinných artefaktech; využijte digitální nástroje pro sběr (OneNote usnadňuje čitelnost) \cite{kb_ai_101_tipu_pdf_04e4a115_page_8_1}. Poznámka: respektujte fakultní pravidla, GDPR a interní politiky VUT.%

%
\subsection{In‑class writings a time‑boxed psaní}%
Krátké, časově omezené písemné úlohy psané v rámci výuky (dozorované nebo digitálně kontrolované), zaměřené na proces tvorby, rychlé opakované verze a artefakty z jejich práce (drafty, poznámky).%


%
Krátké, časově omezené písemné úlohy (time‑boxed writings) sledují pracovní postup studenta — rychlý koncept, minimální draft a krátká metareflexe místo finálního produktu.

Praktický postup pro nasazení
\begin{enumerate}
\item Příprava: úkol 20–40 min; požadujte mezikroky a artefakty (náčrtek, výpočty, krátká reflexe).
\item Technické odevzdání: sken/fotka nebo upload do OneNote/LMS; OneNote usnadňuje převod rukopisu a sběr postupů \cite{kb_ai_101_tipu_pdf_04e4a115_page_8_1}.
\item Vedení: zopakujte pravidla (včetně AI), sledujte průběh; na konci 1–3vět. metareflexi; artefakty: první draft, opravený náčrt, reflexe.
\item Hodnocení: zaměřte se na proces — srozumitelnost mezikroků, logiku, identifikaci chyb a kvalitu reflexe; použijte jednoduchou rubriku.
\end{enumerate}

Praktické tipy a šablona
\begin{itemize}
\item Před hodinou sdělte povolené pomůcky a formát; OneNote zjednodušuje práci s rukopisem a převodem rovnic \cite{kb_ai_101_tipu_pdf_04e4a115_page_8_1}. Zvažte krátké opakované verze pro sledování procesu.
\item Šablona: „30 min — napište postup k úloze X; odevzdejte draft, upravte hlavní krok a přiložte 2–3vět. reflexi.“
\end{itemize}

Poznámka k ověřování a integritě: Požadujte artefakty dokazující proces (drafty, fotky, časová razítka); OneNote/LMS pomáhají se sběrem a ověřením \cite{kb_ai_101_tipu_pdf_04e4a115_page_8_1}.%

%
\subsection{Reflexivní portfolia a metareflexe}%
Sady úkolů založené na portfoliích, kde studenti předkládají výběr prací doplněný povinnou reflexí o použití AI, záznamy o pracovním postupu a sebehodnocení; obsahuje návrh hodnoticích kritérií pro reflexi.%


%
Reflexivní portfolia kladou důraz na výběr 3--5 reprezentativních prací doplněných o metareflexi a dokumentaci postupu; cílem je hodnotit schopnost vysvětlit postup, rozhodnutí a případné použití AI.

Praktický postup: 
\begin{enumerate}
\item Návrh: povinné komponenty (3--5 artefaktů), dokumentace postupu (drafty, mezikroky, screenshoty, časová razítka) a metareflexe.
\item Technické odevzdání: sběr v prostředí zachovávajícím mezikroky (např. OneNote) nebo upload do LMS; OneNote usnadňuje převod rukopisu a sběr postupů \cite{kb_ai_101_tipu_pdf_04e4a115_page_8_1}.
\item Instrukce a kontrola: přiložte šablonu metareflexe, pravidla pro deklaraci AI (typ, prompt, úpravy) a použijte spot‑check pro ověření autenticity.
\end{enumerate}

Minimum portfolia: 3--5 prací; u každé první draft, revidovaná verze/poznámky a chronologický záznam; povinná metareflexe a explicitní deklarace AI.

Rubrika (procesně orientovaná): dokumentace postupu 35\,\% \cite{kb_ai_101_tipu_pdf_04e4a115_page_8_1}, kvalita metareflexe 35\,\%, učební pokrok 20\,\%, formální náležitosti 10\,\%.

Šablona metareflexe: 
\begin{enumerate}
\item Co jsem dělal(a) a proč byl artefakt vybrán? (1--2 odstavce) \item Které AI nástroje jsem použil(a) a jak (prompt/funkce)? Jak jsem ověřil(a) výstupy?
\item Co jsem se naučil(a) mezi verzemi — hlavní zlepšení a nezodpovězené otázky. \item Krátký plán dalšího kroku (1--2 věty).
\end{enumerate}

Poznámka: požadujte artefakty dokazující proces (drafty, screenshoty, časová razítka); OneNote a LMS usnadní sběr a ověření \cite{kb_ai_101_tipu_pdf_04e4a115_page_8_1}.%

%
\subsection{Procesně orientované a lokálně specifické úkoly}%
Úkoly vázané na lokální data, kontext nebo fáze procesu (inkrementální odevzdání, laboratorní zápisníky, analytické kroky), které kladou důraz na sledovatelný postup a jedinečnost řešení.%


%
Procesně orientované a lokálně specifické úkoly kladou důraz na sledovatelný postup (inkrementální odevzdání, mezi‑verze, laboratorní zápisníky) a na vazbu úlohy na lokální data; hodnotí se postup i výsledek, čímž se snižuje motivace ke „komoditnímu“ používání generativní AI bez dokumentace.

Praktický workflow pro návrh a provoz úlohy
\begin{enumerate}
  \item Definujte milníky: např. výběr dat, analýza, první draft, revidovaná verze, závěrečná reflexe; u každého milníku požadujte konkrétní artefakty (kód, krátký log rozhodnutí, screenshoty).
  \item Požadavek na sledovatelnost: vyžadujte meziverze, poznámky o postupu a časová razítka (exporty z OneNote nebo LMS). OneNote usnadňuje sběr mezikroků a převod rukopisu na digitální artefakty \cite{kb_ai_101_tipu_pdf_04e4a115_page_8_1}.
  \item Technické sběrné místo a validace: určete platformu/formáty (LMS, OneNote, repo; notebooky, PDF), využijte integrace Copilot/Notepad pro přípravu dílčích textů a provádějte náhodné kontroly procesních artefaktů včetně krátké metareflexe \cite{kb_ai_101_tipu_pdf_04e4a115_page_14_1, kb_ai_101_tipu_pdf_04e4a115_page_8_1}.
\end{enumerate}

Konkrétní požadované artefakty (vzorně)
\begin{itemize}
  \item První draft (kód/postup) + krátký log rozhodnutí (1--3 věty).
  \item Mezikrok: screenshot běhu, chybový výstup nebo ukázka experimentu.
  \item Revidovaný výstup s komentáři, co se změnilo a proč.
  \item Závěrečná metareflexe: použité nástroje, způsob ověření, co by se dělalo jinak (šablona viz kap. Reflexivní portfolia) \cite{kb_ai_101_tipu_pdf_04e4a115_page_8_1}.
\end{itemize}

Krátká rubrika (zkr.): dokumentace postupu a mezikroků — vysoká váha; kvalita metareflexe — vysoká váha; technická kvalita a formální náležitosti. Pozn.: zahrňte lokální data tam, kde to zvýší jedinečnost řešení (snižuje možnost přímého kopírování z veřejných zdrojů).%

%
\subsection{AI‑enhanced assignments a příklady nástrojů}%
Úkoly, které explicitně povolují použití AI jako nástroje za podmínky povinné dokumentace a reflexe o jeho použití; obsahuje krátké šablony a ilustrativní nástroje (např. Copilot v Poznámkovém bloku, Minecraft Education jako edukační příklad z praxe).%


%
Krátké shrnutí a princip: Úlohy „AI‑enhanced“ povolují generativní AI za podmínky povinné dokumentace použitých nástrojů, promptů a krátké reflexe o ověření výstupu.

\begin{enumerate}
\item Uveďte povolené nástroje a podmínky;
\item Přiložte artefakty: prompt/konverzaci, export odpovědi, krátký log a 2–3větá reflexe ověření;
\item Hodnoťte i proces (dokumentace, kvalita reflexe) a provádějte náhodné spot‑checky (např. 5–10\,\%);
\end{enumerate}

\begin{itemize}\item Příklady: Copilot v Notepadu (přepis, shrnutí) \cite{kb_ai_101_tipu_pdf_04e4a115_page_14_1}; Minecraft Education jako edukační prostředí \cite{kb_ai_101_tipu_pdf_04e4a115_page_15_2}; nástroje pro procvičování a sběr postupů (OneNote, „Pokrok v matematice“) \cite{kb_ai_101_tipu_pdf_04e4a115_page_8_1}.\end{itemize}

Šablona deklarace (stručně): použité nástroje a verze; prompt/export; co převzato; jak ověřeno; reflexe (2–3 věty).

Poznámka: uvedené příklady ilustrují možnosti; konkrétní volbu a podmínky schvalují garanti v souladu s etickými a GDPR zásadami.%

%
\section{Rubriky, transparentnost a milestone‑based hodnocení}%
Praktické rubriky, které zahrnují explicitní pole pro deklaraci použití AI; šablony pro hodnocení po milnících (milestone‑based assessment) a příklady, jak propojit průběžné milníky s konečným hodnocením.%


%
Tato sekce shrnuje praktické rubriky a postupy pro transparentní hodnocení v prostředí s možností použití generativní AI. Cílem je hodnotit nejen konečný produkt, ale i proces, ověřování výsledků a reflexi použití AI.

Principy a požadavky
\begin{itemize}
  \item Každé zadání s možností použití AI musí obsahovat pole „Deklarace použití AI“ s povinnými položkami (viz níže).
  \item Hodnocení zahrnuje procesní artefakty (milníky, drafty, logy) vedle výsledku; u vybraných výstupů požadujte stručnou verifikaci.
  \item Formativní podpora nástroji (např. Copilot) může učiteli šetřit čas, vždy však vyžadujte lidské přezkoumání \cite{kb_ai_101_tipu_pdf_04e4a115_page_21_2}.
\end{itemize}

Povinné pole pro deklaraci AI (šablona)
\begin{enumerate}
  \item Použité nástroje a verze.
  \item Prompt nebo export konverzace.
  \item Co bylo převzato a co upraveno.
  \item Jak bylo ověřeno (stručně).
  \item Reflexe: 2–3 věty o vlivu AI na postup a výsledek.
\end{enumerate}

Rubrika — klíčové položky
\begin{description}
  \item[Proces a dokumentace] Kvalita milníků, draftů a logů; sledovatelnost postupu.
  \item[Deklarace AI a verifikace] Kompletní deklarace + stručné doklady ověření.
  \item[Technická kvalita] Správnost řešení, argumentace, metodika.
  \item[Reflexe a etika] Sebereflexe role AI, etické úvahy.
  \item[Originalita a aplikace] Originalita a lokální uplatnění poznatků.
\end{description}

Návrh vah (ilustrační)
\begin{itemize}
  \item Proces a dokumentace: 25\,\%
  \item Deklarace AI a verifikace: 10\,\%
  \item Technická kvalita: 45\,\%
  \item Reflexe a etika: 10\,\%
  \item Originalita a aplikace: 10\,\%
\end{itemize}

Milestone‑based hodnocení — implementační vzor
\begin{enumerate}
  \item Definujte 3–5 milníků s konkrétními výstupy (návrh, prototyp, testy, závěrečná zpráva + reflexe).
  \item U každého milníku určete váhu a požadované artefakty (draft, log, export konverzace).
  \item Požadujte pro každý milník „Deklaraci AI“ a krátkou verifikaci (testy, zdroje, porovnání).
  \item Zaveďte formativní kontrolní body s komentáři učitele/TA a provádějte náhodné spot‑checky (např. 5–10\,\%).
\end{enumerate}

Příklad rozložení milníků (ilustrativně)
\begin{itemize}
  \item A: Návrh a plán verifikace (10\,\%).
  \item B: Prototyp a logy použití (30\,\%).
  \item C: Testy a verifikace výsledků (30\,\%).
  \item D: Závěrečná zpráva + deklarace AI + reflexe (30\,\%).
\end{itemize}

Praktická šablona pro formativní zpětnou vazbu (kopírovatelný prompt)
\begin{verbatim}
Jsi asistent pro zpětnou vazbu. Analyzuj studentský postup a:
1) vypiš max. 3 konkrétní chyby nebo nejasnosti;
2) pro každou chybu navrhni krátké cvičení nebo kontrolní otázku (max. 2 věty);
3) pokud chybí krok, požádej o doplnění.
Označ tagy: [Chyba], [Cvičení], [Požadavek].
\end{verbatim}

Kontrolní seznam pro zavedení rubriky
\begin{enumerate}
  \item Do sylabu přidejte pole „Deklarace AI“ a očekávání ke každému milníku.
  \item Připravte rubriku pro proces, deklaraci AI a verifikaci.
  \item Pilotujte postup na malé skupině a sbírejte zpětnou vazbu od studentů i TA.
  \item Stanovte mechanismus spot‑checků a eskalace podezření na neférové použití.
  \item Dokumentujte pilot a upravte rubriku před plošným nasazením.
\end{enumerate}

Krátké doporučení k transparentnosti
\begin{itemize}
  \item Studentům jasně sdělte, kde je AI povolena a jak bude její použití hodnoceno.
  \item Povinné deklarace a dokumentace promptů/konverzací zvyšují ověřitelnost a podporují zodpovědné používání nástrojů.
\end{itemize}

Poznámka: Hromadné generování formulací formativní zpětné vazby nástroji typu Copilot může ušetřit čas; vždy však vyžadujte lidskou validaci generovaných komentářů \cite{kb_ai_101_tipu_pdf_04e4a115_page_21_2}.%

%
\section{Formativní hodnocení a nástroje pro podporu hodnocení}%
Jak AI nástroje efektivně podporují formativní zpětnou vazbu (např. použití asistentů k hromadnému generování komentářů); přehled podporovaných nástrojů pro rychlejší a konzistentnější zpětnou vazbu (vč. příkladů implementace jako Gradescope a ukázka využití Copilota pro tvorbu textových hodnocení).%


%
% Text těla sekce (bez nadpisu)
Formativní hodnocení s podporou AI — praktický přístup a doporučení pro pedagogy.

Úvodní poznámka: AI může ulehčit hromadné generování textových komentářů a formativní zpětné vazby, pokud je proces správně strukturován a vždy obsahuje lidskou validaci. Praktický princip ilustruje postup, kdy má učitel připravenou jednoduchou tabulku s názvy studentů a základními charakteristikami; tuto tabulku lze použít jako vstupní data pro asistent typu Copilot, který na jejím základě vygeneruje detailní formulace zpětné vazby \cite{kb_ai_101_tipu_pdf_04e4a115_page_21_2}.

Konkrétní workflow (krátce, krok po kroku)
\begin{enumerate}
  \item Sestavte strukturovaný zdroj dat: tabulka (např. CSV/Excel) se sloupci jako Jméno, Silné stránky (body), Oblasti ke zlepšení (body), krátká poznámka učitele (volitelné).
  \item Vytvořte přesnou instrukci pro asistent (prompt) s požadovaným stylem a délkou komentářů (např. „stručně, pozitivně, 2–3 věty“). Viz příklad níže.
  \item Požádejte asistenta o vygenerování komentářů pro všechny záznamy a export výsledku zpět do tabulky.
  \item Proveďte lidskou revizi: opravte faktické chyby, upravte tón a zajistěte konzistenci. AI návrhy jsou východiskem, nikoli hotovým produktem \cite{kb_ai_101_tipu_pdf_04e4a115_page_21_2}.
  \item Distribuujte studentům formativní komentáře spolu s případnými krátkými doporučeními pro další kroky.
\end{enumerate}

Praktický příklad (inspirováno ukázkou v KB)
\begin{quote}
V příkladu učitel vytvoří tabulku s jmény a základními charakteristikami žáků a následně pošle Copilotovi požadavek: „Můžeš mi pro následující žáky vygenerovat detailní formativní hodnocení?“ — výsledky se potom dolaďují podle potřeby \cite{kb_ai_101_tipu_pdf_04e4a115_page_21_2}.
\end{quote}

Ukázkový prompt pro generování zpětné vazby (kopírujte a upravte podle potřeby)
\begin{verbatim}
Jsi asistent pro učitele. Pro každého žáka vytvoř 2-3 věty formativní zpětné vazby:
1) jedna věta shrne silnou stránku;
2) jedna věta navrhne konkrétní krok pro zlepšení;
3) udržuj tón pozitivní a konkrétní.
Formát: Jméno: <komentář>
\end{verbatim}

Nástroje a scénáře (ověřené poznatky z KB)
\begin{itemize}
  \item Použití Copilotu pro hromadné generování textové zpětné vazby je prakticky realizovatelné pomocí strukturovaného vstupu (tabulka) a následné lidské kontroly \cite{kb_ai_101_tipu_pdf_04e4a115_page_21_2}.
  \item Pro předměty zaměřené na procvičování matematiky existují specializované funkce (např. „Pokrok v matematice“), které generují cvičné příklady, umožňují vyžadovat nahrání postupu a poskytují učiteli pohledy na oblasti, kde žák chyboval; tyto nástroje mohou být zdrojem artefaktů pro formativní hodnocení (postupy žáků jako důkaz) \cite{kb_ai_101_tipu_pdf_04e4a115_page_8_1}.
  \item Digitální nástroje jako OneNote umožňují převod rukopisných rovnic na digitální formu včetně zobrazení postupu; to usnadňuje sběr procesních artefaktů u matematických úloh \cite{kb_ai_101_tipu_pdf_04e4a115_page_8_1}.
\end{itemize}

Doporučené zásady bezpečného a efektivního nasazení
\begin{itemize}
  \item Human‑in‑the‑loop: vždy osvojte rutinu lidské revize generovaných komentářů před jejich odesláním studentům \cite{kb_ai_101_tipu_pdf_04e4a115_page_21_2}.
  \item Struktura vstupů: čím přesněji a více strukturovaně vstup připravíte, tím konzistentnější a užitečnější budou výstupy.
  \item Jasné instrukce pro formátování: nastavte očekávaný tón, délku a konkrétní komponenty (silná stránka / návrh zlepšení).
  \item Ochrana soukromí: při sdílení tabulek s údaji studentů dbejte pravidel GDPR a interních pravidel; minimalizujte citlivé údaje.
\end{itemize}

Krátký checklist pro pilot na katedře (praktické kroky)
\begin{enumerate}
  \item Vyzkoušejte workflow na malé skupině (10–30 studentů) a sbírejte zpětnou vazbu učitele i studentů.
  \item Definujte standardizovaný formát vstupní tabulky a šablonu promptu.
  \item Stanovte postupy lidské kontroly a zodpovědnosti (kdo finalizuje komentáře).
  \item Monitorujte čas úspory a kvalitu zpětné vazby (srovnání s manuálními komentáři).
\end{enumerate}

Poznámka o dalších nástrojích (general background knowledge)
\begin{itemize}
  \item Některé platformy pro hromadné hodnocení a anotaci (např. Gradescope) rovněž podporují automatizaci částí hodnocení a lze je kombinovat s AI‑podporou pro generování komentářů; toto je obecné pozorování a záleží na konkrétní integraci a licencích (general background knowledge).
\end{itemize}

Závěrem: AI je v oblasti formativního hodnocení užitečný akcelerátor rutinních textových úkonů, zejména generování individualizovaných komentářů; klíčem k úspěchu je kvalitní strukturovaný vstup, přesné instrukce pro asistenta a povinná lidská validace konečných výstupů \cite{kb_ai_101_tipu_pdf_04e4a115_page_21_2, kb_ai_101_tipu_pdf_04e4a115_page_8_1}.%

%
\section{Implementační tipy pro technické obory}%
Konkrétní doporučení pro laboratoře, programovací úlohy a projektové předměty: nastavení zadání, ověřitelnost práce (lokální data, repo‑specifické artefakty), workflow pro korekci kódu a hodnocení procesu vývoje, integrace s LMS/SSO a enterprise řešení.%


%
Úvodní poznámka: Technické obory (laboratoře, programovací a projektové předměty) kladou zvláštní důraz na ověřitelnost výsledků a sledovatelnost procesu. Níže jsou praktické doporučení a krátké checklisty, které lze přímo adaptovat při redesignu zadání a nastavení hodnocení.

Nastavení zadání — zásady
\begin{itemize}
  \item Požadujte procesní artefakty: mezi‑verze, krátké zápisky, ukázky mezikroků nebo logy spuštění — u matematických/technických úloh lze vyžadovat nahrání postupu řešení (step‑by‑step) jako podmínku odevzdání \cite{kb_ai_101_tipu_pdf_04e4a115_page_8_1}.
  \item Časové a místní omezení: tam, kde je to pedagogicky oprávněné, volte časově omeřené in‑class či supervised úlohy kombinované s následnou reflexí studenta (general background knowledge).
  \item Jasné metriky pro hodnocení procesu: definujte milestone‑based artefakty (např. návrh, prototyp, testy, závěrečná zpráva) a přiřaďte váhy v rubrice (general background knowledge).
\end{itemize}

Ověřitelnost práce (praktická doporučení)
\begin{itemize}
  \item U programovacích úloh preferujte odevzdání v repozitářích nebo prostředí, která zachovávají historii změn (commit log, issue/ticket artefakty) — tyto artefakty zvyšují důvěru v originalitu řešení a usnadňují hodnocení procesu (general background knowledge).
  \item Pro numerické a experimentální úlohy upřednostněte použití lokálních dat nebo repozitářově identifikovatelných datasetů/konfigurací, aby výsledky byly reprodukovatelné (general background knowledge).
  \item Pokud je v kurzu využito automatizované generování úloh či predikce chyb (např. platformy generující cvičné příklady), nastavte povinnost nahrání postupu tak, aby byl kladen důraz na proces nad pouhým výsledkem \cite{kb_ai_101_tipu_pdf_04e4a115_page_8_1}.
\end{itemize}

Workflow pro korekci kódu a hodnocení procesu
\begin{enumerate}
  \item Automatizované testy a statická analýza: nakonfigurujte CI pipeline, která spustí základní testy, linting a metriky pokrytí (general background knowledge).
  \item AI jako před‑filter: použijte asistenta k vygenerování návrhů komentářů k běžným chybám nebo k shrnutí výstupu testů; finální rozhodnutí a hodnocení vždy provádí člověk (princip „teacher finalizes grading“ platí i pro generativní nástroje) \cite{kb_ai_101_tipu_pdf_04e4a115_page_8_1}.
  \item Lidská validace a škálování: nastavte pravidla, kdy lidský hodnotitel provede plnou revizi (náhodné spot‑checks, projekty nad určitou hranicí složitosti apod.) — tím se kombinuje efektivita automatizace s odpovědností učitele (general background knowledge).
  \item Dokumentace procesu v hodnocení: v rubrice měřte kvalitu procesu (dostatečnost mezikroků, testy, komentáře v kódu), schopnost odůvodnit rozhodnutí a metakognici (general background knowledge).
\end{enumerate}

Integrace s LMS, SSO a enterprise nástroji — praktické poznámky
\begin{itemize}
  \item Pokud využíváte nástroje, které generují úlohy či předpřipravené materiály (např. některé vzdělávací akcelerátory pro matematiku), nastavení toku odevzdání by mělo vyžadovat nahrání procesních artefaktů a jasnou odpovědnost vyučujícího za finální hodnocení \cite{kb_ai_101_tipu_pdf_04e4a115_page_8_1}.
  \item Volba enterprise/integračních řešení a autentizace (SSO) usnadní řízení přístupů a auditování používání nástrojů; detaily integrace závisí na místních IT procesech a licencech (general background knowledge).
  \item Uživatelům doporučujeme využít existujících oficiálních vzdělávacích kanálů a dokumentace poskytovatelů (např. Microsoft Learn pro Microsoft‑stack) při zavádění nových funkcionalit a školení učitelů \cite{kb_ai_101_tipu_pdf_04e4a115_page_54_1}.
\end{itemize}

Krátký checklist pro zavedení do praxe (katedra / garant předmětu)
\begin{enumerate}
  \item Definujte požadované procesní artefakty pro každé odevzdání (meziverze, postupy, testy).
  \item Nasadťe automatické testy a základní CI kroky pro okamžitou kontrolu funkčnosti (general background knowledge).
  \item Určete pravidla pro AI‑assisted pre‑review a pro lidskou finální validaci (kdo, kdy, jak).
  \item Aktualizujte sylabus: pole pro deklaraci použití AI a pokyny pro odevzdání procesních artefaktů.
  \item Pilotujte workflow na malé skupině a vyhodnoťte jak kvalitu, tak provozní zatížení vyučujících (general background knowledge).
\end{enumerate}

Vzory promptů a krátká ukázka (general background knowledge)
\begin{verbatim}
Jsi asistent učitele pro předměty programování. Pro každé odevzdání vygeneruj:
1) Shrnutí výsledků CI (passed/failed tests).
2) Tři nejpravděpodobnější příčiny případných selhání.
3) Doporučení pro studenta, jak opravit chybu (konkrétní kroky).
Formát: Student: <shrnutí> | Chyby: <1,2,3> | Doporučení: <body>
\end{verbatim}

Závěrem: v technických oborech je klíčové kombinovat automatizaci (testy, CI, před‑review nástroji) s povinnou lidskou validací a se zaměřením na procesní důkazy. Pro matematické a technické cvičení využijte funkcionality platforem, které umožňují požadovat nahrání postupu řešení; učitel pak finalizuje hodnocení a rozhoduje o pedagogické akceptaci automatizovaných návrhů \cite{kb_ai_101_tipu_pdf_04e4a115_page_8_1, kb_ai_101_tipu_pdf_04e4a115_page_54_1}.%

%
\section{Etika, GDPR a akademická integrita v kontextu hodnocení}%
Zásady transparentního oznamování použití AI, požadavky na souhlas a zpracování osobních údajů při sběru a zpracování studentských výstupů, pravidla akademické integrity při povolení/omezení AI nástrojů v zadáních.%


%
Úvodní poznámka: při hodnocení v době dostupné generativní AI se musí řešit současně tři propojené oblasti: transparentnost vůči studentům, ochrana osobních údajů při sběru a zpracování artefaktů a jasná pravidla akademické integrity. Tyto oblasti spolu úzce souvisejí a rozhodnutí v jedné oblasti mají dopad na ostatní; je proto vhodné navrhovat opatření komplexně, nikoli izolovaně. Níže najdete stručné zásady a praktické postupy k okamžitému použití na úrovni předmětu nebo katedry, které lze adaptovat podle konkrétního oboru a typu zadání.

Zásady (konkrétní doporučení)
\begin{itemize}
  \item Požadavek na procesní artefakty: vždy, kde je to pedagogicky oprávněné, vyžadujte od studentů nahrání meziverzí, zápisků postupů, commit logů nebo screenshotů z průběhu práce — procesní důkazy snižují motivaci k neoprávněnému použití AI, protože hodnocení váží postup stejně nebo více než finální výsledek. Praktický příklad: v některých nástrojích lze nastavit povinnost nahrání postupu řešení; u „Pokroku v matematice“ je explicitně možné požadovat nahrání postupu jako součást zadání \cite{kb_ai_101_tipu_pdf_04e4a115_page_8_1}. Procesní artefakty také pomáhají učiteli při ověřování porozumění a identifikaci zdrojů chyb.
  \item Transparentnost vůči studentům: do sylabu uveďte, jaké AI nástroje jsou v kurzu povoleny, za jakých podmínek (povinná dokumentace, omezení generování hotových řešení) a jak mají studenti deklarovat použití AI (general background knowledge). Zveřejněná pravidla by měla být konkrétní a srozumitelná, aby studenti věděli, čeho se mají vyvarovat, a jak dokumentovat oprávněné použití nástrojů.
  \item Minimální sběr osobních údajů: sbírejte jen ta data, která jsou nezbytná pro hodnocení (general background knowledge). Pokud používáte nástroje třetích stran, ověřte smluvní podmínky zpracování (DPA) a možnosti data residency (general background knowledge). Zaznamenávejte, proč jsou jednotlivá data potřeba, a jak dlouho budou uchovávána.
  \item Upřednostněte enterprise/školní licence tam, kde je to možné — integrace do školních balíků (např. Microsoft 365) často nabízí lepší kontrolu nad daty a nastavením soukromí; některé školní licence zahrnují i specifické vzdělávací aplikace jako Minecraft Education nebo funkce OneNote pro převod rukopisu \cite{kb_ai_101_tipu_pdf_04e4a115_page_15_2, kb_ai_101_tipu_pdf_04e4a115_page_14_1}. Tyto licence navíc obvykle umožňují centrální správu přístupů, auditování a případnou integraci do školních systémů LMS.
  \item Jasná pravidla akademické integrity: rozhodněte, které typy úloh AI povolíte (např. povoleno s povinnou reflexí a deklarací) a které jsou zakázané; pro podezření z neférového použití preferujte procesní ověření před automatickými detektory (general background knowledge). Doporučuje se také definovat důsledky porušení pravidel a postupy pro jejich vynucení, včetně postupného stupňování sankcí.
\end{itemize}

Krátký operační checklist pro vyučující (krok za krokem)
\begin{enumerate}
  \item Aktualizujte sylabus: vložte krátký odstavec o pravidlech používání AI, požadavku na deklaraci a popisu požadovaných procesních artefaktů (general background knowledge). Doporučujeme ukázkový text vložit přímo do elektronického sylabu i do první přednášky.
  \item Informujte studenty při zahájení kurzu: vysvětlete prakticky, jak deklarovat použití AI, jaké artefakty nahrávat a proč jsou tyto kroky důležité pro hodnocení a důvěryhodnost. Pořaďte krátký demonstrační návod nebo FAQ v LMS.
  \item Před odevzdáním nastavte požadavek na postup: v LMS vyžadujte nahrání meziverzí, logu nebo exportu z nástrojů (commit history, OneNote stránky apod.) — viz příklady praxe v ostatních vzdělávacích nástrojích, kde je nahrání postupu součástí zadání \cite{kb_ai_101_tipu_pdf_04e4a115_page_8_1}. Explicitní nastavení v LMS usnadní konzistenci hodnocení mezi asistenty a učiteli.
  \item Použijte školní/enterprise řešení pro generování či asistenci (pokud je dostupné), abyste omezili nechtěné sdílení osobních dat do veřejných cloudů (general background knowledge). V rámci Microsoft‑stacku lze využít integrace do Poznámkového bloku / Copilot funkcí pro bezpečnější workflow \cite{kb_ai_101_tipu_pdf_04e4a115_page_14_1}. Ověřte také, zda použitá řešení podporují export auditních záznamů.
  \item Při podezření postupujte systematicky: zkontrolujte procesní artefakty → pohovor se studentem (ověření porozumění) → případná disciplinární opatření podle místních pravidel (general background knowledge). Dokumentujte průběh šetření a uchovejte záznamy pro případné další kroky.
  \item Zpětná vazba a iterace: na konci semestru vyhodnoťte, jak opatření fungovala, sesbírejte názory studentů a asistentů a upravte postupy pro další běh kurzu. Takto udržíte pravidla aktuální vůči rychle se měnícím technologiím.
\end{enumerate}

Vzorový text do sylabu (kratká formulace — upravte dle potřeby)
\begin{quote}
Studenti mohou v této části kurzu využívat generativní AI pouze za podmínky, že: (1) každé použití AI bude v odevzdání explicitně deklarováno, (2) k odevzdání budou přiloženy procesní artefakty (meziverze, log změn, krátké vysvětlení postupu) a (3) konečné hodnocení zahrnuje ověření osobního porozumění učitelem. Neoprávněné použití AI bude posuzováno podle pravidel akademické integrity fakulty. V případě nejasností kontaktujte vyučujícího nebo garanta předmětu. (general background knowledge)
\end{quote}

Krátké technické doporučení (bezpečnost dat a nástroje)
\begin{itemize}
  \item Pokud ve výuce využíváte nástroje, které umí převádět rukopis nebo pomáhat s řešením (např. OneNote či specializované „education“ aplikace), nastavte to tak, aby citlivá data studentů nesdílela více než nutné a aby bylo jasné, kde jsou data uchovávána \cite{kb_ai_101_tipu_pdf_04e4a115_page_8_1, kb_ai_101_tipu_pdf_04e4a115_page_14_1}. Zvažte anonymizaci vzorků tam, kde to hodnocení umožňuje.
  \item U didaktických světů a simulací (např. Minecraft Education) zkontrolujte licenční podmínky a dostupnost bezplatných hodin kódu jako pomůcky pro výuku AI principů; tyto prostředí mohou být užitečné pro demonstraci etických aspektů a principů práce s daty \cite{kb_ai_101_tipu_pdf_04e4a115_page_15_2}. Pro simulace také ověřte, zda prostředí neposkytuje zbytečně široké sdílení studentových údajů.
  \item Správa přístupů a šablony: vytvořte jednoduché šablony pro nahrávání procesních artefaktů a instrukce pro pojmenování souborů, aby bylo možné rychle vyhledávat a kontrolovat odevzdání v LMS.
\end{itemize}

Poznámka o právních a GDPR aspektech: konkrétní povinnosti (např. provedení DPIA, dohody o zpracování osobních údajů, role správce/zpracovatele) závisí na rozsahu zpracování dat a na lokálních pravidlech VUT — tyto kroky považujeme za nezbytné, ale jejich podrobný postup je uveden v kapitole o ochraně osobních údajů a v přiložené šabloně DPIA (viz Přílohy). (general background knowledge) V praxi je vhodné konzultovat právní oddělení instituce před zavedením nových nástrojů do výuky a uchovávat záznamy o provedených posouzeních.

Shrnutí: kombinujte pedagogicky odůvodněné požadavky na procesní artefakty s jasnou komunikací pravidel do sylabu a s preferencí školních/enterprise řešení tam, kde to zajišťuje lepší kontrolu nad daty. Tam, kde to technicky umožňují nástroje, využijte možnosti nastavení povinného nahrání postupu řešení jako standardního požadavku při hodnocení \cite{kb_ai_101_tipu_pdf_04e4a115_page_8_1}. Pravidelně revidujte pravidla a technická nastavení podle zkušeností z praxe a změn v legislativě či nabídce nástrojů.%

%
\section{Proč se vyvarovat detektorů AI a alternativy}%
Argumenty proti spoléhaní se na detektory AI a doporučení alternativních postupů (design úloh, procesní ověřování, reflexe studentů, odborné interview), včetně praktických kroků, jak správně řešit podezření z neférového použití AI.%


%
Krátké shrnutí problému
Detektory „AI‑textu“ mohou pomoci, ale nejsou spolehlivé ani právně/eticky bezproblémové; níže stručně proč a praktické alternativy.
Proč detektorům nevěřit \begin{itemize}
  \item Nejsou definitivním důkazem: falešné pozitivy/negativy; neřeší autorství ani kontext.
  \item Jako hlavní nástroj mohou vést k nespravedlnosti a poškození důvěry.
\end{itemize}
Praktické alternativy a doporučené postupy
1) Požadavek na procesní artefakty \begin{itemize}
  \item Požadujte meziverze, commit logy, exporty nebo screenshoty; hodnocení se zaměří na proces \cite{kb_ai_101_tipu_pdf_04e4a115_page_8_1}.
  \item V LMS mějte povinné pole pro průběžné verze a šablony.
\end{itemize}
2) Úlohy odolné vůči triviálnímu použití AI \begin{itemize}
  \item Dejte přednost kontextovým úlohám, místním zdrojům a iteracím; používejte časově omezená in‑class writings, ústní zkoušky či portfolia.
\end{itemize}
3) Transparentní deklarace a reflexe použití AI \begin{itemize}
  \item Pokud je AI povolena, vyžadujte deklaraci a krátkou sebereflexi o roli nástroje a ověření výstupů.
\end{itemize}
4) Ověřovací rozhovor \begin{itemize}
  \item Při podezření doplňte krátký pohovor, kde student vysvětlí postup a rozhodnutí.
\end{itemize}
5) Školní/enterprise řešení \begin{itemize}
  \item Licencované balíky nabízejí lepší kontrolu dat, auditní logy a integraci do výuky \cite{kb_ai_101_tipu_pdf_04e4a115_page_36_1, kb_ai_101_tipu_pdf_04e4a115_page_15_2}.
\end{itemize}
Postup při podezření na nevhodné použití AI \begin{enumerate}
  \item Zkontrolujte povinné procesní artefakty; chybějící vyžádejte \cite{kb_ai_101_tipu_pdf_04e4a115_page_8_1}.
  \item V pohovoru se zaměřte na vysvětlení postupu, zdrojů a rozhodnutí.
  \item Dokumentujte zjištění podle interních pravidel.
  \item Pravidelně vyhodnocujte a upravujte opatření.
\end{enumerate}
Závěrečná poznámka
Místo spoléhání se na automaty doporučujeme kombinaci dobře navržených úloh, povinných artefaktů, transparentních pravidel a krátkých ověřovacích pohovorů; tyto přístupy lépe chrání akademickou integritu a respektují právní a etické aspekty práce se studentskými daty \cite{kb_ai_101_tipu_pdf_04e4a115_page_8_1, kb_ai_101_tipu_pdf_04e4a115_page_36_1, kb_ai_101_tipu_pdf_04e4a115_page_15_2}.%

%
\section{Šablony, checklisty a prompt‑knihovna pro hodnocení}%
Sada připravených šablon pro milestone‑based hodnocení, rubriky obsahující pole pro deklaraci AI, checklisty pro garanty a prompt‑šablony pro zadávání úloh i pro asistenty hodnotitelů při generování formativní zpětné vazby.%


%
Úvodní poznámka. Níže jsou zkrácené šablony, checklisty a prompty pro zavedení milestone‑based hodnocení a transparentní dokumentaci použití AI (požadavek na procesní artefakty) \cite{kb_ai_101_tipu_pdf_04e4a115_page_8_1}. Preferujte školní/enterprise řešení pro kontrolu dat a auditní stopy, pokud to technicky a právně jde \cite{kb_ai_101_tipu_pdf_04e4a115_page_15_2}.

1) Rychlý checklist pro garanty předmětů (před spuštěním)
\begin{itemize}
  \item Uveďte v sylabu pravidla AI: povoleno/zakázáno, povinná deklarace a požadované artefakty (meziverze, commit logy, screenshoty) \cite{kb_ai_101_tipu_pdf_04e4a115_page_8_1}.
  \item Rozdělte závěrečnou úlohu na 1–3 milníky s jasnými výstupy a tím, co hodnotí AI vs. člověk.
  \item Prověřte právní/technické podmínky ukládání dat (enterprise licence, DPA) \cite{kb_ai_101_tipu_pdf_04e4a115_page_15_2}.
  \item Připravte studentské instrukce pro deklaraci AI a nahrávání artefaktů.
\end{itemize}

2) Vzorová položka do rubriky
\begin{verbatim}
Kritérium: Dokumentace postupu
Popis: Odevzdejte meziverze / commit log / screenshoty.
Max. bodů: 15
Hodnocení: úplnost, srozumitelnost, návaznost verzí
Deklarace AI: [Ano/Ne] Nástroj: [název] Role AI: [nápověda/generování/ověření]
\end{verbatim}

3) Jednoduchá rubrika (kopírovat do LMS)
\begin{center}
\begin{tabular}{p{6cm}p{1.5cm}p{6cm}}
\hline
Kritérium & Max. bodů & Poznámky \\
\hline
Kvalita řešení & 40 & Ověřitelné kroky v artefaktech \\
Proces a transparentnost & 20 & Povinné artefakty \\
Argumentace a reflexe & 20 & Krátká reflexe o roli AI \\
Originalita & 10 & Lokální/datumové prvky oceňovány \\
Formální náležitosti & 10 & Pole deklarace AI musí být vyplněno \\
\hline
\end{tabular}
\end{center}

4) Šablona deklarace AI (do sylabu/formuláře)
\begin{verbatim}
Deklarace použití AI: Uveďte, zda jste použili AI, nástroj a roli (nápověda/editace/generování).
Popište ověření správnosti výsledků. Přiložte krátkou sebereflexi (max. 200 slov).
\end{verbatim}

5) Prompt‑knihovna pro asistenty hodnotitelů
Pozn.: anonymizujte data nebo provozujte v enterprise prostředí před generováním komentářů \cite{kb_ai_101_tipu_pdf_04e4a115_page_15_2}.
\begin{verbatim}
You are an assistant producing formative feedback.
Input: (student_submission) (rubric) (declared_AI_use).
Task: Up to 250 words:
 - 2 sentences praising strongest parts (linked to rubric),
 - 2 specific improvement suggestions (linked to rubric),
 - 1 suggestion for next milestone,
 - If AI was used: "Please verify [specific claim] & cite sources."
Tone: constructive, actionable.
\end{verbatim}

6) Vzor workflow pro milestone‑based hodnocení
\begin{enumerate}
  \item Student odevzdá milník + artefakty + deklaraci AI.
  \item Učitel / AI‑asistent v enterprise prostředí vygeneruje formativní komentáře dle rubriky.
  \item Student upraví práci a odevzdá další milník; učitel validuje klíčové kroky.
  \item Finální hodnocení zohlední proces, reflexi a kvalitu podle rubriky.
\end{enumerate}
Pozn.: Procesní artefakty snižují motivaci k nevhodnému použití AI a zvyšují ověřitelnost postupu \cite{kb_ai_101_tipu_pdf_04e4a115_page_8_1}.

7) Praktické doporučení pro multimodální úlohy
Zaznamenávejte zdroje a postupy při generování/úpravě obrázků; některé nástroje umožňují nastavit poměry a stylizaci, což pomáhá reprodukovatelnosti. Požadujte logy práce s obrazovými výstupy \cite{kb_ai_101_tipu_pdf_04e4a115_page_51_1}.

8) Jak začít rychle (shrnutí)
\begin{itemize}
  \item Přidejte do sylabu formulaci o deklaraci AI.
  \item Zvolte 1–2 jasné milníky s definovanými výstupy a artefakty.
  \item Připravte jednoduchou rubriku s polem pro AI deklaraci.
  \item Ověřte možnosti LMS pro sběr meziverzí/commit logů.
  \item Pilotujte s jednou sekcí a sbírejte zpětnou vazbu.
\end{itemize}

Odkazy: viz Přílohy 9.1 (Prompt Library) a 9.2–9.4 (šablony, DPIA, checklist); doporučení pro enterprise nasazení a volbu licencí viz kap. o institucionálním nasazení \cite{kb_ai_101_tipu_pdf_04e4a115_page_15_2, kb_ai_101_tipu_pdf_04e4a115_page_8_1}.%

%
\section{Případové studie a ověřené postupy z praxe}%
Krátké, ověřené příklady implementace z univerzit a kurzů (modely adaptovatelné pro VUT), klíčové poučení a rychlý plán zavedení vybraného vzoru v lokálním předmětu.%


%
Krátké, ověřené příklady nasazení a doporučené postupy, které lze adaptovat na předmět nebo pilot na úrovni katedry. Níže uvedené příklady vycházejí z praktických doporučení a nástrojů zmíněných ve zdrojovém materiálu; u každého příkladu je uvedena klíčová poznámka k přenositelnosti.

Příklady z praxe
\begin{itemize}
  \item Minecraft Education — výukové světy zaměřené na seznámení se s principy AI (Hour of Code) jsou volně dostupné a škola může použít připravené lekce k demonstraci principů strojového učení a praktických aktivit ve vyučování. Část obsahu je dostupná zdarma; rozšiřující světy jsou součástí školní licence Microsoft 365 A3/A5 \cite{kb_ai_101_tipu_pdf_04e4a115_page_15_2}.
  \item Microsoft výukové nástroje pro matematiku a OneNote — nástroje jako „Pokrok v matematice“ a funkce OneNote pro převod rukopisu na digitální formu usnadňují generování cvičných příkladů, sběr postupů a diagnostiku chyb u studentů; lze je využít pro adaptivní procvičování a sledování pokroku \cite{kb_ai_101_tipu_pdf_04e4a115_page_8_1}.
  \item Profesní rozvoj učitelů — doporučení zaměřit se na vybranou oblast/technologii a učit se v hloubce namísto snahy „stíhat vše“; pro praktické školení lze využít oficiální materiály (např. Microsoft Learn) jako systematický zdroj kurzů a modulků pro pedagogy i IT podporu \cite{kb_ai_101_tipu_pdf_04e4a115_page_54_1}.
\end{itemize}

Klíčové poučení (co si odnést)
\begin{itemize}
  \item Využijte hotové didaktické artefakty (např. připravené světy v Minecraft Education) pro rychlý start s praktickými ukázkami ve výuce \cite{kb_ai_101_tipu_pdf_04e4a115_page_15_2}.
  \item Integrace nástrojů, které zaznamenávají postupy (commit/artefakty, nahrávání postupu), zvyšuje ověřitelnost učení a usnadňuje hodnocení procesu — při výběru nástrojů preferujte školní/enterprise řešení, pokud to právně/technicky dává smysl \cite{kb_ai_101_tipu_pdf_04e4a115_page_15_2, kb_ai_101_tipu_pdf_04e4a115_page_8_1}.
  \item Systematické školení pedagogů a podpůrného personálu (vybraná tematická oblast + oficiální learning moduly) zrychlí adaptaci a sníží riziko nekonzistentního nasazení \cite{kb_ai_101_tipu_pdf_04e4a115_page_54_1}.
\end{itemize}

Rychlý plán zavedení vybraného vzoru (1–2měsíční pilot)
\begin{enumerate}
  \item Vyberte jeden kurz nebo modul a konkrétní didaktický cíl (např. demonstrace principů AI v 90minutové hodině pomocí Minecraft Education). Podle doporučení začněte s úzkou oblastí \cite{kb_ai_101_tipu_pdf_04e4a115_page_54_1}.
  \item Ověřte licenční a datovou situaci (je-li potřeba školní licence Microsoft 365 A3/A5 pro rozšířené světy nebo enterprise řešení pro zpracování studentských dat) a zajistěte DPA / souhlas tam, kde je to nutné \cite{kb_ai_101_tipu_pdf_04e4a115_page_15_2}.
  \item Připravte konkrétní hodinu/aktivitu: stáhněte volnou lekci Hour of Code nebo připravený svět, upravte ji pro lokální kontext, připravte krátký doprovodný materiál pro studenty (cíle, očekávané výstupy) \cite{kb_ai_101_tipu_pdf_04e4a115_page_15_2}.
  \item Nasazení a sběr dat o průběhu: zaznamenávejte artefakty studentské práce a krátkou reflexi (pomůže při vyhodnocení a integritě). Pokud používáte nástroje pro procvičování (math tools, OneNote), zapojte požadavek na nahrání postupu řešení \cite{kb_ai_101_tipu_pdf_04e4a115_page_8_1}.
  \item Vyhodnocení pilotu a škálování: zhodnoťte dosažení cílů, zkušenosti učitele i studentů a rozhodněte o rozšíření (úpravy obsahu, školení dalších učitelů pomocí dostupných modulů, např. Microsoft Learn) \cite{kb_ai_101_tipu_pdf_04e4a115_page_54_1}.
\end{enumerate}

General background knowledge (stručně)
\begin{itemize}
  \item General background knowledge: Ve veřejně známých případech existují širší institucionální nasazení AI (např. Georgia Tech, Harvard CS50, Arizona State, UCL, TU Delft) — tyto kazuistiky slouží jako inspirace pro design patterny a škálování, avšak konkrétní podrobnosti a výsledky je nutné konzultovat v originálních zdrojích (uváděno jako všeobecná reference, není detailně citováno v tomto zdroji).
\end{itemize}%

%
\chapter{AI PRO PODPORU STUDENTŮ}%
Kapitola mapuje AI jako personalizovaného tutora, nástroje pro zpětnou vazbu na písemné práce, podporu jazykových dovedností a přístupnosti. Praktické scénáře zahrnují nastavení pravidel pro AI tutora (sokratovský přístup, hintování, zákaz hotových řešení) a ukázky promptů. Nabízí tierovaný seznam nástrojů pro feedback (Grammarly, Hemingway, Turnitin Draft Coach), doporučení pro mezinárodní studenty (DeepL, ChatGPT pro konverzaci, ELSA Speak) a řešení pro accessibility (Microsoft Immersive Reader, Speechify). Vychází z praxe univerzit, kde AI tutori a asistivní technologie prokazatelně zlepšují dostupnost podpory a studijní zkušenost.%
\section{Role AI jako personalizovaný tutor}%
Přehled pedagogických rolí AI (personalizace učení, adaptivní nápověda) a doporučené režimy interakce (sokratovský přístup, hintování, zákaz poskytování hotových řešení). Zmíněny principy pro nastavení pravidel pro tutora a očekávané přínosy pro individuální podporu studentů.%


%
Krátké shrnutí role
AI může v kurzu sloužit jako personalizovaný tutor, který poskytuje adaptivní cvičení, cílené nápovědy a rychlou formativní zpětnou vazbu přizpůsobenou potřebám jednotlivého studenta.
Praktické ukázky komerčních nástrojů potvrzují, že rozhraní pro učitele umožňují sledovat, v jakých oblastech studenti chybují, a generovat pro ně upravené úlohy či doporučení (např. funkce „Pokrok v matematice“) \cite{kb_ai_101_tipu_pdf_04e4a115_page_8_1}.
Tyto funkce učitelům nabízí přehledy výkonu, trendy v čase a konkrétní doporučení pro intervenci.
Pro demonstrační a didaktické účely mohou posloužit i interaktivní světy (např. Minecraft Education) vhodné pro vizualizaci principů učení strojů a první praktické aktivity s žáky \cite{kb_ai_101_tipu_pdf_04e4a115_page_15_2}.
Interaktivní prostředí pomáhají studentům porozumět abstraktním konceptům prostřednictvím vizuálních a herních metafor.
Nástroje typu Microsoft 365 Copilot pak mohou usnadnit hromadnou generaci testů a kvízů, což učiteli uvolní kapacitu pro individuální podporu studentů \cite{kb_ai_101_tipu_pdf_04e4a115_page_36_1}.
Automatizace rutinních úkonů umožňuje věnovat více času reflexi, doplňujícím vysvětlením a práci s rizikovými studenty.

Režimy interakce (praktické varianty)
- Adaptivní cvičení a diagnostika: systém nabízí úlohy podle dosavadního výkonu studenta a učiteli zobrazuje přehled oblastí s chybami (viz „Pokrok v matematice“) \cite{kb_ai_101_tipu_pdf_04e4a115_page_8_1}.
- Adaptive režim může být nastaven tak, aby zvýšil opakování klíčových konceptů u studentů s nižší úspěšností.
- Hintování místo hotových řešení (general background knowledge): tutor dává postupné nápovědy, které vedou studenta k vlastnímu řešení místo poskytnutí kompletní odpovědi.
- Tento přístup podporuje rozvoj metakognitivních dovedností a samostatného myšlení.
- Sokratovský styl dotazování (general background knowledge): tutor klade vedoucí otázky zaměřené na vysvětlení myšlenkového postupu studenta.
- Otázky mohou být konstruovány tak, aby diagnostikovaly nesprávné předpoklady a pomohly je odhalit.
- Generování cvičných testů a materiálů pro opakování: Copilot/Forms lze použít pro rychlé vytvoření testů a jejich variant pro diferenciaci \cite{kb_ai_101_tipu_pdf_04e4a115_page_36_1}.
- Hromadné generování šablon zrychluje přípravu a umožňuje snadno vytvářet varianty pro různé úrovně.
- Demonstrace konceptů a motivace v herním prostředí: připravené světy v Minecraft Education lze využít pro praktické ukázky a motivaci studentů při seznámení s problematikou AI \cite{kb_ai_101_tipu_pdf_04e4a115_page_15_2}.
- Hry a simulace zvyšují angažovanost a dávají možnost bezpečného experimentování.

Praktická pravidla pro nastavení AI tutora (doporučené)
- Požadujte důkazy o postupu: při úkolech, kde je důležitý proces (např. matematika), zařaďte povinnost nahrát nebo odevzdat postup řešení; to zvyšuje ověřitelnost učení a umožní cílenou zpětnou vazbu \cite{kb_ai_101_tipu_pdf_04e4a115_page_8_1}.
- Požadavky na postup pomáhají odhalit pouhé kopírování odpovědí bez pochopení.
- Omezte generování kompletních řešení tam, kde je cílem rozvíjet postupy a dovednosti (general background knowledge).
- Místo toho nastavte systém, aby poskytoval nápovědu po krocích a kontrolní otázky.
- Transparentnost vůči studentům: jasně formulujte pravidla používání AI v kurzu (co je povoleno, jak deklarovat použití AI) — viz příslušné šablony v příručce (general background knowledge).
- Komunikujte také důsledky porušení pravidel a požadavky na citování použití nástrojů.
- Využívejte enterprise/školní varianty nástrojů tam, kde se zpracovávají osobní údaje či sledované artefakty studentů (general background knowledge); pro demonstrační aktivity je možné použít i volně dostupné učební světy (Minecraft Education) \cite{kb_ai_101_tipu_pdf_04e4a115_page_15_2}.
- Zvažte nastavení retence dat a přístupu ke záznamům, aby byla zajištěna ochrana soukromí.

Praktický workflow pro zavedení tutora na úrovni předmětu (krátký návod)
1. Definujte cíle tutorování: které dovednosti má tutor podporovat (např. procedurální řešení úloh, jazyková podpora, opakování pojmů). (general background knowledge)
2. Vyberte nástroje a režimy: pro procvičování matematiky zvažte nástroje s možností sběru postupu (viz „Pokrok v matematice“), pro tvorbu testů použijte Copilot/Forms pro hromadné generování otázek \cite{kb_ai_101_tipu_pdf_04e4a115_page_8_1, kb_ai_101_tipu_pdf_04e4a115_page_36_1}.
3. Nastavte pravidla interakce a omezení (hinty ano, kompletní řešení ne; požadavek na postup). Dokumentujte tato pravidla v sylabu (general background knowledge).
4. Pilotujte malou skupinu studentů: zaznamenejte artefakty práce, sbírejte zpětnou vazbu a upravte chování tutora podle zjištěných potřeb.
5. Pro ukázky principů použijte připravené hodiny v Minecraft Education tam, kde je vhodné vizuální či herní zpracování \cite{kb_ai_101_tipu_pdf_04e4a115_page_15_2}.
6. Škálování: rozšiřte adaptivní banku úloh a připravte šablony testů/kvízů pro různé úrovně pomocí Copilot/Forms \cite{kb_ai_101_tipu_pdf_04e4a115_page_36_1}.
- Během pilotu sledujte metriky zapojení, přesnosti a času stráveného u úloh a iterujte podle dat.

Příklad krátkých pravidel pro studenta (možno vložit do sylabu) — general background knowledge:
- Tutor může poskytovat nápovědy a kontrolovat postupy; kompletní řešení je zakázáno u zadaných úloh, kde se hodnotí postup.
- Při použití AI uvolněného typu (např. pro stylistickou pomoc) uveďte stručnou deklaraci použití a přiložte vlastní reflexi (max. 200 slov).
- Dodržujte zásady akademické poctivosti při využívání AI a konzultujte pochybné případy s vyučujícím.

Očekávané přínosy (konkrétní + obecné)
- Konkrétní: nástroje pro diagnostiku umožní učiteli snadno identifikovat oblasti, kde studenti chybují, a nasadit cílené intervence (viz „Pokrok v matematice“) \cite{kb_ai_101_tipu_pdf_04e4a115_page_8_1}.
- Obecné (general background knowledge): lepší personalizace, rychlejší formativní zpětná vazba a zvýšená angažovanost studentů při adaptivních aktivitách a interaktivních demonstracích (např. herní světy).
- Dlouhodobě může systém přispět ke zefektivnění výuky a lepšímu sledování pokroku meziročními daty.

Kontrolní seznam pro učitele (rychle)
- Má nastavené pravidlo pro sběr postupu u úloh? (viz \cite{kb_ai_101_tipu_pdf_04e4a115_page_8_1})
- Je vymezena politika pro povolené/zakázané chování tutora a je komunikována studentům? (general background knowledge)
- Je připraven demonstrační materiál (např. Minecraft Education world) pro první hodinu seznámení? (viz \cite{kb_ai_101_tipu_pdf_04e4a115_page_15_2})
- Jsou připraveny šablony pro rychlé generování testů a jejich varianty (Copilot/Forms)? (viz \cite{kb_ai_101_tipu_pdf_04e4a115_page_36_1})

Poznámka k bezpečnosti a ochraně dat
- Pokud tutor zpracovává osobní údaje nebo artefakty studentů, zvažte použití školní/enterprise varianty služby a postupy pro anonymizaci či omezení retence dat (general background knowledge).
- Vhodné demonstrace bez citlivých dat lze realizovat i pomocí bezplatných výukových světů a veřejných materiálů \cite{kb_ai_101_tipu_pdf_04e4a115_page_15_2}.
- Pravidelně kontrolujte nastavení sdílení a záloh, aby nedocházelo k nežádoucímu zpřístupnění citlivých informací.%

%
\section{Praktické nastavení AI tutora: pravidla, workflow a ukázky promptů}%
Konkrétní postupy pro konfiguraci AI tutora v kurzu: pravidla chování, příklady promptů (pro vedení otázkami, kontrolu postupů, omezení generování hotových řešení) a doporučené pracovní toky mezi studentem, AI a vyučujícím.%


%
\begin{flushleft}
Krátké shrnutí a cíle této sekce: prakticky nastavit chování a pracovní toky AI tutora v kurzu tak, aby podporoval učení (hinty, diagnostika, adaptivní cvičení), zároveň minimalizoval riziko „hotových řešení“ a respektoval ochranu osobních údajů. Níže naleznete doporučená pravidla, navržený workflow pro zavedení a ukázkové prompty (označené jako general background knowledge tam, kde nejsou přímo podloženy KB).
\end{flushleft}

\paragraph{Klíčová pravidla chování tutora (nutné nastavit před pilotem)}
\begin{itemize}
  \item Požadovat důkazy o postupu u úloh, kde je důležitý proces (např. matematické výpočty): nastavit povinnost nahrát/sklíčit postupy, které tutor kontroluje a k nimž generuje nápovědy, nikoli kompletní řešení \cite{kb_ai_101_tipu_pdf_04e4a115_page_8_1}.
  \item Upřesnit režim nápověd (hintování) místo poskytování hotových řešení; formulovat to ve srozumitelné politice pro studenty a vložit do sylabu (viz šablony v příručce) — toto pravidlo ozřejmuje očekávané chování tutora (general background knowledge).
  \item Transparentnost: studenty informovat, kdy a jak tutor zpracovává jejich artefakty a jak mají deklarovat použití AI (general background knowledge).
  \item Pro demonstrace a non‑sensitive cvičení využít výukové světy (např. Minecraft Education) — vhodné i pro vizualizaci principů strojového učení; některé lekce jsou zdarma, širší obsah může být součástí školní licence Microsoft 365 A3/A5 \cite{kb_ai_101_tipu_pdf_04e4a115_page_15_2}.
  \item Při zpracování osobních údajů či sledování artefaktů preferovat školní/enterprise varianty nástrojů a ověřené datové dohody (general background knowledge). Pro generování a hromadné vytváření testů lze využít Microsoft 365 Copilot + Forms \cite{kb_ai_101_tipu_pdf_04e4a115_page_36_1}.
\end{itemize}

\paragraph{Doporučený praktický workflow pro nasazení tutora (krátký návod)}
\begin{enumerate}
  \item Definujte pedagogické cíle tutorování: které dovednosti má tutor podporovat (procedurální řešení, jazyková podpora, opakování pojmů). (general background knowledge)
  \item Vyberte nástroje podle účelu: pro procvičování matematiky zvažte řešení, které umožní sběr postupu řešení (viz funkce „Pokrok v matematice“); pro tvorbu kvízů zvažte Copilot+Forms \cite{kb_ai_101_tipu_pdf_04e4a115_page_8_1, kb_ai_101_tipu_pdf_04e4a115_page_36_1}.
  \item Nastavte pravidla: explicitně omezte generování kompletních řešení tam, kde požadujete proces; definujte formát odevzdávaných důkazů (snímky postupu, mezi‑verze, commit‑history). Požadavek na postup zvyšuje ověřitelnost výkonu \cite{kb_ai_101_tipu_pdf_04e4a115_page_8_1}.
  \item Pilotujte malou skupinu: sbírejte artefakty práce a krátké reflexe studentů; iterujte nastavení nápověd a omezení podle zpětné vazby.
  \item Škálování: vytvořte banku adaptivních úloh a šablon testů; uvažujte o enterprise licencích tam, kde se budou zpracovávat citlivá data nebo kde je vyžadována centrální správa \cite{kb_ai_101_tipu_pdf_04e4a115_page_36_1}.
\end{enumerate}

\paragraph{Praktické provozní nastavení (konkrétní doporučení)}
\begin{itemize}
  \item Konfigurace ukládání a retence: rozhodněte, jak dlouho se ukládají konverzace a artefakty, a informujte studenty. (general background knowledge)
  \item Auditní logy a přístup: zajistěte, kdo má přístup k datům tutora (učitelé, asistenti) a jak se provádí případná revize.
  \item Metody validace: propojte automatické kontroly s lidskou validací pro klíčové milníky; při podezření na nesprávné použití použijte verifikační rozhovor (ověřovací ústní fáze) jako doplněk k artefaktům. (general background knowledge)
\end{itemize}

\paragraph{Ukázkové prompty pro běžné scénáře (příklady, označeno jako general background knowledge)}
Poznámka: následující prompty jsou šablony k okamžitému použití; upravte je pro konkrétní kurz a pravidla omezování generování řešení.

\begin{itemize}
  \item „Sokratovský“ průvodce (vedení otázkami):
    \begin{quote}
      (general background knowledge) \\
      Jsi studentský tutor. Student posílá krátký výňatek svého řešení. Neodpovídej kompletním řešením. Polož 3 kontrolní otázky, které ověří klíčové předpoklady studenta a nabídni jednu drobnou nápovědu vedoucí ke správnému kroku.
    \end{quote}

  \item Kontrola kroků u matematického příkladu:
    \begin{quote}
      (general background knowledge) \\
      Student poslal postup řešení. Projděte každý krok a označte: (1) správný/částečně správný/nesprávný; (2) u nesprávného kroku popište, v čem je chyba (max. 30 slov); (3) navrhni krátkou nápovědu, která vede ke správnému dalšímu kroku, bez hotového řešení.
    \end{quote}

  \item Generování adaptivních cvičení (inspirace funkcí „Pokrok v matematice“):
    \begin{quote}
      (general background knowledge) \\
      Vygeneruj 5 krátkých úloh zaměřených na lineární rovnice se zadanou obtížností (stupeň: základní/střední/pokročilý). U každé úlohy uveď očekávaný typ chyby a návrh nápovědy, pokud student odpoví chybně.
    \end{quote}

  \item Vytvoření testu v učitelském workflow (inspirace Copilot+Forms):
    \begin{quote}
      (general background knowledge) \\
      Vygeneruj test na téma „lineární rovnice“ pro 1. ročník VŠ, 10 otázek: 6 ABC, 2 pravda/nepravda, 2 krátké otevřené. Přidej klíč odpovědí a krátké vysvětlení správných odpovědí (pro učitele).
    \end{quote}
\end{itemize}

\paragraph{Kontrolní checklist pro učitele před spuštěním tutora}
\begin{itemize}
  \item Je definován požadavek na sběr postupu tam, kde je to potřeba? Viz doporučení pro matematiku a funkci „Pokrok v matematice“ \cite{kb_ai_101_tipu_pdf_04e4a115_page_8_1}.
  \item Jsou stanovena pravidla pro omezení generování kompletních řešení a je to komunikováno studentům? (general background knowledge)
  \item Máte připravený demonstrační scénář (např. v Minecraft Education) pro první seznámení s principy a chováním tutora? Některé lekce jsou zdarma, další jsou v rámci školní licence \cite{kb_ai_101_tipu_pdf_04e4a115_page_15_2}.
  \item Pro generování testů a hromadné úlohy: máte ověřenu integraci s Copilot+Forms nebo jiným nástrojem a nastavenou kontrolu kvality otázek \cite{kb_ai_101_tipu_pdf_04e4a115_page_36_1}.
\end{itemize}

\paragraph{Závěr sekce}
Nejlepší praxe kombinuje jasná pravidla (transparentnost, požadavek na postup), malý pilotní běh s iterací a technologii přizpůsobenou účelu (nástroje podporující sběr postupu pro matematiku; Copilot+Forms pro hromadné kvízy). Pro demonstrační a motivační aktivity uvažujte o připravených výukových světech (Minecraft Education) — to usnadní pochopení principů AI studentům i učitelům \cite{kb_ai_101_tipu_pdf_04e4a115_page_15_2, kb_ai_101_tipu_pdf_04e4a115_page_36_1, kb_ai_101_tipu_pdf_04e4a115_page_8_1}.%

%
\section{Nástroje pro zpětnou vazbu na písemné práce}%
Tierovaný seznam a krátké srovnání nástrojů pro textovou zpětnou vazbu a korekturu v praxi (např. Grammarly, Hemingway, Turnitin Draft Coach) a doporučené postupy jejich využití v procesu hodnocení a rozvoje studenta.%


%
Úvodní poznámka: tato sekce shrnuje praktické nástroje a postupy pro poskytování zpětné vazby na písemné práce v kurzech VUT. Kde to dovolují zdroje, uvedeme konkrétní příklady integrovaných řešení; obecné doporučené postupy jsou označeny jako (general background knowledge).

Krátké rozdělení nástrojů (tierovaný přístup)
\begin{itemize}
  \item Tier 1 — integrované / školní enterprise nástroje: nástroje, které jsou dostupné přímo v rámci školního Office/M365 ekosystému a které lze spravovat centrálně (např. Microsoft 365 Copilot v aplikacích Word, Forms, OneNote apod.). Microsoft 365 Copilot je poskytován jako součást školních Office služeb (ve variantách) a placený doplněk přidává generativní funkce přímo do Word/Excel/PowerPoint/Forms/OneNote atd.; pro školní prostředí existují i webové verze chatbota Copilot Chat \cite{kb_ai_101_tipu_pdf_04e4a115_page_3_1}. 
  \item Tier 2 — specializované asistenční nástroje pro psaní (externí služby): nástroje typu kontrola gramatiky, stylu a čitelnosti (např. Grammarly, Hemingway) a nástroje pro zpětnou vazbu nad strukturou textu. Pozn.: konkrétní popisy funkcí těchto komerčních nástrojů nejsou v dostupných KB výpiscích uvedeny; níže jsou jejich využití a omezení popsány jako (general background knowledge).
  \item Tier 3 — nástroje pro detekci podobnosti a draft‑assistance (originality / průvodní nástroje): nástroje pro kontrolu plagiátorství a práce s návrhy (např. Draft Coach v řešeních typu Turnitin) a platformy, které nabízejí workflow pro revizi draftů. Opět: konkrétní vlastnosti a omezení těchto komerčních produktů jsou zde shrnuty jako (general background knowledge).
\end{itemize}

Doporučené scénáře použití (praktický workflow)
\begin{enumerate}
  \item Nastavte očekávání v sylabu: jasně definujte, které nástroje studenti mohou používat a jak mají deklarovat použití AI (viz vlákno deklarace níže). (general background knowledge)
  \item Požadujte meziverze/drafty: namiřte hodnocení na proces — student odevzdá první draft + krátkou sebereflexi + deklaraci AI; učitel nebo AI‑asistent poskytne strukturovanou formativní zpětnou vazbu; student upraví text a odevzdá finální verzi. Tento postup odpovídá principům milestone‑based workflow používaným v příručce (general background knowledge).
  \item Použijte integrované nástroje tam, kde je potřeba škálovat: pokud máte školní/enterprise Copilot nebo nástroj spravovaný fakultou, využijte jej k hromadnému generování komentářů k formálním aspektům (gramatika, čitelnost, kontrolní otázky) a nechte lektora validovat obsahově klíčové části, zvlášť tam, kde je rozhodující faktická přesnost \cite{kb_ai_101_tipu_pdf_04e4a115_page_3_1}.
  \item Pro matematiku a převod rukopisu: u úloh zahrnujících výpočty využijte nástroje, které dokáží zpracovat postupy (OneNote, math.microsoft.com a funkce „Pokrok v matematice“ byly v praxi popisovány jako užitečné k záznamu postupu a analýze chyb) \cite{kb_ai_101_tipu_pdf_04e4a115_page_8_1}.
  \item Pilotujte a dokumentujte: při prvním cyklu zaznamenejte typické chyby AI‑asistentů a nastavte pravidla retence dat a přístupů (kdo může číst konverzace a drafty). Preferujte školní licence a dohodu o zpracování dat pro nástroje, které budou zpracovávat osobní údaje studentů \cite{kb_ai_101_tipu_pdf_04e4a115_page_3_1}.
\end{enumerate}

Vzorový prompt pro asistenta generujícího formativní zpětnou vazbu
\vspace{0.5em}

\noindent (tento prompt odpovídá doporučenému formátu z knihovny promptů — adaptujte podle rubriky)
\begin{quote}
  Vstupy: (student_submission) (rubric) (declared_AI_use). \\
  Výstup (do ~250 slov): 2 věty chvály, 2 konkrétní návrhy ke zlepšení, 1 návrh na další milník; pokud AI bylo použito, přidej instrukci k ověření specifického tvrzení. Tone: konstruktivní, akční.
\end{quote}
\vspace{0.5em}

Praktický checklist pro učitele před nasazením AI‑podpory zpětné vazby
\begin{itemize}
  \item Je ve sylabu jasně formulováno, jak mají studenti deklarovat použití AI (viz níže uvedená krátká šablona deklarace)? (general background knowledge)
  \item Jsou definovány hranice role AI‑asistenta (např. kontrola jazyka a formy vs. obsahová validace)? (general background knowledge)
  \item Máte nastavené pravidlo sběru postupu pro úlohy, kde je důležitý proces (např. matematické postupy nebo programovací commity)? Pro matematiku zvažte funkce, které umožňují nahrání postupu a jeho analýzu \cite{kb_ai_101_tipu_pdf_04e4a115_page_8_1}.
  \item Je dostupná škálovatelná cesta lidské validace klíčových tvrzení (např. u závěrečných milníků)? (general background knowledge)
  \item Jsou řešeny otázky ochrany osobních údajů a je zajištěna odpovídající licence / DPA pro používané nástroje? Preferujte školní/enterprise varianty tam, kde se zpracovávají studentské artefakty \cite{kb_ai_101_tipu_pdf_04e4a115_page_3_1}.
\end{itemize}

Krátká šablona deklarace použití AI (vložit do instrukcí k odevzdání)
\begin{quote}
  Deklarace použití AI: Uveďte, zda jste použili AI, nástroj a roli (nápověda/editace/generování). Popište ověření správnosti výsledků. Přiložte krátkou sebereflexi (max. 200 slov).
\end{quote}
(vzor převzat z šablon uvedených v příručce; použití doporučeno pro každé odevzdání.)

Poznámky k omezením a rizikům (stručně)
\begin{itemize}
  \item AI nástroje dobře škálují korekturu jazyka a návrhy ke struktuře, avšak obsahovou správnost a hodnocení autorství musí validovat lidský garant. (general background knowledge)
  \item U úloh vyžadujících sledovatelný proces (matematika, programování) preferujte workflow, které sbírá mezi‑verze, snímky postupu, nebo repozitářový záznam (commity). Funkce jako převod rukopisu v OneNote a nástroje pro „pokrok v matematice“ mohou zjednodušit analýzu postupu \cite{kb_ai_101_tipu_pdf_04e4a115_page_8_1}.
  \item Pro demonstrační aktivity a motivaci studentů lze využít výukové světy (např. Minecraft Education) — vhodné zejména pro první seznámení s principy AI v učebním kontextu \cite{kb_ai_101_tipu_pdf_04e4a115_page_15_2}.
\end{itemize}

Závěrem: kombinujte integrované nástroje (kde jsou dostupné a spravované univerzitou) pro škálovatelné aspekty zpětné vazby se strukturovanými lidskými revizemi pro klíčové obsahové části. Pilotní cyklus s jasnými pravidly pro deklaraci AI a sběr postupu je efektivní cestou k udržení akademické integrity a podpory rozvoje studenta. (general background knowledge)%

%
\section{Podpora jazykových dovedností a mezinárodní studenti}%
Nástroje a metody pro podporu studentů s různým jazykovým zázemím: strojový překlad a asistence (DeepL, Microsoft Translator), konverzační praxe a simulace (ChatGPT), výcvik výslovnosti (ELSA Speak) a doporučené didaktické postupy.%


%
Krátké shrnutí a cíl sekce: Tato část popisuje praktické nástroje a didaktické postupy pro podporu studentů s různým jazykovým zázemím (přepis mluveného slova, asistované čtení, překlad, konverzační praxe a nácvik výslovnosti). Cílem je nabídnout konkrétní nástroje, scénáře nasazení a jednoduchý checklist pro zajištění přístupnosti a efektivity výuky.

Klíčové nástroje a jejich pedagogické role
\begin{itemize}
  \item Diktování a přepis v reálném čase (speech‑to‑text) — pomáhá studentům se sluchovým postižením nebo s omezenou znalostí jazyka; umožňuje sledovat mluvený projev současně s textovým zápisem, což usnadňuje porozumění a pozdější revizi \cite{kb_ai_101_tipu_pdf_04e4a115_page_11_1}.
  \item Asistivní čtečky a režimy pro lepší čtení (např. čtení nahlas, zvýrazňování, rozdělení textu na slabiky) — zlepšují srozumitelnost a přístupnost studijních materiálů; vhodné pro studenty s dyslexií nebo jinými potížemi se čtením \cite{kb_ai_101_tipu_pdf_04e4a115_page_11_1}.
  \item Microsoft Translator — využití jako „přepisovač“ v reálném čase (např. přepis řeči učitele pro studenty či přepis do jiného jazyka) bylo v praxi zmíněno jako užitečné řešení pro pomůcky ve třídě; hodí se pro multijazykové publikum a při hostujících přednáškách \cite{kb_ai_101_tipu_pdf_04e4a115_page_11_1}.
  \item SeeingAI a podobné asistivní aplikace — nástroje, které dokážou popsat obrazový obsah uživateli, zlepšují inkluzi studentů se zrakovým postižením; lze je využít pro popis pracovních listů, schémat nebo učebních pomůcek \cite{kb_ai_101_tipu_pdf_04e4a115_page_11_1}.
  \item Výuková prostředí hra‑založená (např. Minecraft Education) — pro didaktické přiblížení témat (včetně cizích jazyků) a interaktivní aktivity existují připravené světy; některé moduly jsou dostupné zdarma, plná verze je součástí školních licencí Microsoft 365 A3/A5 \cite{kb_ai_101_tipu_pdf_04e4a115_page_15_2}.
  \item Aplikační doplňky pro nácvik výslovnosti a konverzační praxe — různé aplikace nabízejí opakování frází, rozpoznávání výslovnosti a interaktivní cvičení; při výběru je vhodné ověřit přesnost hodnocení a ochranu osobních údajů.
\end{itemize}

Praktické scénáře a doporučené postupy (stručně)
\begin{enumerate}
  \item Živý přepis přednášky pro multijazykové publikum (podpůrný scénář): zapněte přepis v dostupné aplikaci (či sdílejte odkaz do Microsoft Translator), sdílejte text na plátně nebo do chatu pro studenty. Tento nástroj pomáhá studentům sledovat mluvené informace a zároveň slouží jako zápis při pozdější revizi \cite{kb_ai_101_tipu_pdf_04e4a115_page_11_1}. (general background knowledge: konkrétní kroky zapnutí přepisu závisí na použité platformě a nastavení IT)
  \item Asistované čtení a adaptace materiálů: nasazení režimu čtení nahlas a úprava zobrazení textu (kontrast, velikost, rozdělení vět) pro studenty s různými potřebami zvyšuje srozumitelnost a zapojení; doplňte originální text přepisem nebo poznámkami pro snazší orientaci \cite{kb_ai_101_tipu_pdf_04e4a115_page_11_1}. (general background knowledge: konkrétní nastavení v aplikacích Word/Edge/Immersive Reader)
  \item Interaktivní jazykové aktivity v herním prostředí: využijte předpřipravené světy v Minecraft Education pro praktické konverzační úkoly a simulace; některé hodiny pro Hodinu kódu jsou dostupné zdarma, plná funkcionalita vyžaduje školní licenci Microsoft 365 A3/A5 \cite{kb_ai_101_tipu_pdf_04e4a115_page_15_2}.
  \item Krátké konverzační bloky a nácvik výslovnosti: rozdělte výuku do opakovaných krátkých bloků (5–10 min) zaměřených na opakování frází a zpětnou vazbu; nahrávky a automatické přepisy pomáhají sledovat pokrok.
\end{enumerate}

Krátký pracovní checklist pro nasazení v kurzu
\begin{itemize}
  \item Ověřte dostupnost a licenci nástrojů ve vašem školním prostředí (např. Microsoft 365 školní balíčky obsahují některé didaktické aplikace) \cite{kb_ai_101_tipu_pdf_04e4a115_page_15_2}.
  \item Informujte studenty o cílech použití nástrojů (přístupnost, podpora učení, nikoli náhrada studia) — zařaďte stručnou poznámku do sylabu a při prvním použití vysvětlete účel. (general background knowledge)
  \item Zajistěte souhlas a poučení o ochraně osobních údajů tam, kde se nahrává mluvený projev nebo sdílí přepis (např. při záznamu přednášky); uveďte dobu uchování záznamů. (general background knowledge)
  \item Připravte alternativní cestu přístupu pro studenty, kteří nástroj nemohou nebo nechtějí používat (PDF s kontrolním přepisem, offline materiály nebo osobní konzultace). (general background knowledge)
  \item Otestujte technické nastavení před hodinou (mikrofon, připojení, práva sdílení) a mějte záložní plán pro případ technických potíží.
\end{itemize}

Krátké příklady aktivit (možno vložit do hodiny)
\begin{itemize}
  \item „Live captions + otázky“: Pusťte krátkou ukázku (5–7 min), aktivujte přepis a nechte studenty vybrat klíčová slovesa či fráze; následně diskutujte rozdíly v překladu a v kontextu. (general background knowledge)
  \item „Role play v Minecraftu“: studenti v malých skupinách plní scénář (nákup, orientace ve městě) v cílovém jazyce; učitel hodnotí komunikační strategii a porozumění. (general background knowledge, doplněno z dostupných modulů Minecraft Education) \cite{kb_ai_101_tipu_pdf_04e4a115_page_15_2}
  \item „Výslovnostní kroužek“: krátká individuální cvičení s nahráváním a následnou srovnávací analýzou přepisu; studenti obdrží cílenou zpětnou vazbu na konkrétní fonetické chyby.
\end{itemize}

Poznámka o nástrojích mimo tento výčet: v praxi se často používají i komerční překladače, konverzační agenti a aplikace pro nácvik výslovnosti; tyto nástroje si vyžadují ověření z hlediska přesnosti, ochrany osobních údajů a kompatibility s univerzitními licencemi (general background knowledge).

Doporučení na závěr: začněte s jednoduchým pilotem (např. zapnutí přepisu u jedné přednášky nebo využití jedné hodiny v Minecraft Education), zaznamenejte zkušenosti studentů a upravte postupy podle zpětné vazby. Postupné škálování se doporučuje s ohledem na přístupnost, ochranu soukromí a didaktický přínos \cite{kb_ai_101_tipu_pdf_04e4a115_page_15_2}.%

%
\section{Accessibility a asistivní technologie}%
Řešení pro zlepšení přístupu a inkluze: čtení nahlas a zvýraznění textu (Microsoft Immersive Reader), aplikace pro nevidomé (SeeingAI), převod textu na řeč a diktování (Speechify, hlasové diktování ve Wordu) a scénáře využití ve výuce.%


%
\begin{flushleft}
Krátké shrnutí a cíl sekce: Tato část popisuje praktická asistivní řešení a postupy, které zlepšují přístupnost výuky pro studenty se speciálními vzdělávacími potřebami (přepis mluveného slova, čtení nahlas, popis obrazového obsahu). Uvedené nástroje a scénáře jsou vybrány podle jejich praktické použitelnosti ve výuce a podle zdrojů popisujících reálné nasazení ve školním prostředí \cite{kb_ai_101_tipu_pdf_04e4a115_page_11_1, kb_ai_101_tipu_pdf_04e4a115_page_7_1}.
\end{flushleft}

\bigskip

Klíčové nástroje a jejich pedagogické role:
\begin{itemize}
  \item Režimy čtení nahlas a asistivní čtečky (např. Microsoft Immersive Reader) — zjednodušují čtení tím, že zvýrazňují text, upravují zobrazení a čtou obsah nahlas; jsou dostupné v mnoha školních aplikacích a pomáhají studentům se čtením a porozuměním \cite{kb_ai_101_tipu_pdf_04e4a115_page_11_1}.
  \item Aplikace pro osoby se zrakovým postižením (SeeingAI) — dokáže v reálném čase popsat objekt či scénu a tím pomáhá nevidomým studentům orientovat se v materiálech a prostředí \cite{kb_ai_101_tipu_pdf_04e4a115_page_11_1}.
  \item Přepis řeči (speech‑to‑text) a překlad v reálném čase (např. Microsoft Translator jako „přepisovač“) — umožňují v hodině zobrazit textový přepis mluveného slova pro studenty se sluchovým postižením nebo pro multijazykové publikum \cite{kb_ai_101_tipu_pdf_04e4a115_page_11_1}.
  \item Funkce školních edukačních aplikací pro sledování pokroku čtení — „Pokrok ve čtení“ v rámci školních nástrojů umí sledovat výkonnost čtení a poskytovat personalizovaná cvičení na základě chyb \cite{kb_ai_101_tipu_pdf_04e4a115_page_7_1}.
\end{itemize}

Praktické scénáře a doporučené postupy:
\begin{enumerate}
  \item Živý přepis přednášky (podpůrný scénář): aktivujte přepis v dostupné aplikaci nebo sdílejte odkaz do Microsoft Translator; studenti vidí přepis v reálném čase na svých zařízeních, což usnadňuje porozumění i pozdější revizi \cite{kb_ai_101_tipu_pdf_04e4a115_page_11_1}.
  \item Asistované čtení materiálů: zařaďte režim čtení nahlas (Immersive Reader) při práci s delšími texty; zvýrazněný a čtený text pomáhá studentům s poruchami čtení lépe sledovat obsah a soustředit se \cite{kb_ai_101_tipu_pdf_04e4a115_page_11_1}.
  \item Personalizované čtecí úlohy: použijte nástroje, které vyhodnocují čtení a generují adaptivní cvičení na slabá místa studenta (příklad: Pokrok ve čtení ve školních edicích nástrojů) \cite{kb_ai_101_tipu_pdf_04e4a115_page_7_1}.
\end{enumerate}

Krátký checklist pro rychlé nasazení v kurzu:
\begin{itemize}
  \item Ověřte dostupnost funkcí v univerzitních licencích a jejich zapnutí v používaných aplikacích (např. Teams / školní aplikace) \cite{kb_ai_101_tipu_pdf_04e4a115_page_7_1}. 
  \item Před první hodinou ohlaste studentům účel použití asistivních nástrojů (přístupnost, podpora učení) a získejte případný souhlas pro nahrávání či ukládání přepisů (general background knowledge).
  \item Před hodinou otestujte technické nastavení (mikrofon, sdílení přepisu, přístupová práva) a mějte záložní plán pro případ výpadku technologie (general background knowledge).
  \item Uvažujte o poskytování alternativních materiálů (PDF s vloženým přepisem, audio verze) pro studenty, kteří nástroj nemohou či nechtějí používat (general background knowledge).
\end{itemize}

Krátké příklady aktivit do hodiny:
\begin{itemize}
  \item „Live captions + otázky“: promítněte krátký výklad (5–7 min) s aktivovaným přepisem; požádejte studenty, aby vyznačili klíčová slovesa nebo fráze a následně otevřete diskusi o rozdílech v porozumění mezi zápisem a mluvením \cite{kb_ai_101_tipu_pdf_04e4a115_page_11_1}. (general background knowledge: konkrétní nastavení přepisu závisí na platformě)
  \item „Orientace s SeeingAI“: nechte studenty ve dvojicích popsat objekt a porovnat vlastní popis se zprávou z aplikace SeeingAI; aktivita ukáže limity a silné stránky strojového popisu \cite{kb_ai_101_tipu_pdf_04e4a115_page_11_1}.
  \item „Krátké čtecí bloky s adaptací“: rozdělte text na krátké části, použijte čtecí režim pro každý blok a po každém bloku dejte cvičení zaměřené na slovní zásobu určenou pro individualizaci (podpůrné cvičení dle nástrojů pro pokrok ve čtení) \cite{kb_ai_101_tipu_pdf_04e4a115_page_7_1}.
\end{itemize}

Poznámky o ochraně osobních údajů a etice (stručně):
\begin{itemize}
  \item Pokud nahráváte nebo ukládáte přepisy či audiozáznamy obsahující osobní údaje, postupujte podle interních pravidel VUT a GDPR‑orientovaných postupů; specifikujte dobu uchování a přístupová práva (general background knowledge).
  \item Transparentně informujte studenty o účelu a rozsahu použití asistivních technologií a nabídněte alternativu pro ty, kteří nástroj nechtějí nebo nemohou využívat (general background knowledge).
\end{itemize}

Doporučení na závěr: začněte s jedním jednoduchým pilotem (např. zapnutí přepisu u jedné přednášky nebo použití čtecího režimu u jednoho textu), zaznamenejte zkušenosti studentů a upravte postupy podle zpětné vazby; nástroje popsané výše jsou především doplňkem pro zlepšení přístupnosti a měly by být integrovány do širší podpory studentů \cite{kb_ai_101_tipu_pdf_04e4a115_page_11_1, kb_ai_101_tipu_pdf_04e4a115_page_7_1}.%

%
\section{AI pro cvičení a procvičování (matematika a kontrola postupů)}%
Praktické aplikace AI pro generování cvičných úloh a vyhodnocení postupů (např. funkce Pokrok v matematice, OneNote pro převod rukopisu a math.microsoft.com) a návrhy, jak je integrovat do domácích úkolů a diagnostiky chyb.%


%
Krátké shrnutí a cíl sekce: Popis praktických nástrojů a pracovních postupů pro generování a vyhodnocení cvičných úloh v matematice s důrazem na sběr postupu řešení a diagnostiku chyb pomocí nástrojů z rodiny Microsoft (Pokrok v matematice, OneNote, math.microsoft.com) \cite{kb_ai_101_tipu_pdf_04e4a115_page_8_1}.

Praktický pracovní postup pro jeden cyklus domácího úkolu
\begin{enumerate}
  \item Výběr oblasti a generování úloh: učitel v rozhraní „Pokrok v matematice“ vybere oblast, pro kterou chce vygenerovat příklady; systém nabídne sadu úloh, ze které učitel vybírá \cite{kb_ai_101_tipu_pdf_04e4a115_page_8_1}.
  \item Konfigurace požadavků na odevzdání: při nastavování zadání lze určit, zda studenti musí nahrát i postup řešení a zda mají studenti hodnotit obtížnost příkladů \cite{kb_ai_101_tipu_pdf_04e4a115_page_8_1}.
  \item Odevzdání a předběžné zpracování AI: studenti vypracují úlohy a přiloží postupy; AI nástroj provede předopravování (označí pravděpodobné chyby, doplní oblast, kde došlo k chybě) a připraví přehled pro učitele \cite{kb_ai_101_tipu_pdf_04e4a115_page_8_1}.
  \item Pedagogická revize a diagnostika: učitel rychle projde označené případy, upraví hodnocení a použije výsledky k cílenému doplnění výuky (skupinová remediace, individuální úkoly) \cite{kb_ai_101_tipu_pdf_04e4a115_page_8_1}.
\end{enumerate}

Hlavní nástroje a funkcionality (rychlý přehled)
\begin{itemize}
  \item Pokrok v matematice — generování cvičných sad, možnost vyžadovat nahrání postupu, předopravování a přehledy o silných/slabých stránkách žáků \cite{kb_ai_101_tipu_pdf_04e4a115_page_8_1}.
  \item OneNote — převod rukopisných rovnic do digitální podoby, včetně zobrazení postupu a podpory výpočtu či grafu \cite{kb_ai_101_tipu_pdf_04e4a115_page_8_1}.
  \item math.microsoft.com — webové rozhraní slučující nástroje pro práci s matematickými úlohami a cvičeními \cite{kb_ai_101_tipu_pdf_04e4a115_page_8_1}.
  \item Rodina Microsoft školních aplikací — tyto funkce jsou součástí ekosystému školních licencí (např. Office 365 / školní varianty), ověřte dostupnost pro vaši instituci \cite{kb_ai_101_tipu_pdf_04e4a115_page_3_1}.
\end{itemize}

Návrhy pro zařazení do domácích úkolů a diagnostiky (praktické tipy)
\begin{itemize}
  \item Vyžadujte nahrání postupu řešení tam, kde je cílem hodnotit postupy a porozumění; v takových úlohách je „postup důležitější než výsledek“ — možnost požadovat postup lze nastavit v Pokrok v matematice \cite{kb_ai_101_tipu_pdf_04e4a115_page_8_1}.
  \item Využijte automatické označení chyb AI jako predikci oblastí, kde studenti potřebují podporu; plánujte krátké remediace založené na těchto datech \cite{kb_ai_101_tipu_pdf_04e4a115_page_8_1}.
  \item Kombinujte digitální převod rukopisu (OneNote) s požadavkem na nahrání fáze řešení — usnadní to automatizované zpracování a archivaci postupů \cite{kb_ai_101_tipu_pdf_04e4a115_page_8_1}.
  \item Před nasazením zkontrolujte dostupnost funkcí ve školních licencích (Office 365 A1/A3/A5 apod.) a informujte studenty o pravidlech zpracování odevzdaných dat \cite{kb_ai_101_tipu_pdf_04e4a115_page_3_1}. (general background knowledge: konkrétní nastavení retence a přístupů nastaví správa IT fakulty.)
\end{itemize}

Krátký checklist před prvním použitím
\begin{itemize}
  \item Ověřit dostupnost Pokrok v matematice, OneNote a příslušných funkcí v používané univerzitní licenci \cite{kb_ai_101_tipu_pdf_04e4a115_page_3_1}.
  \item Rozhodnout, kdy požadovat nahrání postupu (procesní úlohy) a jasně to specifikovat v instrukcích pro studenty \cite{kb_ai_101_tipu_pdf_04e4a115_page_8_1}.
  \item Nastavit, kdo má přístup k datům o výkonu studentů a jak budou použita pro diagnostiku (general background knowledge).
  \item Otestovat celý tok (vygenerovat úlohy, odevzdat testovací řešení, projít AI označeními) před nasazením do skutečné třídy \cite{kb_ai_101_tipu_pdf_04e4a115_page_8_1}.
\end{itemize}

Poznámka na závěr: Popisované postupy a nástroje vycházejí z implementací a funkcí rodiny vzdělávacích nástrojů Microsoft (Pokrok v matematice, OneNote, math.microsoft.com) uvedených v dostupné literatuře \cite{kb_ai_101_tipu_pdf_04e4a115_page_8_1, kb_ai_101_tipu_pdf_04e4a115_page_3_1}. Pro konkrétní technickou integraci a nastavení ochrany dat spolupracujte s IT oddělením fakulty (general background knowledge).%

%
\section{Integrace do studijní podpory: workflow, role učitele a etika}%
Doporučení pro zařazení AI do poradenské a podpůrné infrastruktury: role asistenta vs. učitele, transparentnost použití AI a základní zásady akademické integrity a ochrany osobních údajů (stručné směrnice pro praxi).%


%
Krátké shrnutí: tato sekce nabízí praktický, krok‑za‑krok proces pro zařazení AI do služeb studijní podpory (tutoring, zpětná vazba, diagnostika), rozdílení rolí mezi učitelem a AI‑asistentem a základní etické body, které je třeba před nasazením ošetřit.

Doporučený základní workflow pro nasazení AI v poradenství a podpoře studentů
\begin{enumerate}
  \item Definice cíle a hranic role AI: nejprve stanovte, které konkrétní aktivity má AI provádět (např. generování cvičných materiálů, předběžná kontrola odevzdaných prací, odpovídání na FAQ) a kdy je nutné eskalovat k lidskému pedagogovi. (general background knowledge)
  \item Volba nástrojů a ověření dostupnosti ve školních licencích: preferujte nástroje integrované do univerzitního ekosystému, ověřte dostupnost požadovaných funkcí v školních licencích a možnosti správy přístupů; např. některé funkce rodiny Microsoft (OneNote, Pokrok v matematice, math.microsoft.com) jsou součástí školních licencí a je vhodné předem ověřit jejich dostupnost pro instituci \cite{kb_ai_101_tipu_pdf_04e4a115_page_8_1}. 
  \item Nastavení požadavků na odevzdání a zpracování dat: tam, kde je cílem hodnotit postupy studentů, vyžadujte nahrání postupu (umožňuje to např. nástroj Pokrok v matematice) a nakonfigurujte, jaké výstupy AI smí generovat před pedagogickou revizí \cite{kb_ai_101_tipu_pdf_04e4a115_page_8_1}.
  \item Pilot a ověření end‑to‑end toku: spusťte malý pilot (vybraná skupina kurzů), otestujte generování cvičení, odevzdání, automatické označení problémů a následnou pedagogickou revizi; použijte pilot k doladění pravidel přístupu a retence dat \cite{kb_ai_101_tipu_pdf_04e4a115_page_8_1}.
  \item Integrace diagnostiky do pedagogické práce: využijte předběžné označení chyb AI pro rychlou identifikaci oblastí potřeby remediace a plánování cílených intervencí, přičemž finální hodnocení provádí učitel \cite{kb_ai_101_tipu_pdf_04e4a115_page_8_1}.
  \item Škálování a automatizace rutinních úloh: tam, kde je to vhodné, lze využít nástroje typu Microsoft 365 Copilot v kombinaci s Forms pro rychlé generování cvičných testů nebo kvízů; vždy však ponechte kontrolu nad finálními hodnotami a vysvětlením správných odpovědí \cite{kb_ai_101_tipu_pdf_04e4a115_page_36_1}.
\end{enumerate}

Praktické role a odpovědnosti (kratce)
\begin{itemize}
  \item Učitel / metodik: definice didaktických cílů, ověření validity materiálů generovaných AI, pedagogická revize výsledků a rozhodnutí o závažnosti zásahů do hodnocení. (general background knowledge)
  \item AI‑asistent (systém): automatické generování cvičení, předběžné označení chyb, sumarizace častých nedostatků pro učitele; systém by měl být nastaven tak, aby neprováděl konečné hodnocení bez lidské kontroly \cite{kb_ai_101_tipu_pdf_04e4a115_page_8_1}.
  \item IT / správa dat: zajištění licencí, přístupových práv, zálohování a bezpečného zpracování dat studentů (general background knowledge; viz doporučení ověření funkcí v licencích) \cite{kb_ai_101_tipu_pdf_04e4a115_page_8_1}.
\end{itemize}

Krátký checklist před spuštěním služby (praktické položky)
\begin{itemize}
  \item Ověřit dostupnost požadovaných funkcí (např. nahrávání postupu, převod rukopisu, předopravování) ve školních licencích a nastavit odpovídající přístupy \cite{kb_ai_101_tipu_pdf_04e4a115_page_8_1}.
  \item Nastavit jasná pravidla pro to, kdy AI pouze pomáhá (draft/sugesce) a kdy je vyžadována lidská validace (hodnocení, disciplinární rozhodnutí). (general background knowledge)
  \item Otestovat celý tok (vygenerovat úlohy, simulovat odevzdání, projít AI označeními a pedagogickou revizí) v pilotním běhu \cite{kb_ai_101_tipu_pdf_04e4a115_page_8_1}.
  \item Zajistit záznam o tom, kdo má přístup k výsledkům a jak jsou data uchovávána/mazána (general background knowledge).
\end{itemize}

Základní etické a komunikační požadavky (stručně, doporučené)
\begin{itemize}
  \item Transparentnost vůči studentům: informujte je o tom, kdy a jak AI zasahuje do podpory a hodnocení (general background knowledge).
  \item Ochrana soukromí: minimalizujte sběr citlivých údajů a domluvte retenci dat se správcem IT/fakulty (general background knowledge).
  \item Pedagogická odpovědnost: AI poskytuje pomocnou vrstvu — konečné pedagogické rozhodnutí a oprávnění k hodnocení zůstává u učitele \cite{kb_ai_101_tipu_pdf_04e4a115_page_8_1}.
\end{itemize}

Tip na závěr: pro aktivity zaměřené na seznámení a motivaci studentů (např. krátké workshopy nebo hry přibližující principy AI) lze zvážit využití edukativních prostředí jako Minecraft Education, které se v praxi využívají k demonstraci principů AI a Hodiny kódu; ověřte dostupnost v institucionálních licencích \cite{kb_ai_101_tipu_pdf_04e4a115_page_15_2}.%

%
\section{Checklisty, šablony promptů a příklady použití}%
Praktické checklisty pro nasazení AI v podpoře studentů, sbírka šablon promptů pro běžné situace (zpětná vazba, nápověda, adaptivní úlohy) a krátké ukázky osvědčených scénářů z praxe.%


%
Tato sekce shrnuje krátké checklisty pro rychlé nasazení AI, praktické prompt‑šablony a tři stručné ukázky pracovních scénářů vhodné pro pilot.

\subsubsection*{Krátký checklist před spuštěním služby}
\begin{itemize}
  \item Ověřit dostupnost funkcí a nastavit přístupy/role administrátorů \cite{kb_ai_101_tipu_pdf_04e4a115_page_8_1}.
  \item Definovat rozsah činností AI, pravidla lidské validace a eskalace; nastavit pravidla uchovávání a přístupu k datům.
  \item Spustit malý end‑to‑end pilot (generování → odevzdání → AI předopraví → pedagogická revize) a doladit retenci dat \cite{kb_ai_101_tipu_pdf_04e4a115_page_8_1}.
  \item Ověřit licence pro doplňkové zdroje (např. Minecraft Education, Hour of Code) pokud budou použity \cite{kb_ai_101_tipu_pdf_04e4a115_page_15_2}.
\end{itemize}

\subsubsection*{Rychlé šablony promptů (copy‑paste)}
\paragraph{1) Systémový prompt — formativní zpětná vazba}
\begin{verbatim}
SYSTEM: Asistent pro formativní zpětnou vazbu.
Vstupy: student_submission, rubric (3-5 bodů), declared_AI_use.
Výstup: 200-250 slov:
  1) 2 věty povzbuzení
  2) 2 konkrétní kroky ke zlepšení
  3) 1 doporučený milník
Pokud declared_AI_use == "ano", přidej 1 větu co ověřit.
\end{verbatim}

\paragraph{2) Deklarace použití AI (do instrukcí)}
\begin{verbatim}
Deklarace AI:
 - Použil/a jsem AI: Ano/Ne
 - Pokud Ano: nástroj a role (nápověda/editace/generování)
 - Jak ověřeno (max. 50 slov)
 - Sebereflexe (max. 200 slov)
\end{verbatim}

\paragraph{3) Generování adaptivních cvičení}
\begin{verbatim}
Vygeneruj 10 cvičení: lineární rovnice (úroveň: 2. roč. BS).
 - 6 krátkých výpočtů, 4 slovní úlohy.
 - U každého požaduj nahrání postupu.
 - Označ obtížnost (lehký/standardní/obtížný).
\end{verbatim}
(Možnost vyžadovat postup a AI předopravit odpovědi; funkce Vzdělávacích akcelerátorů \cite{kb_ai_101_tipu_pdf_04e4a115_page_8_1}.)

\paragraph{4) Rychlý kvíz (Copilot + Forms)}
\begin{verbatim}
Vytvoř test 10 otázek "Austrálie" pro 8. roč. ZŠ:
 - 4 ABC, 4 pravda/nepravda, 2 otevřené.
 - U každé napiš správnou odpověď a krátké vysvětlení.
\end{verbatim}
(Praktické použití s Microsoft 365 Copilot ve Forms \cite{kb_ai_101_tipu_pdf_04e4a115_page_36_1}.)

\subsubsection*{Krátké ukázky pracovních scénářů}
\paragraph{A) Adaptivní matematická praxe}
Učitel nastaví generátor (téma, obtížnost, požadavek na postup), studenti odevzdají postup; AI předběžně označí chyby a shrne časté potíže, učitel provede finální kontrolu \cite{kb_ai_101_tipu_pdf_04e4a115_page_8_1}.

\paragraph{B) Rychlá in‑class písemka}
Učitel vygeneruje a upraví písemku, nasdílí v LMS; Copilot přidá vysvětlení a zrychlí přípravu \cite{kb_ai_101_tipu_pdf_04e4a115_page_36_1}.

\paragraph{C) Seznámení s principy AI}
Krátké workshopy využijí Minecraft Education a Hour of Code pro demonstraci principů ML a interaktivní seznámení studentů \cite{kb_ai_101_tipu_pdf_04e4a115_page_15_2}.

\subsubsection*{Závěrečné doporučení}
Přizpůsobte šablony místnímu kurikulu a vždy validujte výstupy AI lidskou kontrolou.%

%
\chapter{PŘÍPADOVÉ STUDIE Z UNIVERZIT (OVĚŘENÉ PŘÍKLADY)}%
Kapitola obsahuje případové studie s jasnou strukturou (instituce, kurz/prostředí, AI řešení, implementace, výsledky, lesson learned, odkazy). Zahrnuje ověřené příklady: Georgia Tech – AI teaching assistant „Jill Watson“ v online kurzu (IBM Watson) s prokázaným zrychlením odpovědí na dotazy; Harvard CS50 – AI bot pro debugging a stylovou zpětnou vazbu s jasnými pravidly použití; Arizona State University – kampusové nasazení ChatGPT Enterprise s fakultním školením a knihovnou promptů; UCL a TU Delft – veřejně dostupné design patterny a příklady pro tvorbu AI‑aware zadání a hodnocení. Případové studie jsou vybrány tak, aby byly snadno přenositelné do kontextu VUT.%
\section{Úvod a metodika případových studií}%
Popisované struktury případových studií (instituce, kurz/prostředí, AI řešení, implementace, výsledky, lessons learned, odkazy) a metodika výběru příkladů tak, aby byly přenositelné na VUT. Pokyny pro čtení kapitoly a jak využít jednotlivé sekce při redesignu předmětu.%


%
% Úvod a metodika případových studií
% Poznámka: níže uvedená metodika obsahuje praktická doporučení; část obsahu je obecné metodické doporučení (general background knowledge).

Tato sekce vysvětluje, jak jsou jednotlivé případové studie strukturovány a jak je číst tak, aby byly použitelné při redesignu předmětů na VUT. Hlavním cílem je poskytnout čtenáři jasný checklist pro posouzení přenositelnosti jednotlivého příkladu do lokálního kontextu a krok‑za‑krokem pracovní postup pro pilotní adaptaci (general background knowledge).

\begin{itemize}
  \item Struktura každé případové studie (pole, která vždy najdete):
    \begin{itemize}
      \item Instituce a kontext kurzu (velikost kurzu, formát: prezenčně/online/hybrid).
      \item Cíl a pedagogické metriky (co se měřilo).
      \item AI řešení (co bylo použito: nástroj, konfigurace, integrace).
      \item Implementace (fáze nasazení, role lidí, technická integrace).
      \item Výsledky a evidence (kvantitativní a kvalitativní poznatky).
      \item Lessons learned (co fungovalo/na co si dát pozor).
      \item Odkazy a zdroje pro další čtení.
    \end{itemize}
  \item Jak číst kazuistiku — rychlý průvodce pro garanta předmětu:
    \begin{enumerate}
      \item Identifikujte, které části řešení jsou technické (model, licence, integrace) a které pedagogické (úkoly, hodnocení, procesní důkazy).
      \item Hledejte konkrétní měřitelné výsledky nebo indikátory úspěchu, které lze reprodukovat ve vašem předmětu (např. čas na opravu, změna engagementu apod.) (general background knowledge).
      \item Zmapujte rizika a mitigace popsané v kazuistice (data, soukromí, akademická integrita).
      \item Zvažte nutné kroky pro lokalizaci: překlad, adaptace zadání na lokální data, ověření licencí a kompatibility s IT prostředím VUT.
    \end{enumerate}
\end{itemize}

Praktické ukázky nástrojů v kazuistikách často obsahují příklady edukačních aplikací a vzdělávacích akcelerátorů; při posuzování jejich přenositelnosti věnujte pozornost, zda řešení vyžaduje školní licence nebo specifické integrační kroky. Jako ilustraci typu nástrojů, které se v praxi používají, lze uvést Minecraft Education pro didaktické světy a Hodinu kódu (viz popis využití v materiálu) \cite{kb_ai_101_tipu_pdf_04e4a115_page_15_2}. Podobně jsou v pedagogických scénářích zmíněny vzdělávací akcelerátory pro procvičování čtení a matematiky (funkce „Pokrok ve čtení“ a „Pokrok v matematice“) — tyto příklady ukazují, jak může systém předzpracovat a vyhodnotit studentovy odpovědi a navrhnout personalizaci \cite{kb_ai_101_tipu_pdf_04e4a115_page_7_1, kb_ai_101_tipu_pdf_04e4a115_page_8_1}.
  
Krátký checklist pro posouzení přenositelnosti případu:
\begin{itemize}
  \item Je cíl pedagogicky kompatibilní s očekávanými výstupy předmětu? (general background knowledge)
  \item Jaké jsou technické požadavky a jsou splnitelné v prostředí VUT (LMS, SSO, enterprise licence)? (general background knowledge)
  \item Jaké údaje se zpracovávají a je potřeba provést DPIA / DPA před nasazením? (general background knowledge)
  \item Existuje jasný plán pilotu (malý rozsah → validace → škálování) a metrika úspěchu? (general background knowledge)
  \item Jsou ve studii popsané konkrétní „lessons learned“, které eliminují známé rizikové kroky? (general background knowledge)
\end{itemize}

Jak využít kazuistiku při redesignu předmětu — doporučený postup (3 kroky, praktické):
\begin{enumerate}
  \item Extrahujte z kazuistiky komponenty, které chcete replikovat: (a) typ úkolu, (b) role AI, (c) požadavek na důkazy studentovy práce.
  \item Navrhněte malý pilot v rámci jednoho tématu/sem. hodiny: definujte vstupy, výstupy, metriky a lidské validační body (např. pedagog prověří 10\% výstupů).
  \item Vyhodnoťte pilot podle předem definovaných metrik a zdokumentujte „lessons learned“ pro širší nasazení nebo úpravu sylabu (general background knowledge).
\end{enumerate}

Na závěr: při čtení kapitol s případovými studiemi hledejte konkrétní technické a pedagogické body, které lze vyzkoušet v malém měřítku; využijte přiložené checklisty této kapitoly jako šablonu pro vlastní pilotní plán (general background knowledge).%

%
\section{Georgia Tech — 'Jill Watson' (AI teaching assistant)}%
Kazuistika zaměřená na použití AI asistenta v online kurzu: popis řešení, implementační kroky, dosažené efekty a klíčové ponaučení. General background knowledge: zahrnuje příklad asistenta 'Jill Watson' používaného v kurzech (AI teaching assistant) a zmínku o pozitivním vlivu na rychlost odpovědí studentům.%


%
% Poznámka: popis této kazuistiky obsahuje obecné shrnutí případu 'Jill Watson' a praktická doporučení pro přenositelnost; konkrétní historické nebo kvantitativní údaje uvedené níže jsou označeny jako general background knowledge, pokud nejsou doloženy v repozitáři příručky.
\medskip

\textbf{Shrnutí případu (stručně)}\\
Příklad „Jill Watson“ ilustruje nasazení AI jako virtuálního výukového asistenta v rozsáhlém online kurzu (general background knowledge). Takoví asistenti obvykle zpracovávají FAQ, nasměrovávají studenty na materiály a snižují administrativní zátěž učitelů, což zkracuje dobu odezvy (general background knowledge).

\bigskip
\textbf{Struktura kazuistiky (pole, která doporučujeme vždy zmapovat)}\\
V souladu s metodikou zahrnuje analýza obvykle tato pole:
\begin{itemize}
  \item Instituce a kontext kurzu (velikost, formát: prezenčně/online/hybrid).
  \item Cíl a pedagogické metriky (jaký konkrétní problém má AI řešit).
  \item AI řešení (typ asistenta, omezení kompetencí, integrační body).
  \item Implementace (fáze nasazení, role lidských aktérů, technická integrace s LMS/SSO).
  \item Výsledky a evidence (měřené ukazatele úspěchu).
  \item Lessons learned (klíčová opatření při provozu).
\end{itemize}

\bigskip
\textbf{Konkrétní body k zamyšlení pro případ „Jill Watson“ (general background knowledge)}\\
\begin{itemize}
  \item Pedagogický cíl: zlepšit dostupnost podpory v asynchronních kurzech a snížit opakující se administrativní dotazy.
  \item Role asistenta: FAQ, odkazy na materiály, eskalace složitých dotazů; nesmí samostatně hodnotit zkoušky ani uzavírat otázky integrity.
  \item Technická integrace: připojení k LMS/diskuzím, logging konverzací pro audit, SSO a posouzení DPA/DPIA při zpracování osobních údajů.
  \item Lidské procesy: plán lidské validace, mechanismus eskalace a proces aktualizace znalostní báze.
\end{itemize}

\bigskip
\textbf{Doporučený postup pro pilotní nasazení na úrovni předmětu (general background knowledge)}\\
Postup přizpůsobitelný pro pilot na VUT:
\begin{enumerate}
  \item Vymezit cíl pilotu a metriky (např. doba odpovědi, počet eskalací, spokojenost).
  \item Definovat kompetence asistenta a scénáře použití; specifikovat otázky, které nesmí odpovídat automaticky.
  \item Připravit znalostní bázi a proces její verze/aktualizací.
  \item Implementovat technickou integraci s LMS/SSO, zajistit logging, přístupy a provést bezpečnostní/GDPR kontrolu.
  \item Spustit pilot v omezeném rozsahu, sbírat data a zpětnou vazbu od studentů a vyučujících.
  \item Pravidelně vyhodnocovat metriky, ladit odpovědi a dokumentovat lessons learned; podle kritérií rozhodnout o škálování.
\end{enumerate}

\bigskip
\textbf{Kontrolní checklist pro garanta předmětu při zvažování replikace takového asistenta (general background knowledge)}\\
\begin{itemize}
  \item Je cíl asistenta kompatibilní s výstupy předmětu?
  \item Jsou hranice kompetencí jasně definovány a sděleny studentům?
  \item Jsou zajištěny právní a GDPR požadavky (DPIA, DPA)?
  \item Existuje plán lidské validace, eskalace a nápravy chyb?
  \item Je nastaven proces logování, auditu a aktualizace znalostní báze?
  \item Má IT kapacity pro bezpečnou integraci a monitoring?
\end{itemize}

\bigskip
\textbf{Klíčová „lessons learned“ k adaptaci pro VUT (general background knowledge)}\\
\begin{itemize}
  \item Jasně stanovte, co asistenta smí a nesmí dělat; uvádějte to v sylabu.
  \item Začněte malým pilotem místo plného nasazení pro omezení rizik.
  \item Buďte transparentní vůči studentům o tom, že komunikují s AI, a o logování dotazů.
  \item Měřte technické i pedagogické metriky (odezva, spokojenost, dopad na zátěž vyučujících).
  \item Pracujte s interními daty opatrně; preferujte anonymizaci nebo pseudonymizaci pro analýzy.
\end{itemize}

\bigskip
\textbf{Krátké doporučení pro další kroky}\\
Pro garanta, který uvažuje přenos řešení na VUT:
\begin{itemize}
  \item Konzultovat plán s místním IT a GDPR/právním oddělením (DPIA/DPA).
  \item Naplánovat 6–8týdenní pilot v jedné sekci s jasnými metrikami.
  \item Použít šablonu pro dokumentaci implementace a lessons learned z repozitáře příručky.
\end{itemize}%

%
\section{Harvard CS50 — AI bot pro debugging a stylovou zpětnou vazbu}%
Popis nasazení AI bota v rozsáhlém kurzu CS50: pravidla použití, omezení, workflow pro poskytování debuggingové a stylistické zpětné vazby a zjištěné dopady na výuku a akademickou integritu. General background knowledge: obsahuje příklad CS50 s jasnými pravidly použití AI bota.%


%
Tato kazuistika shrnuje přístup k nasazení AI bota v rozsáhlém kurzu (např. CS50) zaměřeném na debuggingovou a stylistickou zpětnou vazbu studentům; podrobnosti níže jsou uvedeny jako general background knowledge tam, kde nejsou přímo podloženy interními zdroji nebo daty z repozitáře (general background knowledge).

\medskip
\textbf{Cíl a rozsah nasazení (general background knowledge)}\\
Cílem je:
\begin{itemize}
  \item rychlá prvotní zpětná vazba na chyby v kódu (debugging hints);
  \item stylistická a didaktická doporučení (readability, comments, naming);
  \item snížení administrativní zátěže a zrychlení doby odezvy studentům.
\end{itemize}

\medskip
\textbf{Pravidla použití a omezení (doporučené postupy)}\\
Zásady komunikace studentům:
\begin{itemize}
  \item Bot dává nápovědy a diagnostiku, nesmí generovat kompletní řešení pro hodnocené úlohy (general background knowledge).
  \item Studenti deklarují významné použití bota; interakce se logují pro audit s respektováním GDPR (anonymizace/pseudonymizace).
  \item Eskalace: zásadní opravy prověří učitel nebo asistent před uzavřením hodnocení.
\end{itemize}

\medskip
\textbf{Návrhový workflow pro instruktora (krok‑za‑krokem)}\\
\begin{enumerate}
  \item Definujte scénáře použití (debug hint, style suggestion, explain failing test) a pravidla povolení bota (general background knowledge).
  \item Připravte a otestujte šablony promptů a konverzačních rolí na reprezentativních úkolech.
  \item Integrujte s LMS/rozhraním; zajistěte SSO, logging a auditní záznamy.
  \item Spusťte pilot v jedné sekci; sbírejte metriky (počet interakcí, doba do první nápovědy, eskalace, spokojenost).
  \item Vyhodnoťte dopad na učení i integritu a upravte pravidla/prompting podle zjištění. Pozn.: ukázkové generování artefaktů je popsáno i pro jiné nástroje (např. Microsoft 365 Copilot) \cite{kb_ai_101_tipu_pdf_04e4a115_page_36_1}.
\end{enumerate}

\medskip
\textbf{Praktické šablony promptů (příklady)}\\
Vždy zahrňte kontext zadání a očekávanou úroveň odpovědi:
\begin{itemize}
  \item Debug hint (stručně): ``You are a debugging assistant for beginner CS students. Given this code and the failing test output, identify the most likely cause and suggest one concrete minimal change to fix it. Do not provide the full corrected solution; instead describe the change and why it works.''
  \item Explain failing test: ``Explain in simple steps why the test fails, list up to three hypotheses ordered by likelihood, and suggest one small experiment the student can run to narrow down the cause.''
  \item Style suggestion: ``Review the following function for readability and style. Provide up to five concise suggestions (naming, comments, complexity), and show a one‑line example of an improved variable name or comment.''
\end{itemize}
Šablony se ladí tak, aby bot neprodukoval kompletní řešení, pokud to pravidla kurzu zakazují (general background knowledge).

\medskip
\textbf{Dopady na akademickou integritu a mitigace rizik (general background knowledge)}\\
Minimalizujte zneužití pomocí:
\begin{itemize}
  \item požadavku na procesní důkazy (commit history, postupy, krátké video walkthrough);
  \item preferování časově omezených nebo in‑class hodnocení pro závěrečné části;
  \item vyžadování reflexe a deklarace použití AI v odevzdáních.
\end{itemize}

\medskip
\textbf{Měření úspěchu a doporučené metriky (pilot)}\\
Sledujte operační, pedagogické a integritní metriky:
\begin{itemize}
  \item operační: počet interakcí, prům. délka relace, eskalace k lidem;
  \item pedagogické: doba nápravy chyb u domácích úloh, studentská spokojenost;
  \item integritní: počet porušení pravidel a důvody eskalací.
\end{itemize}
Nastavte cíle a prahy před pilotem a průběžně je revidujte.

\medskip
\textbf{Krátké doporučení pro garanta předmětu}\\
Začněte malým, jasně ohraničeným pilotem s definovanými pravidly, dokumentovanou postupností (pilot → vyhodnocení → revize), transparentností vůči studentům a prověřováním zpracování citlivých údajů přes DPIA/DPA, pokud bot pracuje s osobními údaji (general background knowledge).%

%
\section{Arizona State University — kampusové nasazení ChatGPT Enterprise}%
Případ nasazení na úrovni fakulty/kampusu: proces zavedení enterprise licence, školení pedagogů, knihovna promptů a organizační zajištění provozu. General background knowledge: zahrnuje zmínku o kampusovém nasazení ChatGPT Enterprise a souvisejícím fakultním školení.%


%
\medskip
\textbf{Arizona State University — kampusové nasazení ChatGPT Enterprise (general background knowledge)}\\
Tato kazuistika shrnuje klíčové komponenty kampusového nasazení komerční enterprise služby a slouží jako vstupní šablona pro adaptaci na VUT.
\medskip
\textbf{Cíl a rozsah nasazení (general background knowledge)}\\
Cílem je nabídnout centrálně spravovaný přístup k enterprise verzi (licence, SSO, audit), zrychlit formativní zpětnou vazbu a vytvořit sdílenou knihovnu ověřených promptů.
\medskip
\textbf{Pravidla použití a právně‑etické požadavky (general background knowledge)}\\
Doporučené principy: centrální smlouva a DPA (data residency, zpracování osobních údajů); SSO a role‑based přístup; povinnost oznámit použití AI v sylabu a deklarovat významné využití.
\medskip
\textbf{Návrhový workflow pro fakulty při zavádění (general background knowledge)}\\
Doporučený postup: (1) posoudit potřeby a scénáře (tutoring, grading help, admin‑bot, tvorba materiálů); (2) zajistit smluvní/IT souhlas a SSO; (3) pilotovat v několika kurzech s metrikami a eskalací; (4) školení pedagogů, vytvoření prompt‑knihovny a rozšíření podle výsledků.
\medskip
\textbf{Školení a podpora pedagogů (general background knowledge)}\\
Klíčové prvky: krátké workshopy „how to prompt“ a validace výstupů; ukázky integrace do LMS (logování, anonymizace); FAQ, centrální repozitář promptů a pravidelné webináře.
\medskip
\textbf{Knihovna promptů — ukázkové šablony (general background knowledge)}\\
Přizpůsobitelné příklady: \\
\begin{itemize}
  \item \textbf{Formativní zpětná vazba:} stručné silné stránky, 2–3 akční návrhy ke struktuře a jasnosti, 1 otázka pro instruktora; bez hodnocení.
  \item \textbf{Debugging pro CS:} až 3 pravděpodobné příčiny a jeden minimální experiment pro izolaci problému; neposkytovat kompletní řešení.
  \item \textbf{In‑class kvíz:} 10 krátkých MCQ s jednou správnou odpovědí a jedno‑větým vysvětlením.
\end{itemize}
\medskip
\textbf{Monitoring, metriky a vyhodnocení pilotu (general background knowledge)}\\
Doporučené metriky: operační (aktivní uživatelé, interakce, doba odezvy, eskalace); pedagogické (spokojenost, doba do první zpětné vazby, kvalita dle lidské kontroly); compliance (shoda s DPA/DPIA, incidenty úniku dat nebo nevhodného obsahu).
\medskip
\textbf{Krátké doporučení pro garanta předmětu (general background knowledge)}\\
Začněte s jasně definovaným scénářem použití (např. pouze formativní zpětná vazba), zdokumentujte pravidla pro studenty (sylabus, deklarace) a nastavte ověřování výstupů (losovací kontrola, lidská validace u kritických momentů).
\medskip
\textit{Poznámka:} Tento text shrnuje praktická doporučení jako general background knowledge — konkrétní právní, technické a smluvní kroky musí vždy projít interním ověřením právního oddělení a IT na úrovni VUT před nasazením.%

%
\section{UCL — design patterns pro AI‑aware zadání a hodnocení}%
Shrnutí veřejně dostupných design patternů a příkladů pro tvorbu zadání a hodnoticích postupů, které reflektují přítomnost generativní AI ve výuce; konkrétní patterny pro formulaci úloh, transparentnost a validaci výsledků.%


%
\textbf{(general background knowledge)}\par
Tato sekce shrnuje opakovatelné design patterny pro zadání a hodnocení reflektující generativní AI; každý pattern obsahuje stručný postup nasazení a návrh kritérií. (general background knowledge)

\medskip
\textbf{Hlavní design patterny}
\begin{description}
  \item[1) Process‑first assignment (hodnotí se proces):] \hfill \\
    Zdůrazňuje kroky, mezivýstupy a odůvodnění; nasazení: inkrementální odevzdání a komentáře; rubrika: sledovatelnost/odůvodnění/validace/reflexe.
  \item[2) Artifact‑based / local‑data binding:] \hfill \\
    Úlohy vázané na lokální data/experimenty; nasazení: přiložte datový zdroj; validace: metadata/hashy, sanity checks.
  \item[3) AI‑declared + reflective submission:] \hfill \\
    Student deklaruje AI, přiloží prompty a reflexi; nasazení: povinný „AI log“ a hodnocení reflexe.
  \item[4) Milestone‑based assessment (kontrolní body):] \hfill \\
    Úloha rozdělena do milníků s kritérii; nasazení: 3–4 milníky (návrh, prototyp, testy, závěr) a kontrola konzistence.
  \item[5) Oral spot‑check / viva (živá verifikace):] \hfill \\
    Krátké ústní ověření zaměřené na proces; nasazení: náhodný výběr 10–20 % prací pro 10–15min ověření.
  \item[6) Controlled RAG resources (řízené využití interních materiálů):] \hfill \\
    AI povoleno pouze nad omezenou kolekcí; nasazení: poskytněte RAG index a vyžadujte citace.
\end{description}

\medskip
\textbf{Praktické šablony (krátké příklady textů pro sylabus a zadání)}\\
\textbf{Text do sylabu (vzor):} „Použití generativní AI je dovoleno pouze při dokumentování jeho použití. Při odevzdání uveďte: (1) nástroj a verzi, (2) přesné prompty/vstupy, (3) stručnou reflexi, co jste převzali a co upravili. Nepovolené je předkládat práci bez deklarace.“ (general background knowledge)

\textbf{Krátký prompt‑log (šablona pro odevzdání):}
\begin{itemize}
  \item Nástroj / verze; Prompt / vstup (přesný text);
  \item Co AI vygenerovala; Jak jsem výstup upravil/zkontroloval;
  \item Reflexe: v čem AI pomohla, jaké byly limity.
\end{itemize}
(použít jako přílohu k odevzdání) (general background knowledge)

\medskip
\textbf{Rychlý checklist pro garanta předmětu při adaptaci patternů}
\begin{enumerate}
  \item Vyberte 1–2 patterny vhodné pro kurz.
  \item Aktualizujte sylabus o deklaraci AI a pravidlech hodnocení.
  \item Navrhněte rubriky pro proces, validaci a reflexi o AI.
  \item Zaveďte technické požadavky pro ověření artefaktů (metadata, logy, hashes).
  \item Naplánujte verifikaci (ústní kontroly, peer review, human check).
  \item Komunikujte očekávání a poskytněte ukázkové „good practice“.
\end{enumerate}
(vše výše uvedené je general background knowledge) Sekce nabízí šablony k vložení do sylabů a zadání; doporučujeme pilotovat a zohlednit GDPR/právní požadavky (DPIA/DPA) před rozšířením. (general background knowledge)%

%
\section{TU Delft — šablony a ověřené příklady pro redesign předmětů}%
Kompilace konkrétních šablon a ukázek redesignu předmětů s AI‑aware komponentami: zadání, assessment, formative feedback a technické integrace. General background knowledge: obsahuje odkaz na TU Delft jako zdroj veřejně dostupných příkladů a šablon.%


%
Tato sekce nabízí krátkou sbírku praktických šablon a ukázek pro redesign předmětů s AI‑aware komponentami; upravte podle kurzu a pravidel instituce.

\bigskip

\textbf{Krátký text do sylabu (vzor)}
\begin{quote}
Použití generativní AI: AI nástroje povoleny pouze pokud student při odevzdání přiloží „AI log“ a stručnou reflexi. Uveďte: (1) nástroj a verzi, (2) přesné prompty/vstupy, (3) co AI vygenerovala a jak jste výstup upravili, (4) krátkou reflexi o přínosu a limitech. Neúplná deklarace může snížit hodnocení.
\end{quote}

\bigskip

\textbf{Milníky (template)}
\begin{enumerate}
  \item M1 Návrh (týd.2): cíl, postup, zdroje. Hodnocení: jasnost plánu.
  \item M2 Prototyp (týd.4): funkční část, logy včetně AI logu. Hodnocení: reprodukovatelnost.
  \item M3 Testy (týd.8): metriky, ověření dat. M4 Závěr (týd.12): finální artefakt + 1–2 s. reflexe.
\end{enumerate}

\bigskip

\textbf{Prompt‑log (šablona)}\\
\begin{PromptBlock}


- Nástroj/verze; - Přesný prompt; - Co AI vygenerovala; - Jak jsem upravil/ověřil; - Reflexe (2–4 odstavce).
\end{PromptBlock}



\bigskip

\textbf{Rubrika a checklist pro lektora}\\
Rubrika: Proces 40\% (kroky, zdůvodnění, logy), Reprodukovatelnost 25\% (data, metadata, hashe), Reflexe AI 20\%, Konečný artefakt 15\%.\\
Checklist: vymezit proces vs produkt; aktualizovat sylabus; navrhnout 3–4 milníky; zajistit technické požadavky (metadata, formát logu, volitelné hashování); provést pilot.%

%
\section{Přenositelnost na VUT: konkrétní doporučení, checklisty a šablony}%
Praktický přenos poznatků z kazuistik do prostředí VUT: krok‑za‑krokem doporučení pro implementaci, kontrolní seznamy pro garanty předmětů, šablony promptů a scénáře pilotního nasazení v rámci fakulty.%


%
Tato sekce shrnuje konkrétní, opakovatelné kroky pro přenos poznatků z kazuistik do praxe na VUT. Obsahuje krátký pilotní plán, checklist pro garanty předmětů a hotové šablony (text do sylabu, milníky, prompt-log, základní rubrika), které lze okamžitě vložit do kurzu.

\bigskip

Rychlý pilot (6 kroků)
\begin{enumerate}
  \item Definice cíle a rozsahu: vyberte 1 předmět nebo tematický modul, jasně formulujte, co chcete zlepšit (zpětná vazba, tvorba materiálů, diferenciace).
  \item Právní a GDPR check: konzultujte se školním GDPR/PRIV kontaktem; rozhodněte o minimalizaci osobních údajů, anonymizaci dat a získání souhlasů účastníků.
  \item Návrh intervencí: zvolte konkrétní aktivity (nové úlohy, vzorové hodiny, rubrika hodnocení), určete zdroje a odpovědnosti.
  \item Příprava materiálů: připravte sylabus s vloženým textem, milníky kurzu a prompt-log; použijte dodané šablony pro jednotnost.
  \item Pilotní nasazení: spusťte pilot během jednoho semestru, sbírejte průběžnou zpětnou vazbu od studentů a vyučujících, dokumentujte změny.
  \item Vyhodnocení a škálování: analyzujte výsledky, upravte materiály podle zjištění a naplánujte rozšíření na další předměty.
\end{enumerate}

\bigskip

Checklist pro garanty předmětů (rychlé body)
\begin{itemize}
  \item Potvrdit cíl pilotu a získat podporu katedry/fakulty.
  \item Zajistit GDPR/PRIV konzultaci a formuláře souhlasu tam, kde je to nutné.
  \item Zkontrolovat, že sylabus obsahuje stanovené milníky a popis hodnocení.
  \item Nasadit prompt-log pro záznam interakcí a iterací úloh.
  \item Použít základní rubriku pro konzistentní hodnocení, nastavit škálu a kritéria.
  \item Naplánovat sběr dat a termíny pro vyhodnocení pilotu.
  \item Informovat studenty o cílech pilotu a způsobu využití dat.
\end{itemize}

\bigskip

Hotové šablony k okamžitému vložení do kurzu (s krátkým popisem)
\begin{itemize}
  \item Text do sylabu: stručné znění cíle kurzu, očekávané výstupy, milníky a pravidla hodnocení.
  \item Milníky: časová osa s kontrolními body pro odevzdání a zpětnou vazbu.
  \item Prompt-log: formulář pro záznam dotazů, verzí úloh a použitých promptů.
  \item Základní rubrika: jednoduchá tabulka kritérií a popisů úrovní výkonu pro rychlé hodnocení.
  \item Doporučený plán komunikace: vzor e-mailu pro informování studentů a instrukce pro vyučující.
\end{itemize}%

%
\chapter{INSTITUCIONÁLNÍ NASAZENÍ, INTEGRACE A BEZPEČNOST}%
Kapitola nahrazuje obecný přehled nástrojů a zaměřuje se na praktické zavádění AI na úrovni fakulty či univerzity: výběr enterprise variant (ChatGPT Enterprise, Azure OpenAI Service, Claude for Enterprise, Gemini pro Google Workspace, Microsoft Copilot pro M365), SSO, správa oprávnění, auditní logy, DPA a data residency. Popisuje integrační scénáře do LMS (Moodle), knihoven a vývojářského prostředí, architekturu RAG s interními materiály, možnost lokálního/edge nasazení modelů (např. Llama, Mistral) a řízení nákladů (limity, rozpočty, token metriky). Zahrnuje checklist pro IT a katedry, vzory RFP/SLA, governance (Data Stewards, AI Council) a bezpečnostní opatření (guardrails, filtrování obsahu, red‑teaming).%
\section{Strategie výběru enterprise řešení a licencí}%
Kritéria pro výběr enterprise variant (např. ChatGPT Enterprise, Azure OpenAI, Claude for Enterprise, Gemini, Microsoft Copilot), porovnání podle DPA/SLA, integrace do existujícího stacku, business‑case a rozhodovací matice pro fakultu/univerzitu.%


%
Tato sekce přináší praktický, krok‑za‑krokem přístup pro výběr enterprise řešení a licencí pro potřeby fakulty nebo univerzity. Následující doporučení jsou formulována jako pracovní checklist a rozhodovací matice, které lze použít při jednání s poskytovateli i při přípravě business‑case pro vedení školy. Obsahuje jak obecné zásady, tak konkrétní praktické kroky k ověření klíčových smluvních a technických podmínek. (general background knowledge)

\bigskip

\textbf{Rychlý postup výběru (6 kroků)} (general background knowledge)
\begin{enumerate}
  \item Vyjasněte požadavky: rozsah uživatelů, typy zpracovávaných dat (osobní / citlivá / interní), požadované integrační body (LMS, SSO, repozitáře).
  \item Identifikujte kandidáty: sestavte krátký seznam poskytovatelů (interní/on‑premise, cloudové enterprise nabídky, specializovaní dodavatelé).
  \item Požádejte o dokumentaci: DPA/SLA, bezpečnostní whitepapers, certifikace, data residency možnosti a pracovní postupy zálohování/mazání.
  \item Proveďte technické PoC: krátký pilotní test integrace s LMS/SSO a prověření běžných pracovních toků (přihlášení, logování, export auditních stop).
  \item Zhodnoťte obchodní podmínky a náklady: základní TCO, licensing modely (per‑user / enterprise / API), běžné limity a škálovatelnost.
  \item Formulujte doporučení a plán nasazení: včetně rolí (Data Steward, IT), harmonogramu, metrik úspěchu a plánů pro DPIA/legální kontrolu.
\end{enumerate}

\bigskip

\textbf{Praktická rozhodovací matice (návrh)} (general background knowledge)

\noindent Pro každý kandidát ohodnoťte následující kritéria 1–5 (1 = slabé, 5 = výborné) a sečtěte:
\begin{itemize}
  \item Soulad s GDPR a možnosti DPA/SLA (právní záruky, uchovávání, audit).
  \item Data residency a kontrola nad materiály (možnost regionálního uložení).
  \item Integrace se stávajícím stackem (SSO, LMS, repozitáře).
  \item Bezpečnost a provozní vlastnosti (šifrování, logování, incident response).
  \item Podpora a SLA (reakční doby, onboarding, školení).
  \item Funkčnost a vhodnost pro pedagogiku (přizpůsobitelné modely, tooling pro učitele).
  \item Celkové náklady (licence + integrace + provoz).
\end{itemize}

\noindent Doporučení: kandidáty seřaďte podle součtu a pak prověřte top‑3 skrze piloty a právní review. (general background knowledge)

\bigskip

\textbf{Klíčové smluvní body k ověření v DPA / SLA} (general background knowledge, stručný checklist)
\begin{itemize}
  \item Role správce vs. zpracovatele a jasně definované odpovědnosti.
  \item Data residency a právo požadovat mazání dat (retence, vymazání na vyžádání).
  \item Pravidla pro práci s osobními a citlivými údaji, včetně subdodavatelů.
  \item Auditní práva a možnosti bezpečnostního testování (red‑teaming).
  \item Záruky dostupnosti a doby odezvy při incidentu (SLA).
  \item Podmínky ukončení smlouvy a exportu dat v čitelném, strojově zpracovatelném formátu.
\end{itemize}

\bigskip

\textbf{Technické aspekty k ověření v PoC} (general background knowledge)
\begin{itemize}
  \item Autentizace/SSO: lze použít univerzitní identitu (SAML/OAuth) bez duplicitních účtů?
  \item Audit a logy: jsou dostupné záznamy aktivit a lze je dlouhodobě ukládat pro compliance?
  \item Integrace s LMS: existují konektory/plug‑iny nebo API pro synchronizaci rozvrhů, uživatelů a artefaktů?
  \item Schopnost používat interní zdroje bezpečně (RAG pattern): napojení na interní dokumenty bez exfiltrace.
  \item Možnosti šifrování (in transit / at rest) a správa klíčů.
\end{itemize}

\bigskip

\textbf{Poznámka k vendorům a zdrojům pro učení se od dodavatelů} \\
Využijte oficiální dokumentaci a vzdělávací portály poskytovatelů pro rychlé nastudování možností produktů; příkladem je portál Microsoft Learn jako centrální, bezplatný zdroj dokumentace a kurzů pro produkty Microsoftu \cite{kb_ai_101_tipu_pdf_04e4a115_page_54_1}. Podobně u vzdělávacích produktů (např. Minecraft Education) existují připravené edukační světy a materiály, které mohou sloužit jako inspirace pro pilotní scény používání AI ve výuce \cite{kb_ai_101_tipu_pdf_04e4a115_page_15_2}. Pro matematické a didaktické nástroje v rámci ekosystému Microsoft jsou relevantní konkrétní funkce (např. „Pokrok v matematice“, OneNote, math.microsoft.com), které ukazují, jak integrace do školního prostředí může vypadat \cite{kb_ai_101_tipu_pdf_04e4a115_page_8_1}.

\bigskip

\textbf{Komentář ke „on‑premise“ vs. cloudovým enterprise řešením} (general background knowledge)
\begin{itemize}
  \item On‑premise / hybridní: vyšší kontrola nad daty a modely, vyšší provozní nároky (HW, zálohy, aktualizace).
  \item Cloud enterprise: rychlejší nasazení, správa aktualizací dodavatelem, ale je potřeba ověřit DPA/data residency a možnosti izolace dat.
\end{itemize}

\bigskip

\textbf{Doporučené výstupy z výběrového procesu} (general background knowledge)
\begin{itemize}
  \item Krátký souhrn (1–2 str.) s porovnáním top‑3 kandidátů dle matice.
  \item Návrh pilotu (scope, metriky, doba trvání, zdroje).
  \item Právní seznam požadavků pro DPA a technický checklist pro PoC.
  \item Návrh rozpočtu TCO na první 3 roky včetně nákladů na integraci a provoz.
\end{itemize}

\bigskip

Poznamenejte, že praktické nasazení vždy zohledňuje místní požadavky na ochranu osobních údajů, interní bezpečnostní politiky a provozní možnosti fakulty — před konečným rozhodnutím doporučujeme konzultaci s IT, právním oddělením a GDPR/PRIV kontaktem fakulty. (general background knowledge)%

%
\section{Autentizace, autorizace a správa identit}%
SSO integrace, správa oprávnění a rolí, provisioning uživatelů, auditní logy a způsoby zaznamenávání přístupů a akcí pro účely bezpečnosti a compliance.%


%
Krátké shrnutí a účel sekce (general background knowledge): tato část popisuje praktické kroky a kontrolní body pro bezpečné nasazení autentizace, autorizace a správy identit při zavádění enterprise AI řešení na úrovni fakulty/univerzity. Důraz je na využití univerzitní identity (SSO), jasné role‑/právové modely (RBAC), provisioning uživatelů a auditování aktivit pro účely bezpečnosti a souladu s předpisy. V konkrétních krocích níže jsou uvedeny ověřitelné testy pro PoC i provozní doporučení (general background knowledge).

Rychlý kontrolní seznam pro PoC a nasazení (general background knowledge)
\begin{itemize}
  \item Ověřte podporu univerzitních identit (SSO): zkontrolujte, zda lze propojit poskytovatele s autentizačním protokolem, který používáte (SAML, OAuth2/OIDC) a zda není potřeba vytvářet lokální duplicate účty.
  \item Provisioning a deprovisioning: vyzkoušejte automatické provisioning flow (SCIM nebo jiný mechanizmus), včetně okamžitého zablokování a mazání oprávnění při ukončení zaměstnaneckého/student‑vztahu.
  \item Role a oprávnění (RBAC): navrhněte minimální sadu rolí (správce, správce dat / Data Steward, garant předmětu, lektor/asistent, student) a ověřte, že každá role má nejmenší nutná oprávnění.
  \item Auditní logy a tracing: zkontrolujte dostupnost auditních záznamů pro přihlášení, změny rolí, přístupy k citlivým materiálům a volání API; ověřte možnost exportu logů pro dlouhodobé ukládání a forenzní analýzu.
  \item Integrace s LMS / dalšími systémy: proveďte test synchronizace uživatelů a rolí do LMS (např. Moodle), ověřte mapping rolí a oprávnění mezi systémy.
  \item Bezpečnostní limity a metriky: definujte metriky (počet neúspěšných přihlášení, počet změn práv, počet přístupů k interním dokumentům) a nastavte alertování pro podezřelé události.
\end{itemize}

Praktické ověřovací testy pro PoC (general background knowledge)
\begin{enumerate}
  \item End‑to‑end SSO: přihlášení uživatele přes univerzitní SSO, ověření atributů (e‑mail, instituce, role) přenesených do poskytovatele, a návrat bez vytvoření duplicitního lokálního účtu.
  \item SCIM / provisioning: vytvoření testovacího uživatele v IdP → automatické vytvoření účtu u poskytovatele → změna role v IdP → ověření propagace změny → deaktivace v IdP → ověření odebrání přístupu.
  \item Role enforcement: pokus uživatele s nižší rolí o akci vyžadující vyšší oprávnění (např. export interních dokumentů, správa instancí) a kontrola, že je akce zamítnuta a zaznamenána v logu.
  \item Auditní integrita: simulace administrativní změny (změna rolí, mazání zdroje) a ověření, že událost má časové razítko, identifikaci uživatele, originační IP a možnost exportu do centralizovaného SIEMu.
  \item Export a retence logů: ověření, že logy lze exportovat v strojově čitelném formátu (JSON/CSV) a že je možné nastavit retenční politiku odpovídající interním požadavkům na compliance.
\end{enumerate}

Návrh minimálního modelu rolí a oprávnění (general background knowledge)
\begin{itemize}
  \item Super‑admin (IT): správa integrací, auditních nastavení, SLA konfigurací.
  \item Data Steward / správce obsahu: (meta)řízení indexace interních zdrojů pro RAG, zásahy do přístupových politik k citlivým dokumentům.
  \item Garant předmětu / Lektor: vytváření a správa předmětových asistentů (GPTs), delegování rolí asistentům, přístup k analyzám kurzu.
  \item Asistent (TA): provozní úkony (moderace, zpětná vazba), omezené administrativní akce.
  \item Student / koncový uživatel: standardní přístup k povolenému obsahu, nutnost přihlášení přes SSO a povinnost akceptovat podmínky používání.
\end{itemize}

Auditní logy — doporučené položky k zaznamenávání (general background knowledge)
\begin{itemize}
  \item Události přihlášení / odhlášení (user id, čas, zdroj IP, metoda autentizace).
  \item Změny rolí a oprávnění (kdo změnil, původní a nové oprávnění, čas).
  \item Přístupy k interním zdrojům (kdo, co, kdy, zda došlo k exportu nebo sdílení).
  \item Volání rozhraní modelu (kde relevantní: který uživatel/instrument volal API, požadované zdroje); u RAG scén logovat retrieval události pro forenzní kontrolu.
  \item Administrativní akce (konfigurace, nasazení nových asistentů, mazání dat).
\end{itemize}

Doporučení k retenci a exportu logů (general background knowledge)
\begin{itemize}
  \item Definujte retenční politiku v souladu s interními pravidly a GDPR; oddělte krátkodobé operational logy a dlouhodobé auditní stopy.
  \item Zajistěte možnost exportu logů ve strojově čitelném formátu pro archivaci a předání při incidentu.
  \item Zařaďte logy do centralizovaného SIEM pro korelaci a alertování bezpečnostního týmu.
\end{itemize}

Doporučené postupy při navrhování integrace (general background knowledge)
\begin{itemize}
  \item Preferujte univerzitní SSO (SAML/OIDC) před lokálními účty; minimalizuje to správu hesel a usnadní deprovisioning.
  \item Implementujte princip nejmenších oprávnění (least privilege) a pravidelné revize rolí.
  \item Při napojení interních materiálů přes RAG řešení omezte přístup modelu k minimálnímu potřebnému scope a logujte retrieval operace.
  \item Zapojte IT, bezpečnostní tým a GDPR/PRIV kontakty fakulty již do fáze PoC — právní a bezpečnostní požadavky ovlivní volbu nastavení a retence dat.
\end{itemize}

Odkazy a zdroje pro další nastudování
\begin{itemize}
  \item Využijte oficiální dokumentaci poskytovatelů a školicí portály (např. Microsoft Learn) pro detailní návody k nastavení SSO, provisioning a auditních rozhraní \cite{kb_ai_101_tipu_pdf_04e4a115_page_54_1}.
  \item Konzultujte interní IT a právní oddělení před schválením provozu; u větších nasazení doporučujeme navázat na RFP/SLA vzory popsané v kapitole o výběru enterprise řešení (sekce 8.1) (general background knowledge).
\end{itemize}

Krátké shrnutí pro garanta předmětu (general background knowledge): před nasazením požadujte PoC zahrnující SSO, provisioning a auditní testy; zajistěte jasné role a dokumentované postupy pro deprovisioning; do sylabu a komunikačních materiálů uveďte, jak budou přístupy spravovány a kde se studenti mohou obrátit v případě problémů s přístupem.%

%
\section{Ochrana dat, DPA a data residency}%
Požadavky na zpracování osobních a interních dat, smlouvy o zpracování (DPA), možnosti data residency, šifrování v klidu i v přenosu a zásady pro sdílení citlivých materiálů (GDPR‑orientované postupy).%


%
Tato podsekce shrnuje praktické kroky a kontrolní body zaměřené na ochranu osobních a interních dat při nasazení AI služeb na úrovni fakulty/univerzity. Následující doporučení jsou navržena jako pracovní průvodce pro garanty předmětů, IT a právní/GDPR kontakty; konečná právní a procesní rozhodnutí je nutné konzultovat s interním právním oddělením a specialisty na ochranu osobních údajů (general background knowledge).

Poznámka k nástrojům ve výuce — ilustrativní příklady:
\begin{itemize}
  \item Některé edukační nástroje jsou součástí komerčních balíků (např. Minecraft Education v rámci Microsoft 365 A3/A5), což má vliv na smluvní podmínky a zpracování dat \cite{kb_ai_101_tipu_pdf_04e4a115_page_15_2}.
  \item Funkce založené na zpracování obrazu (OCR, extrakce textu z obrázků) mohou konvertovat nahrané skeny či screenshoty do zpracovatelného textu; to je relevantní při sdílení dokumentů obsahujících osobní údaje \cite{kb_ai_101_tipu_pdf_04e4a115_page_53_2}.
  \item Některé edukační AI funkce shromažďují a vyhodnocují výkon/student‑level metriky (např. nástroje generující příklady a analyzující postupy při matematice), což zvyšuje potřebu jasného smluvního ošetření a ochrany dat \cite{kb_ai_101_tipu_pdf_04e4a115_page_8_1}.
\end{itemize}

Co zkontrolovat v DPA (Data Processing Agreement) — praktické body (general background knowledge)
\begin{itemize}
  \item Role a odpovědnosti: jasné vymezení, kdo je správcem (controller) a kdo zpracovatelem (processor), včetně sub‑processorů a práv na audity.
  \item Účel zpracování: omezení na konkrétní účely výuky/administrace, zákaz sekundárního komerčního využití studentských dat bez souhlasu.
  \item Doba retence a pravidla mazání: termíny uchovávání, procesy pro úplné odstranění dat na vyžádání a mechanismy pro obnovu po chybě.
  \item Data residency a lokace zpracování: požadavek na zpracování/uložení v jurisdikci EU/EHP nebo jiná specifika podle interních pravidel fakulty.
  \item Šifrování a zabezpečení: povinnost používat šifrování v přenosu (TLS) a šifrování v klidu; volitelné zákaznické (customer‑managed) klíče jsou výhodou.
  \item Incidentní řízení: SLA pro hlášení úniku dat, povinnost notifikace do interních kontaktních bodů a lhůty pro oznámení.
  \item Podmínky sub‑processorů a auditu: seznam povolených sub‑processorů, právo na bezpečnostní audity a penále/odpovědnost za porušení.
\end{itemize}

Data residency — možnosti a implikace (general background knowledge)
\begin{itemize}
  \item Možné varianty: 1) úložiště a zpracování v EU/EHP, 2) regionální cloud s možnostmi smluvního omezení lokace, 3) globální cloud bez záruk residency. Volba ovlivní právní compliance a provozní latenci.
  \item Praktické pravidlo: pro citlivá studentská data preferovat EU‑based residency nebo on‑prem/hybridní řešení; u veřejně dostupných studijních materiálů lze zvážit širší varianty s odpovídajícími zárukami.
\end{itemize}

Šifrování a technická opatření (general background knowledge)
\begin{itemize}
  \item Požadujte šifrování v přenosu (TLS 1.2+/TLS 1.3) a šifrování v klidu (AES‑256 nebo ekvivalent).
  \item Zvažte customer‑managed keys (CMK) pro klíčovou kontrolu, rotaci klíčů a možnost odnětí přístupu v případě incidentu.
  \item U RAG scén a indexace interních zdrojů logujte retrieval události a volání modelu pro forenzní sledovatelnost (general background knowledge).
\end{itemize}

Praktický workflow pro sdílení citlivých materiálů v kurzu — krok po kroku (general background knowledge)
\begin{enumerate}
  \item Klasifikujte data používaná v kurzu (veřejné, interní, citlivé/osobní).
  \item Minimalizujte sběr: sbírejte pouze nezbytná pole (pseudonymizujte tam, kde lze).
  \item Proveďte DPIA (posouzení dopadů na ochranu osobních údajů) před zahájením rozsáhlejšího nasazení AI nástroje.
  \item Uzavřete nebo ověřte DPA se službou, zkontrolujte residency, šifrování a sub‑processor pravidla.
  \item Technicky nastavte přístup: SSO/IdP integrace, role‑based access, omezení exportů dat a auditní logování.
  \item Testujte PoC: ověřte mazání dat, export logů, chování při incidentu a chování extrakčních funkcí (OCR atp.).
  \item Do sylabu a komunikace s studenty jasně specifikujte pravidla sdílení dat a jak budou jejich výstupy zpracovávány (general background knowledge).
\end{enumerate}

Kontrolní checklist pro rychlý PoC (general background knowledge)
\begin{itemize}
  \item Existuje podepsaná DPA pokrývající všechny plánované datové toky?
  \item Je lokace uložených dat (data residency) v souladu s interní politikou?
  \item Jsou zapnutá šifrování v přenosu i v klidu a je dostupná možnost CMK?
  \item Jsou nastavena auditní logy pro přístupy k interním dokumentům a volání modelu?
  \item Proběhla DPIA a byla schválena příslušným orgánem/fakultním GDPR kontaktem?
\end{itemize}

Závěrečné doporučení pro garanty a IT (general background knowledge):
\begin{itemize}
  \item Zapojte právní/GDPR a IT/security týmy již ve fázi výběru řešení; smluvní zajištění a technické konfigurace ovlivní možnosti využití nástroje ve výuce.
  \item Piloty navrhujte s omezeným dosahem dat (pseudonymizace, anonymizace) a jasným plánem rozšíření až po ověření bezpečnosti a souladu.
  \item Dokumentujte všechna rozhodnutí o residency, DPA a retenčních politikách a zahrňte je do sylabů a komunikace ke studentům.
\end{itemize}%

%
\section{Integrace do univerzitní infrastruktury}%
Scénáře a technické vzory integrace s LMS (Moodle), knihovnami, vývojářskými prostředími a dalšími službami VUT; connectors, webhooky, SSO‑flow a doporučené implementační postupy pro katedry a IT.%


%
Krátké shrnutí a cíl sekce

Tato část popisuje praktické scénáře a technické vzory, jak napojit AI‑funkce do univerzitní infrastruktury (LMS, knihovny, repozitáře kódu, tooling pro vývoj) tak, aby byly splněny požadavky pedagogů i pravidla ochrany dat. Uvádíme konkrétní příklady, doporučené kroky pro PoC (pilot) a kontrolní seznamy pro spolupráci mezi katedrou a IT oddělením.

Praktické příklady nástrojů a integrací (ilustrační, podpořené zdroji)
\begin{itemize}
  \item Využití herních edukačních světů (např. Minecraft Education) jako zdroje interaktivních materiálů — poznámka k licenci: některé edukační nástroje jsou součástí komerčních balíků (Minecraft Education je součástí Microsoft 365 A3/A5), proto je důležité ověřit smluvní podmínky a zpracování dat u poskytovatele \cite{kb_ai_101_tipu_pdf_04e4a115_page_15_2}.
  \item Integrace nástrojů pro podporu výuky matematiky (např. OneNote s převodem rukopisu na rovnice, math.microsoft.com) do LMS pro sběr postupů a automatickou analýzu řešení \cite{kb_ai_101_tipu_pdf_04e4a115_page_8_1}.
  \item Zdroje pro technické školení a průvodce integrací (např. oficiální vzdělávací materiály Microsoft Learn) jako podpora pro IT týmy a učitele při zavádění nových funkcí \cite{kb_ai_101_tipu_pdf_04e4a115_page_54_1}.
\end{itemize}

Architektury a vzory integrace (general background knowledge)
\begin{itemize}
  \item LMS connector pattern (general background knowledge): vytvoření vrstvy, která synchronizuje kurzy, uživatele a odevzdání mezi LMS (např. Moodle) a AI službou. Obvyklé komponenty: autentizace (SSO), API gateway, transformace dat (anonymizace/pseudonymizace) a auditní logy.
  \item RAG / dokumentový indexer (general background knowledge): bezpečné indexování interních studijních materiálů (PDF, HTML, repository) do embeddings indexu s řízením přístupu a logováním retrieval událostí.
  \item Webhook / grade sync (general background knowledge): volání webů pro synchronizaci výsledků a metadat (např. automatické importy hodnocení z AI‑asistentů do gradebooku LMS) s retry a idempotentní logikou.
  \item Developer workflow a verzování promptů (general background knowledge): správa prompt šablon, verzí modelu a prompt‑logů v repozitáři (git), včetně review procesu a audit trailu pro reproducibilitu.
\end{itemize}

Doporučený PoC workflow pro katedru + IT (general background knowledge)
\begin{enumerate}
  \item Definujte pedagogický cíl pilotu: jasně popište, co má integrace přinést studentům a vyučujícím (např. rychlejší formativní zpětná vazba, interaktivní cvičení s Minecraft Education \cite{kb_ai_101_tipu_pdf_04e4a115_page_15_2}).
  \item Vyberte datové zdroje a rozsah: identifikujte, které materiály budou indexovány/nahrávány, a proveďte klasifikaci dat (veřejné vs. interní vs. osobní).
  \item Smluvně zajistěte podmínky (DPA) a technické požadavky: data residency, šifrování, možnosti mazání logů (viz předchozí kapitola o ochraně dat — obsahuje checklist pro DPA) (general background knowledge).
  \item Implementujte minimální integraci: nastavení SSO/IdP pro bezpečný přístup (general background knowledge), implementace connectoru do LMS (export/import uživatelů a kurzu), a jednoduchý webhook pro notifikace o dokončení zpracování.
  \item Testování s omezenými daty: prověřte anonymizaci, audit logy a chování při chybách; ověřte, že nástroj nevrací citlivé informace mimo určený kontext.
  \item Zapojení pedagogů a pilotních studentů: sběr zpětné vazby na použitelnost a kvalitu výsledků; iterativní úpravy.
\end{enumerate}

Konkrétní technické kontrolní body pro nasazení connectoru do LMS (general background knowledge)
\begin{itemize}
  \item Autentizace: preferovat univerzitní SSO (SAML/OAuth/OIDC) pro jednotné řízení identit a odvolání přístupu.
  \item Minimální privilegia: service‑accounts a scopes omezené pouze na nezbytné API volání.
  \item Anonymizace/pseudonymizace: před odesláním do externích služeb odstranit identifikátory, pokud to pedagogika dovolí.
  \item Auditování: logovat kdo/ kdy/ které dokumenty byly poslány, které retrievaly byly provedeny a jaké odpovědi model vrátil (pro účely vyšetřování a GDPR).
  \item Fail‑safe chování: pokud externí AI služba vrací chybu nebo dojde k incidentu, zajistit, aby do LMS neunikla neúplná nebo nesprávná data.
\end{itemize}

Příklady uplatnění v praxi (ilustrativní)
\begin{itemize}
  \item Integrace Minecraft Education do kurzu: ověřte, zda je školní licence (Microsoft 365 A3/A5) k dispozici a upravte přístup studentů podle smluvních podmínek; použijte připravené světy pro výuku AI a Hour of Code aktivity \cite{kb_ai_101_tipu_pdf_04e4a115_page_15_2}.
  \item Sběr postupů z OneNote pro automatickou kontrolu řešení: nasměrujte studenty, aby odevzdávali postupy v OneNote (který umí konvertovat rukopis na digitální rovnice) a propojte export s analyzátorem chyb ve vyučovacím workflowu \cite{kb_ai_101_tipu_pdf_04e4a115_page_8_1}.
\end{itemize}

Doporučené role a odpovědnosti (general background knowledge)
\begin{itemize}
  \item Katedra / garant předmětu: definice pedagogických požadavků, ověření obsahu a schválení testovacích dat.
  \item IT / integrátor: implementace connectorů, SSO, auditních logů a provozní podpory PoC.
  \item GDPR / právní kontakty: posouzení DPA, data residency a DPIA (pokud potřebné).
  \item Data Steward / AI‑governance: správa přístupů k interním indexům, pravidel pro retenci a audit.
\end{itemize}

Kde hledat další technické návody a školení (podložené zdroji)
\begin{itemize}
  \item Oficiální platformy s technickými kurzy a průvodci (např. Microsoft Learn) jsou vhodným zdrojem pro IT týmy i vyučující, kteří potřebují rychlé seznámení s konkrétními nástroji a možnostmi integrace \cite{kb_ai_101_tipu_pdf_04e4a115_page_54_1}.
  \item Dokumentace konkrétních nástrojů (LMS, cloud provider, enterprise AI) by měla být vždy konzultována před nasazením — zejména části týkající se DPA, residency a možností zákaznické správy klíčů.
\end{itemize}

Rychlý checklist pro spuštění pilotu (general background knowledge)
\begin{enumerate}
  \item Definovaný pedagogický cíl a metriky úspěchu.
  \item Ověřená smlouva/DPA a klasifikace dat.
  \item Implementovaný connector s SSO a auditováním.
  \item Testy anonymizace a chování při incidentu.
  \item Zapojená skupina pedagogů a pilotních studentů pro sběr zpětné vazby.
\end{enumerate}

Poznámka k rizikům a dalším krokům

Při každé integraci dbejte na to, aby pedagogický přínos převýšil rizika spojená s ochranou dat a spolehlivostí výstupů. Doporučíme zahájit s omezeným PoC rozsahem (malá třída, anonymizovaná data), zdokumentovat rozhodnutí o residency a DPA a iterovat podle výsledků pilotu (general background knowledge). Pro konkrétní návody k nástrojům Microsoft (Minecraft Education, OneNote, math.microsoft.com) použijte odkazy a školení uvedené v příručce \cite{kb_ai_101_tipu_pdf_04e4a115_page_15_2, kb_ai_101_tipu_pdf_04e4a115_page_8_1, kb_ai_101_tipu_pdf_04e4a115_page_54_1}.%

%
\section{Architektura RAG s interními materiály}%
Návrhy pro Retrieval‑Augmented Generation: indexace interních dokumentů, embeddings, retrieval vrstvy, bezpečné napojení na LLM, varianty cloud vs. hybrid vs. on‑prem a mitigace rizik úniku informací.%


%
% Poznámka pro čtenáře:
% Položky a technické kroky v této sekci jsou převážně praktickými návrhy pro návrh a provoz Retrieval‑Augmented Generation (RAG) systémů nad interními materiály.
% Pokud není uveden explicitní citát na konkrétní zdroj, považujte tvrzení za "general background knowledge" (praktické doporučení spojené s běžnými integračními vzory).

Úvodní shrnutí (general background knowledge).
RAG (Retrieval‑Augmented Generation) architektura kombinuje indexaci interních dokumentů a rychlé vyhledávání relevantních pasáží s generativním jazykovým modelem tak, aby model mohl „doplňovat“ své odpovědi o kontext z interních zdrojů. Níže uvádíme doporučené komponenty, bezpečnostní opatření a praktický PoC workflow pro nasazení na úrovni katedry nebo fakulty (general background knowledge).

Klíčové komponenty RAG (general background knowledge)
\begin{itemize}
  \item Ingestní vrstva: sběr zdrojů (PDF, HTML, prezentace, OneNote stránky, repozitáře kódu, exporty z LMS).
  \item Preprocessing a čištění: konverze formátů, odstranění metadat, segmentace dokumentů do jednotlivých chunků.
  \item Anonymizace / Pseudonymizace: odstranění nebo zástupné znaky pro citlivé identifikátory před indexací.
  \item Embeddings \\ \& index: výpočet vektorových embeddingů pro jednotlivé chunk‑y a uložení do vyhledávacího vektorového úložiště (vector store).
  \item Retrieval vrstva: dotazy vůči indexu, relevance‑scoring a filtrování výsledků podle oprávnění.
  \item Prompt orchestration / answer synthesis: skládání kontextu (retrieved passages) do promptu pro LLM, pravidla pro ukázání provenance a omezení délky.
  \item Audit \\ \& logging: zaznamenávání retrieval událostí (kdo/co/ kdy), uchování provenance a verzování indexů pro forenzní účely.
  \item Governance \\ \& access control: SSO/RBAC pro řízení přístupu k indexu a retencím dat.
\end{itemize}

Praktický PoC workflow pro katedru (general background knowledge)
\begin{enumerate}
  \item Definice rozsahu a cílů: vyberte sadu dokumentů a konkrétní pedagogický scénář (FAQ pro kurz, sumarizace laboratorních protokolů, doplňování studijních materiálů).
  \item Klasifikace dat: označte, které zdroje jsou interní, veřejné nebo obsahují osobní údaje; podle toho naplánujte anonymizaci a retenci.
  \item Implementace ingestu: např. export OneNote stránek nebo PowerPoint prezentací z repozitáře kurzu a konverze do textového formátu \cite{kb_ai_101_tipu_pdf_04e4a115_page_8_1, kb_ai_101_tipu_pdf_04e4a115_page_44_2}.
  \item Vytvoření embeddings a indexu: zvolte embedding model kompatibilní s provozním režimem (cloud/hybrid/on‑prem).
  \item Nastavení retrieval pravidel: omezení výsledků podle kurzu/katedry, filtrování citlivých entit a nastavení maxima number_of_chunks pro prompt.
  \item Testování přesnosti a bezpečnosti: kontrola zda retrieval nevrací citlivé údaje mimo kontext; validace návratů lidským recenzentem.
  \item Nasazení v omezeném pilotu: malá skupina pedagogů a studentů, sběr zpětné vazby a metrik (relevance, chybovost, incidenty).
  \item Rozšíření a provoz: zavedení auditních procedur, DPA/DPIA kroky dle potřeby a plán retencí.
\end{enumerate}

Ukázkový, zjednodušený příklad indexace výukových materiálů
\begin{enumerate}
  \item Sběr zdrojů: exportujte učební prezentace (PowerPoint), OneNote stránky s řešením úloh a průvodní dokumenty kurzu. (OneNote obsahuje nástroje pro převod rukopisu do digitální podoby vhodné k indexaci) \cite{kb_ai_101_tipu_pdf_04e4a115_page_8_1}.
  \item Normalizace: převod PPT/OneNote do textu, rozdělení na odstavce/sekce, odstranění připomínek s osobními údaji.
  \item Embeddings: spočítejte embedding pro každý chunk a uložte do zabezpečeného vector store (s přístupovým řízením).
  \item Retrieval + LLM: při dotazu nejprve proveďte retrieval, přiložte několik nejrelevantnějších excerptů s jasnou poznámkou o provenance do promptu a spusťte generaci v LLM.
  \item Validace výsledku: u pedagogů/nebo automatizovaného validačního kroku ověřte, zda odpověď nekonkretizuje citlivé interní údaje mimo kontext.
  \item Příklady zdrojů vhodných pro indexaci: výukové slidy, zadání cvičení, lab instrukce, FAQ kurzu; edukační světy (např. Minecraft Education) mohou sloužit jako obsahová složka pro didaktické aktivity, které je možné indexovat jako metadata kurzu \cite{kb_ai_101_tipu_pdf_04e4a115_page_15_2}.
\end{enumerate}

Bezpečnostní mitigace rizika úniku informací (general background knowledge)
\begin{itemize}
  \item Minimalizace ingestu: indexovat pouze to, co je nezbytné pro pedagogický cíl.
  \item Pseudonymizace před indexací: nahradit identifikátory, pokud to pedagogika dovolí.
  \item Přístupové filtry při retrieval: výsledky omezit podle rolí a kontextu (kurs, rok studia).
  \item Red‑teaming a testovací retrievaly: simulovat útoky a nežádoucí dotazy, abyste odhalili nečekané leakage scénáře.
  \item Provenance v odpovědích: vždy zobrazit, ze kterých dokumentů byly úryvky použity.
  \item Human‑in‑the‑loop pro citlivé dotazy: například žádosti týkající se osobních údajů nebo bezpečnostních postupů musí vracet pouze předem schválené kanály.
\end{itemize}

Volba nasazení: cloud vs. hybrid vs. on‑prem (general background knowledge)
\begin{itemize}
  \item Cloud (rychlé nasazení, škálovatelnost): vhodné pro rychlé PoC a menší rizika s daty; vyžaduje silný DPA a data residency kontrolu.
  \item Hybrid: index běží lokálně, inference v cloud nebo opačně; kompromis mezi bezpečností a výkonem.
  \item On‑prem / lokální modely: vyšší kontrola nad daty a nižší dependency na třetích stranách; vyžaduje vyšší provozní kapacity a licenční posouzení (modely, HW).
\end{itemize}

Integrace s existující infrastrukturou a užitečné poznámky
\begin{itemize}
  \item Napojení na LMS: použijte connector pattern pro synchronizaci kurzů a oprávnění (general background knowledge).
  \item SSO a audit: používejte univerzitní SSO (SAML/OIDC) a logujte retrieval akce pro forenzní účely (general background knowledge).
  \item Školení pedagogů a dokumentace: zajistěte krátké návody pro učitele jak interpretovat RAG výstupy a jak vyžadovat provenance; v případě nástrojů Microsoft lze využít dostupné školení a dokumentaci (např. průvodce a školicí materiály pro Copilot a Microsoft 365) \cite{kb_ai_101_tipu_pdf_04e4a115_page_44_2, kb_ai_101_tipu_pdf_04e4a115_page_8_1}.
\end{itemize}

Krátký checklist před spuštěním pilotu RAG (general background knowledge)
\begin{itemize}
  \item Definovaný pedagogický cíl a seznam indexovaných zdrojů.
  \item Prověřená DPA/DPIA cesta pro zpracování interních dat.
  \item Mechanismy anonymizace/pseudonymizace aktivní před indexací.
  \item Přístupové řízení (SSO/RBAC) a auditování retrieval událostí.
  \item Testovací scénáře pro odhalení leakage a plán rollbacku.
  \item Zapojení pilotních pedagogů a jasný plán sběru metrik a zpětné vazby.
\end{itemize}

Závěrem: RAG přináší silný nástroj pro zpřístupnění interních znalostí v konverzačním prostředí, ale jeho bezpečné nasazení vyžaduje plánování ingestu, anonymizaci, řízení přístupu a auditování. Pro konkrétní technické kroky doporučujeme koordinaci mezi garantem předmětu, IT oddělením a GDPR/právním oddělením fakulty (general background knowledge).%

%
\section{Lokální a edge nasazení modelů}%
Možnosti lokálního/edge provozu (např. menší modely jako Llama, Mistral), provozní požadavky, licensing‑considerations, zálohy, aktualizace modelů a provozní bezpečnostní doporučení.%


%
Tato část shrnuje praktické zásady pro lokální (on‑prem) a edge nasazení menších modelů v univerzitním prostředí. Obsahuje doporučené kroky před nasazením, provozní požadavky, bezpečnostní body a provozní rutiny. Pokud není uvedeno jinak, níže uvedená doporučení jsou general background knowledge.

Krátké vymezení a kdy uvažovat o lokálním/edge provozu (general background knowledge).
Lokální nebo edge nasazení zvyšuje kontrolu nad citlivými daty a může snížit závislost na externích cloudových poskytovatelích; na druhou stranu obvykle vyžaduje vyšší provozní kapacity, správu HW a pečlivé licenční posouzení (modely, knihovny, runtime) (general background knowledge).

Předběžné kroky před PoC
\begin{itemize}
  \item Definujte pedagogický cíl nasazení (např. lokální tutor pro laboratoře, inference nad interními materiály) a rozsah dat, která budou modelu zpřístupněna (general background knowledge).
  \item Zapojte IT, právní/GDPR a garanty předmětu ještě v návrhové fázi — rozhodnutí o on‑prem vs. hybrid musí zahrnovat posouzení DPA/DPIA, požadavky na data residency a auditní stopu (general background knowledge).
  \item Zvolte rozsah funkcionality: offline inference na klientském zařízení (edge), inference na lokálním serveru, nebo hybridní režim (lokální index + cloud inference) (general background knowledge).
\end{itemize}

Technické požadavky a doporučení (general background knowledge)
\begin{itemize}
  \item Hardware: odhady nároků závisí na velikosti modelu a režimu (CPU inference vs. GPU/TPU). Pro edge/embedded scénáře volte kvantizované a optimalizované verze modelů; pro server‑side inference plánujte dostatek paměti, diskového úložiště pro indexy a způsob řízení GPU času.
  \item Modely a optimalizace: u menších open‑source modelů (příklady modelových rodin se často používají v praxi) volte varianty optimalizované pro inference, případně použijte kvantizaci a distilaci ke snížení nároků (general background knowledge).
  \item Správa závislostí: používejte kontejnerizaci (Docker/OCI), verziování modelů a reproducibilní CI/CD pipeline pro nasazení modelu a runtime (general background knowledge).
  \item Licensing‑considerations: ověřte licenční podmínky modelu a knihoven (komerční použití, distribuce upravených modelů, požadavky na attribution). Licence mohou ovlivnit možnost nasazení na univerzitním HW nebo sdílení modelu mezi útvary (general background knowledge).
\end{itemize}

Bezpečnost, soukromí a provozní opatření (general background knowledge)
\begin{itemize}
  \item Omezení ingestu: indexujte a zpřístupňujte lokálně pouze nezbytné dokumenty. Pseudonymizujte nebo anonymizujte osobní údaje před použitím v tréninku nebo indexaci.
  \item Přístup a audit: integrujte SSO/SSO‑flow a RBAC pro řízení přístupu k inference endpointům a vektorovým indexům; logujte requesty a uchovávejte auditní trail pro forenzní účely.
  \item Red‑teaming a testování leakage: simulujte dotazy, které by mohly způsobit únik interních informací, a ověřte, že model nevrací citlivé části dokumentů mimo kontext.
  \item Human‑in‑the‑loop: pro dotazy obsahující citlivá data nebo rozhodnutí s vysokou mírou dopadu zajistěte schvalovací procesy nebo eskalaci na lidského pracovníka.
\end{itemize}

Provozní rutiny: zálohy, aktualizace a monitoring (general background knowledge)
\begin{enumerate}
  \item Monitoring: sledujte latenci, chybovost, využití HW a anomálie v inference (neobvyklé odpovědi, náhlý nárůst chyb).
  \item Aktualizace modelu a indexů: verziujte změny (modely i prompt šablony), testujte nové verze v izolovaném prostředí a plánujte rollback strategie.
  \item Zálohy a retence: pravidelně zálohujte modely, indexy a metadata (včetně provenance). Definujte retenční politiku pro staré indexy a logy.
  \item Incident response: připravte postup pro bezpečnostní incidenty včetně kontaktů IT, právního oddělení a odpovědné osoby na katedře.
\end{enumerate}

Praktický checklist pro pilotní nasazení (general background knowledge)
\begin{itemize}
  \item Jasně definovaný pedagogický případ použití a seznam typů dokumentů, které budou použity.
  \item Posouzení licencí modelu a souvisejících knihoven.
  \item Schválení DPIA/DPA v případě zpracování osobních údajů.
  \item Testovací prostředí s monitoringem a možností rollbacku.
  \item Dokumentace provozních postupů, školení pro pedagogy a kontaktní body pro technickou podporu.
\end{itemize}

Závěrem: lokální a edge nasazení přináší výhodu zvýšené kontroly nad daty a nižší závislosti na cloudových operátorech, ale s sebou nese vyšší provozní nároky, potřebu řízení licencí a robustní provozní postupy. Doporučujeme plánovat PoC s úzkou spoluprací mezi garantem předmětu, IT oddělením a právním/GDPR týmem a začít s omezeným pilotem před rozšířením do produkce (general background knowledge).%

%
\section{Governance, řízení nákladů, checklisty a bezpečnostní opatření}%
Role a procesy (Data Stewards, AI Council), řízení rozpočtů a token/metrické limity, vzory RFP/SLA, IT+katederní checklisty pro nasazení a provoz, plus guardrails, filtrování obsahu a red‑teaming postupy.%


%
% Governance, řízení nákladů, checklisty a bezpečnostní opatření
% (Obsah této části následuje po nadpisu sekce; nadpis není součástí tohoto výstupu.)

Krátké shrnutí účelu sekce. Tato část popisuje doporučené organizační role, základní principy řízení nákladů a metrik při provozu AI služeb na úrovni fakulty/univerzity a poskytuje praktické check‑listy pro IT i katedry včetně návrhů základních bezpečnostních opatření a red‑teamingu.

\medskip

\textbf{Klíčové organizační role (stručně)} \\
(Doplňující popis níže je uveden jako general background knowledge.)
\begin{itemize}
  \item Data Steward — odpovědná osoba za klasifikaci a správu interních dat, rozhoduje o anonymizaci/pseudonymizaci a o rozsahu ingestu. (general background knowledge)
  \item AI Council / Governance Board — meziodborová skupina garantující pravidla pro nasazení (DPA/DPIA, SLA, bezpečnostní požadavky) a schvalování pilotů. (general background knowledge)
  \item IT Security Owner + Incident Response kontakt — provozní zodpovědnosti za monitoring, logy a reakce na bezpečnostní incidenty. (general background knowledge)
\end{itemize}

\medskip

\textbf{Bezpečnostní a riziková priorita} \\
Analýza rizik by měla být nedílnou součástí aplikací AI ve vzdělávání; při návrhu řešení je třeba explicitně posoudit bezpečnostní, etická a ochrany‑osobních‑údajů rizika a zapojit příslušné odborné týmy (IT, právní/GDPR, garanti předmětů) do návrhové fáze \cite{kb_malinka_chatgpt_ve_vyuce_dva_roky_pote_2024_pdf_90bc3aaf_page_12_1}. 

\medskip

\textbf{Řízení nákladů a metriky provozu} \\
(Zde uvedené body jsou uvedeny jako general background knowledge — slouží k rychlému nasazení rozhodovacích mechanismů.)
\begin{itemize}
  \item Rozpočet a chargeback: alokovat náklady podle kateder / pilotních projektů; definovat limitní rozpočty pro PoC fáze. (general background knowledge)
  \item Tokeny / metrické limity: u cloudových LLM sledovat spotřebu tokenů, počet requestů a latency; nastavit varovné hranice a automatické throttlingy. (general background knowledge)
  \item Měření dopadu: metriky spokojenosti uživatelů, poměr lidské validace vs. automatických odpovědí, chybovost/halucinace (sampling a auditování). (general background knowledge)
\end{itemize}

\medskip

\textbf{RFP / SLA – klíčové body pro požadavky na dodavatele} \\
(Následující body jsou doporučenou šablonou; pokud jsou použity v zadávacích dokumentech, doplňte právní prověrku.)
\begin{itemize}
  \item Zásady zpracování dat (DPA): jasně definované role správce a zpracovatele, popis typů zpracovávaných dat a uložení dat (data residency). (general background knowledge)
  \item Auditní schopnosti: požadavek na auditní logy, možnost exportu logů, a dostupnost traceability/provenance pro retrieval události. (general background knowledge)
  \item SLA metriky: dostupnost služby, latence, garantované časy obnovy, podmínky pro rollback a bezpečnostní incidenty. (general background knowledge)
  \item Bezpečnostní certifikace a red‑teaming: doložení bezpečnostních testů, výsledky red‑team cvičení nebo povinnost umožnit třetí straně audit. (general background knowledge)
\end{itemize}

\medskip

\textbf{Praktické checklisty pro IT a katedry před nasazením (rychlý přehled)} \\
Následující checklist je uveden jako general background knowledge pro rychlé ověření připravenosti:
\begin{enumerate}
  \item Vymezení pedagogického cíle a seznam dokumentů/data, která budou použita (scope). (general background knowledge)
  \item Posouzení nutnosti DPIA a zapojení právního/GDPR týmu (analýza dopadu na ochranu dat je součástí návrhu) \cite{kb_malinka_chatgpt_ve_vyuce_dva_roky_pote_2024_pdf_90bc3aaf_page_12_1}.
  \item Klasifikace dat (interní / veřejné / obsahující PII) a rozhodnutí o anonymizaci/pseudonymizaci. (general background knowledge)
  \item Definice přístupových práv (SSO + RBAC) a auditních požadavků. (general background knowledge)
  \item Testovací fáze (pilot s malou skupinou, měření metrik a lidská validace výsledků). (general background knowledge)
  \item Nasazení monitoringu, alertů na překročení token/metrických limitů a plán pro rollback. (general background knowledge)
\end{enumerate}

\medskip

\textbf{Guardrails, filtrování obsahu a red‑teaming} \\
Doporučené bezpečnostní postupy a testy — shrnutí jako general background knowledge:
\begin{itemize}
  \item Guardrails: definujte limity pro typ dotazů (zákaz generování citlivých dat), předfiltraci vstupů a blacklisty/whitelisty pro externí zdroje. (general background knowledge)
  \item Filtrování retrieval vrstev (u RAG): omezení retrieve na schválené zdroje, filtrování podle role/kurzu a maximální počet chunků na odpověď. (general background knowledge)
  \item Red‑teaming: pravidelné interní testy zaměřené na leakage interních dokumentů, zkoušení nebezpečných promptů a simulace chybových stavů; výstupy dokumentovat a implementovat nápravná opatření. (general background knowledge)
  \item Human‑in‑the‑loop: pro dotazy s vysokým rizikem zajistit eskalaci k lidskému recenzentovi. (general background knowledge)
\end{itemize}

\medskip

\textbf{Školení a znalostní podpora} \\
Pro kapacitní zajištění a zvyšování kvality provozu doporučujeme zřídit školení pro IT a pedagogy — u produktů Microsoft jsou užitečné oficiální moduly na Microsoft Learn jako zdroj pro technologie Microsoft/365 a integrace \cite{kb_ai_101_tipu_pdf_04e4a115_page_54_1}. V rámci interní governance je vhodné vést centrální repozitář dokumentace (SOP), incidentních postupů a šablon pro DPA/DPIA.

\medskip

\textbf{Rychlý návrh eskalačního postupu při incidentu} \\
(Uvedené kroky jsou general background knowledge a měly by být adaptovány podle interních pravidel VUT.)
\begin{enumerate}
  \item Izolace incidentu a omezení přístupu k postiženým službám.
  \item Notifikace Data Stewarda, IT Security Ownera a právního/GDPR kontaktu.
  \item Interní forenzní analýza (logy, provenance retrieval událostí) a vyhodnocení rozsahu uniklých informací.
  \item Opatření nápravy (revoke keys, rotace modelů/indexů, aktualizace filtrovacích pravidel).
  \item Zpráva o incidentu a aktualizace DPIA / provozní dokumentace. 
\end{enumerate}

\medskip

Závěrem: implementace governance a řízení nákladů pro AI na úrovni fakulty kombinuje technické (monitoring, throttling, RBAC), právní (DPA/DPIA) a organizační (role, board) prvky. Analýza dopadů a zapojení relevantních týmů musí proběhnout již v předběžné fázi návrhu řešení \cite{kb_malinka_chatgpt_ve_vyuce_dva_roky_pote_2024_pdf_90bc3aaf_page_12_1}. Pro praktické nasazení použijte tento text jako výchozí checklist a adaptujte body podle interních pravidel a smluvních požadavků dodavatelů. \cite{kb_ai_101_tipu_pdf_04e4a115_page_54_1}%

%
\chapter{PŘÍLOHY}%
Sada praktických příloh: Prompt Library s 20+ šablonami k okamžitému použití (lesson planning, assessment, feedback), Checklist pro redesign předmětu, AI Policy Template do sylabů, vzor DPIA pro použití AI ve výuce, základní RAG blueprint pro kurzy, Glossář pojmů a kontakty/podpora VUT. Dále kompletní bibliografie a anotované odkazy na univerzitní guideline a patterny (UCL, TU Delft aj.). Přílohy jsou navrženy jako pracovní nástroje připravené ke kopírování a úpravám.%
\section{Prompt Library (20+ šablon)}%
Knihovna 20+ praktických prompt šablon k okamžitému použití v přípravě výuky, tvorbě zadání a hodnocení (lesson planning, assessment, feedback, generování testů, rubrik, zpětné vazby). Šablony jsou copy-\&-paste připravené s krátkým návodem k přizpůsobení pro různé předměty a úrovně.%


%
\medskip

Krátké vysvětlení: níže je kolekce 20+ copy‑\\ \&‑paste promptů (šablon) rozdělených podle účelu: plánování lekcí, generování cvičení a kvízů, formativní zpětné vazby, rubrik a scénářů pro interaktivní hodiny. Prompty jsou navrženy pro rychlé upravení podle oboru, úrovně studentů a výukového cíle. Některé příklady odkazují na nástroje zmíněné v této příručce \cite{kb_ai_101_tipu_pdf_04e4a115_page_15_2, kb_ai_101_tipu_pdf_04e4a115_page_8_1, kb_ai_101_tipu_pdf_04e4a115_page_7_1}.

\medskip

\textbf{A. Plánování lekcí a sylabů}

\begin{enumerate}
  \item \textbf{90minutová přednáška — timeboxed plán}
  \begin{verbatim}
You are an instructional designer. Create a 90-min plan on: [TÉMA].
Split into 4–6 segments with: time, 1–3 objectives, activity, materials.
Add 3 short homework tasks with expected times.
Audience: [Bakalář/Magistr/PhD], prior knowledge: [..]
  \end{verbatim}
  \item \textbf{Generování sylabu (ECTS‑kompatibilní)}
  \begin{verbatim}
Generate a course syllabus: [NÁZEV].
Include: 3–6 learning outcomes, 12 weekly topics, assessment weights,
recommended literature, short AI policy.
Target audience: [level].
  \end{verbatim}
  \item \textbf{Differenciace cvičení podle úrovně}
  \begin{verbatim}
Create Easy/Standard/Advanced variants of an exercise on [TÉMA]:
For each: task, time, hints, answer outline.
  \end{verbatim}
\end{enumerate}

\medskip

\textbf{B. Generování cvičení, kvízů a testovacích případů}

\begin{enumerate}
  \setcounter{enumi}{3}
  \item \textbf{MCQ banka — variabilita obtížnosti}
  \begin{verbatim}
Generate 10 MCQs on [TÉMA].
For each: question, 4 options (1 correct), short explanation, difficulty.
Output JSON: id, question, options[], correct_index, explanation, difficulty.
  \end{verbatim}
  \item \textbf{Time‑boxed in‑class writing}
  \begin{verbatim}
Create a 20-min in-class writing prompt on [TÉMA].
Include 1–2 directed questions, a short rubric (<=5 items), and 3 sample responses.
  \end{verbatim}
  \item \textbf{Programátorský úkol + testovací případy}
  \begin{verbatim}
Design a programming assignment: [NÁZEV].
Provide: statement, I/O spec, 6 unit tests (including edges), complexity limits, solution outline.
Language: [Python/Java/C++].
  \end{verbatim}
\end{enumerate}

\medskip

\textbf{C. Formativní zpětná vazba a komentáře}

\begin{enumerate}
  \setcounter{enumi}{6}
  \item \textbf{Rychlá konzistentní zpětná vazba ke krátkému textu}
  \begin{verbatim}
You are a feedback assistant. Given student's answer provide:
1) 2-sentence summary; 2) 3 strengths; 3) 3 actionable suggestions;
4) suggested grade band w/ justification.
Criteria: [clarity, argumentation, sources]
  \end{verbatim}
  \item \textbf{Komentáře ke kódu (code review)}
  \begin{verbatim}
Review code: list critical bugs/security, perf improvements, style notes,
and 3 unit tests to add. Include short remediation suggestions.
Code: [VLOŽIT]
  \end{verbatim}
  \item \textbf{Šablona pro studentovu sebereflexi o použití AI}
  \begin{verbatim}
Student reflection template:
1) AI tools used; 2) what AI generated (paste if possible);
3) how modified/verified; 4) lessons learned; self-rating 1–5 on reliance.
  \end{verbatim}
\end{enumerate}

\medskip

\textbf{D. Rubriky a hodnoticí šablony}

\begin{enumerate}
  \setcounter{enumi}{9}
  \item \textbf{Rubrika pro procesní hodnocení}
  \begin{verbatim}
Create rubric for: [NÁZEV].
Dimensions: Process documentation, Correctness, Reproducibility, Reflection on AI.
For each: Excellent/Good/Satisfactory/Needs improvement descriptors and weights.
  \end{verbatim}
  \item \textbf{Milestone‑based assessment template}
  \begin{verbatim}
Define 3 milestones for [NÁZEV]: deliverable, acceptance criteria, reviewer, date window.
Include checklist for evidence of process.
  \end{verbatim}
\end{enumerate}

\medskip

\textbf{E. Interaktivní hodiny a speciální nástroje}

\begin{enumerate}
  \setcounter{enumi}{11}
  \item \textbf{Scénář pro Minecraft Education — úvod do AI}
  \begin{verbatim}
Design a 45-min Minecraft Education activity demonstrating an AI concept.
Include: learning objective, step-by-step tasks with times, discussion Qs, homework.
  \end{verbatim}
  \textit{Pozn.: využijte připravené světy a školní licence \cite{kb_ai_101_tipu_pdf_04e4a115_page_15_2}.}
  \item \textbf{Matematické procvičování — integrace s Microsoft nástroji}
  \begin{verbatim}
Prepare 20 practice problems in [oblast matematiky].
For each: problem, answer, hint for common error, and import notes for "Progress in Math".
  \end{verbatim}
  \textit{Funkce umožňují nahrávání postupu pro hodnocení \cite{kb_ai_101_tipu_pdf_04e4a115_page_8_1}.}
  \item \textbf{Převod rukopisu na digitální cvičení — OneNote}
  \begin{verbatim}
Student instructions: write solution by hand in OneNote, convert to digital, submit,
and attach short reflection. Include expected file format/screenshot placeholders.
  \end{verbatim}
\end{enumerate}

\medskip

\textbf{F. RAG / indexing — tvorba krátkých context chunks}

\begin{enumerate}
  \setcounter{enumi}{15}
  \item \textbf{Chunking interního dokumentu pro RAG}
  \begin{verbatim}
Split text into context chunks (max 250 words).
For each: chunk_id, title, 2 keypoints, 1-sentence summary for embedding.
Text: [VLOŽIT DLOUHÝ DOKUMENT]
  \end{verbatim}
  \item \textbf{Prompt pro bezpečnostní filtrování před ingestem}
  \begin{verbatim}
Scan document for PII/confidential passages to redact before indexing.
Return: line numbers, reason, suggested redaction.
  \end{verbatim}
\end{enumerate}

\medskip

\textbf{G. Šablony pro pravidla a komunikaci se studenty}

\begin{enumerate}
  \setcounter{enumi}{17}
  \item \textbf{Text do sylabu — AI policy}
  \begin{verbatim}
Insert into syllabus: short policy (allowed/forbidden), requirement to declare AI use,
and consequences for misuse. Keep 100–150 words, formal tone.
  \end{verbatim}
  \item \textbf{Oznámení o nasazení AI asistenta v kurzu}
  \begin{verbatim}
Draft email announcing AI assistant: purpose, capabilities, limits, how to report issues, privacy rules.
  \end{verbatim}
\end{enumerate}

\medskip

\textbf{H. Rychlý checklist pro adaptaci šablon}

\begin{itemize}
  \item Uveďte úroveň (BSc/MSc/PhD), cílové kompetence a formát odevzdání před použitím.
  \item Rozhodněte, zda je AI povolena, deklarovaná nebo zakázaná — upravte prompty.
  \item Při práci s interními daty nejprve spustit filtrovací prompt pro PII.
  \item Pro workshopy s Minecraft/Microsoft nástroji využijte školní licence a odkazy v příručce \cite{kb_ai_101_tipu_pdf_04e4a115_page_15_2, kb_ai_101_tipu_pdf_04e4a115_page_8_1}.
\end{itemize}

\medskip

Závěr: šablony jsou vstupní sadou — otestujte je v pilotu, upravte dle lokálních pravidel (DPA/DPIA) a uložte verzi v repozitáři pro opakované použití. Pro ukázky viz online repozitář příručky.%

%
\section{Checklist pro redesign předmětu}%
Pracovní kontrolní seznam kroků při redesignu kurzu s AI (cíle, výstupy, metody hodnocení, akademická integrita, transparentnost použití AI, ověření výsledků), připravený k přímému použití a úpravě garanty předmětů.%


%
Krátké vysvětlení: níže uvedený checklist slouží jako praktický pracovní nástroj pro garanty předmětů při redesignu kurzu s využitím generativní AI.

\begin{enumerate}
  \item[1.] [\ ] Vymezit pedagogické cíle a měřitelné výstupy: formulujte 3–5 klíčových kompetencí.
  \item[2.] [\ ] Zvolit typy hodnocení a jejich vztah k výstupům: preferujte procesní důkazy a u technických úloh kombinaci automatizace a kontroly postupu.
  \item[3.] [\ ] Definovat politiku používání AI v sylabu: specifikujte, co je povoleno, kdy deklarovat AI a případné sankce.
  \item[4.] [\ ] Navrhnout artefakty a metriky hodnocení: zařaďte pole pro deklaraci AI a požadavky na dokumentaci postupu (logy, drafty, verze).
  \item[5.] [\ ] Požadavek na nahrání postupu u cvičení (kde relevantní):
    \begin{itemize}
      \item Vyžadujte nahrání postupu u matematických/výpočetních úloh; některé nástroje to podporují \cite{kb_ai_101_tipu_pdf_04e4a115_page_8_1}.
      \item U rukopisných řešení využijte převod rukopisu (OneNote) pro sběr a automatizaci kontroly \cite{kb_ai_101_tipu_pdf_04e4a115_page_8_1}.
    \end{itemize}
  \item[6.] [\ ] Zvážit inkluzivní a interaktivní aktivity: použijte demonstrace a edukační světy (např. Minecraft Education) \cite{kb_ai_101_tipu_pdf_04e4a115_page_15_2}.
  \item[7.] [\ ] Posoudit rizika ochrany osobních údajů a plán datového toku: identifikujte exportovaná data, potřebu (pseudo)anonymizace a připravte DPA/DPIA podklady.
  \item[8.] [\ ] Pilot a validační plán: pilotujte s malou skupinou; sbírejte data o použitelnosti, přesnosti nástrojů a dopadu na zátěž vyučujících.
  \item[9.] [\ ] Technické a provozní kroky: uložte verze, repozitář, šablony a doporučená nastavení nástrojů a přístupů.
  \item[10.] [\ ] Komunikace se studenty a školení týmu: připravte krátkou informaci o roli AI v kurzu, očekávaných artefaktech a způsobu deklarace.
\end{enumerate}

Doporučení pro rychlé nasazení: před spuštěním pilotu projděte body 5–9 a uložte šablony do online repozitáře (living document). Při práci s interními daty nejprve spusťte kontrolu PII/redakci podle lokálních pravidel a DPA (viz poznámky v Prompt Library) \cite{kb_ai_101_tipu_pdf_04e4a115_page_15_2, kb_ai_101_tipu_pdf_04e4a115_page_8_1}.%

%
\section{AI Policy Template do sylabů}%
Vzor textu pro zařazení do sylabů s doporučenou formulací o povoleném/zakázaném použití generativní AI, požadavcích na citování, a doporučených postupech pro studenty i vyučující.%


%
Krátké vysvětlení: návrh textu do sylabu s instrukcemi pro garanty a vzorem citace AI.
Doporučené znění do sylabu (krátká verze)
\begin{quote}
V tomto předmětu je povoleno využití generativní AI pro podporu učení. Studentům je povinné:
\begin{enumerate}
\item uvést, kde a jakou AI použili (krátká deklarace v metadatech nebo poznámce),
\item doplnit odkazy/citace k relevantním automatizovaným výstupům,
\item nezveřejňovat osobní údaje třetích stran; dodržujte pokyny vyučujícího a fakultní GDPR.
\end{enumerate}
Používání AI k obcházení originální práce (plagiátorství) je zakázáno.
\end{quote}

Rozšířená verze (body pro detailní sylabus): \begin{itemize}
\item Rozsah: specifikujte povolené/nedovolené použití; Deklarace: stručná poznámka s názvem nástroje a popisem ovlivněných částí \cite{kb_ai_101_tipu_pdf_04e4a115_page_36_1}.
\item Dokumentace: u procesních úkolů vyžadujte meziverze/drafty/logy; Ochrana dat: zákaz vkládání citlivých údajů, uveďte repozitář a kontakty.
\end{itemize}

Praktický vzor deklarace a citace: \begin{quote}Deklarace: nástroj: <název>; použití: <stručný popis>; generované části upraveny a ověřeny. Citace: Generovaný text: <popis>. Vytvořeno: <název nástroje>, datum: <DD.MM.RRRR>. Ověřeno autorem.\end{quote}

Krátké poznámky pro garanty: \begin{itemize}
\item Při povolení AI v testech uveďte pravidla validace a lidskou kontrolu \cite{kb_ai_101_tipu_pdf_04e4a115_page_36_1}; pro edukační nástroje uveďte licenční poznámky \cite{kb_ai_101_tipu_pdf_04e4a115_page_15_2}.
\item Uveďte kontakty pro podporu a hlášení incidentů.
\end{itemize}

Poznámka k ochraně osobních údajů a DPIA: \begin{quote}Zvažte minimalizaci dat, pseudonymizaci a konzultaci DPIA/DPA před nasazením nástrojů zpracovávajících osobní údaje.\end{quote}

Rychlé nasazení do sylabu: \begin{enumerate}\item Zkopírujte krátkou verzi do "Organizační informace" a vložte pole pro deklaraci AI do hodnocení; \item Uložte příklady povoleného/zakázaného použití v repozitáři (ai.vut.cz) a při sběru citlivých dat kontaktujte GDPR/IT tým.\end{enumerate}%

%
\section{Vzor DPIA (posouzení dopadů na ochranu osobních údajů)}%
Šablona DPIA připravená pro nasazení AI ve výuce: popis zpracování, identifikace rizik (GDPR), navrhovaná opatření a doporučený postup schvalování na úrovni katedry/FS/IT.%


%
\paragraph{Krátké shrnutí (executive summary)}
\begin{itemize}
\item Název: \underline{\hspace{8cm}}; Garant: \underline{\hspace{6cm}}; Cíl: \underline{\hspace{6cm}}; Riziko: \underline{\hspace{3cm}}
\end{itemize}
\paragraph{General background knowledge:} DPIA je proces identifikace a hodnocení rizik pro práva a svobody subjektů údajů při plánovaném zpracování.
\paragraph{1) Popis zpracování (co, kdo, proč)}
\begin{itemize}
\item Funkce nástroje ve výuce (asistent, RAG); kategorie údajů (jméno, email, výsledky, nahrávky apod.); účel a právní základ; tok dat: sběr→ukládání→přístup→retence→mazání (lokace: VUT cloud/externí).
\end{itemize}
\paragraph{2) Identifikace a hodnocení rizik}
Navržené záznamy (u každého uveďte příčinu, opatření a zodpovědnou roli):
\begin{itemize}
\item Neoprávněný přístup — Pravděpodobnost: \underline{\hspace{2cm}} — Dopad: \underline{\hspace{2cm}}
\item Únik citlivých údajů při indexaci — Pravděpodobnost: \underline{\hspace{2cm}} — Dopad: \underline{\hspace{2cm}}
\item Nesprávné rozhodnutí modelu ovlivňující hodnocení — Pravděpodobnost: \underline{\hspace{2cm}} — Dopad: \underline{\hspace{2cm}}
\end{itemize}
\paragraph{3) Navrhovaná opatření (kontroly)}\begin{itemize}
\item Minimalizace dat; pseudonymizace/anonymizace před ingestem. \textit{(general background knowledge)}
\item IAM: SSO + role‑based access; auditní logy; šifrování v klidu i přenosu; uvést data residency; smluvní zajištění (DPA).
\item Lidská validace AI u rozhodnutí ovlivňujících práva studentů; retence/mazání; testy leakage a bezpečnostní testy před pilotem.
\end{itemize}
\paragraph{4) Závěr a doporučení}\begin{itemize}
\item Shrnutí residualního rizika: \underline{\hspace{6cm}}; Doporučení: povolit pilot / upravit návrh / neimplementovat (uveďte podmínky)
\end{itemize}
\paragraph{5) Schvalovací workflow}\begin{enumerate}
\item Konzultace: garant + metodik + GDPR; IT technická kontrola; právní review; schválení vedoucím; pilot, monitoring a revize DPIA.
\end{enumerate}
\paragraph{6) Kontrolní checklist (před pilotem)}\begin{itemize}
\item Určen GDPR/IT kontakt? \underline{\hspace{4cm}}; Kategorie a minimalizace definovány? \underline{\hspace{2cm}}; Uzavřena DPA? \underline{\hspace{2cm}}; Auditní logy, retence a plán lidské validace? \underline{\hspace{2cm}}
\end{itemize}
\paragraph{Příloha: vzorový záznam rizika}\begin{itemize}
\item Název: indexace osobních poznámek; Popis: ingest studentských poznámek; Dopad: možné zveřejnění citlivých informací; Opatření: pseudonymizace, zákaz uploadu osobních údajů, revize před publikací; Stav: návrh/implementováno/ověřeno
\end{itemize}%

%
\section{Základní RAG blueprint pro kurzy}%
Jednostránkový plán (blueprint) pro nasazení Retrieval-Augmented Generation ve výuce: architektura, datové zdroje, bezpečnostní a etické body, role vyučujícího a doporučené scénáře použití.%


%
\paragraph{Stručný přehled} Krátký, opakovatelný PoC blueprint pro nasazení Retrieval‑Augmented Generation (RAG) v kurzu; cíl: rychle ověřit pedagogický přínos při řízení rizik a ochraně dat.

\paragraph{Krokový PoC workflow (katedra)} \begin{enumerate}
  \item Definovat pedagogický cíl, typy dotazů, výstupy a metriky.
  \item Klasifikovat zdroje (interní / veřejné / PII) a vyloučit citlivé soubory.
  \item Navrhnout anonymizaci, pseudonymizaci a pravidla retence/mazání.
  \item Ingest/index: export → preprocessing/chunking → embeddings → index (on‑prem/hybrid/cloud).
  \item Nastavit retrieval (filtrování podle kurzu/role, max. kontext, re‑ranking) a vybrat LLM.
  \item Testovat přesnost a bezpečnost: lidská validace, leakage testy; iterovat prompt/retrieval.
  \item Pilot s malou skupinou; sběr zpětné vazby a úpravy politik.
  \item Zavést provozní procedury: audit, DPA/DPIA, zodpovědnosti, rollback a monitoring.
\end{enumerate}

\paragraph{Základní architektonické komponenty} \begin{itemize}
  \item Zdrojové soubory: slides, skripta, zadání, FAQ, lab návody.
  \item Preprocessing/chunking, embeddings (volba podle bezpečnosti) a index/retrieval (filtrování, re‑ranking).
  \item LLM + prompt layer s guardraily a monitorování/audit logy.
\end{itemize}

\paragraph{Bezpečnostní a etické body} \begin{itemize}
  \item Proveďte DPIA; zajistěte DPA a jasné data‑residency podmínky.
  \item Minimalizujte ingest: nevkládat student. poznámky bez pseudonymizace; vyloučit fotografie/medical.
  \item IAM a audit: SSO + role‑based přístup; auditní logy pro IT/garanta.
  \item Lidská validace u rozhodnutí s dopadem; leakage testing/red‑teaming před nasazením.
\end{itemize}

\paragraph{Role vyučujícího} Kurátor zdrojů, ověřovatel kvality indexu, definovatel retrievalu a poskytovatel lidské verifikace; zodpovědný za komunikaci se studenty.

\paragraph{Doporučené scénáře použití} FAQ asistent, shrnutí přednášek, vyhledávání ve skriptech, nápověda k labům s odkazy.

\paragraph{Rychlé checklisty před pilotem} \begin{itemize}
  \item Pedagogický cíl a metriky? \underline{\hspace{4cm}} \item Klasifikace zdrojů a anonymizace? \underline{\hspace{3cm}}
  \item DPA/DPIA schváleny? \underline{\hspace{3cm}} \item Auditní logy a role‑based přístup nastaveny? \underline{\hspace{3cm}}
  \item Plán lidské validace existuje? \underline{\hspace{3cm}}
\end{itemize}

\noindent Poznámka: blueprint navržen jako krátký, opakovatelný plán přizpůsobitelný předmětu; po pilotu doplnit technickou dokumentaci a SLA.%

%
\section{Glossář pojmů}%
Krátký slovník klíčových pojmů (generativní AI, LLM, tokeny, embeddings, RAG, DPIA apod.) jako rychlá reference pro pedagogy při přípravě materiálů a komunikaci se studenty.%


%
Krátký slovník klíčových pojmů určený pro pedagogy — rychlá reference při přípravě materiálů a komunikaci se studenty, s krátkými definicemi a praktickými poznámkami.

\begin{description}
  \item[Generativní AI] (general background knowledge) Algoritmy a modely, které vytvářejí nový obsah (text, obrázky, kód apod.) na základě naučených vzorů. V kontextu výuky jde především o nástroje, které studenti i vyučující mohou využít jako asistenty při tvorbě materiálů, vysvětlování či generování cvičení; je dobré rozlišovat tvorbu pod vedením učitele a plně automatické generování. Doporučuje se určit jasná pravidla použití a ověřování výsledků.
  \item[LLM — velké jazykové modely] (general background knowledge) Modely trénované na rozsáhlých korpusech textu, které předpovídají další tokeny a lze je použít pro generování textu, sumarizace či konverzační asistenci. V pedagogickém nasazení sledujte limity ve faktické přesnosti a vhodnosti obsahu. Učitelé by měli nastavit explicitní instrukce a validační kroky.
  \item[Tokeny] (general background knowledge) Základní jednotky textu, které model zpracovává (mohou být část slova, celé slovo nebo více slov v závislosti na tokenizéru). Důležité pro řízení délky vstupu/výstupu modelu a pro účtování spotřeby; při přípravě úloh plánujte rozpočet tokenů. V praxi to ovlivňuje délku promptů a rozsah poskytovaných kontextů.
  \item[Embeddings] (general background knowledge) Vektorové reprezentace textu/slov, které zachycují jejich význam a umožňují měřit podobnost mezi dokumenty; základní komponenta pro vyhledávání relevantních zdrojů nebo clustering odpovědí. Prakticky se používají pro vyhledávání v databázi učebních materiálů a pro návrh personalizovaných doporučení.
  \item[RAG — Retrieval‑Augmented Generation] (general background knowledge) Vzor, kdy se před generováním textu vyhledají relevantní dokumenty (retrieval) a tyto dokumenty se použijí jako kontext pro LLM, čímž se zvyšuje přesnost a faktická podpora odpovědí. (Viz také „Základní RAG blueprint pro kurzy“ v předchozí kapitole.) Vhodný při tvorbě fact‑checked materiálů a při omezování halucinací.
  \item[DPIA — posouzení dopadů na ochranu osobních údajů] (general background knowledge) Dokumentovaný postup hodnocení rizik zpracování osobních údajů (GDPR), doporučený před nasazením AI řešení, která pracují se studentskými daty. Obsahuje identifikaci rizik, návrh mitigací a záznamy o rozhodnutích. Výsledek DPIA pomáhá definovat technické a organizační opatření.
  \item[Minecraft Education] Výuková verze Minecraftu s připravenými světy a aktivitami vhodnými i pro ilustraci principů strojového učení a AI; některé světy pro „Hour of Code“ jsou zdarma, další jsou dostupné v rámci školních licencí Microsoft 365 A3/A5 \cite{kb_ai_101_tipu_pdf_04e4a115_page_15_2}. Lze jej využít pro projektové vyučování a interaktivní demonstrace, které zvyšují zapojení studentů.
  \item[Microsoft 365 Copilot / Copilot Chat] Chatbot založený na modelech GPT dostupný ve školních variantách Office 365; ve školním prostředí je k dispozici webová i mobilní verze Copilot Chat a placený doplněk Microsoft 365 Copilot rozšiřuje funkce přímo v aplikacích Word, Excel, PowerPoint apod. \cite{kb_ai_101_tipu_pdf_04e4a115_page_3_1}. Využití zahrnuje asistenční psaní, návrh prezentací a analýzu dat; při zavádění kontrolujte nastavení ochrany dat.
  \item["Pokrok v matematice"] Příklad nástroje pro generování a nasazení cvičných úloh, který umožňuje učiteli vybrat oblast, generovat příklady, nechat studenty odeslat postupy řešení; systém předopraví zadání a učitel finalizuje hodnocení, přičemž jsou dostupné přehledy výkonu jednotlivých žáků \cite{kb_ai_101_tipu_pdf_04e4a115_page_8_1}. Hodí se pro individualizovanou práci i zpětnou vazbu a pro monitoring pokroku třídy.
  \item[OneNote / math.microsoft.com] Praktické nástroje pro technické obory: OneNote umí převádět rukopisné rovnice do digitální podoby (včetně výpočtů a grafů) a web math.microsoft.com agreguje nástroje pro řešení matematických úloh \cite{kb_ai_101_tipu_pdf_04e4a115_page_8_1}. Tyto nástroje usnadňují sběr zadání a automatickou kontrolu postupů, což šetří čas učitele.
  \item[Halucinace (hallucination)] (general background knowledge) Situace, kdy model vygeneruje fakticky nepřesné nebo vymyšlené informace; riziko vyžaduje ověřování výstupů a lidskou validaci ve výuce. Doporučuje se zadávat ověřitelné zdroje a kontrolní otázky a učit studenty kritickému ověřování.
  \item[Bias / zkreslení] (general background knowledge) Sklony v datech nebo chování modelu, které mohou vést k nespravedlivým nebo nevhodným závěrům; při použití v kurzech je třeba pozornost při výběru a revizi obsahu. Zapojení různé perspektivy a revizní proces pomáhá odhalit a opravit zkreslení, a je vhodné mít transparentní auditní postupy.
  \item[GDPR a ochrana soukromí] (general background knowledge) Základní požadavky na zpracování osobních údajů (minimalizace, účelové omezení, anonymizace/pseudonymizace a dokumentace zpracování) — důležitá součást plánování nasazení AI ve výuce. Zahrnuje i informovaný souhlas studentů a bezpečné ukládání dat, včetně omezení přístupu a pravidelného vymazávání nepotřebných údajů.
\end{description}

Poznámka: položky označené jako "general background knowledge" jsou záměrně stručné; podrobné vysvětlení (včetně ilustrací tokenizace, vizualizace embeddings nebo RAG‑pipeline) najdete v odpovídajících kapitolách příručky. Tento slovník slouží jako startovací bod pro rychlou orientaci a plánování výuky a pro diskusi o praktických pravidlech nasazení v konkrétních třídních scénářích.%

%
\section{Kontakty a podpora VUT}%
Seznam interních kontaktů pro podporu (metodické oddělení, IT, právní/GDPR, e-learning, národní kontakty), doporučené SLA pro žádosti o pomoc a postup při hlášení incidentů souvisejících s AI.%


%
Krátký přehled kontaktů a postupů pro získání podpory při nasazení a provozu AI‑nástrojů na VUT; šablona pro doplnění konkrétních interních kontaktů a odkazů do fakultních ticketovacích systémů. (general background knowledge)
\begin{itemize}
\item Doporučené kontaktní role (doplnit e‑maily/aliasy a tel.): \begin{itemize}
  \item Metodické oddělení — didaktika, šablony, prompty.
  \item IT / DevOps — integrace (LMS, SSO), přístupy, incidenty.
  \item Právní / GDPR — DPIA, DPA, ochrana osobních údajů.
  \item E‑learning / LMS a Data Steward / IT Security Owner — Moodle, pluginy, bezpečnost dat, logy, incident response. (general background knowledge)
\end{itemize}
\item Template pro ticket — identifikace kurzu/katedry; popis požadavku a dopad; informace o datech (osobní/citlivé); přílohy; preferovaný termín. (general background knowledge)
\item Doporučené SLA (příklad): Kritické: potvrzení do 4 h, komunikace každých 4–8 h; Vážné: potvrzení do 24 h, plán do 3 pracovních dnů; Běžné: potvrzení do 3 pracovních dnů, řešení do 10 dnů.
\item Hlášení bezpečnostního incidentu: izolovat systémy; nahlásit ticketem jako bezpečnostní incident; informovat IT Security Owner a právní; dokumentovat timeline a logy; provést post‑mortem. (general background knowledge)
\item Praktické poznámky: doplnit kontakty do sylabu; používat jednotný ticketovací systém; konzultovat požadavky se studentskými daty s GDPR kontaktem. (general background knowledge)
\end{itemize}
Další zdroje, šablony ticketů a vzor DPIA jsou v online repozitáři příručky — viz repozitář a video tutoriály (doplnit interní odkaz, např. ai.vut.cz). (general background knowledge)%

%
\section{Bibliografie a anotované odkazy na univerzitní guideline}%
Kompaktní bibliografie a anotované odkazy na relevantní univerzitní guideline a patterny (UCL, TU Delft, další příklady dobré praxe), připraveno jako výchozí seznam pro další studium a odkazy v online repozitáři.%


%
Tento souhrnný seznam uvádí vybrané zdroje a univerzitní guideliney užitečné pro další studium a adaptaci postupů do praxe na VUT. U zdrojů, které jsou přímo citovány z interního materiálu „AI – 101 tipů“, uvádíme poznámky a doporučené využití v kurzech.

\begin{itemize}
  \item Microsoft 365 / Office 365 a Copilot — přehled nástrojů (Chat, integrované funkce v Word/Excel/PowerPoint) a poznámka, že školní varianty (A1 / A3/A5) obsahují specifické edukační funkce; Copilot Chat je dostupný ve školních verzích a Copilot může rozšířit běžné aplikace o generativní funkce \cite{kb_ai_101_tipu_pdf_04e4a115_page_3_1}.
    \begin{itemize}
      \item Praktické využití: integrace do přípravy materiálů, rychlé sumarizace, generování kvízů a nápovědy pro studenty (využít v pilotu s jasnými pravidly pro akademickou integritu). % (podrobnosti viz repozitář).
    \end{itemize}

  \item Pokrok v matematice (součást Vzdělávacích akcelerátorů Microsoft) — nástroj pro generování a správu procvičovacích úloh, sběr postupů žáků a analýzu silných / slabých stránek ve třídě; v rámci rodiny Microsoft nástroje doplňuje OneNote (převod rukopisu a výpočty) a webová služba math.microsoft.com \cite{kb_ai_101_tipu_pdf_04e4a115_page_8_1}.
    \begin{itemize}
      \item Praktické využití: vhodné pro automatizované cvičení doma i diagnostiku chyb; při nasazení v kurzech Zajistit, aby studenti odevzdávali i postup (procesní hodnocení).
    \end{itemize}

  \item Minecraft Education — edukační světy a „Hour of Code“ aktivity vhodné pro názorné seznámení studentů s principy strojového učení a AI; některé lekce jsou zdarma, širší obsah je dostupný v rámci školních licencí Microsoft 365 A3/A5 \cite{kb_ai_101_tipu_pdf_04e4a115_page_15_2}.
    \begin{itemize}
      \item Praktické využití: hands‑on ukázky pro základní i mezioborové kurzy (viz dostupné světy zaměřené na AI a Hodinu kódu).
    \end{itemize}
\end{itemize}

General background knowledge: Pro inspiraci a implementační patterny je užitečné prostudovat guideliney a design patterny z jiných univerzit (např. UCL, TU Delft, Harvard CS50, Georgia Tech) — tyto zdroje a přenositelné příklady jsou shromážděny a anotovány v online repozitáři příručky (viz ai.vut.cz jako living repository). (general background knowledge)%

%
\section{Repozitář, video tutoriály a mechanismus aktualizací}%
Instrukce pro přístup k online repozitáři s šablonami a video tutoriály, popis mechanismu pro kvartální aktualizace a sběr zpětné vazby od učitelů (živý dokument).%


%
Repozitář a video tutoriály (přístup a provoz)
Online repozitář slouží jako „living document“ s šablonami, video‑návody a verzovanými materiály pro rychlé nasazení.
Přístup a základní orientace (krátký návod)
\begin{itemize}
\item Kde: centrální vstup ai.vut.cz; doporučená struktura: \texttt{/prompts}, \texttt{/templates}, \texttt{/videos}, \texttt{/case‑studies}, \texttt{/blueprints}.
\item Rychlý start: stáhněte ZIP nebo klonujte, prostudujte \texttt{README.md} a otevřete \texttt{/videos} s titulky a časovými razítky.
\end{itemize}
Video tutoriály — doporučené postupy pro učitele
\begin{itemize}
\item Každé video má abstrakt, délku a klíčové časové body.
\item Témata zahrnují praktické návody (např. Minecraft Education) \cite{kb_ai_101_tipu_pdf_04e4a115_page_15_2} a integrace Microsoft (OneNote, Copilot, živé titulky) \cite{kb_ai_101_tipu_pdf_04e4a115_page_8_1, kb_ai_101_tipu_pdf_04e4a115_page_16_1}.
\end{itemize}
Mechanismus kvartálních aktualizací a sběru zpětné vazby
\begin{itemize}
\item Příspěvky průběžně; formální vydání kvartálně. Použijte issue šablonu nebo pull request s popisem pedagogického dopadu a testováním.
\item Ediční tým kontroluje etiku, GDPR a kvalitu; každé vydání obsahuje release notes.
\end{itemize}
Podpora a eskalace
\begin{itemize}
\item Pro dotazy a chyby použijte issues nebo formulář „Kontakt / Podpora“; technické/institucionální otázky směřujte na lokální IT/metodické kontakty (kap. 9.7).
\end{itemize}
Poznámka o ověřeném obsahu
\begin{itemize}
\item Videa a šablony odkazují na ověřené příklady a v popisu vždy uveďte doporučené licence a omezení z hlediska ochrany osobních údajů \cite{kb_ai_101_tipu_pdf_04e4a115_page_15_2, kb_ai_101_tipu_pdf_04e4a115_page_8_1, kb_ai_101_tipu_pdf_04e4a115_page_16_1}.
\end{itemize}%

%
\chapter{ZÁVĚR}%
Shrnutí klíčových myšlenek příručky, důraz na výhody AI pro kvalitu a efektivitu výuky a výzva k odpovědnému experimentování. Nabízí další kroky, pozvánku do online repository a formulář pro zpětnou vazbu. Motivuje k průběžné evaluaci dopadů a sdílení praxe napříč VUT.%


%
Tato příručka shrnuje principy a konkrétní nástroje pro zodpovědné a efektivní začlenění generativní AI do výuky na VUT. Jejím cílem není poskytovat vyčerpávající technický manuál, ale nabídnout praktický rámec kroků, které lze okamžitě aplikovat v kurzech i na úrovni kateder. Zaměřuje se na principy bezpečného nasazení, ověřitelnosti výsledků a minimalizace rizik při zachování didaktické hodnoty výuky. Hlavní sdělení, která si ponechte jako praktický rámec pro další kroky, jsou tato:

- AI je prostředek pro zlepšení kvality výuky — především pro personalizaci, rychlejší zpětnou vazbu a tvorbu variabilních cvičení — nikoli náhrada didaktické práce učitele. Role pedagoga zůstává v navrhování výukových cílů, ověřování výsledků učení a rozhodování o etickém rámci použití nástrojů (general background knowledge). Pedagogové také určují, které části úlohy mohou být delegovány na nástroj a které vždy musí projít lidskou kontrolou.
- Začněte malými, měřitelnými piloty a ověřujte dopady na učení i integritu hodnocení. Krokové piloty minimalizují rizika a umožňují iterativní zlepšování na základě konkrétních metrik a empirických pozorování (general background knowledge). Pro piloty doporučujeme stanovit kontrolní skupinu nebo předpilotní baseline, aby byly výsledky porovnatelné a interpretovatelné.

Konkrétní doporučení a ukázky, které najdete v příručce, jsou zaměřené na rychlé nasazení v hodině i na úrovni katedry. Například použití vzdělávacího prostředí Minecraft Education jako demonstračního nástroje pro vysvětlení principů strojového učení a základních konceptů AI je ověřený praktický přístup, který lze rychle implementovat ve výuce \cite{kb_ai_101_tipu_pdf_04e4a115_page_15_2}. Podobně nástroje z rodiny Microsoft (funkce „Pokrok v matematice“ či „Pokrok ve čtení“) ilustrují, jak lze AI využít pro generování cvičení, sběr postupů a sledování individuálního pokroku žáků/studentů \cite{kb_ai_101_tipu_pdf_04e4a115_page_8_1, kb_ai_101_tipu_pdf_04e4a115_page_7_1}. V příručce jsou k těmto případům doplněny konkrétní lekce, prompt‑šablony a návrhy metrik ke sběru dat.

Další kroky — doporučené krátké akce (praktický checklist)
\begin{itemize}
\item Vyberte malý pilotní kurz (1–2 týmy/sekce) a definujte jasné metriky úspěchu (zapojení studentů, kvalita zpětné vazby, čas vyučujících). Stanovte také časový rámec pilotu a zodpovědné osoby za sběr dat a vyhodnocení. (general background knowledge) Doporučujeme vymezit minimální a cílové hodnoty metrik, aby bylo možné rozhodnout o rozšíření projektu.
\item Použijte šablony a prompt‑knihovnu z repozitáře pro rychlé vytvoření testovacích aktivit a rubrik; validujte každý generovaný materiál lidskou kontrolou před nasazením. Doporučujeme záznam kontroly (kdo, kdy, jaké změny) pro zajištění reprodukovatelnosti a odpovědnosti. (general background knowledge) Přidejte do repozitáře verze šablon a log změn pro auditní stopu.
\item Zajistěte transparentnost: upravte sylaby tak, aby obsahovaly pravidla pro používání AI a povinnost deklarace (viz kap. 2 a přílohy). Do sylabu uveďte také konkrétní příklady povoleného a zakázaného použití, aby studenti přesně věděli, jak jednat. (general background knowledge) Doporučujeme rovněž formulovat sankce a kroky v případě porušení těchto pravidel.
\item Pro piloty využijte ilustrativní nástroje (např. Minecraft Education pro demonstrace, Microsoft funkce pro procvičování) podle doporučení v kapitolách a videonávodech \cite{kb_ai_101_tipu_pdf_04e4a115_page_15_2, kb_ai_101_tipu_pdf_04e4a115_page_8_1}. Využijte ukázkové lekce a připravené scénáře z repozitáře pro rychlejší start a srovnání výsledků. Přizpůsobte scénáře konkrétním učebním cílům a věkové skupinám.
\item Dokumentujte DPIA/DPA a bezpečnostní opatření tam, kde jsou zpracovávána osobní data studentů; zapojte IT/právní oddělení k revizi. Ujistěte se o minimalizaci sdílených osobních údajů a o šifrování či jiných technických opatřeních tam, kde to je nutné. (general background knowledge) V případě potřeby zařiďte informovaný souhlas studentů nebo jejich zákonných zástupců.
\item Monitorujte etické aspekty a bias v datech: zaznamenávejte známé omezení nástrojů a informujte o nich studenty i kolegy, aby bylo používání transparentní.
\end{itemize}

Praktické provozní doporučení pro piloty
\begin{itemize}
\item Zaveďte krátké školení pro vyučující a asistenty o fungování používaných nástrojů, správném promptování a základních rizicích (bias, hallucinations). Součástí školení by měla být i praktická část s ukázkami ověřování výstupů.
\item Určete jasné postupy pro ověřování autenticity studentské práce a pro řešení podezření z porušení akademických pravidel při použití AI. Doporučujeme kombinovat techniky (plagiátorský software, ústní obhajoby, laboratorní záznamy) pro robustní ověření.
\item Sledujte časovou náročnost: měřte, zda AI řešení skutečně šetří čas vyučujících, a zvažte celkovou režii související s validací výstupů. Zaznamenávejte skutečné časy na přípravu, korekci a správu aktivit.
\item Zaznamenávejte zpětnou vazbu studentů systematicky, aby bylo možné upravovat prompt šablony a aktivity na základě uživatelských zkušeností. Využijte krátké anonymní ankety i fokusové rozhovory pro hlubší vhled.
\item Zaveďte malé kontrolní mechanismy kvality: náhodné audity generovaných materiálů, peer review mezi vyučujícími a běžné aktualizace prompt‑knihovny.
\end{itemize}

Pozvánka k participaci a udržování praxe
(general background knowledge) Tento materiál je navržen jako „living document“ — sdílejte svoji zkušenost, přispívejte šablonami a případovými studiemi, a účastněte se kvartálních aktualizací repozitáře a video‑tutoriálů. Zapojení širší komunity vyučujících urychlí vznik ověřených vzorů a pomůže rozšířit osvědčené postupy napříč fakultami. Navrhujeme také pořádat interní workshopy pro předávání znalostí mezi katedrami a pravidelné evaluační setkání po ukončení každého pilotu. Do repozitáře ukládejte nejen úspěšné případy, ale i neúspěchy a klíčová rizika, aby ostatní mohli těžit z vaší zkušenosti.

Závěrem: experimentujte zodpovědně, měřte výsledky a sdílejte poznatky. AI může posunout kvalitu výuky, pokud zůstane nástrojem pedagoga — nikoli náhradou. Dodržováním zásad transparentnosti, ochrany dat a pečlivého ověřování výstupů lze minimalizovat rizika a maximalizovat přínos pro studenty i vyučující. Udržujte dialog v rámci komunity, aktualizujte praktiky podle nových poznatků a považujte piloty za iterativní proces učení pro všechny zúčastněné.%

%
\bibliographystyle{unsrt}%
\bibliography{refs}%
\end{document}